%%%%%%%%%%%%%%%%%%%%%%%%%%%%%%%%%%%%%%%%%%%%%%%%%%%%%%%%%%%%%
%% CATEGORY THEORY — Preamble + Chapters 1–2
%% Categories, Functors, Natural Transformations / Duality
%%%%%%%%%%%%%%%%%%%%%%%%%%%%%%%%%%%%%%%%%%%%%%%%%%%%%%%%%%%%%

\documentclass[a4paper,12pt]{book}

%% ---- Encoding and language ----
\usepackage[utf8]{inputenc}
\usepackage[T1]{fontenc}
\usepackage[english]{babel}

%% ---- Mathematics ----
\usepackage{amsmath}
\usepackage{amssymb}
\usepackage{amsthm}
\usepackage{mathtools}
\usepackage{stmaryrd}

%% ---- Coloured boxes ----
\usepackage[theorems,skins,breakable]{tcolorbox}

%% ---- Diagrams and graphics ----
\usepackage{tikz}
\usetikzlibrary{calc}
\usepackage{tikz-cd}
\usepackage{pgfplots}
\pgfplotsset{compat=1.18}
\usepackage{graphicx}
\usepackage{float}
\usepackage{caption}

%% ---- Lists, tables ----
\usepackage{enumitem}
\usepackage{booktabs}
\usepackage{array}
\usepackage{longtable}

%% ---- Cross-references and indexing ----
\usepackage{hyperref}
\hypersetup{
    colorlinks=true,
    linkcolor=blue!70!black,
    citecolor=green!50!black,
    urlcolor=blue!60!black
}
\usepackage{cleveref}
\usepackage{makeidx}
\usepackage{minitoc}

\makeindex

%% ============================================================
%%  Math operators
%% ============================================================
\DeclareMathOperator{\Hom}{Hom}
\DeclareMathOperator{\End}{End}
\DeclareMathOperator{\Aut}{Aut}
\DeclareMathOperator{\im}{im}
\DeclareMathOperator{\Id}{Id}
\DeclareMathOperator{\Ker}{Ker}
\DeclareMathOperator{\coker}{coker}
\DeclareMathOperator{\Ob}{Ob}
\DeclareMathOperator{\Mor}{Mor}
\DeclareMathOperator{\Fun}{Fun}
\DeclareMathOperator{\Nat}{Nat}
\DeclareMathOperator{\colim}{colim}
\DeclareMathOperator{\eq}{eq}
\DeclareMathOperator{\coeq}{coeq}

%% ---- Named categories ----
\newcommand{\Set}{\mathbf{Set}}
\newcommand{\Grp}{\mathbf{Grp}}
\newcommand{\Ab}{\mathbf{Ab}}
\newcommand{\Ring}{\mathbf{Ring}}
\newcommand{\Top}{\mathbf{Top}}
\newcommand{\Vect}{\mathbf{Vect}}
\newcommand{\Mod}{\mathbf{Mod}}
\newcommand{\Cat}{\mathbf{Cat}}

%% ---- Convenience commands ----
\newcommand{\op}{\mathrm{op}}
\newcommand{\id}{\mathrm{id}}
\newcommand{\N}{\mathbb{N}}
\newcommand{\Z}{\mathbb{Z}}
\newcommand{\Q}{\mathbb{Q}}
\newcommand{\R}{\mathbb{R}}
\newcommand{\C}{\mathbb{C}}
\newcommand{\K}{\mathbb{K}}

%% ============================================================
%%  Theorem environments (amsthm + tcolorbox styling)
%% ============================================================

%% --- amsthm declarations ---
\theoremstyle{definition}
\newtheorem{definition}{Definition}[section]
\newtheorem{example}[definition]{Example}
\newtheorem{exercise}[definition]{Exercise}

\theoremstyle{plain}
\newtheorem{theorem}[definition]{Theorem}
\newtheorem{proposition}[definition]{Proposition}
\newtheorem{lemma}[definition]{Lemma}
\newtheorem{corollary}[definition]{Corollary}

\theoremstyle{remark}
\newtheorem{remark}[definition]{Remark}
\newtheorem{notation}[definition]{Notation}

%% --- tcolorbox styling ---
\tcolorboxenvironment{definition}{
    enhanced, breakable,
    colback=blue!3, colframe=blue!50!black,
    boxrule=0.6pt, arc=2pt,
    left=6pt, right=6pt, top=4pt, bottom=4pt,
    fonttitle=\bfseries
}
\tcolorboxenvironment{example}{
    enhanced, breakable,
    colback=green!3, colframe=green!40!black,
    boxrule=0.5pt, arc=2pt,
    left=6pt, right=6pt, top=4pt, bottom=4pt,
    fonttitle=\bfseries
}
\tcolorboxenvironment{exercise}{
    enhanced, breakable,
    colback=orange!4, colframe=orange!60!black,
    boxrule=0.5pt, arc=2pt,
    left=6pt, right=6pt, top=4pt, bottom=4pt,
    fonttitle=\bfseries
}
\tcolorboxenvironment{theorem}{
    enhanced, breakable,
    colback=red!3, colframe=red!60!black,
    boxrule=0.7pt, arc=2pt,
    left=6pt, right=6pt, top=4pt, bottom=4pt,
    fonttitle=\bfseries
}
\tcolorboxenvironment{proposition}{
    enhanced, breakable,
    colback=red!2, colframe=red!45!black,
    boxrule=0.6pt, arc=2pt,
    left=6pt, right=6pt, top=4pt, bottom=4pt,
    fonttitle=\bfseries
}
\tcolorboxenvironment{lemma}{
    enhanced, breakable,
    colback=yellow!4, colframe=yellow!50!black,
    boxrule=0.5pt, arc=2pt,
    left=6pt, right=6pt, top=4pt, bottom=4pt,
    fonttitle=\bfseries
}
\tcolorboxenvironment{corollary}{
    enhanced, breakable,
    colback=purple!3, colframe=purple!40!black,
    boxrule=0.5pt, arc=2pt,
    left=6pt, right=6pt, top=4pt, bottom=4pt,
    fonttitle=\bfseries
}
\tcolorboxenvironment{remark}{
    enhanced, breakable,
    colback=gray!4, colframe=gray!50!black,
    boxrule=0.4pt, arc=2pt,
    left=6pt, right=6pt, top=4pt, bottom=4pt,
    fonttitle=\bfseries
}
\tcolorboxenvironment{notation}{
    enhanced, breakable,
    colback=cyan!3, colframe=cyan!40!black,
    boxrule=0.4pt, arc=2pt,
    left=6pt, right=6pt, top=4pt, bottom=4pt,
    fonttitle=\bfseries
}

%% ============================================================
%%  Title
%% ============================================================
\title{\Huge\textbf{Category Theory}\\[0.5em]
       \Large A Graduate Course}
\author{}
\date{}

%% ============================================================
\begin{document}
%% ============================================================

\begin{titlepage}
\begin{tikzpicture}[remember picture, overlay]
  \fill[left color=blue!15!white, right color=blue!5!white]
    (current page.south west) rectangle (current page.north east);
  \fill[color=orange!70!yellow]
    ([yshift=-2cm]current page.north west) rectangle ([yshift=-2.3cm]current page.north east);
  \fill[color=blue!80!black, rounded corners=8pt]
    ([xshift=3cm, yshift=-4cm]current page.north west) rectangle
    ([xshift=-3cm, yshift=-10cm]current page.north east);
  \node[anchor=center, text=white, font=\Huge\bfseries]
    at ([yshift=-6cm]current page.north) {Category Theory};
  \node[anchor=center, text=orange!80!yellow, font=\Large]
    at ([yshift=-7.5cm]current page.north) {Lecture Notes};
  \node[anchor=center, text=white!80, font=\large]
    at ([yshift=-8.8cm]current page.north) {M1--PhD --- 2025--2026};
  \node[anchor=center, text=blue!80!black, font=\Large\itshape]
    at ([yshift=-12cm]current page.north) {Ya\'e Ulrich Gaba};
  \draw[orange!70!yellow, line width=2pt]
    ([xshift=5cm, yshift=-13.5cm]current page.north west) --
    ([xshift=-5cm, yshift=-13.5cm]current page.north east);
  \node[anchor=center, text width=12cm, align=center, text=blue!60!black, font=\itshape]
    at ([yshift=-16cm]current page.north) {``I thought the notion of category was needed to axiomatise natural transformations.'' --- Saunders Mac\,Lane};
  \node[anchor=south, text=gray, font=\small]
    at ([yshift=1.5cm]current page.south) {\today};
\end{tikzpicture}
\end{titlepage}

\dominitoc
\tableofcontents

%% ============================================================
%%  PREFACE
%% ============================================================
\chapter*{Preface}\label{ch:preface}
\addcontentsline{toc}{chapter}{Preface}

Category theory provides a unifying language for mathematics.  Originally
introduced by Eilenberg and Mac\,Lane in the 1940s to formalise natural
transformations in algebraic topology, it has since become an
indispensable tool in algebra, geometry, logic, computer science, and
mathematical physics.

This text is designed for a one-semester graduate course.  We assume
familiarity with basic algebra (groups, rings, modules) and point-set
topology, but no prior knowledge of category theory itself.  The
exposition emphasises both the abstract framework and its concrete
manifestations: every definition is illustrated by examples drawn from
algebra, topology, and analysis.

\medskip
\noindent\textbf{Organisation.}
Chapters~1 and~2 lay the foundations: categories, functors, natural
transformations, and the duality principle.
Chapters~3 and~4 develop limits, colimits, and adjunctions.
Later chapters treat representability, the Yoneda lemma, Kan extensions,
abelian categories, and monoidal categories.

\medskip
\noindent\textbf{Exercises.}
Each chapter ends with a graded exercise set.  Working through these
is essential for mastering the material.

%% ============================================================
%%  NOTATION
%% ============================================================
\chapter*{Notation and Conventions}\label{ch:notation}
\addcontentsline{toc}{chapter}{Notation and Conventions}

\begin{longtable}{>{\raggedright}p{4cm} p{9.5cm}}
\toprule
\textbf{Symbol} & \textbf{Meaning} \\
\midrule
\endhead
$\N,\;\Z,\;\Q,\;\R,\;\C$ &
    Natural numbers (including~$0$), integers, rationals, reals,
    complex numbers \\[4pt]
$\K$ &
    An arbitrary field \\[4pt]
$\Set,\;\Grp,\;\Ab$ &
    Categories of sets, groups, abelian groups \\[4pt]
$\Ring,\;\Top,\;\Vect_\K$ &
    Categories of rings, topological spaces,
    $\K$-vector spaces \\[4pt]
$\Mod_R$ &
    Category of left $R$-modules \\[4pt]
$\Cat$ &
    Category of small categories \\[4pt]
$\Hom_\mathcal{C}(A,B)$ &
    Set of morphisms $A\to B$ in~$\mathcal{C}$ \\[4pt]
$\mathcal{C}^{\op}$ &
    Opposite (dual) category of~$\mathcal{C}$ \\[4pt]
$F\colon\mathcal{C}\to\mathcal{D}$ &
    Functor from $\mathcal{C}$ to $\mathcal{D}$ \\[4pt]
$\alpha\colon F\Rightarrow G$ &
    Natural transformation from $F$ to $G$ \\[4pt]
$\id_A$ &
    Identity morphism on~$A$ \\[4pt]
$g\circ f$ &
    Composition: first~$f$, then~$g$ \\[4pt]
$[\mathcal{C},\mathcal{D}]$ or $\Fun(\mathcal{C},\mathcal{D})$ &
    Functor category \\[4pt]
\bottomrule
\end{longtable}

\noindent
Throughout, ``category'' means a locally small category unless otherwise
stated.  We use the terms ``morphism'', ``arrow'', and ``map''
interchangeably.  All rings are assumed to have a multiplicative
identity.


%%%%%%%%%%%%%%%%%%%%%%%%%%%%%%%%%%%%%%%%%%%%%%%%%%%%%%%%%%%%
%%       CHAPTER 1: CATEGORIES, FUNCTORS, NATURAL         %%
%%                  TRANSFORMATIONS                        %%
%%%%%%%%%%%%%%%%%%%%%%%%%%%%%%%%%%%%%%%%%%%%%%%%%%%%%%%%%%%%

\chapter{Categories, Functors, and Natural Transformations}
\label{ch:cat-fun-nat}
\minitoc

\vspace{1em}

Categories, functors, and natural transformations constitute the three
fundamental notions of category theory.  A category axiomatises the idea
of ``objects with structure-preserving maps between them''.  A functor is
a morphism of categories, and a natural transformation is a morphism of
functors.  Already at this first level of abstraction one finds a
striking pattern: the passage from objects to morphisms between objects
repeats at every stage.

In this chapter we introduce these three notions carefully, give
copious examples, and establish basic terminology---monomorphisms,
epimorphisms, isomorphisms, initial and terminal objects---that will be
used throughout the course.

\index{category|(}


%% ============================================================
\section{Categories}\label{sec:categories}
%% ============================================================

\begin{definition}[Category]\label{def:category}
\index{category!definition}
\index{object}
\index{morphism}
\index{composition}
\index{identity morphism}
A \textbf{category}~$\mathcal{C}$ consists of the following data:
\begin{enumerate}[label=(\roman*)]
    \item A collection $\Ob(\mathcal{C})$ of \emph{objects}.
    \item For every ordered pair of objects $A,B\in\Ob(\mathcal{C})$,
          a set $\Hom_\mathcal{C}(A,B)$ of \emph{morphisms}
          (or \emph{arrows}) from~$A$ to~$B$.
          We write $f\colon A\to B$ to indicate
          $f\in\Hom_\mathcal{C}(A,B)$.
    \item For every triple of objects $A,B,C$, a \emph{composition law}
          \[
              \circ\colon \Hom_\mathcal{C}(B,C)\times
                          \Hom_\mathcal{C}(A,B)
                          \longrightarrow \Hom_\mathcal{C}(A,C),
              \qquad (g,f)\longmapsto g\circ f.
          \]
    \item For every object $A$, an \emph{identity morphism}
          $\id_A\in\Hom_\mathcal{C}(A,A)$.
\end{enumerate}
These data are subject to two axioms:
\begin{enumerate}[label=\textbf{(C\arabic*)}]
    \item \textbf{Associativity.}\index{associativity}
          For all $f\colon A\to B$, $g\colon B\to C$,
          $h\colon C\to D$,
          \[
              h\circ(g\circ f) = (h\circ g)\circ f.
          \]
    \item \textbf{Identity.}
          For all $f\colon A\to B$,
          \[
              f\circ\id_A = f = \id_B\circ f.
          \]
\end{enumerate}
\end{definition}

\begin{remark}\label{rem:hom-disjoint}
\index{hom-set}
We require the hom-sets to be pairwise disjoint: every morphism~$f$ has a
uniquely determined \emph{domain} (or \emph{source}) and \emph{codomain}
(or \emph{target}).  Some authors encode this by defining a morphism as a
triple $(f,A,B)$.
\end{remark}

\begin{notation}\label{not:hom-notation}
\index{hom-set!notation}
We write $\Hom_\mathcal{C}(A,B)$, or $\mathcal{C}(A,B)$, or
$\Mor_\mathcal{C}(A,B)$ interchangeably.
When the ambient category is clear we simply write $\Hom(A,B)$.
\end{notation}


%% ============================================================
\section{First examples of categories}\label{sec:first-examples}
%% ============================================================

\begin{example}[Algebraic categories]\label{ex:algebraic-cats}
\index{category!Set@$\Set$}
\index{category!Grp@$\Grp$}
\index{category!Ab@$\Ab$}
\index{category!Ring@$\Ring$}
\index{category!Mod@$\Mod_R$}
\index{category!Vect@$\Vect_\K$}
The following are categories:
\begin{enumerate}[label=(\roman*)]
    \item $\Set$: objects are sets, morphisms are functions.
    \item $\Grp$: objects are groups, morphisms are group homomorphisms.
    \item $\Ab$: objects are abelian groups, morphisms are group
          homomorphisms.
    \item $\Ring$: objects are (unital) rings, morphisms are ring
          homomorphisms (preserving~$1$).
    \item $\Mod_R$: for a ring~$R$, objects are left $R$-modules,
          morphisms are $R$-linear maps.
    \item $\Vect_\K$: objects are $\K$-vector spaces, morphisms are
          $\K$-linear maps.  This is the special case
          $\Mod_\K$.
\end{enumerate}
In each case, composition is the usual composition of functions and the
identity morphism on an object~$A$ is the identity function
$\id_A\colon A\to A$.
\end{example}

\begin{example}[Topological category]\label{ex:top-cat}
\index{category!Top@$\Top$}
$\Top$: objects are topological spaces, morphisms are continuous maps.
\end{example}

\begin{example}[Ordered sets as categories]\label{ex:poset-cat}
\index{category!poset}
\index{poset!as category}
\index{preorder!as category}
Let $(P,\le)$ be a preorder (a set equipped with a reflexive, transitive
relation).  Define a category~$\mathcal{C}_P$ by:
\begin{itemize}
    \item $\Ob(\mathcal{C}_P) = P$.
    \item For $a,b\in P$:
    \[
        \Hom(a,b) =
        \begin{cases}
            \{\ast\} & \text{if } a\le b, \\
            \varnothing & \text{otherwise.}
        \end{cases}
    \]
\end{itemize}
Transitivity gives composition; reflexivity gives identities.
If $(P,\le)$ is a partial order, distinct objects are never isomorphic.
\end{example}

\begin{example}[Monoids as categories]\label{ex:monoid-cat}
\index{monoid!as category}
\index{category!one-object}
A monoid $(M,\cdot,e)$ may be viewed as a category with a single
object~$\ast$ and $\Hom(\ast,\ast)=M$.
Composition is the monoid operation and the identity morphism is~$e$.

Conversely, every category with exactly one object is a monoid.
A group is a one-object category in which every morphism is invertible.
\end{example}

\begin{example}[The category $\Cat$]\label{ex:cat-of-cats}
\index{category!Cat@$\Cat$}
The category $\Cat$ has small categories as objects and functors
(defined in 
ef{sec:functors}) as morphisms.
\end{example}


%% ============================================================
\section{Small and locally small categories}\label{sec:size}
%% ============================================================

\begin{definition}[Size conditions]\label{def:size}
\index{small category}
\index{locally small category}
\index{large category}
A category~$\mathcal{C}$ is called:
\begin{enumerate}[label=(\roman*)]
    \item \textbf{small} if $\Ob(\mathcal{C})$ is a set
          (not a proper class);
    \item \textbf{locally small} if for every pair of objects
          $A,B$, the collection $\Hom_\mathcal{C}(A,B)$ is a set.
\end{enumerate}
A category that is not small is called \textbf{large}.
\end{definition}

\begin{remark}\label{rem:size-remarks}
Every small category is locally small.
The categories $\Set$, $\Grp$, $\Top$, etc.\ are locally small but not
small (their collection of objects is a proper class).
Every preorder viewed as a category is small provided the underlying
set is a set.
Throughout this text, ``category'' means ``locally small category''
unless otherwise stated.
\end{remark}


%% ============================================================
\section{Monomorphisms, epimorphisms, and isomorphisms}
\label{sec:mono-epi-iso}
%% ============================================================

\begin{definition}[Monomorphism]\label{def:mono}
\index{monomorphism}
\index{monic|see{monomorphism}}
A morphism $f\colon A\to B$ in a category~$\mathcal{C}$ is a
\textbf{monomorphism} (or is \emph{monic}) if for every pair of
morphisms $g_1,g_2\colon C\to A$,
\[
    f\circ g_1 = f\circ g_2 \implies g_1 = g_2.
\]
We denote a monomorphism by $f\colon A\rightarrowtail B$.
\end{definition}

\begin{definition}[Epimorphism]\label{def:epi}
\index{epimorphism}
\index{epic|see{epimorphism}}
A morphism $f\colon A\to B$ is an \textbf{epimorphism} (or is
\emph{epic}) if for every pair $h_1,h_2\colon B\to D$,
\[
    h_1\circ f = h_2\circ f \implies h_1 = h_2.
\]
We denote an epimorphism by $f\colon A\twoheadrightarrow B$.
\end{definition}

\begin{definition}[Isomorphism]\label{def:iso}
\index{isomorphism}
\index{inverse morphism}
A morphism $f\colon A\to B$ is an \textbf{isomorphism} if there exists a
morphism $g\colon B\to A$ such that
\[
    g\circ f = \id_A \qquad\text{and}\qquad f\circ g = \id_B.
\]
The morphism~$g$ is unique and is called the \emph{inverse} of~$f$,
written $f^{-1}$.  Two objects are \emph{isomorphic},
$A\cong B$\index{isomorphic objects}, if there exists an isomorphism
between them.
\end{definition}

\begin{proposition}\label{prop:iso-mono-epi}
\index{isomorphism!versus monomorphism and epimorphism}
Every isomorphism is both a monomorphism and an epimorphism.
\end{proposition}

\begin{proof}
Let $f\colon A\to B$ be an isomorphism with inverse~$g$.
If $f\circ g_1 = f\circ g_2$, apply~$g$ on the left:
$g_1 = g\circ f\circ g_1 = g\circ f\circ g_2 = g_2$,
so~$f$ is monic.  Dually, $f$ is epic.
\end{proof}

\begin{remark}\label{rem:mono-epi-not-iso}
\index{epimorphism!not surjective}
The converse fails in general.  In~$\Ring$, the inclusion
$\Z\hookrightarrow\Q$ is both monic and epic, but it is not an
isomorphism.
\end{remark}

\begin{example}[Mono and epi in concrete categories]\label{ex:mono-epi-concrete}
\index{monomorphism!in Set@in $\Set$}
\index{epimorphism!in Set@in $\Set$}
\begin{enumerate}[label=(\roman*)]
    \item In $\Set$: monomorphisms are injections, epimorphisms are
          surjections, isomorphisms are bijections.
    \item In $\Grp$: monomorphisms are injective homomorphisms,
          epimorphisms are surjective homomorphisms.
    \item In $\Top$: monomorphisms are injective continuous maps;
          epimorphisms are surjective continuous maps
          (but \emph{not} necessarily quotient maps).
\end{enumerate}
\end{example}

\begin{lemma}\label{lem:mono-epi-composition}
\index{monomorphism!composition}
\index{epimorphism!composition}
Let $f\colon A\to B$ and $g\colon B\to C$ be morphisms.
\begin{enumerate}[label=(\roman*)]
    \item If $g\circ f$ is monic, then $f$ is monic.
    \item If $g\circ f$ is epic, then $g$ is epic.
    \item If both $f$ and $g$ are monic (resp.\ epic), then so is
          $g\circ f$.
\end{enumerate}
\end{lemma}

\begin{proof}
(i) Suppose $g\circ f$ is monic and $f\circ h_1 = f\circ h_2$.
Then $g\circ f\circ h_1 = g\circ f\circ h_2$, hence $h_1=h_2$.
Parts~(ii) and~(iii) are similar.
\end{proof}


%% ============================================================
\section{Initial and terminal objects}\label{sec:initial-terminal}
%% ============================================================

\begin{definition}[Initial object]\label{def:initial}
\index{initial object}
\index{object!initial}
An object $I\in\mathcal{C}$ is \textbf{initial} if for every object
$A\in\mathcal{C}$ there exists a unique morphism $I\to A$.
\end{definition}

\begin{definition}[Terminal object]\label{def:terminal}
\index{terminal object}
\index{object!terminal}
An object $T\in\mathcal{C}$ is \textbf{terminal} if for every object
$A\in\mathcal{C}$ there exists a unique morphism $A\to T$.
\end{definition}

\begin{definition}[Zero object]\label{def:zero-object}
\index{zero object}
An object that is both initial and terminal is called a
\textbf{zero object}.
\end{definition}

\begin{proposition}\label{prop:initial-terminal-unique}
\index{initial object!uniqueness}
\index{terminal object!uniqueness}
If an initial (resp.\ terminal, resp.\ zero) object exists, it is unique
up to unique isomorphism.
\end{proposition}

\begin{proof}
Let $I$ and $I'$ be initial.  By the universal property there exist
unique morphisms $f\colon I\to I'$ and $g\colon I'\to I$.
Then $g\circ f\colon I\to I$ must equal $\id_I$ (uniqueness of the
morphism $I\to I$), and similarly $f\circ g = \id_{I'}$.
The argument for terminal objects is dual.
\end{proof}

\begin{example}[Initial and terminal objects]\label{ex:initial-terminal}
\index{empty set!as initial object}
\begin{enumerate}[label=(\roman*)]
    \item In $\Set$: the empty set $\varnothing$ is initial; any
          singleton $\{*\}$ is terminal.
    \item In $\Grp$ and $\Ab$: the trivial group $\{e\}$ is a
          zero object.
    \item In $\Ring$: $\Z$ is initial (unique ring homomorphism
          $\Z\to R$ sending $1\mapsto 1_R$); the zero ring~$\{0\}$
          is terminal.
    \item In $\Top$: the empty space is initial; any one-point space
          is terminal.
    \item In $\Vect_\K$: the zero space $\{0\}$ is a zero object.
\end{enumerate}
\end{example}

\begin{remark}\label{rem:zero-morphism}
\index{zero morphism}
If $\mathcal{C}$ has a zero object~$0$, then for any objects $A,B$ there
is a unique morphism $A\to 0\to B$, called the \emph{zero morphism}
and denoted $0_{AB}$ (or simply~$0$).
\end{remark}


%% ============================================================
\section{Functors}\label{sec:functors}
%% ============================================================

\index{functor|(}

\begin{definition}[Functor]\label{def:functor}
\index{functor!covariant}
\index{covariant functor|see{functor}}
Let $\mathcal{C}$ and $\mathcal{D}$ be categories.
A \textbf{(covariant) functor} $F\colon\mathcal{C}\to\mathcal{D}$
consists of:
\begin{enumerate}[label=(\roman*)]
    \item A map on objects:
          $A\mapsto F(A)$ for each $A\in\Ob(\mathcal{C})$.
    \item A map on morphisms: for each $f\colon A\to B$ in~$\mathcal{C}$,
          a morphism $F(f)\colon F(A)\to F(B)$ in~$\mathcal{D}$.
\end{enumerate}
These must satisfy:
\begin{enumerate}[label=\textbf{(F\arabic*)}]
    \item $F(\id_A) = \id_{F(A)}$ for every object $A$.
    \item $F(g\circ f) = F(g)\circ F(f)$ for every composable pair
          $f,g$.
\end{enumerate}
\end{definition}

\begin{definition}[Contravariant functor]\label{def:contravariant}
\index{functor!contravariant}
\index{contravariant functor}
A \textbf{contravariant functor} $F\colon\mathcal{C}\to\mathcal{D}$ is
a (covariant) functor $F\colon\mathcal{C}^{\op}\to\mathcal{D}$.
Equivalently, $F$ reverses the direction of morphisms:
$F(f)\colon F(B)\to F(A)$ when $f\colon A\to B$, and
$F(g\circ f) = F(f)\circ F(g)$.
\end{definition}

\begin{example}[Forgetful functors]\label{ex:forgetful}
\index{functor!forgetful}
\index{forgetful functor}
\begin{enumerate}[label=(\roman*)]
    \item $U\colon\Grp\to\Set$: sends a group to its underlying set
          and a homomorphism to the underlying function.
    \item $U\colon\Top\to\Set$: forgets the topology.
    \item $U\colon\Vect_\K\to\Ab$: forgets scalar multiplication,
          retaining only the additive group.
    \item $U\colon\Ring\to\Ab$: sends a ring to its underlying
          additive group.
\end{enumerate}
\end{example}

\begin{example}[Free functors]\label{ex:free-functor}
\index{functor!free}
\index{free functor}
\begin{enumerate}[label=(\roman*)]
    \item $F\colon\Set\to\Grp$: sends a set~$S$ to the free group on~$S$.
    \item $F\colon\Set\to\Vect_\K$: sends a set~$S$ to the vector
          space with basis~$S$.
    \item $F\colon\Set\to\Ab$: sends a set~$S$ to the free abelian
          group $\Z^{(S)} = \bigoplus_{s\in S}\Z$.
\end{enumerate}
\end{example}

\begin{example}[Power-set functor]\label{ex:powerset-functor}
\index{functor!power-set}
\index{power-set functor}
There are two natural functors $\Set\to\Set$ associated with the power
set:
\begin{enumerate}[label=(\roman*)]
    \item \textbf{Covariant:} $\mathcal{P}\colon\Set\to\Set$ sends a
          set~$S$ to $\mathcal{P}(S)$ and a function $f\colon S\to T$
          to the direct image $f_*\colon \mathcal{P}(S)\to\mathcal{P}(T)$,
          $A\mapsto f(A)$.
    \item \textbf{Contravariant:} $\mathcal{P}^{\op}\colon\Set^{\op}\to\Set$
          sends $f\colon S\to T$ to the pre-image
          $f^*\colon\mathcal{P}(T)\to\mathcal{P}(S)$,
          $B\mapsto f^{-1}(B)$.
\end{enumerate}
\end{example}

\begin{example}[Hom-functors]\label{ex:hom-functors}
\index{functor!hom}
\index{hom-functor}
For any locally small category~$\mathcal{C}$ and object $A\in\mathcal{C}$:
\begin{enumerate}[label=(\roman*)]
    \item The \textbf{covariant hom-functor}
          $\Hom_\mathcal{C}(A,-)\colon\mathcal{C}\to\Set$
          sends $B\mapsto\Hom(A,B)$ and $f\colon B\to C$ to
          $f_*\colon\Hom(A,B)\to\Hom(A,C)$, $g\mapsto f\circ g$.
    \item The \textbf{contravariant hom-functor}
          $\Hom_\mathcal{C}(-,B)\colon\mathcal{C}^{\op}\to\Set$
          sends $A\mapsto\Hom(A,B)$ and $f\colon A'\to A$ to
          $f^*\colon\Hom(A,B)\to\Hom(A',B)$, $g\mapsto g\circ f$.
\end{enumerate}
\end{example}

\begin{example}[Fundamental group functor]\label{ex:pi1-functor}
\index{functor!fundamental group}
\index{fundamental group!functor}
Let $\Top_*$ denote the category of pointed topological spaces and
base-point-preserving continuous maps.  The fundamental group construction
defines a functor
\[
    \pi_1\colon\Top_*\longrightarrow\Grp,
    \qquad (X,x_0)\longmapsto\pi_1(X,x_0).
\]
A continuous map $f\colon(X,x_0)\to(Y,y_0)$ induces the group
homomorphism $f_*\colon\pi_1(X,x_0)\to\pi_1(Y,y_0)$ given by
$[\gamma]\mapsto[f\circ\gamma]$.
\end{example}


%% ============================================================
\section{Faithful, full, and essentially surjective functors}
\label{sec:faithful-full}
%% ============================================================

\begin{definition}[Faithful, full, fully faithful]\label{def:faithful-full}
\index{functor!faithful}
\index{functor!full}
\index{functor!fully faithful}
\index{faithful functor}
\index{full functor}
A functor $F\colon\mathcal{C}\to\mathcal{D}$ is called:
\begin{enumerate}[label=(\roman*)]
    \item \textbf{faithful} if for every pair $A,B\in\mathcal{C}$
          the map
          \[
              F_{A,B}\colon\Hom_\mathcal{C}(A,B)
              \longrightarrow\Hom_\mathcal{D}(F(A),F(B))
          \]
          is injective;
    \item \textbf{full} if every $F_{A,B}$ is surjective;
    \item \textbf{fully faithful} if every $F_{A,B}$ is bijective.
\end{enumerate}
\end{definition}

\begin{definition}[Essentially surjective]\label{def:ess-surj}
\index{functor!essentially surjective}
\index{essentially surjective}
A functor $F\colon\mathcal{C}\to\mathcal{D}$ is
\textbf{essentially surjective} (or \emph{essentially surjective on
objects}) if for every $D\in\mathcal{D}$ there exists $C\in\mathcal{C}$
with $F(C)\cong D$.
\end{definition}

\begin{example}\label{ex:faithful-full-examples}
\begin{enumerate}[label=(\roman*)]
    \item Every forgetful functor (e.g.\ $U\colon\Grp\to\Set$) is
          faithful.  It is not full in general: not every function
          between the underlying sets of two groups is a homomorphism.
    \item The inclusion functor $\Ab\hookrightarrow\Grp$ is fully
          faithful: a homomorphism between abelian groups is the same
          thing whether viewed in~$\Ab$ or~$\Grp$.
    \item A functor is an \emph{equivalence of categories}
          if and only if it is fully faithful and essentially surjective
          (assuming the axiom of choice).
\end{enumerate}
\end{example}

\begin{proposition}\label{prop:ff-reflects-iso}
\index{functor!fully faithful!reflects isomorphisms}
A fully faithful functor reflects isomorphisms: if
$F(f)$ is an isomorphism in~$\mathcal{D}$, then $f$ is an isomorphism
in~$\mathcal{C}$.
\end{proposition}

\begin{proof}
Let $F(f)\colon F(A)\to F(B)$ be an isomorphism with inverse~$h$.
Since $F$ is full, $h = F(g)$ for some $g\colon B\to A$.
Then $F(g\circ f) = F(g)\circ F(f) = \id_{F(A)} = F(\id_A)$.
Since $F$ is faithful, $g\circ f = \id_A$.  Similarly
$f\circ g = \id_B$.
\end{proof}


%% ============================================================
\section{Natural transformations}\label{sec:nat-trans}
%% ============================================================

\index{natural transformation|(}

\begin{definition}[Natural transformation]\label{def:nat-trans}
\index{natural transformation!definition}
\index{component!of a natural transformation}
Let $F,G\colon\mathcal{C}\to\mathcal{D}$ be functors.
A \textbf{natural transformation} $\alpha\colon F\Rightarrow G$ is a
family of morphisms in~$\mathcal{D}$,
\[
    \bigl\{\alpha_A\colon F(A)\to G(A)\bigr\}_{A\in\Ob(\mathcal{C})},
\]
called the \emph{components} of~$\alpha$, such that for every morphism
$f\colon A\to B$ in~$\mathcal{C}$ the following \emph{naturality square}
commutes:
\[
\begin{tikzcd}
    F(A) \arrow[r, "\alpha_A"] \arrow[d, "F(f)"'] &
    G(A) \arrow[d, "G(f)"] \\
    F(B) \arrow[r, "\alpha_B"'] &
    G(B)
\end{tikzcd}
\]
That is, $G(f)\circ\alpha_A = \alpha_B\circ F(f)$ for all
$f\colon A\to B$.
\end{definition}

\begin{definition}[Natural isomorphism]\label{def:nat-iso}
\index{natural isomorphism}
\index{natural transformation!isomorphism}
A natural transformation $\alpha\colon F\Rightarrow G$ is a
\textbf{natural isomorphism} if every component $\alpha_A$ is an
isomorphism.  We then write $F\cong G$ and say $F$ and $G$ are
\emph{naturally isomorphic}.
\end{definition}

\begin{example}[Double dual]\label{ex:double-dual}
\index{natural transformation!double dual}
\index{double dual}
Let $\Vect_\K^{\mathrm{fd}}$ denote the category of finite-dimensional
$\K$-vector spaces.  Define $\eta\colon\Id\Rightarrow(-)^{**}$ by
\[
    \eta_V\colon V\longrightarrow V^{**}, \qquad
    v\longmapsto\bigl(\varphi\mapsto\varphi(v)\bigr).
\]
This is a natural isomorphism (each $\eta_V$ is an isomorphism by
dimension counting).  Crucially, the isomorphism $V\cong V^{**}$ is
\emph{natural} in~$V$, whereas the isomorphism $V\cong V^*$ (which
requires choosing a basis) is not.
\end{example}

\begin{example}[Determinant]\label{ex:determinant-nat-trans}
\index{natural transformation!determinant}
\index{determinant!as natural transformation}
Let $\Ring$ be the category of commutative rings.  Consider the functors
$\mathrm{GL}_n,\;(-)^\times\colon\Ring\to\Grp$ sending a
commutative ring~$R$ to the group of invertible $n\times n$ matrices over~$R$
and to the group of units~$R^\times$, respectively.
The determinant gives a natural transformation
\[
    \det\colon\mathrm{GL}_n\Longrightarrow(-)^\times.
\]
Naturality says: for every ring homomorphism $\varphi\colon R\to S$ and
every $M\in\mathrm{GL}_n(R)$,
$\varphi(\det M) = \det(\varphi(M))$.
\end{example}

\begin{definition}[Vertical composition]\label{def:vert-comp}
\index{natural transformation!vertical composition}
\index{vertical composition}
Let $\alpha\colon F\Rightarrow G$ and $\beta\colon G\Rightarrow H$ be
natural transformations between functors
$F,G,H\colon\mathcal{C}\to\mathcal{D}$.
Their \textbf{vertical composition}
$\beta\circ\alpha\colon F\Rightarrow H$ has components
\[
    (\beta\circ\alpha)_A = \beta_A\circ\alpha_A.
\]
\end{definition}

This is depicted schematically as:
\[
\begin{tikzcd}[column sep=huge]
    \mathcal{C}
    \arrow[r, bend left=50, "F"{name=F}]
    \arrow[r, "G"{name=G, description}]
    \arrow[r, bend right=50, "H"{name=H, below}]
    &
    \mathcal{D}
    \arrow[Rightarrow, from=F, to=G, "\alpha"]
    \arrow[Rightarrow, from=G, to=H, "\beta"]
\end{tikzcd}
\]

\begin{lemma}\label{lem:vert-comp-nat}
Vertical composition is associative, and the identity natural
transformation $\id_F\colon F\Rightarrow F$ (with components
$(\id_F)_A = \id_{F(A)}$) serves as identity.
\end{lemma}

\begin{proof}
Both statements follow immediately from the corresponding properties
of morphism composition in~$\mathcal{D}$.
\end{proof}

\begin{definition}[Horizontal composition]\label{def:horiz-comp}
\index{natural transformation!horizontal composition}
\index{horizontal composition}
\index{Godement product|see{horizontal composition}}
Let $\alpha\colon F\Rightarrow G$ between functors
$\mathcal{C}\to\mathcal{D}$ and
$\beta\colon H\Rightarrow K$ between functors
$\mathcal{D}\to\mathcal{E}$.
The \textbf{horizontal composition} (or \emph{Godement product})
$\beta*\alpha\colon H\circ F\Rightarrow K\circ G$ has components
\[
    (\beta*\alpha)_A = \beta_{G(A)}\circ H(\alpha_A)
                     = K(\alpha_A)\circ\beta_{F(A)}.
\]
The two expressions are equal by naturality of~$\beta$:
\[
\begin{tikzcd}
    HF(A) \arrow[r, "H(\alpha_A)"] \arrow[d, "\beta_{F(A)}"'] &
    HG(A) \arrow[d, "\beta_{G(A)}"] \\
    KF(A) \arrow[r, "K(\alpha_A)"'] &
    KG(A)
\end{tikzcd}
\]
\end{definition}


%% ============================================================
\section{Functor categories}\label{sec:functor-categories}
%% ============================================================

\begin{definition}[Functor category]\label{def:functor-category}
\index{functor category}
\index{category!functor}
Let $\mathcal{C}$ and $\mathcal{D}$ be categories with $\mathcal{C}$
small.  The \textbf{functor category}
$[\mathcal{C},\mathcal{D}]$ (also written
$\Fun(\mathcal{C},\mathcal{D})$ or $\mathcal{D}^{\mathcal{C}}$) is
the category whose:
\begin{itemize}
    \item objects are functors $F\colon\mathcal{C}\to\mathcal{D}$;
    \item morphisms are natural transformations;
    \item composition is vertical composition of natural
          transformations.
\end{itemize}
\end{definition}

\begin{remark}\label{rem:functor-cat-size}
\index{functor category!size}
If $\mathcal{C}$ is small and $\mathcal{D}$ is locally small, then
$[\mathcal{C},\mathcal{D}]$ is locally small.
\end{remark}

\begin{example}[Presheaf categories]\label{ex:presheaf}
\index{presheaf}
\index{presheaf category}
A particularly important special case: let $\mathcal{C}$ be a small
category.  The functor category
\[
    \widehat{\mathcal{C}} \;=\; [\mathcal{C}^{\op},\Set]
\]
is called the \textbf{category of presheaves} on~$\mathcal{C}$.  Its
objects are contravariant functors $\mathcal{C}\to\Set$, called
\emph{presheaves}.  This category has remarkable properties (it is
complete, cocomplete, and cartesian closed) that we shall explore later.
\end{example}

\begin{proposition}\label{prop:nat-iso-pointwise}
\index{natural isomorphism!pointwise criterion}
A natural transformation $\alpha\colon F\Rightarrow G$ is an isomorphism
in $[\mathcal{C},\mathcal{D}]$ if and only if each component
$\alpha_A$ is an isomorphism in~$\mathcal{D}$.
\end{proposition}

\begin{proof}
If $\alpha$ is a natural isomorphism, its inverse $\alpha^{-1}$ has
components $(\alpha^{-1})_A = (\alpha_A)^{-1}$.  One checks naturality
of $\alpha^{-1}$ by applying $(-\,)^{-1}$ to the naturality squares
of~$\alpha$.
Conversely, if each $\alpha_A$ is an isomorphism, define
$\beta_A = (\alpha_A)^{-1}$.  The naturality of~$\beta$ follows from
that of~$\alpha$: for $f\colon A\to B$,
\[
    \beta_B\circ G(f) = \beta_B\circ G(f)\circ\alpha_A\circ\beta_A
    = \beta_B\circ\alpha_B\circ F(f)\circ\beta_A
    = F(f)\circ\beta_A. \qedhere
\]
\end{proof}

\begin{definition}[Equivalence of categories]\label{def:equivalence}
\index{equivalence of categories}
\index{category!equivalence}
An \textbf{equivalence of categories} between $\mathcal{C}$ and
$\mathcal{D}$ consists of functors
$F\colon\mathcal{C}\to\mathcal{D}$ and
$G\colon\mathcal{D}\to\mathcal{C}$ together with natural isomorphisms
\[
    \eta\colon\Id_\mathcal{C}\xRightarrow{\;\sim\;}G\circ F
    \qquad\text{and}\qquad
    \varepsilon\colon F\circ G\xRightarrow{\;\sim\;}\Id_\mathcal{D}.
\]
We write $\mathcal{C}\simeq\mathcal{D}$ and say $\mathcal{C}$ and
$\mathcal{D}$ are \emph{equivalent}.
\end{definition}

\begin{theorem}\label{thm:equiv-ff-eso}
\index{equivalence of categories!characterisation}
A functor $F\colon\mathcal{C}\to\mathcal{D}$ is part of an equivalence
of categories if and only if $F$ is fully faithful and essentially
surjective.
\end{theorem}

\begin{proof}
($\Rightarrow$) Suppose $F$ is an equivalence with quasi-inverse~$G$ and
natural isomorphisms $\eta\colon\Id\xRightarrow{\sim} GF$ and
$\varepsilon\colon FG\xRightarrow{\sim}\Id$.

\emph{Essentially surjective:} For $D\in\mathcal{D}$,
$F(G(D))\cong D$ via $\varepsilon_D$.

\emph{Faithful:} Suppose $F(f)=F(g)$ for $f,g\colon A\to B$.
Then $GF(f)=GF(g)$.  Naturality of~$\eta$ gives
$\eta_B\circ f = GF(f)\circ\eta_A = GF(g)\circ\eta_A = \eta_B\circ g$,
and since $\eta_B$ is an isomorphism, $f=g$.

\emph{Full:} Let $h\colon F(A)\to F(B)$.  Consider
$f = \eta_B^{-1}\circ G(h)\circ\eta_A\colon A\to B$.
Then $F(f) = \varepsilon_{F(B)}\circ FG(h)\circ(\varepsilon_{F(A)})^{-1}$.
By naturality of~$\varepsilon$,
$\varepsilon_{F(B)}\circ FG(h) = h\circ\varepsilon_{F(A)}$,
so $F(f) = h$.

($\Leftarrow$)  Assume $F$ is fully faithful and essentially surjective.
For each $D\in\mathcal{D}$, choose $G(D)\in\mathcal{C}$ with an
isomorphism $\varepsilon_D\colon FG(D)\xrightarrow{\sim} D$.
For $h\colon D\to D'$, define $G(h)$ to be the unique morphism
(existing by full faithfulness) satisfying
$F(G(h)) = \varepsilon_{D'}^{-1}\circ h\circ\varepsilon_D$.
One verifies $G$ is a functor and that $\varepsilon$ and the induced
$\eta$ are natural isomorphisms.
\end{proof}


%% ============================================================
\section{Exercises for Chapter~1}\label{sec:exercises-ch1}
%% ============================================================

\begin{exercise}\label{exo:cat-rels}
\index{category!of relations}
Define a category $\mathbf{Rel}$ whose objects are sets and whose
morphisms $A\to B$ are relations $R\subseteq A\times B$.
What is composition?  What are the identity morphisms?
Show that $\mathbf{Rel}$ is indeed a category.
\end{exercise}

\begin{exercise}\label{exo:groupoid}
\index{groupoid}
A \emph{groupoid} is a category in which every morphism is an
isomorphism.  Show that the following are groupoids:
\begin{enumerate}[label=(\roman*)]
    \item A group, viewed as a one-object category.
    \item The \emph{fundamental groupoid} $\Pi_1(X)$ of a topological
          space~$X$: objects are points of~$X$, morphisms $x\to y$
          are homotopy classes of paths from~$x$ to~$y$.
    \item The \emph{core} $\mathrm{core}(\mathcal{C})$ of any category
          $\mathcal{C}$: same objects, but only isomorphisms.
\end{enumerate}
\end{exercise}

\begin{exercise}\label{exo:slice-category}
\index{slice category}
\index{category!slice}
Let $\mathcal{C}$ be a category and $A\in\mathcal{C}$ an object.
The \emph{slice category} (or \emph{over category}) $\mathcal{C}/A$
has:
\begin{itemize}
    \item Objects: morphisms $f\colon X\to A$ in~$\mathcal{C}$.
    \item Morphisms: from $(f\colon X\to A)$ to $(g\colon Y\to A)$,
          a morphism $h\colon X\to Y$ in~$\mathcal{C}$ such that
          $g\circ h = f$:
\[
\begin{tikzcd}
    X \arrow[rr, "h"] \arrow[dr, "f"'] && Y \arrow[dl, "g"] \\
    & A &
\end{tikzcd}
\]
\end{itemize}
Verify that $\mathcal{C}/A$ is a category.  What is the terminal object?
Describe $\Set/\{0,1\}$ explicitly.
\end{exercise}

\begin{exercise}\label{exo:comma-category}
\index{comma category}
Let $F\colon\mathcal{A}\to\mathcal{C}$ and
$G\colon\mathcal{B}\to\mathcal{C}$ be functors.
Define the \emph{comma category} $(F\downarrow G)$ and show that
slice categories and coslice categories are special cases.
\end{exercise}

\begin{exercise}\label{exo:mono-epi-exercises}
\begin{enumerate}[label=(\roman*)]
    \item Show that in any category, the composition of two
          monomorphisms is a monomorphism.
    \item Show that in $\Set$, a morphism is an epimorphism if and
          only if it is surjective.
    \item Find an example of a morphism in a concrete category that is
          both monic and epic but not an isomorphism
          (beyond the example of $\Z\hookrightarrow\Q$ in $\Ring$).
\end{enumerate}
\end{exercise}

\begin{exercise}\label{exo:endomorphism-monoid}
\index{endomorphism!monoid}
For any object $A$ in a category~$\mathcal{C}$, show that
$\End_\mathcal{C}(A)=\Hom_\mathcal{C}(A,A)$ forms a monoid under
composition.  When is this monoid a group?
\end{exercise}

\begin{exercise}\label{exo:functor-composition}
\index{functor!composition}
Let $F\colon\mathcal{A}\to\mathcal{B}$ and
$G\colon\mathcal{B}\to\mathcal{C}$ be functors.
Show that the composition $G\circ F\colon\mathcal{A}\to\mathcal{C}$
is again a functor.  Verify that functor composition is associative
and that the identity functor $\Id_\mathcal{C}$ is a two-sided identity.
\end{exercise}

\begin{exercise}\label{exo:nat-trans-id-functor}
\index{natural transformation!from identity functor}
Let $F\colon\mathcal{C}\to\mathcal{C}$ be a functor.
Show that the following are equivalent:
\begin{enumerate}[label=(\roman*)]
    \item $\alpha\colon\Id_\mathcal{C}\Rightarrow F$ is a natural
          transformation;
    \item For every morphism $f\colon A\to B$, one has
          $F(f)\circ\alpha_A = \alpha_B\circ f$.
\end{enumerate}
\end{exercise}

\begin{exercise}\label{exo:interchange-law}
\index{interchange law}
\index{natural transformation!interchange law}
Prove the \emph{interchange law} for horizontal and vertical
composition of natural transformations.  Specifically, given
\[
\begin{tikzcd}[column sep=large]
    \mathcal{A}
    \arrow[r, bend left=40, "F_1"{name=F1}]
    \arrow[r, bend right=40, "F_2"{name=F2, below}]
    &
    \mathcal{B}
    \arrow[r, bend left=40, "G_1"{name=G1}]
    \arrow[r, bend right=40, "G_2"{name=G2, below}]
    &
    \mathcal{C}
    \arrow[Rightarrow, from=F1, to=F2, "\alpha"]
    \arrow[Rightarrow, from=G1, to=G2, "\beta"]
\end{tikzcd}
\]
and natural transformations $\alpha'\colon F_2\Rightarrow F_3$,
$\beta'\colon G_2\Rightarrow G_3$, show that
\[
    (\beta'\circ\beta)*(\alpha'\circ\alpha)
    = (\beta'*\alpha')\circ(\beta*\alpha).
\]
\end{exercise}

\begin{exercise}\label{exo:skeleton}
\index{skeleton of a category}
A category~$\mathcal{C}$ is \emph{skeletal} if isomorphic objects are
equal.  A \emph{skeleton} of~$\mathcal{C}$ is a full subcategory that
is skeletal and contains exactly one object from each isomorphism class.
Show that every category has a skeleton and that inclusion of the
skeleton is an equivalence of categories.
\end{exercise}

\begin{exercise}\label{exo:nat-trans-precomposition}
Let $F,G\colon\mathcal{C}\to\mathcal{D}$ be functors,
$\alpha\colon F\Rightarrow G$ a natural transformation, and
$H\colon\mathcal{B}\to\mathcal{C}$ a functor.
Define the \emph{whiskering} $\alpha H\colon FH\Rightarrow GH$ by
$(\alpha H)_B = \alpha_{H(B)}$.
Show that $\alpha H$ is a natural transformation.
Similarly define whiskering on the other side.
\end{exercise}

\index{functor|)}
\index{natural transformation|)}
\index{category|)}


%%%%%%%%%%%%%%%%%%%%%%%%%%%%%%%%%%%%%%%%%%%%%%%%%%%%%%%%%%%%
%%       CHAPTER 2: DUALITY AND THE DUAL PRINCIPLE        %%
%%%%%%%%%%%%%%%%%%%%%%%%%%%%%%%%%%%%%%%%%%%%%%%%%%%%%%%%%%%%

\chapter{Duality and the Dual Principle}\label{ch:duality}
\minitoc

\vspace{1em}

One of the most powerful features of category theory is the
\emph{duality principle}: every categorical statement has a dual,
obtained by reversing all morphisms.  This single observation doubles the
theorems one can prove ``for free''.  In this chapter we make this
precise by introducing opposite categories, formulating the duality
principle, and examining several concrete dualities.  We also study the
two-variable hom-functor, which will be essential in later chapters.

\index{duality|(}


%% ============================================================
\section{The opposite category}\label{sec:opposite-cat}
%% ============================================================

\begin{definition}[Opposite category]\label{def:opposite-cat}
\index{opposite category}
\index{category!opposite}
\index{dual category|see{opposite category}}
Let $\mathcal{C}$ be a category.  The \textbf{opposite category} (or
\emph{dual category}) $\mathcal{C}^{\op}$ is defined by:
\begin{enumerate}[label=(\roman*)]
    \item $\Ob(\mathcal{C}^{\op}) = \Ob(\mathcal{C})$.
    \item For objects $A,B$:
          $\Hom_{\mathcal{C}^{\op}}(A,B) = \Hom_\mathcal{C}(B,A)$.
    \item Composition in $\mathcal{C}^{\op}$:
          given $f^{\op}\colon A\to B$ and $g^{\op}\colon B\to C$
          in~$\mathcal{C}^{\op}$ (corresponding to $f\colon B\to A$ and
          $g\colon C\to B$ in~$\mathcal{C}$), we set
          $g^{\op}\circ f^{\op} = (f\circ g)^{\op}$.
    \item The identity on~$A$ in $\mathcal{C}^{\op}$ is $(\id_A)^{\op}
          = \id_A$.
\end{enumerate}
\end{definition}

\begin{proposition}\label{prop:op-involution}
\index{opposite category!involution}
For any category~$\mathcal{C}$, $(\mathcal{C}^{\op})^{\op}=\mathcal{C}$.
\end{proposition}

\begin{proof}
The objects coincide.  For morphisms,
$\Hom_{(\mathcal{C}^{\op})^{\op}}(A,B)
= \Hom_{\mathcal{C}^{\op}}(B,A)
= \Hom_\mathcal{C}(A,B)$.
Composition: $(f^{\op})^{\op}\circ(g^{\op})^{\op}
= ((g^{\op}\circ f^{\op})^{\op})
= ((f\circ g)^{\op})^{\op}
= f\circ g$.
Everything reduces to the original data of~$\mathcal{C}$.
\end{proof}

\begin{example}\label{ex:op-examples}
\begin{enumerate}[label=(\roman*)]
    \item If $(P,\le)$ is a poset viewed as a category, then
          $P^{\op}$ is the poset $(P,\ge)$.
          \index{poset!opposite}
    \item If $G$ is a group viewed as a one-object category, then
          $G^{\op}$ is the \emph{opposite group}: same elements,
          with multiplication $a\cdot^{\op}b = b\cdot a$.
          For abelian groups, $G^{\op}\cong G$.
          \index{group!opposite}
    \item $\Set^{\op}$ is a perfectly valid category, but it is not
          equivalent to any ``familiar'' category of structured sets
          (it is, however, equivalent to the category of complete
          atomic Boolean algebras, by a non-trivial theorem).
          \index{category!Set op@$\Set^{\op}$}
\end{enumerate}
\end{example}


%% ============================================================
\section{The duality principle}\label{sec:duality-principle}
%% ============================================================

\begin{definition}[Dual statement]\label{def:dual-statement}
\index{dual statement}
\index{duality principle}
Let $\Sigma$ be a statement formulated purely in the language of
category theory (objects, morphisms, composition, identities).
The \textbf{dual statement}~$\Sigma^{\op}$ is obtained from~$\Sigma$
by:
\begin{itemize}
    \item reversing the direction of every morphism;
    \item replacing each composition $g\circ f$ by $f\circ g$;
    \item interchanging domain and codomain.
\end{itemize}
\end{definition}

\begin{theorem}[Duality principle]\label{thm:duality-principle}
\index{duality principle!statement}
If a statement~$\Sigma$ holds in every category, then so does its
dual~$\Sigma^{\op}$.
More generally, if $\Sigma$ holds in a category~$\mathcal{C}$, then
$\Sigma^{\op}$ holds in~$\mathcal{C}^{\op}$.
\end{theorem}

\begin{proof}
A statement about~$\mathcal{C}^{\op}$ is, by definition, a statement
about~$\mathcal{C}$ with all arrows reversed.
Since $\mathcal{C}^{\op}$ is a legitimate category, any theorem valid
for all categories applies to it.
Unwinding the definitions in $\mathcal{C}^{\op}$ recovers the dual
statement in~$\mathcal{C}$.
\end{proof}

\begin{remark}\label{rem:duality-examples}
\index{duality principle!examples}
The duality principle immediately gives us dual pairs of concepts:
\begin{center}
\begin{tabular}{ll}
\toprule
\textbf{Concept} & \textbf{Dual concept} \\
\midrule
monomorphism & epimorphism \\
initial object & terminal object \\
product & coproduct \\
limit & colimit \\
kernel & cokernel \\
left adjoint & right adjoint \\
\bottomrule
\end{tabular}
\end{center}
Once we prove a theorem about monomorphisms, the dual theorem about
epimorphisms follows automatically.
\end{remark}

\begin{example}\label{ex:duality-mono-epi}
\index{monomorphism!dual of epimorphism}
\index{epimorphism!dual of monomorphism}
We showed in 
ef{lem:mono-epi-composition} that if $g\circ f$ is monic
then $f$ is monic.  By duality: if $g\circ f$ is epic, then $g$ is
epic.  No additional proof is needed---it follows from the proof of the
original statement applied to $\mathcal{C}^{\op}$.
\end{example}


%% ============================================================
\section{Functors and duality}\label{sec:functors-duality}
%% ============================================================

\begin{proposition}\label{prop:contravariant-op}
\index{functor!contravariant!via opposite category}
A contravariant functor $F\colon\mathcal{C}\to\mathcal{D}$ is the same
thing as a covariant functor $\mathcal{C}^{\op}\to\mathcal{D}$, and
also the same as a covariant functor $\mathcal{C}\to\mathcal{D}^{\op}$.
\end{proposition}

\begin{proof}
A covariant functor $\mathcal{C}^{\op}\to\mathcal{D}$ assigns to each
morphism $f^{\op}\colon B\to A$ in $\mathcal{C}^{\op}$ (i.e.\
$f\colon A\to B$ in $\mathcal{C}$) a morphism
$F(f^{\op})\colon F(B)\to F(A)$ in~$\mathcal{D}$,
which is precisely a contravariant assignment.
The preservation of composition and identities translates
directly.  The argument for $\mathcal{C}\to\mathcal{D}^{\op}$ is
analogous.
\end{proof}

\begin{remark}\label{rem:op-functor}
\index{opposite functor}
Every functor $F\colon\mathcal{C}\to\mathcal{D}$ induces a functor
$F^{\op}\colon\mathcal{C}^{\op}\to\mathcal{D}^{\op}$ defined by
$F^{\op}(A)=F(A)$ on objects and $F^{\op}(f^{\op})=(F(f))^{\op}$ on
morphisms.  In fact, $(-)^{\op}$ is a functor $\Cat\to\Cat$.
\end{remark}


%% ============================================================
\section{Concrete dualities}\label{sec:concrete-dualities}
%% ============================================================

\index{duality!concrete|(}

The abstract duality principle tells us about dual \emph{concepts}, but
there also exist concrete \emph{equivalences}
$\mathcal{C}^{\op}\simeq\mathcal{D}$ for specific categories.  These are
called \emph{concrete dualities} or \emph{dualities of categories}.

\begin{example}[Stone duality]\label{ex:stone-duality}
\index{Stone duality}
\index{Boolean algebra}
\index{Stone space}
The category $\mathbf{Stone}$ of Stone spaces (compact, Hausdorff,
totally disconnected topological spaces) and continuous maps is
equivalent to $\mathbf{Bool}^{\op}$, where $\mathbf{Bool}$ is the
category of Boolean algebras and their homomorphisms:
\[
    \mathbf{Stone} \;\simeq\; \mathbf{Bool}^{\op}.
\]
The equivalence is implemented by the functor sending a Stone space to
its Boolean algebra of clopen subsets, and conversely sending a
Boolean algebra to its space of ultrafilters.
\end{example}

\begin{example}[Pontryagin duality]\label{ex:pontryagin-duality}
\index{Pontryagin duality}
\index{locally compact abelian group}
Let $\mathbf{LCA}$ be the category of locally compact abelian groups
and continuous homomorphisms.  Pontryagin duality provides a
contravariant equivalence
\[
    (-)^\wedge\colon\mathbf{LCA}^{\op}
    \xrightarrow{\;\sim\;}\mathbf{LCA},
\]
where $G^\wedge = \Hom_{\mathrm{cts}}(G,\R/\Z)$ is the
\emph{Pontryagin dual}.
The natural map $G\to G^{\wedge\wedge}$, $g\mapsto(\chi\mapsto\chi(g))$,
is an isomorphism.
\end{example}

\begin{example}[Gelfand duality]\label{ex:gelfand-duality}
\index{Gelfand duality}
\index{C*-algebra@$C^*$-algebra}
There is an equivalence
\[
    \mathbf{CHaus}^{\op} \;\simeq\; \mathbf{cC^*Alg}_1,
\]
where $\mathbf{CHaus}$ is the category of compact Hausdorff spaces and
$\mathbf{cC^*Alg}_1$ is the category of commutative unital
$C^*$-algebras.  The functor sends $X\mapsto C(X,\C)$
(continuous functions) and the quasi-inverse sends a $C^*$-algebra to
its maximal ideal space (Gelfand spectrum).
\end{example}

\begin{example}[Finite-dimensional vector space duality]\label{ex:fd-vect-duality}
\index{duality!finite-dimensional vector spaces}
\index{vector space!duality}
The dual space functor $(-)^*\colon\Vect_\K^{\mathrm{fd}}\to
(\Vect_\K^{\mathrm{fd}})^{\op}$ is an equivalence of categories.
Combined with $(\Vect_\K^{\mathrm{fd}})^{\op}\simeq
\Vect_\K^{\mathrm{fd}}$
(via the double dual), we see that finite-dimensional vector spaces
are self-dual.
\end{example}

\index{duality!concrete|)}


%% ============================================================
\section{The hom-functor in both variables}\label{sec:hom-bifunctor}
%% ============================================================

\begin{definition}[Bifunctor]\label{def:bifunctor}
\index{bifunctor}
\index{functor!of two variables}
Let $\mathcal{C}$, $\mathcal{D}$, $\mathcal{E}$ be categories.
A \textbf{bifunctor} is a functor
$F\colon\mathcal{C}\times\mathcal{D}\to\mathcal{E}$.
\end{definition}

\begin{remark}\label{rem:bifunctor-partial}
A bifunctor $F\colon\mathcal{C}\times\mathcal{D}\to\mathcal{E}$
determines, for each $C\in\mathcal{C}$, a functor
$F(C,-)\colon\mathcal{D}\to\mathcal{E}$ and, for each
$D\in\mathcal{D}$, a functor
$F(-,D)\colon\mathcal{C}\to\mathcal{E}$.
Conversely, a family of functors in each variable that is
``jointly functorial'' determines a bifunctor.
\end{remark}

\begin{proposition}\label{prop:hom-bifunctor}
\index{hom-functor!bifunctor}
\index{hom-bifunctor}
For any locally small category~$\mathcal{C}$, the assignment
\[
    \Hom_\mathcal{C}(-,-)\colon\mathcal{C}^{\op}\times\mathcal{C}
    \longrightarrow\Set
\]
defined on objects by $(A,B)\mapsto\Hom_\mathcal{C}(A,B)$ and on
morphisms by
\[
    (f\colon A'\to A,\; g\colon B\to B')\;\longmapsto\;
    \bigl(h\mapsto g\circ h\circ f\bigr)
    \colon\Hom(A,B)\to\Hom(A',B'),
\]
is a bifunctor (i.e.\ a functor from the product category
$\mathcal{C}^{\op}\times\mathcal{C}$ to $\Set$).
\end{proposition}

\begin{proof}
We verify the functor axioms.  On identities:
\[
    \Hom(\id_A,\id_B)(h) = \id_B\circ h\circ\id_A = h,
\]
so $\Hom(\id_A,\id_B) = \id_{\Hom(A,B)}$.

On composition: let $f'\colon A''\to A'$, $f\colon A'\to A$,
$g\colon B\to B'$, $g'\colon B'\to B''$.
We need
$\Hom(f\circ f',\,g'\circ g) = \Hom(f',g')\circ\Hom(f,g)$.
The left side maps $h\mapsto (g'\circ g)\circ h\circ(f\circ f')$.
The right side maps
$h\mapsto g\circ h\circ f\mapsto g'\circ(g\circ h\circ f)\circ f'$,
which is the same by associativity.
\end{proof}

\begin{remark}\label{rem:hom-functor-partial-functors}
\index{hom-functor!covariant}
\index{hom-functor!contravariant}
The bifunctor $\Hom(-,-)$ encodes the two ``partial'' hom-functors
from 
ef{ex:hom-functors}:
\begin{itemize}
    \item Fixing the first argument: $\Hom(A,-)$ is covariant.
    \item Fixing the second argument: $\Hom(-,B)$ is contravariant
          (i.e.\ covariant from $\mathcal{C}^{\op}$).
\end{itemize}
The two partial functors determine the bifunctor (they agree on pairs
of identities and satisfy a compatibility condition that is automatic
from the bifunctor structure).
\end{remark}

\begin{proposition}\label{prop:hom-preserves-limits}
\index{hom-functor!preserves limits}
For any object $A$ in a locally small category~$\mathcal{C}$, the
covariant hom-functor $\Hom(A,-)\colon\mathcal{C}\to\Set$ preserves
all limits that exist in~$\mathcal{C}$.
\end{proposition}

\begin{proof}
This is a consequence of the universal property of limits and will be
proved rigorously in 
ef{ch:limits-colimits}.  The key idea is that a
natural bijection $\Hom(A,\lim F)\cong\lim\Hom(A,F-)$ follows directly
from the defining universal property of the limit.
\end{proof}

\begin{definition}[Representable functor (preview)]\label{def:representable-preview}
\index{representable functor!preview}
\index{functor!representable}
A functor $F\colon\mathcal{C}\to\Set$ is \textbf{representable} if it
is naturally isomorphic to $\Hom_\mathcal{C}(A,-)$ for some object
$A\in\mathcal{C}$.  The object~$A$ is called a \emph{representing
object}.

A contravariant functor $F\colon\mathcal{C}^{\op}\to\Set$ is
representable if $F\cong\Hom_\mathcal{C}(-,B)$ for some~$B$.

Representable functors are central to category theory; we shall study
them in depth in connection with the Yoneda lemma.
\end{definition}


%% ============================================================
\section{Exercises for Chapter~2}\label{sec:exercises-ch2}
%% ============================================================

\begin{exercise}\label{exo:op-op}
Let $\mathcal{C}$ be a category.
\begin{enumerate}[label=(\roman*)]
    \item Verify carefully that $\mathcal{C}^{\op}$ as defined in
          
ef{def:opposite-cat} satisfies the axioms of a category.
    \item Show that $(\mathcal{C}^{\op})^{\op}=\mathcal{C}$ (equality,
          not merely isomorphism).
\end{enumerate}
\end{exercise}

\begin{exercise}\label{exo:initial-terminal-dual}
\index{initial object!dual of terminal}
Show that $A$ is an initial object in $\mathcal{C}$ if and only if
$A$ is a terminal object in $\mathcal{C}^{\op}$.
Deduce the uniqueness (up to unique isomorphism) of terminal objects
from the corresponding result for initial objects via duality.
\end{exercise}

\begin{exercise}\label{exo:mono-epi-dual}
Show that $f$ is a monomorphism in $\mathcal{C}$ if and only if
$f^{\op}$ is an epimorphism in $\mathcal{C}^{\op}$.
\end{exercise}

\begin{exercise}\label{exo:op-functor-category}
\index{functor category!opposite}
Let $\mathcal{C}$ be small and $\mathcal{D}$ any category.
Construct a canonical isomorphism of categories
\[
    [\mathcal{C},\mathcal{D}]^{\op}
    \;\cong\;
    [\mathcal{C},\mathcal{D}^{\op}].
\]
(Hint: what happens to naturality squares when you reverse arrows in the
target?)
\end{exercise}

\begin{exercise}\label{exo:self-dual-category}
\index{self-dual category}
A category~$\mathcal{C}$ is called \emph{self-dual} if
$\mathcal{C}\simeq\mathcal{C}^{\op}$.
\begin{enumerate}[label=(\roman*)]
    \item Show that $\Vect_\K^{\mathrm{fd}}$ is self-dual.
    \item Show that any groupoid is self-dual.
    \item Is $\Set$ self-dual?  Justify.
\end{enumerate}
\end{exercise}

\begin{exercise}\label{exo:product-category}
\index{product category}
Let $\mathcal{C}$ and $\mathcal{D}$ be categories.
Define the product category $\mathcal{C}\times\mathcal{D}$ and show that
\[
    (\mathcal{C}\times\mathcal{D})^{\op}
    \;\cong\;
    \mathcal{C}^{\op}\times\mathcal{D}^{\op}.
\]
\end{exercise}

\begin{exercise}\label{exo:hom-bifunctor-verification}
\index{hom-bifunctor!detailed verification}
Let $\mathcal{C}$ be a locally small category.
Give a detailed verification that
$\Hom_\mathcal{C}(-,-)\colon\mathcal{C}^{\op}\times\mathcal{C}\to\Set$
is a functor by checking the functor axioms (preservation of identities
and composition) explicitly.
\end{exercise}

\begin{exercise}\label{exo:contravariant-nat-trans}
\index{natural transformation!between contravariant functors}
Let $F,G\colon\mathcal{C}^{\op}\to\Set$ be contravariant functors.
Spell out explicitly what a natural transformation
$\alpha\colon F\Rightarrow G$ looks like: what are the components, and
what does the naturality condition say?
Draw the naturality square and compare it with the covariant case.
\end{exercise}

\begin{exercise}\label{exo:op-preserves-equivalence}
\index{equivalence of categories!and opposite}
Show that if $\mathcal{C}\simeq\mathcal{D}$, then
$\mathcal{C}^{\op}\simeq\mathcal{D}^{\op}$.
\end{exercise}

\begin{exercise}\label{exo:concrete-duality-exploration}
\index{duality!in lattices}
Let $L$ be a bounded lattice, viewed as a category (poset).
Show that $L^{\op}$ is again a bounded lattice, with joins and meets
interchanged.  Describe the dual of the lattice of open sets of a
topological space~$X$.
\end{exercise}

\begin{exercise}\label{exo:op-slice}
\index{slice category!opposite}
\index{coslice category}
Let $\mathcal{C}$ be a category and $A\in\mathcal{C}$ an object.
The \emph{coslice category} (or \emph{under category}) $A/\mathcal{C}$
has objects $f\colon A\to X$ and morphisms the evident commutative
triangles.
Show that $(A/\mathcal{C})^{\op}\cong\mathcal{C}^{\op}/A$.
\end{exercise}

\begin{exercise}\label{exo:representable-dual}
\index{representable functor!and duality}
Let $\mathcal{C}$ be a locally small category and $A\in\mathcal{C}$.
Show that the contravariant hom-functor $\Hom(-,A)$ is representable
as a functor $\mathcal{C}^{\op}\to\Set$, with representing object~$A$
(viewed as an object of $\mathcal{C}^{\op}$).
\end{exercise}

\index{duality|)}
%%%%%%%%%%%%%%%%%%%%%%%%%%%%%%%%%%%%%%%%%%%%%%%%%%%%%%%%%%%%%
%% CATEGORY THEORY — Chapters 3–4
%% Limits and Colimits / Adjunctions
%%%%%%%%%%%%%%%%%%%%%%%%%%%%%%%%%%%%%%%%%%%%%%%%%%%%%%%%%%%%%

%%%%%%%%%%%%%%%%%%%%%%%%%%%%%%%%%%%%%%%%%%%%%%%%%%%%%%%%%%%%
%%           CHAPTER 3: LIMITS AND COLIMITS               %%
%%%%%%%%%%%%%%%%%%%%%%%%%%%%%%%%%%%%%%%%%%%%%%%%%%%%%%%%%%%%

\chapter{Limits and Colimits}\label{ch:limits-colimits}
\minitoc

\vspace{1em}

Limits and colimits are the categorical generalisations of constructions
that pervade all of mathematics: products, intersections, kernels,
pullbacks, inverse limits, and their duals.  Understanding them in full
generality is one of the main rewards of learning category theory.
In this chapter we develop the theory systematically, beginning with the
language of diagrams and cones, proceeding to the universal property that
defines a limit, and then examining the most important special cases.
We prove the fundamental existence theorem---that products and equalisers
suffice for all finite limits---and establish the crucial relationship
between limits and representable functors.

\index{limit|(}
\index{colimit|(}

%% ============================================================
\section{Diagrams and index categories}\label{sec:diagrams}
%% ============================================================

\begin{definition}[Diagram]\label{def:diagram}
\index{diagram}
\index{index category}
\index{functor!diagram}
Let $\mathcal{C}$ be a category and let $\mathcal{J}$ be a small category
(the \emph{index category} or \emph{shape category}).
A \textbf{diagram of shape~$\mathcal{J}$ in~$\mathcal{C}$} is a functor
\[
    F \colon \mathcal{J} \longrightarrow \mathcal{C}.
\]
We write $F_j$ for the image of an object $j \in \Ob(\mathcal{J})$
and $F_\alpha \colon F_j \to F_k$ for the image of a morphism
$\alpha \colon j \to k$ in~$\mathcal{J}$.
\end{definition}

\begin{example}\label{ex:diagrams-shapes}
\index{discrete category}
\index{parallel pair}
\index{span}
\index{cospan}
The following small categories appear constantly as index categories.
\begin{enumerate}[label=(\roman*)]
\item \textbf{Discrete categories.}
    Let $\mathcal{J}$ be a discrete category on a set~$I$
    (objects~$I$, only identity morphisms).
    A diagram $F\colon \mathcal{J}\to\mathcal{C}$ is simply a family
    of objects $\{F_i\}_{i\in I}$.

\item \textbf{The parallel pair.}
    Let $\mathcal{J} = (\bullet \rightrightarrows \bullet)$,
    i.e.\ two objects $0,1$ with two non-identity morphisms
    $\alpha,\beta\colon 0\to 1$.
    A diagram of this shape is a pair of parallel morphisms
    $f,g\colon A\to B$ in~$\mathcal{C}$.

\item \textbf{The span.}
    $\mathcal{J} = (\bullet \leftarrow \bullet \rightarrow \bullet)$.
    A diagram is a span $A \xleftarrow{f} C \xrightarrow{g} B$.

\item \textbf{The cospan.}
    $\mathcal{J} = (\bullet \rightarrow \bullet \leftarrow \bullet)$.
    A diagram is a cospan $A \xrightarrow{f} C \xleftarrow{g} B$.
\end{enumerate}
\end{example}

\begin{remark}\label{rem:diagram-notation}
\index{diagram!as functor}
We shall sometimes write a diagram simply as
$\{F_j,\,F_\alpha\}_{j\in\mathcal{J}}$
when the index category is understood.
A diagram is nothing more than a functor; the terminology ``diagram''
merely signals that we are about to take its limit or colimit.
\end{remark}


%% ============================================================
\section{Cones and cocones}\label{sec:cones}
%% ============================================================

\begin{definition}[Cone]\label{def:cone}
\index{cone}
\index{cone!vertex}
\index{cone!leg}
Let $F\colon \mathcal{J}\to\mathcal{C}$ be a diagram.
A \textbf{cone over~$F$} is a pair $(N,\,\{\pi_j\}_{j\in\mathcal{J}})$
consisting of an object $N\in\mathcal{C}$ (the \emph{vertex} of the cone)
and a family of morphisms $\pi_j\colon N\to F_j$ (the \emph{legs}),
one for each object $j\in\mathcal{J}$,
such that for every morphism $\alpha\colon j\to k$ in~$\mathcal{J}$
the following triangle commutes:
\[
\begin{tikzcd}
& N \arrow[dl, "\pi_j"'] \arrow[dr, "\pi_k"] & \\
F_j \arrow[rr, "F_\alpha"'] & & F_k
\end{tikzcd}
\]
That is, $F_\alpha \circ \pi_j = \pi_k$ for every $\alpha\colon j\to k$.
\end{definition}

\begin{definition}[Morphism of cones]\label{def:cone-morphism}
\index{cone!morphism}
Let $(N,\pi_j)$ and $(N',\pi'_j)$ be two cones over $F$.
A \textbf{morphism of cones} from $(N,\pi_j)$ to $(N',\pi'_j)$
is a morphism $u\colon N\to N'$ in~$\mathcal{C}$ such that
$\pi'_j \circ u = \pi_j$ for every $j\in\mathcal{J}$:
\[
\begin{tikzcd}
N \arrow[rr, "u"] \arrow[dr, "\pi_j"'] && N' \arrow[dl, "\pi'_j"] \\
& F_j &
\end{tikzcd}
\]
The cones over~$F$ and their morphisms form a category,
which we denote $\mathsf{Cone}(F)$.
\index{cone!category of cones}
\end{definition}

\begin{definition}[Cocone]\label{def:cocone}
\index{cocone}
\index{cocone!vertex}
\index{cocone!leg}
Dually, a \textbf{cocone under~$F$} is a pair
$(N,\,\{\iota_j\}_{j\in\mathcal{J}})$ where $\iota_j\colon F_j\to N$
and $\iota_k \circ F_\alpha = \iota_j$ for every $\alpha\colon j\to k$:
\[
\begin{tikzcd}
F_j \arrow[rr, "F_\alpha"] \arrow[dr, "\iota_j"'] & & F_k \arrow[dl, "\iota_k"] \\
& N &
\end{tikzcd}
\]
Cocones under~$F$ form a category $\mathsf{Cocone}(F)$.
\end{definition}

\begin{remark}\label{rem:cone-as-nat}
\index{constant functor}
\index{natural transformation!and cones}
Let $\Delta_N \colon \mathcal{J}\to\mathcal{C}$ denote the constant functor
sending every object to~$N$ and every morphism to~$\id_N$.
Then a cone $(N,\pi_j)$ over~$F$ is exactly a natural transformation
$\pi\colon \Delta_N \Rightarrow F$.
Dually, a cocone is a natural transformation $\iota\colon F\Rightarrow\Delta_N$.
This reformulation will be important when we study the relationship
between limits and the functor category $\Fun(\mathcal{J},\mathcal{C})$.
\end{remark}


%% ============================================================
\section{Limits}\label{sec:limits}
%% ============================================================

\begin{definition}[Limit]\label{def:limit}
\index{limit!definition}
\index{universal cone}
\index{limit!universal property}
Let $F\colon\mathcal{J}\to\mathcal{C}$ be a diagram.
A \textbf{limit} of~$F$ is a cone $(\varprojlim F,\,\{\pi_j\}_{j})$
that is \emph{terminal} in $\mathsf{Cone}(F)$.
Explicitly, this means:
\begin{enumerate}[label=(\alph*)]
\item $(\varprojlim F,\,\pi_j)$ is a cone over~$F$; and
\item for every cone $(N,\,\psi_j)$ over~$F$, there exists a
      \emph{unique} morphism $u\colon N\to\varprojlim F$
      such that $\pi_j\circ u = \psi_j$ for every $j\in\mathcal{J}$:
\end{enumerate}
\[
\begin{tikzcd}
N \arrow[dr, "\psi_j"'] \arrow[rr, dashed, "\exists!\, u"] & & \varprojlim F \arrow[dl, "\pi_j"] \\
& F_j &
\end{tikzcd}
\]
We also write $\lim_{j\in\mathcal{J}} F_j$ or simply $\lim F$.
The morphisms $\pi_j$ are called the \textbf{projection morphisms}
(or \textbf{canonical projections}).
\index{projection morphisms}
\end{definition}

\begin{proposition}[Uniqueness of limits]\label{prop:limit-unique}
\index{limit!uniqueness}
If $F\colon\mathcal{J}\to\mathcal{C}$ admits a limit, then the limit
is unique up to unique isomorphism.
\end{proposition}

\begin{proof}
This is a standard argument for terminal objects.
Suppose $(L,\pi_j)$ and $(L',\pi'_j)$ are both limits of~$F$.
By the universal property of~$L$, there is a unique
$u\colon L'\to L$ with $\pi_j\circ u = \pi'_j$ for all~$j$.
By the universal property of~$L'$, there is a unique
$v\colon L\to L'$ with $\pi'_j\circ v = \pi_j$ for all~$j$.
Then $\pi_j\circ(u\circ v) = \pi'_j\circ v = \pi_j$
for all~$j$, and by uniqueness applied to $L$ with the cone $(L,\pi_j)$,
we get $u\circ v = \id_L$.  Similarly $v\circ u = \id_{L'}$.
Hence $u$ is an isomorphism, and it is the unique isomorphism of cones.
\[
\begin{tikzcd}
L \arrow[rr, bend left, "v"] & & L' \arrow[ll, bend left, "u"] \\
& F_j \arrow[ul, leftarrow, "\pi_j"] \arrow[ur, leftarrow, "\pi'_j"'] &
\end{tikzcd}
\]
\end{proof}


%% ============================================================
\section{Colimits}\label{sec:colimits}
%% ============================================================

\begin{definition}[Colimit]\label{def:colimit}
\index{colimit!definition}
\index{universal cocone}
\index{colimit!universal property}
Let $F\colon\mathcal{J}\to\mathcal{C}$ be a diagram.
A \textbf{colimit} of~$F$ is a cocone $(\varinjlim F,\,\{\iota_j\}_{j})$
that is \emph{initial} in $\mathsf{Cocone}(F)$.
That is, for every cocone $(N,\,\psi_j)$ under~$F$, there exists a
unique morphism $u\colon \varinjlim F\to N$
with $u\circ\iota_j = \psi_j$ for all~$j$:
\[
\begin{tikzcd}
F_j \arrow[rr, "\iota_j"] \arrow[dr, "\psi_j"'] & & \varinjlim F \arrow[dl, dashed, "\exists!\, u"] \\
& N &
\end{tikzcd}
\]
The morphisms $\iota_j$ are the \textbf{coprojection morphisms}
(or \textbf{canonical injections}).
\index{coprojection morphisms}
\end{definition}

\begin{remark}\label{rem:colimit-dual}
\index{duality!limits and colimits}
A colimit of $F\colon\mathcal{J}\to\mathcal{C}$
is the same as a limit of $F^{\op}\colon\mathcal{J}^{\op}\to\mathcal{C}^{\op}$.
Hence every theorem about limits has a dual statement about colimits.
We will exploit this duality systematically.
\end{remark}


%% ============================================================
\section{Products and coproducts}\label{sec:products-coproducts}
%% ============================================================

Products and coproducts are limits and colimits indexed by a
discrete category.

\begin{definition}[Product]\label{def:product}
\index{product}
\index{product!universal property}
Let $\{A_i\}_{i\in I}$ be a family of objects in~$\mathcal{C}$,
viewed as a diagram $F\colon I_{\mathrm{disc}}\to\mathcal{C}$.
The \textbf{product} $\prod_{i\in I} A_i$ is the limit of this diagram.
Concretely, it is an object $\prod_{i} A_i$ equipped with projections
$\pi_i\colon \prod_j A_j\to A_i$ such that for every object~$N$
with morphisms $f_i\colon N\to A_i$, there exists a unique
$\langle f_i \rangle\colon N\to\prod_j A_j$ with
$\pi_i\circ\langle f_i \rangle = f_i$:
\[
\begin{tikzcd}
& N \arrow[dl, "f_1"'] \arrow[d, "f_2"] \arrow[dr, "f_3"]
      \arrow[dd, dashed, "\langle f_i\rangle" description] & \\
A_1 & A_2 & A_3 \\
& \prod_j A_j \arrow[ul, "\pi_1"] \arrow[u, "\pi_2"'] \arrow[ur, "\pi_3"'] &
\end{tikzcd}
\]
For a binary product we write $A\times B$ with projections
$\pi_1\colon A\times B\to A$ and $\pi_2\colon A\times B\to B$.
\index{product!binary}
\end{definition}

\begin{definition}[Coproduct]\label{def:coproduct}
\index{coproduct}
\index{coproduct!universal property}
Dually, the \textbf{coproduct} $\coprod_{i\in I} A_i$ is the
colimit of a family indexed by a discrete category.
It is an object with injections $\iota_i\colon A_i\to\coprod_j A_j$
such that for every object~$N$ with morphisms $g_i\colon A_i\to N$,
there is a unique $[g_i]\colon\coprod_j A_j\to N$
with $[g_i]\circ\iota_i = g_i$:
\[
\begin{tikzcd}
A_1 \arrow[dr, "\iota_1"] \arrow[ddr, bend right, "g_1"'] & &
A_2 \arrow[dl, "\iota_2"'] \arrow[ddl, bend left, "g_2"] \\
& \coprod_j A_j \arrow[d, dashed, "{[g_j]}" description] & \\
& N &
\end{tikzcd}
\]
For a binary coproduct we write $A\amalg B$ or $A + B$.
\index{coproduct!binary}
\end{definition}

\begin{example}[Products and coproducts in concrete categories]%
\label{ex:products-concrete}
\index{product!in Set@in $\Set$}
\index{product!in Grp@in $\Grp$}
\index{product!in Top@in $\Top$}
\index{product!in Ab@in $\Ab$}
\index{coproduct!in Set@in $\Set$}
\index{coproduct!in Grp@in $\Grp$}
\index{coproduct!in Top@in $\Top$}
\index{coproduct!in Ab@in $\Ab$}
\begin{enumerate}[label=(\roman*)]
\item \textbf{$\Set$.} The product is the Cartesian product
    $\prod_{i} A_i = \{\,(a_i)_{i\in I} : a_i\in A_i\,\}$
    with the usual projections.
    The coproduct is the disjoint union $\coprod_{i} A_i$.

\item \textbf{$\Grp$.} The product is the direct product of groups
    (Cartesian product with component-wise operations).
    The coproduct is the \emph{free product} $\ast_{i} G_i$,
    which is considerably more complicated than the direct product.
    \index{free product}

\item \textbf{$\Ab$.} Both the product and the finite coproduct
    coincide with the direct sum $\bigoplus_{i} A_i$
    (for finite families).  For infinite families the product
    is the full direct product and the coproduct is the direct sum
    (elements with finitely many nonzero components).
    \index{direct sum}

\item \textbf{$\Top$.} The product is the Cartesian product with the
    product topology (the coarsest topology making all projections
    continuous).  The coproduct is the disjoint union with the
    coproduct (disjoint union) topology.
    \index{product topology}

\item \textbf{$\Vect_\K$.} The product is the direct product
    of vector spaces and the finite coproduct is the direct sum.
    For infinite families, $\prod_i V_i \neq \bigoplus_i V_i$
    in general.
\end{enumerate}
\end{example}

\begin{exercise}\label{exer:product-comm-assoc}
\index{product!commutativity}
\index{product!associativity}
Let $\mathcal{C}$ be a category with binary products.
Show that:
\begin{enumerate}[label=(\alph*)]
\item $A \times B \cong B \times A$ (commutativity);
\item $(A \times B) \times C \cong A \times (B \times C)$ (associativity);
\item if $\mathcal{C}$ has a terminal object~$1$, then $A \times 1 \cong A$.
\end{enumerate}
\end{exercise}


%% ============================================================
\section{Equalisers and coequalisers}\label{sec:equalisers}
%% ============================================================

\begin{definition}[Equaliser]\label{def:equaliser}
\index{equaliser}
\index{equaliser!universal property}
Let $f,g\colon A\rightrightarrows B$ be a parallel pair.
The \textbf{equaliser} of $f$ and~$g$ is the limit of the diagram
$(\bullet\rightrightarrows\bullet)$: an object $E$ with a morphism
$e\colon E\to A$ such that $f\circ e = g\circ e$,
and universal with this property.
That is, for every $h\colon N\to A$ with $f\circ h = g\circ h$,
there exists a unique $u\colon N\to E$ with $e\circ u = h$:
\[
\begin{tikzcd}
N \arrow[dr, dashed, "\exists!\,u"'] \arrow[drr, bend left, "h"] & & \\
& E \arrow[r, "e"] & A \arrow[r, shift left, "f"] \arrow[r, shift right, "g"'] & B
\end{tikzcd}
\]
We write $E = \eq(f,g)$.
\end{definition}

\begin{example}\label{ex:equaliser-set}
\index{equaliser!in Set@in $\Set$}
In~$\Set$, the equaliser of $f,g\colon A\rightrightarrows B$ is
$\eq(f,g) = \{a\in A : f(a) = g(a)\}$
with the inclusion $e\colon\eq(f,g)\hookrightarrow A$.
\end{example}

\begin{proposition}\label{prop:equaliser-mono}
\index{equaliser!is monic}
\index{monomorphism!and equalisers}
Every equaliser is a monomorphism.
\end{proposition}

\begin{proof}
Let $e\colon E\to A$ be the equaliser of $f,g$.
Suppose $e\circ u = e\circ v$ for $u,v\colon N\to E$.
Then $f\circ(e\circ u) = f\circ(e\circ v)$ and
$g\circ(e\circ u) = g\circ(e\circ v)$.
Since $f\circ e = g\circ e$, the morphism $e\circ u = e\circ v$
satisfies $f\circ(e\circ u) = g\circ(e\circ u)$.
By the universal property of the equaliser applied to $h = e\circ u$,
there is a \emph{unique} morphism $N\to E$ making the triangle commute;
both $u$ and~$v$ satisfy this, so $u = v$.
\end{proof}

\begin{definition}[Coequaliser]\label{def:coequaliser}
\index{coequaliser}
\index{coequaliser!universal property}
Dually, the \textbf{coequaliser} of $f,g\colon A\rightrightarrows B$
is an object $Q$ with $q\colon B\to Q$ such that $q\circ f = q\circ g$,
universal from below:
\[
\begin{tikzcd}
A \arrow[r, shift left, "f"] \arrow[r, shift right, "g"'] &
B \arrow[r, "q"] \arrow[rrd, bend right, "h"'] &
Q \arrow[rd, dashed, "\exists!\,u"] & \\
& & & N
\end{tikzcd}
\]
We write $Q = \coeq(f,g)$.
\end{definition}

\begin{example}\label{ex:coequaliser-set}
\index{coequaliser!in Set@in $\Set$}
In~$\Set$, the coequaliser of $f,g\colon A\rightrightarrows B$ is the
quotient $B/{\sim}$ where $\sim$ is the equivalence relation generated
by $f(a)\sim g(a)$ for all $a\in A$.
\end{example}

\begin{proposition}\label{prop:coequaliser-epi}
\index{coequaliser!is epic}
\index{epimorphism!and coequalisers}
Every coequaliser is an epimorphism.
\end{proposition}

\begin{proof}
Dual to the proof that equalisers are monic (Proposition~\ref{prop:equaliser-mono}).
\end{proof}


%% ============================================================
\section{Kernels and cokernels}\label{sec:kernels-cokernels}
%% ============================================================

\begin{definition}[Kernel]\label{def:kernel-categorical}
\index{kernel!categorical}
\index{zero morphism}
Let $\mathcal{C}$ be a category with a zero object~$0$.
For a morphism $f\colon A\to B$, the \textbf{kernel} of~$f$ is the
equaliser of~$f$ and the zero morphism $0_{AB}\colon A\to 0\to B$:
\[
\Ker(f) = \eq(f,\,0_{AB}).
\]
Concretely, it is a morphism $k\colon K\to A$ such that $f\circ k = 0$
and every $h\colon N\to A$ with $f\circ h = 0$ factors uniquely through~$k$:
\[
\begin{tikzcd}
N \arrow[dr, dashed, "\exists!\,u"'] \arrow[drr, bend left, "h"] & & \\
& K \arrow[r, "k"] & A \arrow[r, "f"] & B
\end{tikzcd}
\]
\end{definition}

\begin{definition}[Cokernel]\label{def:cokernel-categorical}
\index{cokernel!categorical}
The \textbf{cokernel} of $f\colon A\to B$ is the coequaliser of~$f$
and~$0_{AB}$:
\[
\coker(f) = \coeq(f,\,0_{AB}).
\]
It is a morphism $c\colon B\to Q$ with $c\circ f = 0$, universal
with this property.
\end{definition}

\begin{example}\label{ex:ker-coker-ab}
\index{kernel!in Ab@in $\Ab$}
\index{cokernel!in Ab@in $\Ab$}
In~$\Ab$, the kernel of a homomorphism $f\colon A\to B$ is
$\Ker(f) = \{a\in A : f(a) = 0\}$ with the inclusion into~$A$,
and the cokernel is $\coker(f) = B/\im(f)$ with the projection.
\end{example}

\begin{exercise}\label{exer:kernel-mono}
\index{kernel!is monic}
Show that in any category with a zero object, a kernel is a monomorphism.
(This follows from the fact that equalisers are monic.)
\end{exercise}


%% ============================================================
\section{Pullbacks and pushouts}\label{sec:pullbacks-pushouts}
%% ============================================================

\begin{definition}[Pullback]\label{def:pullback}
\index{pullback}
\index{pullback!universal property}
\index{fibre product}
Given a cospan $A \xrightarrow{f} C \xleftarrow{g} B$, the
\textbf{pullback} (or \textbf{fibre product}) is the limit of this
diagram.  It is an object $A\times_C B$ equipped with morphisms
$p_1\colon A\times_C B\to A$ and $p_2\colon A\times_C B\to B$
such that $f\circ p_1 = g\circ p_2$, and universal with this property:
\[
\begin{tikzcd}
N \arrow[drr, bend left, "\psi_2"]
  \arrow[ddr, bend right, "\psi_1"']
  \arrow[dr, dashed, "\exists!\,u" description] & & \\
& A\times_C B \arrow[r, "p_2"] \arrow[d, "p_1"']
    \arrow[dr, phantom, "\lrcorner" very near start]
  & B \arrow[d, "g"] \\
& A \arrow[r, "f"'] & C
\end{tikzcd}
\]
The symbol $\lrcorner$ in the square indicates that it is a pullback square.
\end{definition}

\begin{example}[Pullback in $\Set$]\label{ex:pullback-set}
\index{pullback!in Set@in $\Set$}
Given $f\colon A\to C$ and $g\colon B\to C$ in $\Set$,
\[
    A\times_C B = \{(a,b)\in A\times B : f(a) = g(b)\}.
\]
The projections are $(a,b)\mapsto a$ and $(a,b)\mapsto b$.
\end{example}

\begin{example}[Pullback in $\Top$]\label{ex:pullback-top}
\index{pullback!in Top@in $\Top$}
In $\Top$, the pullback is the set-theoretic fibre product
$\{(a,b)\in A\times B : f(a) = g(b)\}$ equipped with the
subspace topology inherited from $A\times B$.
\end{example}

\begin{example}[Pullback as inverse image]\label{ex:pullback-preimage}
\index{pullback!as inverse image}
In~$\Set$, consider $f\colon A\to C$ and the inclusion
$\iota\colon S\hookrightarrow C$ of a subset.
The pullback $A\times_C S$ is the preimage
$f^{-1}(S) = \{a\in A : f(a)\in S\}$.
\[
\begin{tikzcd}
f^{-1}(S) \arrow[r, hook] \arrow[d]
  \arrow[dr, phantom, "\lrcorner" very near start]
& A \arrow[d, "f"] \\
S \arrow[r, hook, "\iota"'] & C
\end{tikzcd}
\]
\end{example}

\begin{proposition}\label{prop:pullback-of-mono}
\index{pullback!of monomorphism}
\index{monomorphism!stable under pullback}
If $g\colon B\to C$ is a monomorphism, then $p_1\colon A\times_C B\to A$
is also a monomorphism.  That is, monomorphisms are stable under pullback.
\end{proposition}

\begin{proof}
Suppose $p_1\circ u = p_1\circ v$ for $u,v\colon N\to A\times_C B$.
Then $f\circ p_1\circ u = f\circ p_1\circ v$, hence
$g\circ p_2\circ u = g\circ p_2\circ v$, hence
$p_2\circ u = p_2\circ v$ since $g$ is monic.
Now $u$ and $v$ are both morphisms $N\to A\times_C B$ satisfying
$p_1\circ u = p_1\circ v$ and $p_2\circ u = p_2\circ v$.
By the uniqueness part of the pullback universal property, $u = v$.
\end{proof}

\begin{definition}[Pushout]\label{def:pushout}
\index{pushout}
\index{pushout!universal property}
\index{amalgamated sum}
Given a span $A \xleftarrow{f} C \xrightarrow{g} B$, the
\textbf{pushout} (or \textbf{amalgamated sum}) is the colimit of
this diagram.  It is an object $A\amalg_C B$ with morphisms
$q_1\colon A\to A\amalg_C B$ and $q_2\colon B\to A\amalg_C B$
such that $q_1\circ f = q_2\circ g$, universal with this property:
\[
\begin{tikzcd}
C \arrow[r, "g"] \arrow[d, "f"'] & B \arrow[d, "q_2"]
  \arrow[ddr, bend left, "\psi_2"] & \\
A \arrow[r, "q_1"'] \arrow[rrd, bend right, "\psi_1"']
& A\amalg_C B \arrow[rd, dashed, "\exists!\,u" description]
  \arrow[from=1-1, phantom, "\ulcorner" very near end, pos=0.9] & \\
& & N
\end{tikzcd}
\]
\end{definition}

\begin{example}[Pushout in $\Set$]\label{ex:pushout-set}
\index{pushout!in Set@in $\Set$}
Given a span $A\xleftarrow{f} C \xrightarrow{g} B$ in~$\Set$,
\[
    A\amalg_C B = (A\amalg B)/{\sim}
\]
where $\sim$ is generated by $f(c)\sim g(c)$ for all $c\in C$.
\end{example}

\begin{example}[Pushout in $\Top$]\label{ex:pushout-top}
\index{pushout!in Top@in $\Top$}
\index{attaching space}
In~$\Top$, pushouts correspond to gluing constructions.
If $f\colon A\hookrightarrow X$ is a subspace inclusion and
$g\colon A\to Y$ is continuous, then $X\amalg_A Y$ is the space
obtained by gluing $Y$ to~$X$ along~$A$ via~$g$.
A fundamental instance: if $D^n$ is the $n$-disc,
$S^{n-1}$ its boundary, and $\varphi\colon S^{n-1}\to X$ an attaching map,
then
\[
\begin{tikzcd}
S^{n-1} \arrow[r, hook] \arrow[d, "\varphi"'] & D^n \arrow[d] \\
X \arrow[r] & X\cup_\varphi D^n
\end{tikzcd}
\]
is a pushout square giving the space $X$ with an $n$-cell attached.
\end{example}

\begin{exercise}\label{exer:pushout-epi}
\index{pushout!of epimorphism}
Show that epimorphisms are stable under pushout.
(This is dual to Proposition~\ref{prop:pullback-of-mono}.)
\end{exercise}

\begin{exercise}\label{exer:pullback-pasting}
\index{pullback!pasting lemma}
\textbf{(Pullback pasting lemma.)}
Consider a diagram
\[
\begin{tikzcd}
A \arrow[r] \arrow[d] & B \arrow[r] \arrow[d] & C \arrow[d] \\
D \arrow[r] & E \arrow[r] & F
\end{tikzcd}
\]
Show that if both squares are pullbacks, then the outer rectangle
is a pullback.  Show also that if the right square and the outer
rectangle are pullbacks, then the left square is a pullback.
\end{exercise}


%% ============================================================
\section{General existence theorem}\label{sec:limit-existence}
%% ============================================================

The following fundamental theorem reduces the existence of all finite
limits to two basic building blocks.

\begin{theorem}[Products and equalisers give all finite limits]%
\label{thm:prod-eq-limits}
\index{limit!existence theorem}
\index{finite limits}
\index{completeness!finite}
A category $\mathcal{C}$ has all finite limits if and only if it has all
finite products and all equalisers.
\end{theorem}

\begin{proof}
The ``only if'' direction is clear, since products and equalisers are
special cases of finite limits.

For the ``if'' direction, let $F\colon\mathcal{J}\to\mathcal{C}$
be a diagram where $\mathcal{J}$ is a finite category.
We construct $\lim F$ as an equaliser of two maps between products.

\medskip\noindent
\textbf{Step 1.} Form the products
\[
    P = \prod_{j\in\Ob(\mathcal{J})} F_j,
    \qquad
    Q = \prod_{\alpha\in\Mor(\mathcal{J})} F_{\mathrm{cod}(\alpha)},
\]
where for each morphism $\alpha\colon j\to k$ in~$\mathcal{J}$,
the factor of~$Q$ indexed by~$\alpha$ is $F_k = F_{\mathrm{cod}(\alpha)}$.

\medskip\noindent
\textbf{Step 2.} Define two morphisms $s,t\colon P\to Q$ as follows.
For each morphism $\alpha\colon j\to k$ in~$\mathcal{J}$, the
$\alpha$-component of~$s$ is the composite
\[
    s_\alpha\colon P \xrightarrow{\pi_j} F_j \xrightarrow{F_\alpha} F_k,
\]
and the $\alpha$-component of~$t$ is
\[
    t_\alpha\colon P \xrightarrow{\pi_k} F_k.
\]

\[
\begin{tikzcd}
\eq(s,t) \arrow[r, "e"] & P \arrow[r, shift left, "s"] \arrow[r, shift right, "t"'] & Q
\end{tikzcd}
\]

\medskip\noindent
\textbf{Step 3.} Let $e\colon E\to P$ be the equaliser of $s$ and~$t$.
We claim $E = \lim F$ with projections
$\pi_j\circ e\colon E\to F_j$.

The condition $s\circ e = t\circ e$ means that for every
$\alpha\colon j\to k$,
\[
    F_\alpha \circ \pi_j \circ e = \pi_k \circ e,
\]
which is exactly the cone condition.

Conversely, if $(N,\psi_j)$ is any cone over~$F$, the universal property
of~$P$ gives a unique $h\colon N\to P$ with $\pi_j\circ h = \psi_j$.
The cone condition ensures $s\circ h = t\circ h$, so the universal
property of the equaliser gives a unique $u\colon N\to E$ with
$e\circ u = h$.
Then $\pi_j\circ e\circ u = \pi_j\circ h = \psi_j$, as required.
\end{proof}

\begin{corollary}\label{cor:complete-category}
\index{complete category}
\index{cocomplete category}
A category $\mathcal{C}$ is \textbf{(finitely) complete} if and only if
it has (finite) products and equalisers.  Dually, $\mathcal{C}$ is
\textbf{(finitely) cocomplete} if and only if it has (finite) coproducts
and coequalisers.
\end{corollary}

\begin{example}\label{ex:complete-categories}
\index{complete category!examples}
The categories $\Set$, $\Grp$, $\Ab$, $\Ring$, $\Top$, $\Vect_\K$,
and $\Mod_R$ (for any ring~$R$) are all complete and cocomplete.
\end{example}

\begin{remark}\label{rem:limits-general}
\index{limit!small}
The same argument works for arbitrary (small) limits: a category
with all small products and equalisers has all small limits.
A category with all small limits is called \textbf{complete}.
\end{remark}


%% ============================================================
\section{Limits and the Hom functor}\label{sec:limits-hom}
%% ============================================================

\begin{theorem}[Limits commute with Hom]\label{thm:lim-hom}
\index{limit!and Hom functor}
\index{Hom functor!preserves limits}
Let $F\colon\mathcal{J}\to\mathcal{C}$ be a diagram admitting a limit
$L = \lim F$.  For any object $X\in\mathcal{C}$,
\[
    \Hom_{\mathcal{C}}(X,\,\lim_{j} F_j)
    \;\cong\;
    \lim_{j}\, \Hom_{\mathcal{C}}(X,\,F_j),
\]
naturally in~$X$.  Here the limit on the right is computed in $\Set$.
\end{theorem}

\begin{proof}
Let $(L,\pi_j)$ be the limit cone.  Define
\[
    \Phi\colon \Hom(X,L) \longrightarrow \lim_j\Hom(X,F_j)
\]
by $\Phi(h) = (\pi_j\circ h)_{j\in\mathcal{J}}$.  We verify that
$(\pi_j\circ h)_j$ lies in the limit: for $\alpha\colon j\to k$,
$F_\alpha\circ\pi_j\circ h = \pi_k\circ h$ since $(L,\pi_j)$ is a cone.

The map $\Phi$ is \textbf{injective}: if $\Phi(h)=\Phi(h')$,
then $\pi_j\circ h = \pi_j\circ h'$ for all~$j$, so $h = h'$
by the uniqueness in the universal property of~$L$.

The map $\Phi$ is \textbf{surjective}: given
$(f_j)_j\in\lim_j\Hom(X,F_j)$, the family $(f_j\colon X\to F_j)_j$
is a cone over~$F$ with vertex~$X$.
By the universal property of~$L$, there exists a unique
$h\colon X\to L$ with $\pi_j\circ h = f_j$ for all~$j$.
Then $\Phi(h) = (f_j)_j$.

Naturality in~$X$ is straightforward: for $\varphi\colon X'\to X$,
$\Phi(\,h\circ\varphi\,) = (\pi_j\circ h\circ\varphi)_j$,
which is the image of $\Phi(h)$ under pre-composition with~$\varphi$.
\end{proof}

\begin{corollary}\label{cor:representable-preserves-limits}
\index{representable functor!preserves limits}
Every representable functor $\Hom(X,-)$ preserves all limits
that exist in~$\mathcal{C}$.
\end{corollary}

\begin{remark}\label{rem:colim-hom}
\index{colimit!and Hom functor}
Dually, for colimits:
\[
    \Hom_{\mathcal{C}}(\colim_j F_j,\, X)
    \;\cong\;
    \lim_j\, \Hom_{\mathcal{C}}(F_j,\, X).
\]
Note the \emph{contravariance}: the colimit on the left becomes a
limit on the right.
\end{remark}


%% ============================================================
\section{Functors and limits}\label{sec:functors-limits}
%% ============================================================

\begin{definition}[Preservation of limits]\label{def:preserve-limits}
\index{functor!preserving limits}
\index{limit!preservation}
A functor $G\colon\mathcal{C}\to\mathcal{D}$ \textbf{preserves the limit}
of a diagram $F\colon\mathcal{J}\to\mathcal{C}$ if whenever
$(L,\pi_j)$ is a limit cone for~$F$, then $(GL,\,G\pi_j)$ is a limit
cone for $G\circ F$ in~$\mathcal{D}$.

We say $G$ \textbf{preserves all limits} (or is
\textbf{continuous}\index{continuous functor}) if it preserves
the limit of every diagram that admits a limit.
\end{definition}

\begin{definition}[Reflection and creation of limits]%
\label{def:reflect-create-limits}
\index{functor!reflecting limits}
\index{functor!creating limits}
\index{limit!reflection}
\index{limit!creation}
A functor $G\colon\mathcal{C}\to\mathcal{D}$ \textbf{reflects the limit}
of $F\colon\mathcal{J}\to\mathcal{C}$ if whenever $(L,\pi_j)$ is a cone
for~$F$ such that $(GL,G\pi_j)$ is a limit cone for $G\circ F$,
then $(L,\pi_j)$ is already a limit cone for~$F$.

$G$ \textbf{creates the limit} of~$F$ if for every limit cone
$(M,\psi_j)$ for $G\circ F$ in~$\mathcal{D}$, there exists a unique
cone $(L,\pi_j)$ for~$F$ with $GL = M$ and $G\pi_j = \psi_j$,
and moreover this cone is a limit of~$F$.
\end{definition}

\begin{example}\label{ex:forgetful-creates}
\index{forgetful functor!creates limits}
The forgetful functor $U\colon\Grp\to\Set$ creates all limits.
Indeed, given a diagram of groups, one first forms the limit in~$\Set$
(using the Cartesian product construction), then shows this set carries
a unique group structure making all projections group homomorphisms.
The same is true for $\Ab\to\Set$, $\Ring\to\Set$, $\Mod_R\to\Set$, etc.
\end{example}


%% ============================================================
\section{Filtered colimits}\label{sec:filtered-colimits}
%% ============================================================

\begin{definition}[Filtered category]\label{def:filtered-category}
\index{filtered category}
A small category $\mathcal{J}$ is \textbf{filtered} if:
\begin{enumerate}[label=(\roman*)]
\item $\mathcal{J}$ is non-empty;
\item for every pair of objects $j,k\in\mathcal{J}$, there exists
      an object $\ell$ and morphisms $j\to\ell$ and $k\to\ell$;
\item for every pair of parallel morphisms $\alpha,\beta\colon j\rightrightarrows k$,
      there exists a morphism $\gamma\colon k\to\ell$ such that
      $\gamma\circ\alpha = \gamma\circ\beta$.
\end{enumerate}
\end{definition}

\begin{definition}[Filtered colimit]\label{def:filtered-colimit}
\index{filtered colimit}
\index{colimit!filtered}
A \textbf{filtered colimit} (or \textbf{direct limit}) is a colimit
of a diagram $F\colon\mathcal{J}\to\mathcal{C}$ where $\mathcal{J}$
is filtered.  We also write $\varinjlim F$.
\end{definition}

\begin{example}\label{ex:filtered-colimit-directed}
\index{directed system}
\index{directed colimit}
A partially ordered set $(I,\leq)$ that is directed
(every pair has an upper bound) is a filtered category.
A functor $F\colon I\to\mathcal{C}$ is a directed system,
and its colimit is the classical direct limit.
\end{example}

\begin{theorem}\label{thm:filtered-colim-set}
\index{filtered colimit!in Set@in $\Set$}
In~$\Set$, filtered colimits commute with finite limits.
More precisely, if $\mathcal{J}$ is filtered and $\mathcal{K}$ is finite,
then for any functor $F\colon\mathcal{J}\times\mathcal{K}\to\Set$,
\[
    \varinjlim_j\, \varprojlim_k\, F(j,k)
    \;\cong\;
    \varprojlim_k\, \varinjlim_j\, F(j,k).
\]
\end{theorem}

\begin{proof}
We sketch the proof for the key cases.  The filtered colimit
$\varinjlim_j F_j$ in~$\Set$ is computed as
$\bigl(\coprod_j F_j\bigr)/{\sim}$
where $x\in F_j$ and $y\in F_k$ satisfy $x\sim y$ iff there exist
morphisms $\alpha\colon j\to\ell$ and $\beta\colon k\to\ell$ with
$F_\alpha(x) = F_\beta(y)$.

\smallskip\noindent
\emph{Commutation with finite products:}
An element of $\varinjlim_j(F_j\times G_j)$ is an equivalence class
$[(x,y)]$ with $x\in F_j$, $y\in G_j$.  The filteredness condition~(ii)
ensures that any pair of elements from
$\varinjlim_j F_j$ and $\varinjlim_j G_j$ can be represented at a
common index, establishing a bijection with
$(\varinjlim_j F_j)\times(\varinjlim_j G_j)$.

\smallskip\noindent
\emph{Commutation with equalisers:}
Given parallel natural transformations $f,g\colon F\rightrightarrows G$,
the filteredness condition~(iii) ensures that
$\varinjlim_j\eq(f_j,g_j) \cong \eq(\varinjlim_j f_j,\,\varinjlim_j g_j)$.

Since finite products and equalisers generate all finite limits
(Theorem~\ref{thm:prod-eq-limits}), the general statement follows.
\end{proof}

\begin{example}\label{ex:filtered-colim-algebra}
\index{filtered colimit!of groups}
\index{filtered colimit!of rings}
A filtered colimit of groups is a group, a filtered colimit of rings
is a ring, etc.  For instance, $\Q = \varinjlim_{n\geq 1} \frac{1}{n}\Z$
as abelian groups, and algebraic closures can be constructed as
filtered colimits of finite extensions.
\end{example}


%% ============================================================
\section{Exercises}\label{sec:exercises-ch3}
%% ============================================================

\begin{exercise}\label{exer:terminal-limit}
\index{terminal object!as limit}
\index{initial object!as colimit}
Show that a terminal object is the limit of the empty diagram
(i.e.\ $\mathcal{J} = \varnothing$), and an initial object is the
colimit of the empty diagram.
\end{exercise}

\begin{exercise}\label{exer:equaliser-from-pullback}
\index{equaliser!from pullback}
Show that if $\mathcal{C}$ has binary products and pullbacks,
then $\mathcal{C}$ has equalisers.
\emph{Hint:} Given $f,g\colon A\rightrightarrows B$, consider the
pullback of $\langle f,g\rangle\colon A\to B\times B$ along
the diagonal $\Delta\colon B\to B\times B$.
\end{exercise}

\begin{exercise}\label{exer:product-from-pullback}
\index{product!from pullback}
Show that if $\mathcal{C}$ has a terminal object and pullbacks,
then $\mathcal{C}$ has all finite limits.
\end{exercise}

\begin{exercise}\label{exer:colimit-in-set}
\index{colimit!in Set@in $\Set$}
Show that for any small diagram $F\colon\mathcal{J}\to\Set$,
the colimit is
\[
    \varinjlim F \;=\;
    \Bigl(\coprod_{j\in\mathcal{J}} F_j\Bigr)\Big/{\sim}
\]
where $\sim$ is the equivalence relation generated by:
$x\in F_j$ is equivalent to $F_\alpha(x)\in F_k$ for every
$\alpha\colon j\to k$ in~$\mathcal{J}$.
\end{exercise}

\begin{exercise}\label{exer:limit-functor-category}
\index{limit!in functor category}
Let $\mathcal{C}$ be a complete category and $\mathcal{J}$ a small category.
Show that $\Fun(\mathcal{J},\mathcal{C})$ is also complete,
and that limits are computed ``pointwise'': for a diagram
$\Phi\colon\mathcal{K}\to\Fun(\mathcal{J},\mathcal{C})$,
\[
    (\lim_k \Phi_k)(j) \cong \lim_k(\Phi_k(j))
\]
for each $j\in\mathcal{J}$.
\end{exercise}

\begin{exercise}\label{exer:limit-preserves-mono}
\index{limit!preserves monomorphisms}
Let $\mathcal{C}$ be a category with equalisers.
Show that if $\alpha\colon F\Rightarrow G$ is a natural transformation
between diagrams $F,G\colon\mathcal{J}\to\mathcal{C}$ such that each
$\alpha_j$ is a monomorphism, and if $\lim F$ and $\lim G$ both exist,
then the induced morphism $\lim F\to\lim G$ is also a monomorphism.
\end{exercise}

\begin{exercise}\label{exer:inverse-limit-abelian-groups}
\index{inverse limit!of abelian groups}
\index{$p$-adic integers}
Let $p$ be a prime.  Consider the inverse system
$\cdots\to\Z/p^3\Z\to\Z/p^2\Z\to\Z/p\Z$
in~$\Ab$.  Show that the inverse limit is the ring of
$p$-adic integers~$\Z_p$.
\end{exercise}

\begin{exercise}\label{exer:filtered-colim-exact}
\index{filtered colimit!and exactness}
Show that filtered colimits in $\Mod_R$ (for $R$ a ring) are exact.
That is, if $0\to F\to G\to H\to 0$ is a short exact sequence
of directed systems, then
$0\to\varinjlim F\to\varinjlim G\to\varinjlim H\to 0$ is exact.
\end{exercise}

\index{limit|)}
\index{colimit|)}

\newpage

%%%%%%%%%%%%%%%%%%%%%%%%%%%%%%%%%%%%%%%%%%%%%%%%%%%%%%%%%%%%
%%           CHAPTER 4: ADJUNCTIONS                       %%
%%%%%%%%%%%%%%%%%%%%%%%%%%%%%%%%%%%%%%%%%%%%%%%%%%%%%%%%%%%%

\chapter{Adjunctions}\label{ch:adjunctions}
\minitoc

\vspace{1em}

Adjunctions are arguably the most important concept in category theory.
Saunders Mac\,Lane famously wrote that ``adjoint functors arise everywhere.''
\index{Mac Lane, Saunders}
An adjunction captures a universal relationship between two functors
going in opposite directions, generalising at a stroke the notions of
free construction, tensor--hom duality, and a host of other fundamental
correspondences throughout mathematics.

In this chapter we give three equivalent definitions of adjunctions,
prove their equivalence in detail, explore the fundamental examples,
and establish the key theorems---most notably, that left adjoints
preserve colimits and right adjoints preserve limits.  We conclude
with Freyd's adjoint functor theorems.

\index{adjunction|(}

%% ============================================================
\section{Definition via natural bijection}\label{sec:adj-hom-def}
%% ============================================================

\begin{definition}[Adjunction]\label{def:adjunction}
\index{adjunction!definition}
\index{left adjoint}
\index{right adjoint}
\index{adjunction!natural bijection}
Let $\mathcal{C}$ and $\mathcal{D}$ be categories and
$F\colon\mathcal{C}\to\mathcal{D}$,
$G\colon\mathcal{D}\to\mathcal{C}$ be functors.
We say that $F$ is \textbf{left adjoint} to~$G$
(equivalently, $G$ is \textbf{right adjoint} to~$F$),
written $F\dashv G$, if there is a natural bijection
\[
    \Phi_{A,B}\colon
    \Hom_{\mathcal{D}}(FA,\,B)
    \;\xrightarrow{\;\sim\;}\;
    \Hom_{\mathcal{C}}(A,\,GB)
\]
for all $A\in\mathcal{C}$ and $B\in\mathcal{D}$,
natural in both variables.  That is, $\Phi$ is a natural isomorphism
\[
    \Phi\colon \Hom_{\mathcal{D}}(F-,\,-)
    \;\xRightarrow{\;\sim\;}\;
    \Hom_{\mathcal{C}}(-,\,G-)
\]
of bifunctors $\mathcal{C}^{\op}\times\mathcal{D}\to\Set$.
\end{definition}

\begin{remark}\label{rem:adj-notation}
\index{adjunction!notation}
We depict an adjunction as:
\[
\begin{tikzcd}
\mathcal{C}
  \arrow[rr, bend left=25, "F", ""{name=F, below}]
& \bot &
\mathcal{D}
  \arrow[ll, bend left=25, "G", ""{name=G, above}]
  \arrow[from=F, to=G, phantom, "\dashv"]
\end{tikzcd}
\]
The symbol $\bot$ (or $\dashv$) indicates $F\dashv G$.
\end{remark}

Let us spell out what naturality means in this context.

\begin{proposition}[Naturality conditions]\label{prop:adj-naturality}
\index{adjunction!naturality}
The bijection $\Phi_{A,B}$ is natural in $A$ and $B$ if and only if:
\begin{enumerate}[label=(\roman*)]
\item \textbf{Naturality in $B$:}
For every $g\colon B\to B'$ and $f\colon FA\to B$,
\[
    \Phi_{A,B'}(g\circ f) = Gg \circ \Phi_{A,B}(f).
\]
\item \textbf{Naturality in $A$:}
For every $h\colon A'\to A$ and $f\colon FA\to B$,
\[
    \Phi_{A',B}(f\circ Fh) = \Phi_{A,B}(f)\circ h.
\]
\end{enumerate}
\end{proposition}

\begin{proof}
This is a direct translation of the naturality squares for the
bifunctor $\Hom_\mathcal{D}(F-,-) \Rightarrow \Hom_\mathcal{C}(-,G-)$.
For (i), the naturality square in~$B$ is:
\[
\begin{tikzcd}
\Hom(FA,B) \arrow[r, "\Phi_{A,B}"] \arrow[d, "g\circ -"'] &
\Hom(A,GB) \arrow[d, "Gg\circ -"] \\
\Hom(FA,B') \arrow[r, "\Phi_{A,B'}"'] & \Hom(A,GB')
\end{tikzcd}
\]
Commutativity gives $Gg\circ\Phi_{A,B}(f) = \Phi_{A,B'}(g\circ f)$.
The argument for (ii) is analogous.
\end{proof}


%% ============================================================
\section{Unit and counit}\label{sec:unit-counit}
%% ============================================================

\begin{definition}[Unit and counit]\label{def:unit-counit}
\index{unit of adjunction}
\index{counit of adjunction}
\index{adjunction!unit}
\index{adjunction!counit}
Given an adjunction $F\dashv G$ with bijection~$\Phi$, define:
\begin{itemize}
\item The \textbf{unit} $\eta\colon \Id_{\mathcal{C}}\Rightarrow GF$
    by $\eta_A = \Phi_{A,FA}(\id_{FA})$ for each $A\in\mathcal{C}$.
\item The \textbf{counit} $\varepsilon\colon FG\Rightarrow\Id_{\mathcal{D}}$
    by $\varepsilon_B = \Phi^{-1}_{GB,B}(\id_{GB})$ for each
    $B\in\mathcal{D}$.
\end{itemize}
\end{definition}

\begin{proposition}\label{prop:unit-counit-nat}
\index{unit of adjunction!naturality}
\index{counit of adjunction!naturality}
The unit $\eta$ and counit $\varepsilon$ are natural transformations.
Moreover, the bijection $\Phi$ can be recovered from them:
\[
    \Phi_{A,B}(f\colon FA\to B) = Gf\circ\eta_A,
    \qquad
    \Phi^{-1}_{A,B}(g\colon A\to GB) = \varepsilon_B\circ Fg.
\]
\end{proposition}

\begin{proof}
\textbf{Recovery formula.}
For $f\colon FA\to B$, naturality of~$\Phi$ in~$B$ gives
\[
    \Phi_{A,B}(f) = \Phi_{A,B}(f\circ\id_{FA})
    = Gf\circ\Phi_{A,FA}(\id_{FA}) = Gf\circ\eta_A.
\]
Similarly, for $g\colon A\to GB$, naturality in~$A$ gives
\[
    \Phi^{-1}_{A,B}(g)
    = \Phi^{-1}_{A,B}(\id_{GB}\circ g)
    = \Phi^{-1}_{GB,B}(\id_{GB})\circ Fg
    = \varepsilon_B\circ Fg.
\]

\textbf{Naturality of $\eta$.}
We must show that for $h\colon A\to A'$, $\eta_{A'}\circ h = GFh\circ\eta_A$:
\[
\begin{tikzcd}
A \arrow[r, "h"] \arrow[d, "\eta_A"'] & A' \arrow[d, "\eta_{A'}"] \\
GFA \arrow[r, "GFh"'] & GFA'
\end{tikzcd}
\]
We have $\eta_{A'}\circ h = \Phi_{A',FA'}(\id_{FA'})\circ h
= \Phi_{A,FA'}(\id_{FA'}\circ Fh) = \Phi_{A,FA'}(Fh)
= GFh\circ\eta_A$, using naturality in~$A$ for the second
equality and the recovery formula for the last.

\textbf{Naturality of $\varepsilon$} is proved dually.
\end{proof}


%% ============================================================
\section{Triangle identities}\label{sec:triangle-identities}
%% ============================================================

\begin{theorem}[Triangle identities]\label{thm:triangle-identities}
\index{triangle identities}
\index{zig-zag identities}
\index{adjunction!triangle identities}
Let $F\dashv G$ with unit $\eta$ and counit~$\varepsilon$.
Then the following two composites are identities:
\begin{enumerate}[label=(\alph*)]
\item $\varepsilon F \circ F\eta = \id_F$, i.e.\ for every $A$:
\[
\begin{tikzcd}
FA \arrow[r, "F\eta_A"] \arrow[dr, equal] & FGFA \arrow[d, "\varepsilon_{FA}"] \\
& FA
\end{tikzcd}
\]
\item $G\varepsilon \circ \eta G = \id_G$, i.e.\ for every $B$:
\[
\begin{tikzcd}
GB \arrow[r, "\eta_{GB}"] \arrow[dr, equal] & GFGB \arrow[d, "G\varepsilon_B"] \\
& GB
\end{tikzcd}
\]
\end{enumerate}
These are called the \textbf{triangle identities} (or \textbf{zig-zag
identities}) because of their shape in the following diagram:
\[
\begin{tikzcd}[column sep=large]
F \arrow[r, Rightarrow, "F\eta"] \arrow[dr, equal]
  & FGF \arrow[d, Rightarrow, "\varepsilon F"] &
G \arrow[r, Rightarrow, "\eta G"] \arrow[dr, equal]
  & GFG \arrow[d, Rightarrow, "G\varepsilon"] \\
& F & & G
\end{tikzcd}
\]
\end{theorem}

\begin{proof}
\textbf{(a):}
We must show $\varepsilon_{FA}\circ F\eta_A = \id_{FA}$.
By the recovery formula (Proposition~\ref{prop:unit-counit-nat}),
$\Phi^{-1}_{A,FA}(\eta_A) = \varepsilon_{FA}\circ F\eta_A$.
But $\eta_A = \Phi_{A,FA}(\id_{FA})$, so
$\Phi^{-1}_{A,FA}(\eta_A) = \id_{FA}$.
Hence $\varepsilon_{FA}\circ F\eta_A = \id_{FA}$.

\medskip
\textbf{(b):}
We must show $G\varepsilon_B\circ\eta_{GB} = \id_{GB}$.
By the recovery formula,
$\Phi_{GB,B}(\varepsilon_B) = G\varepsilon_B\circ\eta_{GB}$.
But $\varepsilon_B = \Phi^{-1}_{GB,B}(\id_{GB})$, so
$\Phi_{GB,B}(\varepsilon_B) = \id_{GB}$.
Hence $G\varepsilon_B\circ\eta_{GB} = \id_{GB}$.
\end{proof}


%% ============================================================
\section{Equivalence of definitions}\label{sec:adj-equivalence}
%% ============================================================

We now state and prove the equivalence of three different ways
of specifying an adjunction.

\begin{theorem}[Three equivalent definitions of adjunction]%
\label{thm:adj-equiv-defs}
\index{adjunction!equivalent definitions}
Let $F\colon\mathcal{C}\to\mathcal{D}$ and
$G\colon\mathcal{D}\to\mathcal{C}$ be functors.
The following data are in bijective correspondence:
\begin{enumerate}[label=\textbf{(\Alph*)}]
\item\label{item:adj-hom}
    A natural isomorphism
    $\Phi\colon\Hom_\mathcal{D}(F-,-)\xRightarrow{\sim}\Hom_\mathcal{C}(-,G-)$.

\item\label{item:adj-unit-counit}
    Natural transformations $\eta\colon\Id_\mathcal{C}\Rightarrow GF$
    and $\varepsilon\colon FG\Rightarrow\Id_\mathcal{D}$ satisfying
    the triangle identities:
    \[
        \varepsilon F\circ F\eta = \id_F,
        \qquad
        G\varepsilon\circ\eta G = \id_G.
    \]

\item\label{item:adj-universal}
    For each $A\in\mathcal{C}$, an object $FA\in\mathcal{D}$ and
    a morphism $\eta_A\colon A\to GFA$ that is \textbf{universal}:
    for every $B\in\mathcal{D}$ and $g\colon A\to GB$, there exists
    a unique $f\colon FA\to B$ with $Gf\circ\eta_A = g$:
    \[
    \begin{tikzcd}
    A \arrow[r, "\eta_A"] \arrow[dr, "g"'] & GFA \arrow[d, "Gf"] \\
    & GB
    \end{tikzcd}
    \]
\end{enumerate}
\end{theorem}

\begin{proof}
\textbf{\ref{item:adj-hom} $\Rightarrow$ \ref{item:adj-unit-counit}:}
Given $\Phi$, define $\eta_A = \Phi_{A,FA}(\id_{FA})$ and
$\varepsilon_B = \Phi^{-1}_{GB,B}(\id_{GB})$.
By Proposition~\ref{prop:unit-counit-nat}, these are natural transformations.
The triangle identities are proved in Theorem~\ref{thm:triangle-identities}.

\medskip
\textbf{\ref{item:adj-unit-counit} $\Rightarrow$ \ref{item:adj-hom}:}
Given $(\eta,\varepsilon)$ satisfying the triangle identities, define
\[
    \Phi_{A,B}(f) = Gf\circ\eta_A, \qquad
    \Psi_{A,B}(g) = \varepsilon_B\circ Fg.
\]
We verify $\Psi\circ\Phi = \id$ and $\Phi\circ\Psi = \id$.

For $f\colon FA\to B$:
\begin{align*}
    \Psi(\Phi(f)) &= \Psi(Gf\circ\eta_A)
    = \varepsilon_B\circ F(Gf\circ\eta_A) \\
    &= \varepsilon_B\circ FGf\circ F\eta_A
    = f\circ\varepsilon_{FA}\circ F\eta_A \quad\text{(naturality of $\varepsilon$)}\\
    &= f\circ\id_{FA} = f \quad\text{(triangle identity (a))}.
\end{align*}

For $g\colon A\to GB$:
\begin{align*}
    \Phi(\Psi(g)) &= \Phi(\varepsilon_B\circ Fg)
    = G(\varepsilon_B\circ Fg)\circ\eta_A \\
    &= G\varepsilon_B\circ GFg\circ\eta_A
    = G\varepsilon_B\circ\eta_{GB}\circ g \quad\text{(naturality of $\eta$)}\\
    &= \id_{GB}\circ g = g \quad\text{(triangle identity (b))}.
\end{align*}

Naturality of $\Phi$ in both variables follows from the naturality of
$\eta$ and the functoriality of~$G$.

\medskip
\textbf{\ref{item:adj-hom} $\Leftrightarrow$ \ref{item:adj-universal}:}
Given~$\Phi$, setting $\eta_A = \Phi_{A,FA}(\id_{FA})$ provides the
universal morphism.  The universality statement is precisely the
bijectivity of~$\Phi_{A,B}$: given $g\colon A\to GB$, the unique
$f$ with $Gf\circ\eta_A = g$ is $f = \Phi^{-1}_{A,B}(g)$.

Conversely, given universal arrows $(\eta_A)$, define
$\Phi_{A,B}(f) = Gf\circ\eta_A$.  The universality ensures this is
a bijection.  Naturality follows from the naturality of the collection
$(\eta_A)_A$ (which itself follows from universality applied to the
composites $GFh\circ\eta_A$).
\end{proof}


%% ============================================================
\section{Fundamental examples}\label{sec:adj-examples}
%% ============================================================

\subsection{Free--forgetful adjunctions}

\begin{example}[Free group]\label{ex:free-group}
\index{adjunction!free-forgetful}
\index{free group}
\index{free functor}
\index{forgetful functor}
Let $U\colon\Grp\to\Set$ be the forgetful functor and
$F\colon\Set\to\Grp$ the free group functor.
Then $F\dashv U$:
\[
\begin{tikzcd}
\Set
  \arrow[rr, bend left=25, "F", ""{name=F, below}]
& \bot &
\Grp
  \arrow[ll, bend left=25, "U", ""{name=G, above}]
\end{tikzcd}
\]
The adjunction bijection is
\[
    \Hom_\Grp(F(S),\,G) \;\cong\; \Hom_\Set(S,\,U(G)):
\]
a group homomorphism from the free group on~$S$ is determined by
where the generators go, i.e.\ by a function $S\to U(G)$.

The unit $\eta_S\colon S\to UF(S)$ sends each element to the
corresponding generator in the free group.
The counit $\varepsilon_G\colon FU(G)\to G$ sends a word in the
``generators'' $U(G)$ to its evaluation in~$G$.
\end{example}

\begin{example}[Free abelian group]\label{ex:free-abelian}
\index{free abelian group}
The free abelian group functor $\Z[-]\colon\Set\to\Ab$ is left adjoint
to the forgetful functor $U\colon\Ab\to\Set$:
\[
    \Hom_\Ab(\Z[S],\,A) \;\cong\; \Hom_\Set(S,\,U(A)).
\]
Here $\Z[S] = \bigoplus_{s\in S} \Z$ is the free abelian group on~$S$.
\end{example}

\begin{example}[Free $R$-module]\label{ex:free-module}
\index{free module}
For a ring~$R$, the free module functor $R[-]\colon\Set\to\Mod_R$
is left adjoint to the forgetful functor.
\end{example}

\begin{example}[Free--forgetful between $\Ab$ and $\Grp$]%
\label{ex:abelianisation}
\index{abelianisation}
The inclusion functor $\iota\colon\Ab\hookrightarrow\Grp$ has a
left adjoint: the abelianisation $\mathrm{ab}\colon\Grp\to\Ab$
sending $G$ to $G/[G,G]$.
\[
    \Hom_\Ab(G/[G,G],\,A) \;\cong\; \Hom_\Grp(G,\,\iota A).
\]
\end{example}

\subsection{Tensor--Hom adjunction}

\begin{example}[Tensor--Hom]\label{ex:tensor-hom}
\index{adjunction!tensor-Hom}
\index{tensor product!as left adjoint}
\index{Hom!as right adjoint}
Let $R$ be a commutative ring and $M$ an $R$-module.
Then $(-\otimes_R M) \dashv \Hom_R(M,-)$:
\[
    \Hom_R(N\otimes_R M,\, P) \;\cong\; \Hom_R(N,\,\Hom_R(M,P))
\]
naturally in $N$ and~$P$.  This is the universal property of the
tensor product.
\end{example}

\subsection{Product--Internal Hom (Cartesian closed categories)}

\begin{example}[Cartesian closed]\label{ex:cartesian-closed}
\index{adjunction!product-internal Hom}
\index{Cartesian closed category}
\index{exponential object}
In a Cartesian closed category (e.g.\ $\Set$), for each object~$B$
the product functor $(-\times B)$ has a right adjoint, the
\textbf{internal Hom} (or \textbf{exponential}) $B^{(-)}$ or $[B,-]$:
\[
    \Hom(A\times B,\,C) \;\cong\; \Hom(A,\,C^B).
\]
In $\Set$, $C^B = \{f\colon B\to C\}$ is the set of functions.
This is the \textbf{currying} correspondence in computer science.
\index{currying}
\end{example}

\subsection{Direct and inverse image (sheaves)}

\begin{example}[Direct and inverse image]\label{ex:sheaf-adjunction}
\index{adjunction!direct and inverse image}
\index{sheaves!adjunction}
\index{inverse image functor}
\index{direct image functor}
Let $f\colon X\to Y$ be a continuous map of topological spaces.
There is an adjunction
\[
    f^* \dashv f_*
\]
between the inverse image functor $f^*\colon\mathsf{Sh}(Y)\to\mathsf{Sh}(X)$
and the direct image functor $f_*\colon\mathsf{Sh}(X)\to\mathsf{Sh}(Y)$
on categories of sheaves.  This adjunction is fundamental in algebraic
geometry and sheaf theory.
\end{example}

\subsection{Suspension--loops adjunction}

\begin{example}[Suspension and loops]\label{ex:suspension-loops}
\index{adjunction!suspension-loops}
\index{suspension functor}
\index{loop space functor}
In the category $\Top_*$ of pointed topological spaces, the
(reduced) suspension functor $\Sigma$ is left adjoint to the
loops functor~$\Omega$:
\[
    \Sigma \dashv \Omega, \qquad
    [\Sigma X,\, Y]_* \;\cong\; [X,\,\Omega Y]_*
\]
where $[-,-]_*$ denotes pointed homotopy classes.
This adjunction underlies the theory of stable homotopy.
\end{example}

\begin{example}[Galois connections]\label{ex:galois-connection}
\index{Galois connection}
\index{adjunction!as Galois connection}
If $\mathcal{C}$ and $\mathcal{D}$ are posets (viewed as categories),
then an adjunction $F\dashv G$ is the same as a \textbf{Galois connection}:
$F(a)\leq b$ if and only if $a\leq G(b)$.
\end{example}


%% ============================================================
\section{Left adjoints preserve colimits}\label{sec:ladj-colim}
%% ============================================================

\begin{theorem}[Left adjoints preserve colimits]%
\label{thm:ladj-preserve-colim}
\index{left adjoint!preserves colimits}
\index{colimit!preserved by left adjoint}
\index{adjunction!and colimits}
If $F\dashv G$ with $F\colon\mathcal{C}\to\mathcal{D}$,
then $F$ preserves all colimits that exist in~$\mathcal{C}$.
\end{theorem}

\begin{proof}
Let $D\colon\mathcal{J}\to\mathcal{C}$ be a diagram with colimit
$(\varinjlim D,\,\iota_j)$.  We must show that
$(F(\varinjlim D),\,F\iota_j)$ is a colimit of $F\circ D$.

Let $(N,\psi_j)$ be a cocone under $F\circ D$ in~$\mathcal{D}$,
so $\psi_j\colon FD_j\to N$ with compatibility.
For each~$j$, applying the adjunction bijection gives
\[
    \tilde\psi_j = \Phi_{D_j,N}(\psi_j) \colon D_j \to GN.
\]

\textbf{Claim:} $(GN,\tilde\psi_j)$ is a cocone under~$D$ in~$\mathcal{C}$.
Indeed, for $\alpha\colon j\to k$ in~$\mathcal{J}$:
\begin{align*}
\tilde\psi_k\circ D_\alpha
&= \Phi_{D_k,N}(\psi_k)\circ D_\alpha \\
&= \Phi_{D_j,N}(\psi_k\circ FD_\alpha)
    \quad\text{(naturality of $\Phi$ in $A$)} \\
&= \Phi_{D_j,N}(\psi_j) = \tilde\psi_j
    \quad\text{(cocone condition for $\psi$)}.
\end{align*}

By the universal property of $\varinjlim D$, there is a unique
$\tilde u\colon\varinjlim D\to GN$ with
$\tilde u\circ\iota_j = \tilde\psi_j$ for all~$j$.

Now set $u = \Phi^{-1}_{\varinjlim D,\,N}(\tilde u)\colon F(\varinjlim D)\to N$.
Then for each~$j$, using naturality of~$\Phi^{-1}$ in~$A$:
\[
    u\circ F\iota_j
    = \Phi^{-1}_{D_j,N}(\tilde u\circ\iota_j)
    = \Phi^{-1}_{D_j,N}(\tilde\psi_j)
    = \psi_j.
\]
This gives existence.

\textbf{Uniqueness:}
If $u'\colon F(\varinjlim D)\to N$ also satisfies
$u'\circ F\iota_j = \psi_j$, then
$\Phi_{\varinjlim D,N}(u')\circ\iota_j = \tilde\psi_j$ for all~$j$,
so by uniqueness for $\varinjlim D$,
$\Phi(u') = \tilde u = \Phi(u)$, hence $u' = u$.
\[
\begin{tikzcd}
FD_j \arrow[r, "F\iota_j"] \arrow[dr, "\psi_j"']
& F(\varinjlim D) \arrow[d, dashed, "\exists!\, u"] \\
& N
\end{tikzcd}
\]
\end{proof}


%% ============================================================
\section{Right adjoints preserve limits}\label{sec:radj-lim}
%% ============================================================

\begin{theorem}[Right adjoints preserve limits]%
\label{thm:radj-preserve-lim}
\index{right adjoint!preserves limits}
\index{limit!preserved by right adjoint}
\index{adjunction!and limits}
If $F\dashv G$ with $G\colon\mathcal{D}\to\mathcal{C}$,
then $G$ preserves all limits that exist in~$\mathcal{D}$.
\end{theorem}

\begin{proof}
Let $D\colon\mathcal{J}\to\mathcal{D}$ be a diagram with limit
$(L,\pi_j)$, where $\pi_j\colon L\to D_j$.
We show $(GL,G\pi_j)$ is a limit of $G\circ D$.

Let $(N,\psi_j)$ be a cone over $G\circ D$ in~$\mathcal{C}$,
so $\psi_j\colon N\to GD_j$.
By the adjunction, for each~$j$ let
$\tilde\psi_j = \Phi^{-1}_{N,D_j}(\psi_j)\colon FN\to D_j$.

\textbf{Claim:} $(FN,\tilde\psi_j)$ is a cone over~$D$.
For $\alpha\colon j\to k$:
\begin{align*}
D_\alpha\circ\tilde\psi_j
&= D_\alpha\circ\Phi^{-1}_{N,D_j}(\psi_j) \\
&= \Phi^{-1}_{N,D_k}(GD_\alpha\circ\psi_j)
  \quad\text{(naturality of $\Phi^{-1}$ in $B$)} \\
&= \Phi^{-1}_{N,D_k}(\psi_k) = \tilde\psi_k
  \quad\text{(cone condition for $\psi$)}.
\end{align*}

By the universal property of~$L$, there is a unique
$\tilde v\colon FN\to L$ with $\pi_j\circ\tilde v = \tilde\psi_j$.
Set $v = \Phi_{N,L}(\tilde v)\colon N\to GL$.  Then
\[
    G\pi_j\circ v = G\pi_j\circ\Phi_{N,L}(\tilde v)
    = \Phi_{N,D_j}(\pi_j\circ\tilde v)
    = \Phi_{N,D_j}(\tilde\psi_j) = \psi_j,
\]
using naturality of~$\Phi$ in~$B$.

\textbf{Uniqueness} follows as before from the bijectivity of~$\Phi$
and uniqueness for the limit~$L$.
\[
\begin{tikzcd}
N \arrow[dr, "\psi_j"'] \arrow[rr, dashed, "\exists!\, v"]
& & GL \arrow[dl, "G\pi_j"] \\
& GD_j &
\end{tikzcd}
\]
\end{proof}

\begin{corollary}\label{cor:limits-detect-adjoints}
\index{adjunction!and preservation of (co)limits}
If a functor $G\colon\mathcal{D}\to\mathcal{C}$ has a left adjoint,
then $G$ preserves all limits.
If $F\colon\mathcal{C}\to\mathcal{D}$ has a right adjoint,
then $F$ preserves all colimits.
\end{corollary}

\begin{remark}\label{rem:converse-not-true}
The converse does not hold in general: preserving limits is necessary
but not sufficient for having a left adjoint.  The adjoint functor
theorems (below) provide sufficient conditions.
\end{remark}

\begin{example}\label{ex:forgetful-preserves-limits}
\index{forgetful functor!preserves limits}
Since the forgetful functor $U\colon\Grp\to\Set$ has a left adjoint
(the free group functor), it preserves all limits.  This explains why
limits of groups are computed ``on underlying sets'':
the underlying set of a product of groups is the product of their
underlying sets, etc.
\end{example}


%% ============================================================
\section{Uniqueness of adjoints}\label{sec:adj-unique}
%% ============================================================

\begin{theorem}[Uniqueness of adjoints]\label{thm:adj-unique}
\index{adjunction!uniqueness}
\index{right adjoint!uniqueness}
\index{left adjoint!uniqueness}
If $F\dashv G$ and $F\dashv G'$, then $G\cong G'$ via a natural
isomorphism.  Similarly, if $F\dashv G$ and $F'\dashv G$, then
$F\cong F'$.
\end{theorem}

\begin{proof}
Suppose $F\dashv G$ and $F\dashv G'$.  For every $B\in\mathcal{D}$,
\[
    \Hom_\mathcal{C}(-,GB)
    \cong \Hom_\mathcal{D}(F-,B)
    \cong \Hom_\mathcal{C}(-,G'B)
\]
naturally.  By the Yoneda lemma, a natural isomorphism
$\Hom(-,GB)\cong\Hom(-,G'B)$ implies $GB\cong G'B$,
and the naturality in~$B$ gives a natural isomorphism $G\cong G'$.

The uniqueness of left adjoints follows by a dual argument
(using the covariant Yoneda embedding).
\end{proof}


%% ============================================================
\section{Freyd's Adjoint Functor Theorems}\label{sec:GAFT}
%% ============================================================

The adjoint functor theorems give powerful sufficient conditions
for a functor to have an adjoint.  They are among the deepest
results of basic category theory.

\begin{definition}[Solution set condition]\label{def:solution-set}
\index{solution set condition}
A functor $G\colon\mathcal{D}\to\mathcal{C}$ satisfies the
\textbf{solution set condition} if for each object $A\in\mathcal{C}$,
there exists a \emph{set} (not a proper class) of morphisms
$\{f_i\colon A\to GB_i\}_{i\in I}$ such that every morphism
$g\colon A\to GB$ factors as $g = Gh\circ f_i$ for some $i\in I$
and some $h\colon B_i\to B$:
\[
\begin{tikzcd}
A \arrow[r, "f_i"] \arrow[dr, "g"'] & GB_i \arrow[d, "Gh"] \\
& GB
\end{tikzcd}
\]
\end{definition}

\begin{theorem}[General Adjoint Functor Theorem (Freyd)]%
\label{thm:GAFT}
\index{General Adjoint Functor Theorem}
\index{GAFT}
\index{Freyd's Adjoint Functor Theorem}
\index{adjoint functor theorem!general}
Let $\mathcal{D}$ be a complete, locally small category and let
$G\colon\mathcal{D}\to\mathcal{C}$ be a functor.
Then $G$ has a left adjoint if and only if:
\begin{enumerate}[label=(\roman*)]
\item $G$ preserves all (small) limits; and
\item $G$ satisfies the solution set condition.
\end{enumerate}
\end{theorem}

\begin{proof}
\textbf{Necessity:}
If $G$ has a left adjoint~$F$, then $G$ preserves limits by
Theorem~\ref{thm:radj-preserve-lim}.
The solution set condition is satisfied with the singleton
$\{\eta_A\colon A\to GFA\}$ for each~$A$ (where $\eta$ is the unit):
any $g\colon A\to GB$ equals $Gf\circ\eta_A$ where
$f = \Phi^{-1}(g)\colon FA\to B$.

\medskip
\textbf{Sufficiency:}
We construct, for each $A\in\mathcal{C}$, a universal arrow
$\eta_A\colon A\to GFA$ (Definition~\ref{def:adjunction},
form~\ref{item:adj-universal}).

\medskip\noindent
\emph{Step 1: Initial construction.}
Let $\{f_i\colon A\to GB_i\}_{i\in I}$ be a solution set for~$A$.
Form the product $P = \prod_{i\in I} B_i$ in~$\mathcal{D}$
(which exists since $\mathcal{D}$ is complete).
Since $G$ preserves limits, $GP \cong \prod_i GB_i$,
so the $f_i$ assemble into a single morphism
$f\colon A\to GP$.

\medskip\noindent
\emph{Step 2: Equalising.}
Consider the set of all endomorphisms $e$ of~$P$ in~$\mathcal{D}$
such that $Ge\circ f = f$.  Let $E$ be the joint equaliser of all
such endomorphisms (an intersection of equalisers, which exists by
completeness).  Let $m\colon E\to P$ be the canonical morphism
and let $\eta_A = Gm^{-1}\circ f'\colon A\to GE$
be the factored morphism.  More precisely, $f$ factors through $GE$
via a unique $\eta_A\colon A\to GE$ with $Gm\circ\eta_A = f$.

\medskip\noindent
\emph{Step 3: Verification of universality.}
Given $g\colon A\to GB$, the solution set condition gives
$g = Gh\circ f_i$ for some~$i$ and some $h\colon B_i\to B$.
Define $h'\colon P\to B$ extending~$h$ via the product projection.
Then $Gh'\circ f = Gh'\circ Gm\circ\eta_A$,
and one verifies that $h'$ factors through~$E$
(by the equaliser condition and the fact that $G$ preserves the
relevant limits), yielding the desired factorisation
$g = G\bar h\circ\eta_A$ for a unique $\bar h\colon E\to B$.

\medskip\noindent
\emph{Uniqueness} of $\bar h$ requires a further equaliser argument:
if $\bar h,\bar h'\colon E\to B$ both satisfy $G\bar h\circ\eta_A
= G\bar h'\circ\eta_A$, one shows that the equaliser of $\bar h$
and~$\bar h'$ is all of~$E$, using the fact that $G$ preserves
this equaliser.

Setting $FA = E$ for each $A$ and using the universal arrows to define
$F$ on morphisms, one obtains the left adjoint
$F\colon\mathcal{C}\to\mathcal{D}$.
\end{proof}

\begin{theorem}[Special Adjoint Functor Theorem]%
\label{thm:SAFT}
\index{Special Adjoint Functor Theorem}
\index{SAFT}
\index{adjoint functor theorem!special}
\index{cogenerating set}
\index{well-powered category}
Let $\mathcal{D}$ be a complete, locally small, well-powered
category with a cogenerating set.
Then a functor $G\colon\mathcal{D}\to\mathcal{C}$ has a left
adjoint if and only if $G$ preserves all (small) limits.
\end{theorem}

In other words, the solution set condition is automatic when $\mathcal{D}$
has a cogenerating set and is well-powered.

\begin{remark}\label{rem:SAFT-applications}
\index{Special Adjoint Functor Theorem!applications}
The SAFT applies notably to:
\begin{itemize}
\item $\mathcal{D} = \Set$ (cogenerating set: $\{1, 2\}$ where
      $2 = \{0,1\}$);
\item $\mathcal{D} = \Ab$ (cogenerating set: $\{\Q/\Z\}$);
\item $\mathcal{D} = \Mod_R$ for $R$ Noetherian
      (any injective cogenerator suffices).
\end{itemize}
For instance, the SAFT immediately implies that any limit-preserving
functor $G\colon\Ab\to\mathcal{C}$ from the abelian category~$\Ab$
has a left adjoint, since $\Ab$ is complete, locally small,
well-powered, and has $\Q/\Z$ as a cogenerator.
\end{remark}

\begin{example}[Application of GAFT]\label{ex:GAFT-application}
\index{General Adjoint Functor Theorem!application}
\index{Stone--\v{C}ech compactification}
Consider the inclusion functor $G\colon\mathsf{CompHaus}\hookrightarrow\Top$
from compact Hausdorff spaces to all topological spaces.
Since $G$ preserves all limits (closed subspaces of compact Hausdorff
spaces are compact Hausdorff) and the solution set condition is verified
by cardinality arguments, the GAFT gives a left adjoint
$\beta\colon\Top\to\mathsf{CompHaus}$, the
\textbf{Stone--\v{C}ech compactification}.
\end{example}


%% ============================================================
\section{Composition of adjunctions}\label{sec:adj-composition}
%% ============================================================

\begin{proposition}[Composition of adjunctions]\label{prop:adj-compose}
\index{adjunction!composition}
If $F_1\dashv G_1$ between $\mathcal{C}$ and $\mathcal{D}$,
and $F_2\dashv G_2$ between $\mathcal{D}$ and $\mathcal{E}$, then
$F_2 F_1 \dashv G_1 G_2$:
\[
\begin{tikzcd}
\mathcal{C}
  \arrow[r, bend left, "F_1"]
& \mathcal{D}
  \arrow[l, bend left, "G_1"]
  \arrow[r, bend left, "F_2"]
& \mathcal{E}
  \arrow[l, bend left, "G_2"]
\end{tikzcd}
\qquad\Longrightarrow\qquad
\begin{tikzcd}
\mathcal{C}
  \arrow[rr, bend left, "F_2 F_1"]
& &
\mathcal{E}
  \arrow[ll, bend left, "G_1 G_2"]
\end{tikzcd}
\]
\end{proposition}

\begin{proof}
For $A\in\mathcal{C}$ and $E\in\mathcal{E}$:
\[
    \Hom_\mathcal{E}(F_2 F_1 A,\, E)
    \cong \Hom_\mathcal{D}(F_1 A,\, G_2 E)
    \cong \Hom_\mathcal{C}(A,\, G_1 G_2 E),
\]
naturally in $A$ and~$E$.  The composite of natural isomorphisms is
a natural isomorphism.
\end{proof}


%% ============================================================
\section{Adjunctions and equivalences}\label{sec:adj-equiv}
%% ============================================================

\begin{proposition}\label{prop:equiv-is-adjunction}
\index{equivalence of categories!and adjunctions}
\index{adjunction!and equivalence}
Every equivalence of categories gives rise to an adjunction.
Specifically, if $F\colon\mathcal{C}\to\mathcal{D}$ is part of an
equivalence $(F,G,\eta,\varepsilon)$ with $\eta\colon\Id\xRightarrow{\sim}GF$
and $\varepsilon\colon FG\xRightarrow{\sim}\Id$ natural isomorphisms,
then by modifying $\varepsilon$ if necessary, one obtains $F\dashv G$.
\end{proposition}

\begin{proof}
Given the equivalence, set $\varepsilon' = \varepsilon \circ F(\eta^{-1})_{G}
\circ (F\eta G)^{-1}$---more precisely, define
$\varepsilon'_B = \varepsilon_B \circ F(G\varepsilon_B)^{-1} \circ F\eta_{GB}$?
The standard trick is: set
\[
    \varepsilon'_B = \varepsilon_B \circ F(\eta^{-1}_{GB}) \circ F\eta_{GB}
\]
which simplifies, but let us use a cleaner approach.
Define the new counit as
$\varepsilon'_B = \varepsilon_B$ and check whether the triangle
identity holds up to adjusting signs.

In fact, the cleaner argument is: define $\Phi_{A,B}(f) = Gf\circ\eta_A$.
Since $\eta_A$ is an isomorphism, $\Phi_{A,B}$ is a bijection
(given $g\colon A\to GB$, set $f = \varepsilon_B\circ F(\eta_A^{-1}\circ
\cdots)$---the details require care).

Alternatively, observe that $(F,G,\eta,\varepsilon')$ with
$\varepsilon'_B = \varepsilon_B\circ F(G\varepsilon_B^{-1}\circ\eta_{GB})^{-1}
\circ F\eta_{GB}$ satisfies the triangle identities
after suitable simplification using the fact that $\eta$ and $\varepsilon$
are natural isomorphisms.  We omit the (routine but notationally heavy)
verification.
\end{proof}

\begin{remark}\label{rem:adjoint-equivalence}
\index{adjoint equivalence}
An \textbf{adjoint equivalence} is an adjunction $F\dashv G$ in which
both $\eta$ and $\varepsilon$ are natural isomorphisms.
Every equivalence of categories can be promoted to an adjoint equivalence.
\end{remark}


%% ============================================================
\section{Adjunctions via initial objects}\label{sec:adj-comma}
%% ============================================================

\begin{remark}\label{rem:adj-comma-cat}
\index{comma category}
\index{adjunction!via comma category}
Yet another characterisation: $G\colon\mathcal{D}\to\mathcal{C}$ has a
left adjoint if and only if for every $A\in\mathcal{C}$, the
\textbf{comma category} $(A\downarrow G)$ has an initial object.
The initial object is the pair $(FA,\eta_A\colon A\to GFA)$.
Recall that objects of $(A\downarrow G)$ are pairs $(B,\,f\colon A\to GB)$,
and a morphism $(B,f)\to(B',f')$ is an $h\colon B\to B'$ with
$Gh\circ f = f'$.
\end{remark}


%% ============================================================
\section{Exercises}\label{sec:exercises-ch4}
%% ============================================================

\begin{exercise}\label{exer:adj-check-triangle}
\index{triangle identities!verification}
Verify the triangle identities directly for the free--forgetful adjunction
$F\dashv U$ between $\Set$ and $\Grp$ (Example~\ref{ex:free-group}).
\end{exercise}

\begin{exercise}\label{exer:adj-poset}
\index{Galois connection!examples}
Let $f\colon X\to Y$ be a function between sets.  Define
$f_*\colon\mathcal{P}(X)\to\mathcal{P}(Y)$ by $f_*(S)=f(S)$
(direct image) and $f^*\colon\mathcal{P}(Y)\to\mathcal{P}(X)$
by $f^*(T)=f^{-1}(T)$ (inverse image), where $\mathcal{P}$ denotes
power set ordered by inclusion.
Show that $f_*\dashv f^*$ as a Galois connection.
Also show that $f^*$ has a right adjoint $f_!\colon\mathcal{P}(X)\to\mathcal{P}(Y)$
given by $f_!(S) = \{y\in Y : f^{-1}(\{y\})\subseteq S\}$.
\end{exercise}

\begin{exercise}\label{exer:adj-tensor-hom-triangle}
\index{adjunction!tensor-Hom!unit and counit}
For the tensor--Hom adjunction $(-\otimes_R M)\dashv\Hom_R(M,-)$
(Example~\ref{ex:tensor-hom}):
\begin{enumerate}[label=(\alph*)]
\item Describe the unit $\eta_N\colon N\to\Hom_R(M,\,N\otimes_R M)$
    explicitly.
\item Describe the counit
    $\varepsilon_P\colon\Hom_R(M,P)\otimes_R M\to P$ explicitly.
\item Verify the triangle identities.
\end{enumerate}
\end{exercise}

\begin{exercise}\label{exer:ladj-preserves-colim-examples}
\index{left adjoint!preserves colimits!examples}
Use the fact that left adjoints preserve colimits to prove:
\begin{enumerate}[label=(\alph*)]
\item The free group functor $F\colon\Set\to\Grp$ preserves coproducts.
    Conclude that $F(S_1\amalg S_2)\cong F(S_1)*F(S_2)$ (free product).
\item The tensor product $-\otimes_R M$ preserves direct sums:
    $(\bigoplus_i N_i)\otimes_R M \cong \bigoplus_i(N_i\otimes_R M)$.
\end{enumerate}
\end{exercise}

\begin{exercise}\label{exer:radj-preserves-lim-examples}
\index{right adjoint!preserves limits!examples}
Use the fact that right adjoints preserve limits to prove:
\begin{enumerate}[label=(\alph*)]
\item The forgetful functor $U\colon\Grp\to\Set$ preserves products.
\item $\Hom_R(M,\prod_i P_i)\cong\prod_i\Hom_R(M,P_i)$.
\item The loops functor $\Omega\colon\Top_*\to\Top_*$ preserves products.
\end{enumerate}
\end{exercise}

\begin{exercise}\label{exer:adj-iff-representable}
\index{adjunction!and representability}
Show that $G\colon\mathcal{D}\to\mathcal{C}$ has a left adjoint
if and only if for each $A\in\mathcal{C}$, the functor
$\Hom_\mathcal{C}(A,G-)\colon\mathcal{D}\to\Set$ is representable.
\end{exercise}

\begin{exercise}\label{exer:reflective-subcategory}
\index{reflective subcategory}
\index{adjunction!reflective subcategory}
A full subcategory $\mathcal{D}\hookrightarrow\mathcal{C}$ is
\textbf{reflective} if the inclusion functor $\iota$ has a left adjoint
$L$ (the \textbf{reflector}).  Show that:
\begin{enumerate}[label=(\alph*)]
\item The counit $\varepsilon\colon L\iota\Rightarrow\Id_\mathcal{D}$
    is a natural isomorphism.
\item $\Ab\hookrightarrow\Grp$ is reflective (the reflector is abelianisation).
\item The full subcategory of complete metric spaces in $\mathsf{Met}$
    is reflective (the reflector is Cauchy completion).
\end{enumerate}
\end{exercise}

\begin{exercise}\label{exer:fixed-points-adjunction}
\index{adjunction!and fixed points}
Let $F\dashv G$ with unit~$\eta$ and counit~$\varepsilon$.
An object $A\in\mathcal{C}$ is called \textbf{$G$-split} if
$\eta_A$ is a split monomorphism.
\begin{enumerate}[label=(\alph*)]
\item Show that $\eta_A$ is a split mono for every~$A$ if and only if
    $G$ is faithful.
\item Show that $\eta$ is a natural isomorphism if and only if
    $G$ is fully faithful (i.e.\ $\mathcal{D}$ is a reflective
    subcategory of~$\mathcal{C}$ up to equivalence).
\end{enumerate}
\end{exercise}

\begin{exercise}\label{exer:GAFT-apply}
\index{General Adjoint Functor Theorem!exercises}
\begin{enumerate}[label=(\alph*)]
\item Use the GAFT to prove that the forgetful functor
    $U\colon\mathsf{CompHaus}\to\Top$ has a left adjoint
    (the Stone--\v{C}ech compactification~$\beta$).
    Verify the solution set condition.
\item Use the SAFT to show that any continuous (limit-preserving)
    functor $G\colon\Ab\to\Set$ has a left adjoint.
\end{enumerate}
\end{exercise}

\begin{exercise}\label{exer:adj-monad}
\index{monad!from adjunction}
\index{adjunction!associated monad}
Let $F\dashv G$ with unit $\eta$ and counit~$\varepsilon$.
Define $T = GF\colon\mathcal{C}\to\mathcal{C}$.
Show that $(T,\eta,\mu)$ is a \textbf{monad} on~$\mathcal{C}$,
where $\mu = G\varepsilon F\colon T^2 = GFGF\Rightarrow GF = T$.
That is, verify:
\begin{enumerate}[label=(\alph*)]
\item $\mu\circ T\mu = \mu\circ\mu T$ (associativity);
\item $\mu\circ T\eta = \mu\circ\eta T = \id_T$ (unit laws).
\end{enumerate}
\emph{Hint:} Use the triangle identities and naturality.
\[
\begin{tikzcd}
T^3 \arrow[r, "T\mu"] \arrow[d, "\mu T"'] & T^2 \arrow[d, "\mu"] \\
T^2 \arrow[r, "\mu"'] & T
\end{tikzcd}
\qquad
\begin{tikzcd}
T \arrow[r, "T\eta"] \arrow[dr, equal] & T^2 \arrow[d, "\mu"]
  & T \arrow[l, "\eta T"'] \arrow[dl, equal] \\
& T &
\end{tikzcd}
\]
(We will study monads in detail in a later chapter.)
\end{exercise}

\index{adjunction|)}
%%%%%%%%%%%%%%%%%%%%%%%%%%%%%%%%%%%%%%%%%%%%%%%%%%%%%%%%%%%%%
%% CATEGORY THEORY — Chapters 5–6
%% The Yoneda Lemma and Representability / Monads and Algebras
%%%%%%%%%%%%%%%%%%%%%%%%%%%%%%%%%%%%%%%%%%%%%%%%%%%%%%%%%%%%%

%%%%%%%%%%%%%%%%%%%%%%%%%%%%%%%%%%%%%%%%%%%%%%%%%%%%%%%%%%%%
%%    CHAPTER 5: THE YONEDA LEMMA AND REPRESENTABILITY    %%
%%%%%%%%%%%%%%%%%%%%%%%%%%%%%%%%%%%%%%%%%%%%%%%%%%%%%%%%%%%%

\chapter{The Yoneda Lemma and Representability}\label{ch:yoneda}
\minitoc

\vspace{1em}

The Yoneda lemma is, without exaggeration, the single most important
result in category theory.  It tells us that every object of a category
is completely determined by the totality of morphisms into (or out of) it.
More precisely, it establishes a bijection between natural
transformations out of a representable functor and the elements of the
target functor, and it shows that the Yoneda embedding is fully faithful.
This chapter develops the theory of presheaves, states and proves the
Yoneda lemma in its full generality, and explores representable functors
and the fundamental fact that every presheaf is a colimit of
representables.

\index{Yoneda lemma|(}
\index{representable functor|(}
\index{presheaf|(}

%% ============================================================
\section{Presheaves and the functor category}\label{sec:presheaves}
%% ============================================================

\begin{definition}[Presheaf]\label{def:presheaf}
\index{presheaf!definition}
\index{functor category!of presheaves}
Let $\mathcal{C}$ be a locally small category.
A \textbf{presheaf} on~$\mathcal{C}$ is a functor
\[
    F \colon \mathcal{C}^{\op} \longrightarrow \Set.
\]
The \textbf{category of presheaves} on~$\mathcal{C}$ is the functor category
\[
    \widehat{\mathcal{C}} \;=\; \Fun(\mathcal{C}^{\op},\,\Set)
    \;=\; \Set^{\mathcal{C}^{\op}}.
\]
\index{$\widehat{\mathcal{C}}$}
Morphisms in~$\widehat{\mathcal{C}}$ are natural transformations.
\end{definition}

\begin{remark}\label{rem:presheaf-covariant}
\index{copresheaf}
Dually, a \textbf{copresheaf} (or covariant presheaf) on~$\mathcal{C}$
is a functor $\mathcal{C}\to\Set$.  The category of copresheaves is
$\Fun(\mathcal{C},\Set)$.  Both viewpoints are important; the choice of
variance depends on the context.
\end{remark}

\begin{example}\label{ex:presheaf-topology}
\index{presheaf!on a topological space}
\index{open sets!poset of}
Let $X$ be a topological space and let $\mathcal{O}(X)$ be the poset of
open subsets ordered by inclusion, viewed as a category.
A presheaf on~$\mathcal{O}(X)$ in the sense above is a functor
$\mathcal{O}(X)^{\op}\to\Set$: it assigns to each open set~$U$ a
set~$F(U)$ and to each inclusion $V\subseteq U$ a \emph{restriction map}
$F(U)\to F(V)$, contravariantly.  This recovers the classical notion of
presheaf on a topological space.
\end{example}

\begin{example}\label{ex:presheaf-simplicial}
\index{simplicial set}
\index{$\mathbf{\Delta}$}
A \textbf{simplicial set} is a presheaf on the simplex category~$\mathbf{\Delta}$
(whose objects are the finite totally ordered sets $[n]=\{0,1,\ldots,n\}$
and whose morphisms are order-preserving maps).  Thus
$\mathsf{sSet} = \Fun(\mathbf{\Delta}^{\op},\Set)$.
\end{example}


%% ============================================================
\section{Representable functors}\label{sec:representable}
%% ============================================================

\begin{definition}[Covariant representable functor]\label{def:h-lower}
\index{representable functor!covariant}
\index{$h_A$}
\index{Hom functor!covariant}
For each object $A\in\mathcal{C}$, the \textbf{covariant representable
functor} (or \textbf{covariant Hom functor}) is
\[
    h_A \;=\; \Hom_{\mathcal{C}}(A,-) \;\colon\;
    \mathcal{C} \longrightarrow \Set.
\]
On objects: $h_A(B) = \Hom(A,B)$.
On morphisms: for $f\colon B\to B'$, the map
$h_A(f) = f_* = f\circ(-)\colon \Hom(A,B)\to\Hom(A,B')$
is post-composition with~$f$.
\end{definition}

\begin{definition}[Contravariant representable functor]\label{def:h-upper}
\index{representable functor!contravariant}
\index{$h^A$}
\index{Hom functor!contravariant}
For each object $A\in\mathcal{C}$, the \textbf{contravariant representable
functor} is
\[
    h^A \;=\; \Hom_{\mathcal{C}}(-,A) \;\colon\;
    \mathcal{C}^{\op} \longrightarrow \Set.
\]
On objects: $h^A(B) = \Hom(B,A)$.
On morphisms: for $f\colon B'\to B$ in~$\mathcal{C}$ (equivalently
$f\colon B\to B'$ in~$\mathcal{C}^{\op}$), the map
$h^A(f) = f^* = (-)\circ f\colon \Hom(B,A)\to\Hom(B',A)$
is pre-composition with~$f$.
\end{definition}

\begin{remark}\label{rem:representable-conventions}
\index{representable functor!convention}
Convention varies across the literature.  We follow the notation where
$h^A$ denotes the contravariant functor (a presheaf on~$\mathcal{C}$)
and $h_A$ denotes the covariant functor (a copresheaf).
Some authors write $\mathbf{y}(A)$ for $h^A$, where $\mathbf{y}$ is the
Yoneda embedding.
\end{remark}

\begin{definition}[Representable functor]\label{def:representable}
\index{representable functor!definition}
\index{representing object}
A functor $F\colon\mathcal{C}^{\op}\to\Set$ is \textbf{representable}
if there exists an object $A\in\mathcal{C}$ and a natural isomorphism
$F\cong h^A$.  The object~$A$ is called a \textbf{representing object}.
Dually, a functor $G\colon\mathcal{C}\to\Set$ is \textbf{representable}
if $G\cong h_A$ for some~$A$.
\end{definition}

\begin{notation}\label{not:hom-bifunctor}
\index{Hom functor!bifunctor}
\index{bifunctor!Hom}
The assignments $(A,B)\mapsto\Hom(A,B)$ assemble into a bifunctor
\[
    \Hom_{\mathcal{C}}(-,-) \;\colon\;
    \mathcal{C}^{\op}\times\mathcal{C} \longrightarrow \Set.
\]
Fixing the first variable yields $h_A = \Hom(A,-)$; fixing the second
yields $h^A = \Hom(-,A)$.
\end{notation}

\begin{example}\label{ex:representable-forgetful}
\index{forgetful functor!as representable}
\begin{enumerate}[label=(\roman*)]
\item The forgetful functor $U\colon\Grp\to\Set$ is representable:
    $U\cong h_{\mathbb{Z}} = \Hom_{\Grp}(\mathbb{Z},-)$, since a
    group homomorphism $\mathbb{Z}\to G$ is determined by the image
    of~$1$, hence by an element of the underlying set of~$G$.

\item The forgetful functor $U\colon\mathsf{Ring}\to\Set$ is
    representable: $U\cong\Hom_{\mathsf{Ring}}(\mathbb{Z}[x],-)$.

\item The forgetful functor $U\colon\Top\to\Set$ is representable:
    $U\cong\Hom_{\Top}(\{*\},-)$, where $\{*\}$ is the one-point space.

\item The functor $\mathsf{GL}_n\colon\mathsf{CRing}\to\Grp$ sending
    a commutative ring~$R$ to $\mathsf{GL}_n(R)$ is representable, with
    representing object $\mathbb{Z}[x_{ij},\det(x_{ij})^{-1}]$.
    \index{general linear group!as representable functor}
\end{enumerate}
\end{example}


%% ============================================================
\section{The Yoneda embedding}\label{sec:yoneda-embedding}
%% ============================================================

\begin{definition}[Yoneda embedding]\label{def:yoneda-embedding}
\index{Yoneda embedding}
\index{$\mathbf{y}$}
The \textbf{Yoneda embedding} is the functor
\[
    \mathbf{y} \;\colon\; \mathcal{C} \longrightarrow
    \Fun(\mathcal{C}^{\op},\Set) = \widehat{\mathcal{C}}
\]
defined on objects by $\mathbf{y}(A) = h^A = \Hom(-,A)$
and on morphisms by: for $f\colon A\to B$ in~$\mathcal{C}$,
$\mathbf{y}(f) = f_*\colon h^A\Rightarrow h^B$ is the natural transformation
whose component at $C\in\mathcal{C}$ is
\[
    (f_*)_C \;=\; f\circ(-) \;\colon\;
    \Hom(C,A) \longrightarrow \Hom(C,B).
\]
\end{definition}

\begin{lemma}\label{lem:yoneda-well-defined}
\index{Yoneda embedding!well-defined}
The map $\mathbf{y}(f) = f_*$ is indeed a natural transformation
$h^A\Rightarrow h^B$ for each morphism $f\colon A\to B$.
\end{lemma}

\begin{proof}
We must check naturality: for every $g\colon C'\to C$ in~$\mathcal{C}$,
the diagram
\[
\begin{tikzcd}
\Hom(C,A) \arrow[r, "f_*"] \arrow[d, "g^*"'] &
\Hom(C,B) \arrow[d, "g^*"] \\
\Hom(C',A) \arrow[r, "f_*"'] &
\Hom(C',B)
\end{tikzcd}
\]
commutes.  For $\varphi\in\Hom(C,A)$:
$g^*(f_*(\varphi)) = (f\circ\varphi)\circ g = f\circ(\varphi\circ g)
= f_*(g^*(\varphi))$.
\end{proof}

There is a dual version:

\begin{definition}[Co-Yoneda embedding]\label{def:co-yoneda-embedding}
\index{Yoneda embedding!covariant (co-Yoneda)}
The \textbf{co-Yoneda embedding} is the functor
\[
    \mathbf{y}_{\bullet} \;\colon\; \mathcal{C}^{\op}
    \longrightarrow \Fun(\mathcal{C},\Set)
\]
defined by $\mathbf{y}_{\bullet}(A) = h_A = \Hom(A,-)$.
\end{definition}


%% ============================================================
\section{The Yoneda lemma}\label{sec:yoneda-lemma}
%% ============================================================

We now state and prove the central result of this chapter.

\begin{theorem}[Yoneda Lemma]\label{thm:yoneda}
\index{Yoneda lemma!statement}
\index{Yoneda lemma!proof}
Let $\mathcal{C}$ be a locally small category, let $A\in\mathcal{C}$,
and let $F\colon\mathcal{C}^{\op}\to\Set$ be a presheaf.
There is a bijection
\[
    \Phi_{A,F} \;\colon\;
    \Nat(h^A,\,F) \;\xrightarrow{\;\sim\;}\; F(A)
\]
given by $\Phi_{A,F}(\alpha) = \alpha_A(\id_A)$, where
$\alpha\colon h^A\Rightarrow F$ is a natural transformation.
Moreover, this bijection is natural in both $A$ and~$F$.
\end{theorem}

\begin{proof}
The proof proceeds in four steps.

\medskip\noindent
\textbf{Step 1: Definition of the map~$\Phi$.}
Given a natural transformation $\alpha\colon h^A\Rightarrow F$,
we define
\[
    \Phi_{A,F}(\alpha) \;=\; \alpha_A(\id_A) \;\in\; F(A).
\]
This is well-defined since $\id_A\in h^A(A) = \Hom(A,A)$
and $\alpha_A\colon h^A(A)\to F(A)$.

\medskip\noindent
\textbf{Step 2: Definition of the inverse~$\Psi$.}
Given an element $x\in F(A)$, we define a natural transformation
$\Psi_{A,F}(x) = \alpha^x\colon h^A\Rightarrow F$ as follows.
For each $B\in\mathcal{C}$, the component
\[
    \alpha^x_B \;\colon\; h^A(B) = \Hom(B,A) \longrightarrow F(B)
\]
is defined by
\[
    \alpha^x_B(f) \;=\; F(f)(x) \quad\text{for } f\colon B\to A.
\]
Here $F(f)\colon F(A)\to F(B)$ is the action of the presheaf~$F$ on
the morphism $f$ (recall $F$ is contravariant on~$\mathcal{C}$, i.e.\
covariant on~$\mathcal{C}^{\op}$).

We must verify that $\alpha^x$ is a natural transformation.  Let
$g\colon C\to B$ be a morphism in~$\mathcal{C}$.  We need the following
diagram to commute:
\[
\begin{tikzcd}
h^A(B) \arrow[r, "\alpha^x_B"] \arrow[d, "h^A(g) = g^*"'] &
F(B) \arrow[d, "F(g)"] \\
h^A(C) \arrow[r, "\alpha^x_C"'] &
F(C)
\end{tikzcd}
\]
For $f\in h^A(B) = \Hom(B,A)$:
\begin{align*}
    F(g)(\alpha^x_B(f))
    &= F(g)(F(f)(x)) \\
    &= (F(g)\circ F(f))(x) \\
    &= F(f\circ g)(x) \quad\text{(since $F$ is a functor on $\mathcal{C}^{\op}$)}\\
    &= \alpha^x_C(f\circ g) \\
    &= \alpha^x_C(g^*(f)).
\end{align*}
Hence $\alpha^x$ is indeed natural.

\medskip\noindent
\textbf{Step 3: $\Phi$ and $\Psi$ are mutually inverse.}

\emph{$\Phi\circ\Psi = \id$:}
For $x\in F(A)$:
\[
    \Phi(\Psi(x)) = \Phi(\alpha^x) = \alpha^x_A(\id_A)
    = F(\id_A)(x) = \id_{F(A)}(x) = x.
\]

\emph{$\Psi\circ\Phi = \id$:}
For $\alpha\colon h^A\Rightarrow F$, let $x = \Phi(\alpha) = \alpha_A(\id_A)$.
We must show $\alpha^x = \alpha$, i.e.\ $\alpha^x_B = \alpha_B$ for
all $B\in\mathcal{C}$.  Let $f\in\Hom(B,A)$.  Consider the naturality
square of~$\alpha$ at $f\colon B\to A$:
\[
\begin{tikzcd}
h^A(A) \arrow[r, "\alpha_A"] \arrow[d, "f^*"'] &
F(A) \arrow[d, "F(f)"] \\
h^A(B) \arrow[r, "\alpha_B"'] &
F(B)
\end{tikzcd}
\]
Evaluating at $\id_A \in h^A(A)$:
\[
    \alpha_B(f^*(\id_A)) = F(f)(\alpha_A(\id_A)),
\]
i.e.\
\[
    \alpha_B(f) = F(f)(x) = \alpha^x_B(f).
\]
Since this holds for all $f\in\Hom(B,A)$, we have $\alpha_B = \alpha^x_B$
for all~$B$, hence $\alpha = \alpha^x = \Psi(\Phi(\alpha))$.

\medskip\noindent
\textbf{Step 4: Naturality in $A$ and~$F$.}

\emph{Naturality in~$A$:}
Let $g\colon A'\to A$ be a morphism.  We must show that the diagram
\[
\begin{tikzcd}
\Nat(h^A,F) \arrow[r, "\Phi_{A,F}"] \arrow[d, "g^*"'] &
F(A) \arrow[d, "F(g)"] \\
\Nat(h^{A'},F) \arrow[r, "\Phi_{A',F}"'] &
F(A')
\end{tikzcd}
\]
commutes, where $g^*(\alpha) = \alpha\circ\mathbf{y}(g) = \alpha\circ g_*$.
For $\alpha\in\Nat(h^A,F)$:
\begin{align*}
    \Phi_{A',F}(g^*(\alpha))
    &= \Phi_{A',F}(\alpha\circ g_*)
    = (\alpha\circ g_*)_{A'}(\id_{A'})
    = \alpha_{A'}(g\circ\id_{A'})
    = \alpha_{A'}(g).
\end{align*}
On the other hand, by the naturality of~$\alpha$ at~$g\colon A'\to A$:
\[
    F(g)(\alpha_A(\id_A)) = \alpha_{A'}(g) \quad\text{(from the naturality square)}.
\]
Hence $F(g)(\Phi_{A,F}(\alpha)) = \Phi_{A',F}(g^*(\alpha))$.

\emph{Naturality in~$F$:}
Let $\beta\colon F\Rightarrow G$ be a natural transformation of presheaves.
We must show that
\[
\begin{tikzcd}
\Nat(h^A,F) \arrow[r, "\Phi_{A,F}"] \arrow[d, "\beta_*"'] &
F(A) \arrow[d, "\beta_A"] \\
\Nat(h^A,G) \arrow[r, "\Phi_{A,G}"'] &
G(A)
\end{tikzcd}
\]
commutes, where $\beta_*(\alpha) = \beta\circ\alpha$.
For $\alpha\in\Nat(h^A,F)$:
\[
    \Phi_{A,G}(\beta_*(\alpha)) = \Phi_{A,G}(\beta\circ\alpha)
    = (\beta\circ\alpha)_A(\id_A) = \beta_A(\alpha_A(\id_A))
    = \beta_A(\Phi_{A,F}(\alpha)). \qedhere
\]
\end{proof}

\begin{remark}\label{rem:yoneda-key-idea}
\index{Yoneda lemma!key idea}
The essential insight of the Yoneda lemma is that a natural transformation
$\alpha\colon h^A\Rightarrow F$ is \emph{completely determined} by the
single element $\alpha_A(\id_A)\in F(A)$.  Naturality then forces
the value of~$\alpha$ at every other component.  This is sometimes
called the ``Yoneda trick'': evaluate at the identity.
\end{remark}

\begin{remark}[Dual Yoneda lemma]\label{rem:yoneda-dual}
\index{Yoneda lemma!dual}
The dual version states: for $A\in\mathcal{C}$ and a functor
$G\colon\mathcal{C}\to\Set$,
\[
    \Nat(h_A,\,G) \;\cong\; G(A),
\]
naturally in $A$ and~$G$.  The proof is formally dual.
\end{remark}


%% ============================================================
\section{Corollaries of the Yoneda lemma}\label{sec:yoneda-corollaries}
%% ============================================================

\begin{corollary}[Yoneda embedding is fully faithful]%
\label{cor:yoneda-ff}
\index{Yoneda embedding!fully faithful}
The Yoneda embedding $\mathbf{y}\colon\mathcal{C}\to\widehat{\mathcal{C}}$
is fully faithful.
\end{corollary}

\begin{proof}
For objects $A,B\in\mathcal{C}$, we must show that the map
\[
    \mathbf{y}_{A,B} \;\colon\;
    \Hom_{\mathcal{C}}(A,B) \longrightarrow
    \Nat(h^A,\,h^B) = \Hom_{\widehat{\mathcal{C}}}(\mathbf{y}(A),\mathbf{y}(B))
\]
is a bijection.  By the Yoneda lemma applied with $F = h^B$:
\[
    \Nat(h^A,\,h^B) \;\cong\; h^B(A) = \Hom(A,B).
\]
Moreover, the Yoneda bijection sends $\alpha\in\Nat(h^A,h^B)$ to
$\alpha_A(\id_A)\in\Hom(A,B)$.  For $\alpha = \mathbf{y}(f) = f_*$
(where $f\colon A\to B$), we get $(f_*)_A(\id_A) = f\circ\id_A = f$.
Thus the Yoneda bijection $\Nat(h^A,h^B)\xrightarrow{\sim}\Hom(A,B)$
is the inverse of~$\mathbf{y}_{A,B}$, so $\mathbf{y}_{A,B}$ is
a bijection.
\end{proof}

\begin{corollary}[Representables determine objects up to isomorphism]%
\label{cor:representable-iso}
\index{representable functor!determines object}
If $h^A \cong h^B$ as presheaves on~$\mathcal{C}$, then $A\cong B$
in~$\mathcal{C}$.  Dually, $h_A\cong h_B$ implies $A\cong B$.
\end{corollary}

\begin{proof}
Since $\mathbf{y}$ is fully faithful, it reflects isomorphisms
(see Exercise~\ref{exer:ff-reflects-iso}).
If $h^A\cong h^B$ in $\widehat{\mathcal{C}}$, let
$\alpha\colon h^A\xrightarrow{\sim} h^B$ be a natural isomorphism.
Since $\mathbf{y}$ is full, there exists $f\colon A\to B$ with
$\mathbf{y}(f)=\alpha$.  Similarly, the inverse $\alpha^{-1}$ equals
$\mathbf{y}(g)$ for some $g\colon B\to A$.  Since $\mathbf{y}$ is faithful:
\[
    \mathbf{y}(g\circ f) = \mathbf{y}(g)\circ\mathbf{y}(f)
    = \alpha^{-1}\circ\alpha = \id_{h^A} = \mathbf{y}(\id_A)
    \implies g\circ f = \id_A.
\]
Similarly $f\circ g = \id_B$, so $f$ is an isomorphism.
\end{proof}

\begin{corollary}[Uniqueness of representing objects]%
\label{cor:rep-unique}
\index{representing object!uniqueness}
If a functor $F\colon\mathcal{C}^{\op}\to\Set$ is represented by both
$A$ and~$A'$, then $A\cong A'$ in~$\mathcal{C}$.
\end{corollary}

\begin{proof}
If $F\cong h^A$ and $F\cong h^{A'}$, then $h^A\cong h^{A'}$, hence
$A\cong A'$ by Corollary~\ref{cor:representable-iso}.
\end{proof}

\begin{corollary}\label{cor:yoneda-universal-element}
\index{universal element}
Let $F\colon\mathcal{C}^{\op}\to\Set$.  Then $F$ is representable
(i.e.\ $F\cong h^A$ for some~$A$) if and only if there exists an object
$A\in\mathcal{C}$ and an element $u\in F(A)$---called the
\textbf{universal element}---such that for every $B\in\mathcal{C}$, the map
\[
    \Hom(B,A) \longrightarrow F(B), \quad f\longmapsto F(f)(u)
\]
is a bijection.
\end{corollary}

\begin{proof}
By the Yoneda lemma, a natural isomorphism $\alpha\colon h^A\xrightarrow{\sim} F$
corresponds to the element $u = \alpha_A(\id_A)\in F(A)$.
The component $\alpha_B\colon\Hom(B,A)\to F(B)$ sends $f$ to
$F(f)(u)$ (this is the formula from Step~2 of the proof of the
Yoneda lemma).  Conversely, any $u\in F(A)$ such that $f\mapsto F(f)(u)$
is bijective for all~$B$ defines a natural isomorphism
$h^A\xrightarrow{\sim}F$.
\end{proof}


%% ============================================================
\section{Every presheaf is a colimit of representables}%
\label{sec:colim-representables}
%% ============================================================

\begin{definition}[Category of elements]\label{def:category-of-elements}
\index{category of elements}
\index{$\int F$}
\index{Grothendieck construction!for presheaves}
Let $F\colon\mathcal{C}^{\op}\to\Set$ be a presheaf.
The \textbf{category of elements} of~$F$, denoted $\int F$
(or $\mathcal{C}/F$ or $\mathrm{el}(F)$), has:
\begin{itemize}
\item \textbf{Objects:} pairs $(C,x)$ with $C\in\mathcal{C}$ and $x\in F(C)$.
\item \textbf{Morphisms:} a morphism $(C,x)\to(C',x')$ is a morphism
    $f\colon C\to C'$ in~$\mathcal{C}$ such that $F(f)(x') = x$
    (note the contravariance).
\end{itemize}
There is a canonical projection functor
$\pi_F\colon\int F\to\mathcal{C}$, $(C,x)\mapsto C$.
\end{definition}

\begin{example}\label{ex:elements-representable}
\index{category of elements!of representable}
If $F = h^A$, then an object of $\int h^A$ is a pair $(C,f)$ with
$f\in\Hom(C,A)$, i.e.\ a morphism $f\colon C\to A$.
A morphism $(C,f)\to(C',f')$ is $g\colon C\to C'$ with $f = f'\circ g$.
Thus $\int h^A \cong \mathcal{C}/A$, the slice category over~$A$.
\end{example}

\begin{theorem}[Every presheaf is a colimit of representables]%
\label{thm:presheaf-colim-rep}
\index{presheaf!as colimit of representables}
\index{colimit!of representables}
\index{representable functor!presheaf as colimit of}
Let $F\colon\mathcal{C}^{\op}\to\Set$ be a presheaf.  Then $F$ is the
colimit, in the presheaf category~$\widehat{\mathcal{C}}$, of the
diagram of representable presheaves indexed by the category of elements:
\[
    F \;\cong\; \varinjlim_{(C,x)\,\in\,\int F}\; h^C.
\]
More precisely, the composite functor
$\int F \xrightarrow{\pi_F} \mathcal{C} \xrightarrow{\mathbf{y}}
\widehat{\mathcal{C}}$
has colimit~$F$.
\end{theorem}

\begin{proof}
We construct a cocone and verify the universal property.
For each $(C,x)\in\int F$, the element $x\in F(C)$ corresponds, by the
Yoneda lemma, to a natural transformation
$\alpha^x\colon h^C\Rightarrow F$
defined by $\alpha^x_B(f) = F(f)(x)$ for $f\in\Hom(B,C)$.
This gives a cocone $\{\alpha^x\colon h^C\to F\}_{(C,x)\in\int F}$.

We verify compatibility: if $g\colon(C,x)\to(C',x')$ is a morphism in
$\int F$ (so $g\colon C\to C'$ and $F(g)(x') = x$), then for all
$B\in\mathcal{C}$ and $f\in\Hom(B,C)$:
\[
    \alpha^{x'}_B(g\circ f) = F(g\circ f)(x') = F(f)(F(g)(x'))
    = F(f)(x) = \alpha^x_B(f),
\]
so $\alpha^{x'}\circ g_* = \alpha^x$ as required.

Now let $\{\beta^{(C,x)}\colon h^C\to G\}_{(C,x)\in\int F}$ be any
cocone over the diagram.  We define a natural transformation
$\gamma\colon F\Rightarrow G$ by
\[
    \gamma_B(y) \;=\; \beta^{(B,y)}_B(\id_B) \quad\text{for } y\in F(B).
\]
One checks naturality of~$\gamma$ using the cocone compatibility, and
uniqueness follows because any $\gamma$ satisfying
$\gamma\circ\alpha^x = \beta^{(C,x)}$ for all $(C,x)$ must satisfy
$\gamma_C(x) = \gamma_C(\alpha^x_C(\id_C)) = \beta^{(C,x)}_C(\id_C)$.
\end{proof}

\begin{remark}\label{rem:density}
\index{density theorem}
\index{Yoneda embedding!dense}
Theorem~\ref{thm:presheaf-colim-rep} is sometimes called the
\textbf{density theorem} or the statement that the Yoneda embedding
is \textbf{dense}.  It is fundamental in the theory of Kan extensions
and in the study of locally presentable categories.
\end{remark}


%% ============================================================
\section{Applications of the Yoneda lemma}\label{sec:yoneda-applications}
%% ============================================================

\begin{proposition}[Natural transformations via representables]%
\label{prop:nat-trans-representable}
\index{natural transformation!between representable functors}
For any functor $F\colon\mathcal{C}\to\mathcal{D}$ between locally small
categories, and objects $A,B\in\mathcal{C}$, there is a bijection
\[
    \Nat(h^{F(A)},\, h^{F(B)}) \;\cong\; \Hom_{\mathcal{D}}(F(A),F(B)).
\]
In particular, if $F$ is fully faithful, then it induces a fully faithful
functor on presheaf categories.
\end{proposition}

\begin{proof}
This is the Yoneda lemma applied to~$\mathcal{D}$ with
$G = h^{F(B)}$:
$\Nat(h^{F(A)},h^{F(B)})\cong h^{F(B)}(F(A)) = \Hom(F(A),F(B))$.
\end{proof}

\begin{proposition}[Characterisation of limits via representables]%
\label{prop:limit-representable}
\index{limit!via representable functor}
Let $D\colon\mathcal{J}\to\mathcal{C}$ be a diagram.  An object
$L\in\mathcal{C}$ together with a cone $(\pi_j\colon L\to D_j)$ is a
limit of~$D$ if and only if for every $C\in\mathcal{C}$, the canonical map
\[
    \Hom(C,L) \longrightarrow \varprojlim_{j\in\mathcal{J}}\,\Hom(C,D_j)
\]
is a bijection.  Equivalently, $h^L\cong\varprojlim_j\, h^{D_j}$
as presheaves, i.e.\ limits in~$\mathcal{C}$ are computed pointwise
by the Yoneda embedding.
\end{proposition}

\begin{proof}
By definition, $L = \varprojlim D$ means
$\Hom(C,L)\cong\mathrm{Cone}(C,D)$ naturally in~$C$.
The set of cones from~$C$ to~$D$ is precisely
$\varprojlim_j\Hom(C,D_j)$ (where the limit is in~$\Set$).
Naturality in~$C$ means $h^L\cong\varprojlim_j h^{D_j}$ as presheaves.
Since the Yoneda embedding is fully faithful, this limit in the
presheaf category reflects to a limit in~$\mathcal{C}$.
\end{proof}


%% ============================================================
\section{Exercises}\label{sec:yoneda-exercises}
%% ============================================================

\begin{exercise}\label{exer:ff-reflects-iso}
\index{fully faithful!reflects isomorphisms}
Show that a fully faithful functor reflects isomorphisms: if
$F\colon\mathcal{C}\to\mathcal{D}$ is fully faithful and $F(f)$
is an isomorphism, then $f$ is an isomorphism.
\end{exercise}

\begin{exercise}\label{exer:yoneda-explicit}
\index{Yoneda lemma!explicit example}
Let $\mathcal{C} = \mathsf{Vect}_k$ (vector spaces over a field~$k$).
\begin{enumerate}[label=(\alph*)]
\item Show that $\Nat(h^V,h^W)\cong\Hom(V,W)$ for all vector
    spaces $V,W$ by tracing through the Yoneda bijection.
\item Let $F(V) = V\otimes V$ (a covariant functor).  Compute
    $\Nat(h_k,\,F)$ and identify it with~$F(k) = k$.
\end{enumerate}
\end{exercise}

\begin{exercise}\label{exer:group-as-presheaf}
\index{group!as presheaf}
Let $\mathcal{C}$ be a category with finite products.  A \textbf{group
object} in~$\mathcal{C}$ is an object~$G$ with morphisms
$m\colon G\times G\to G$, $e\colon 1\to G$, $\iota\colon G\to G$
satisfying the usual diagrams.  Show that~$G$ is a group object if and
only if $h^G = \Hom(-,G)\colon\mathcal{C}^{\op}\to\Set$ factors through
$\Grp\to\Set$.
\emph{Hint:} Use the Yoneda lemma to translate element-free diagrams
to element-wise conditions.
\end{exercise}

\begin{exercise}\label{exer:yoneda-monoid}
\index{Yoneda lemma!for monoids}
Let $M$ be a monoid viewed as a one-object category~$\mathcal{B}_M$.
\begin{enumerate}[label=(\alph*)]
\item Show that a presheaf on~$\mathcal{B}_M$ is the same as a
    right $M$-set.
\item Identify the representable presheaf $h^{*}$ (where $*$ is
    the unique object) with $M$ acting on itself by right multiplication.
\item Formulate and verify the Yoneda lemma in this case: for every
    right $M$-set~$S$, the set of $M$-equivariant maps $M\to S$
    is in bijection with~$S$.
\end{enumerate}
\end{exercise}

\begin{exercise}\label{exer:subfunctor}
\index{subfunctor}
\index{sieve}
Let $F\colon\mathcal{C}^{\op}\to\Set$ be a presheaf and let
$S\subseteq F$ be a subfunctor ($S(C)\subseteq F(C)$ for all~$C$,
compatible with restriction maps).  Show that every subfunctor
of~$h^A$ corresponds to a \textbf{sieve} on~$A$, i.e.\ a set~$S$ of
morphisms with codomain~$A$ such that if $f\in S$ and $g$ is composable,
then $f\circ g\in S$.
\end{exercise}

\begin{exercise}\label{exer:cayley-yoneda}
\index{Cayley's theorem!via Yoneda}
\index{Yoneda lemma!Cayley's theorem as special case}
Show that Cayley's theorem in group theory (every group embeds in a
symmetric group) is a special case of the Yoneda lemma.
\emph{Hint:} Apply the Yoneda embedding to the one-object category
associated to a group.
\end{exercise}

\begin{exercise}\label{exer:presheaf-colim-rep}
\index{presheaf!as colimit of representables}
Fill in the remaining details of the proof of
Theorem~\ref{thm:presheaf-colim-rep}: verify that $\gamma$ as defined
there is indeed natural.
\end{exercise}

\begin{exercise}\label{exer:yoneda-enriched}
\index{Yoneda lemma!enriched}
\index{enriched category}
Let $\mathcal{V}$ be a complete, cocomplete, symmetric monoidal closed
category.  Formulate (without proof) the enriched Yoneda lemma for a
$\mathcal{V}$-enriched category~$\mathcal{C}$: for a
$\mathcal{V}$-functor $F\colon\mathcal{C}^{\op}\to\mathcal{V}$ and
$A\in\mathcal{C}$,
\[
    [\mathcal{C}^{\op},\mathcal{V}](h^A,F) \;\cong\; F(A)
\]
in~$\mathcal{V}$, where the left-hand side is the $\mathcal{V}$-object
of $\mathcal{V}$-natural transformations.
\end{exercise}

\index{Yoneda lemma|)}
\index{representable functor|)}
\index{presheaf|)}


%%%%%%%%%%%%%%%%%%%%%%%%%%%%%%%%%%%%%%%%%%%%%%%%%%%%%%%%%%%%
%%        CHAPTER 6: MONADS AND ALGEBRAS                  %%
%%%%%%%%%%%%%%%%%%%%%%%%%%%%%%%%%%%%%%%%%%%%%%%%%%%%%%%%%%%%

\chapter{Monads and Algebras}\label{ch:monads}
\minitoc

\vspace{1em}

Monads are one of the most versatile concepts in category theory,
appearing in algebra, topology, logic, and computer science.
A monad on a category captures the essence of algebraic structure:
it packages together a ``free construction'' and the ``equations''
that the resulting algebraic objects must satisfy.
In this chapter we define monads, show how they arise from adjunctions,
study the categories of algebras over a monad (the Eilenberg--Moore and
Kleisli categories), and prove Beck's monadicity theorem, which
characterises those adjunctions that arise from monads.

\index{monad|(}

%% ============================================================
\section{Monads}\label{sec:monad-def}
%% ============================================================

\begin{definition}[Monad]\label{def:monad}
\index{monad!definition}
\index{unit!of a monad}
\index{multiplication!of a monad}
A \textbf{monad} on a category~$\mathcal{C}$ is a triple
$(T,\eta,\mu)$ where:
\begin{itemize}
\item $T\colon\mathcal{C}\to\mathcal{C}$ is an endofunctor,
\item $\eta\colon\Id_{\mathcal{C}}\Rightarrow T$ is a natural
    transformation called the \textbf{unit},
\item $\mu\colon T^2\Rightarrow T$ is a natural transformation called
    the \textbf{multiplication},
\end{itemize}
subject to the following diagrams commuting for every object
$A\in\mathcal{C}$:

\textbf{Associativity:}
\[
\begin{tikzcd}
T^3A \arrow[r, "T\mu_A"] \arrow[d, "\mu_{TA}"'] &
T^2A \arrow[d, "\mu_A"] \\
T^2A \arrow[r, "\mu_A"'] &
TA
\end{tikzcd}
\]

\textbf{Unit laws:}
\[
\begin{tikzcd}
TA \arrow[r, "T\eta_A"] \arrow[dr, equal] &
T^2A \arrow[d, "\mu_A"] &
TA \arrow[l, "\eta_{TA}"'] \arrow[dl, equal] \\
& TA &
\end{tikzcd}
\]
Equivalently, $\mu\circ T\mu = \mu\circ\mu T$ and
$\mu\circ T\eta = \mu\circ\eta T = \id_T$.
\end{definition}

\begin{remark}\label{rem:monad-as-monoid}
\index{monad!as monoid in endofunctor category}
\index{monoid!in monoidal category}
A monad is a \textbf{monoid} in the monoidal category
$(\mathrm{End}(\mathcal{C}),\circ,\Id)$ of endofunctors with
composition as the monoidal product.  The unit~$\eta$ is the monoid
unit and the multiplication~$\mu$ is the monoid multiplication.
This perspective is extremely useful.
\end{remark}


%% ============================================================
\section{Monads from adjunctions}\label{sec:monad-adj}
%% ============================================================

\begin{proposition}[Every adjunction gives a monad]%
\label{prop:adj-monad}
\index{monad!from adjunction}
\index{adjunction!associated monad}
Let $F\dashv G$ be an adjunction with
$F\colon\mathcal{C}\to\mathcal{D}$, $G\colon\mathcal{D}\to\mathcal{C}$,
unit $\eta\colon\Id_{\mathcal{C}}\Rightarrow GF$, and counit
$\varepsilon\colon FG\Rightarrow\Id_{\mathcal{D}}$.
Then $(T,\eta,\mu)$ is a monad on~$\mathcal{C}$, where $T = GF$ and
$\mu = G\varepsilon F\colon GFGF\Rightarrow GF$.
\end{proposition}

\begin{proof}
\textit{Associativity:}
We must show $\mu\circ T\mu = \mu\circ\mu T$, i.e.\
$G\varepsilon F\circ GF\cdot G\varepsilon F
= G\varepsilon F\circ G\varepsilon F\cdot GF$.
Equivalently, for each~$A$:
\begin{align*}
    \mu_A \circ T\mu_A
    &= G\varepsilon_{FA}\circ GFG\varepsilon_{FA}
    = G(\varepsilon_{FA}\circ FG\varepsilon_{FA}), \\
    \mu_A \circ \mu_{TA}
    &= G\varepsilon_{FA}\circ G\varepsilon_{FGF A}
    = G(\varepsilon_{FA}\circ\varepsilon_{FGFA}).
\end{align*}
By naturality of~$\varepsilon$ at $\varepsilon_{FA}\colon FGFA\to FA$:
$\varepsilon_{FA}\circ FG\varepsilon_{FA} = \varepsilon_{FA}\circ\varepsilon_{FGFA}$.
Hence the two expressions agree.

\textit{Left unit law:} We show $\mu\circ\eta T = \id_T$.
For each~$A$:
$\mu_A\circ\eta_{TA} = G\varepsilon_{FA}\circ\eta_{GFA}
= \id_{GFA} = \id_{TA}$,
by one of the triangle identities ($G\varepsilon\circ\eta G = \id_G$
applied to $FA$).

\textit{Right unit law:} We show $\mu\circ T\eta = \id_T$.
For each~$A$:
$\mu_A\circ T\eta_A = G\varepsilon_{FA}\circ GF\eta_A
= G(\varepsilon_{FA}\circ F\eta_A) = G(\id_{FA}) = \id_{TA}$,
by the other triangle identity ($\varepsilon F\circ F\eta = \id_F$
applied to~$A$).
\end{proof}

\begin{remark}\label{rem:comonad}
\index{comonad}
\index{comonad!from adjunction}
Dually, the adjunction $F\dashv G$ also gives a \textbf{comonad}
$(S,\varepsilon,\delta)$ on~$\mathcal{D}$, where $S = FG$,
the counit is $\varepsilon\colon FG\Rightarrow\Id_{\mathcal{D}}$,
and the comultiplication is $\delta = F\eta G\colon FG\Rightarrow FGFG$.
\end{remark}


%% ============================================================
\section{Examples of monads}\label{sec:monad-examples}
%% ============================================================

\begin{example}[Free monoid monad]\label{ex:monad-free-monoid}
\index{monad!free monoid}
\index{free monoid}
The free-forgetful adjunction $F\dashv U\colon\mathsf{Mon}\to\Set$
gives rise to the monad $T = UF$ on~$\Set$.  Explicitly:
\[
    T(X) \;=\; \coprod_{n\geq 0} X^n \;=\;
    \{(x_1,\ldots,x_n) : n\geq 0,\, x_i\in X\},
\]
the set of finite words on the alphabet~$X$.
The unit $\eta_X\colon X\to T(X)$ sends $x$ to the one-letter
word~$(x)$.  The multiplication $\mu_X\colon T^2(X)\to T(X)$ is
``flattening'' (concatenation): it sends a word of words to a single
word by removing parentheses.  The monad axioms correspond to the
associativity of concatenation and the fact that the empty word is
a unit.
\end{example}

\begin{example}[Free group monad]\label{ex:monad-free-group}
\index{monad!free group}
\index{free group}
The adjunction $F\dashv U\colon\Grp\to\Set$ gives a monad~$T$
where $T(X)$ is the underlying set of the free group on~$X$, i.e.\
the set of reduced words in $X\cup X^{-1}$.
\end{example}

\begin{example}[Free $R$-module monad]\label{ex:monad-free-module}
\index{monad!free module}
\index{free module}
For a ring~$R$, the adjunction $F\dashv U\colon R\text{-}\mathsf{Mod}\to\Set$
gives a monad with
\[
    T(X) \;=\; \bigoplus_{x\in X} R \;=\;
    \Bigl\{\sum_{x\in X} r_x\cdot x :
    r_x\in R,\;\text{finitely many nonzero}\Bigr\}.
\]
\end{example}

\begin{example}[Power set monad]\label{ex:monad-powerset}
\index{monad!power set}
\index{power set monad}
The \textbf{power set monad} $\mathcal{P}$ on $\Set$ is defined by
$\mathcal{P}(X) = 2^X$ (the power set), with
$\eta_X(x) = \{x\}$ and $\mu_X(\mathcal{A}) = \bigcup_{A\in\mathcal{A}} A$
for $\mathcal{A}\in\mathcal{P}(\mathcal{P}(X))$.
This monad arises from the adjunction between $\Set$ and the category
of complete sup-lattices.
\end{example}

\begin{example}[Maybe monad]\label{ex:monad-maybe}
\index{monad!maybe}
\index{maybe monad}
The functor $T(X) = X\sqcup\{*\}$ on $\Set$, with $\eta_X$ the
inclusion $X\hookrightarrow X\sqcup\{*\}$ and $\mu_X$ collapsing
the two copies of~$*$, is a monad.  In computer science, this is the
\textbf{Maybe} (or \textbf{option}) monad, modelling computations that
may fail.
\end{example}

\begin{example}[Continuation monad]\label{ex:monad-continuation}
\index{monad!continuation}
\index{continuation monad}
Fix a set~$R$.  The functor $T(X) = R^{R^X}$ on $\Set$ carries a
monad structure called the \textbf{continuation monad}.
In functional programming, it arises from the self-adjunction
$(-)^R\dashv(-)^R\colon\Set^{\op}\to\Set$.
\end{example}

\begin{example}[Monad from a closure operator]\label{ex:monad-closure}
\index{monad!closure operator}
\index{closure operator}
Let $(P,\leq)$ be a poset viewed as a category.  A monad on~$P$ is a
\textbf{closure operator}: a monotone map $c\colon P\to P$ with
$x\leq c(x)$ (unit) and $c(c(x))\leq c(x)$ (multiplication), which
combined with monotonicity gives $c(c(x)) = c(x)$ (idempotence).
\end{example}


%% ============================================================
\section{Eilenberg--Moore algebras}\label{sec:em-algebras}
%% ============================================================

\begin{definition}[Algebra over a monad]\label{def:em-algebra}
\index{Eilenberg--Moore algebra}
\index{monad!algebra}
\index{$T$-algebra}
Let $(T,\eta,\mu)$ be a monad on~$\mathcal{C}$.
A \textbf{$T$-algebra} (or \textbf{Eilenberg--Moore algebra}) is a pair
$(A,a)$ where $A\in\mathcal{C}$ is an object and
$a\colon TA\to A$ is a morphism (the \textbf{structure map}) such that:

\textbf{Unit:}
$a\circ\eta_A = \id_A$.
\[
\begin{tikzcd}
A \arrow[r, "\eta_A"] \arrow[dr, equal] &
TA \arrow[d, "a"] \\
& A
\end{tikzcd}
\]

\textbf{Associativity:}
$a\circ\mu_A = a\circ Ta$.
\[
\begin{tikzcd}
T^2A \arrow[r, "\mu_A"] \arrow[d, "Ta"'] &
TA \arrow[d, "a"] \\
TA \arrow[r, "a"'] &
A
\end{tikzcd}
\]
\end{definition}

\begin{definition}[Morphism of algebras]\label{def:em-morphism}
\index{Eilenberg--Moore algebra!morphism}
\index{$T$-algebra!morphism}
A \textbf{morphism of $T$-algebras} $f\colon(A,a)\to(B,b)$ is a
morphism $f\colon A\to B$ in~$\mathcal{C}$ such that $f\circ a = b\circ Tf$:
\[
\begin{tikzcd}
TA \arrow[r, "Tf"] \arrow[d, "a"'] &
TB \arrow[d, "b"] \\
A \arrow[r, "f"'] &
B
\end{tikzcd}
\]
\end{definition}

\begin{definition}[Eilenberg--Moore category]\label{def:em-category}
\index{Eilenberg--Moore category}
\index{$\mathcal{C}^T$}
The \textbf{Eilenberg--Moore category} $\mathcal{C}^T$ has:
\begin{itemize}
\item \textbf{Objects:} $T$-algebras $(A,a)$,
\item \textbf{Morphisms:} morphisms of $T$-algebras,
\item \textbf{Composition and identities:} inherited from~$\mathcal{C}$.
\end{itemize}
There is a forgetful functor $U^T\colon\mathcal{C}^T\to\mathcal{C}$,
$(A,a)\mapsto A$, which has a left adjoint
$F^T\colon\mathcal{C}\to\mathcal{C}^T$ defined by
$F^T(A) = (TA,\mu_A)$.
\end{definition}

\begin{proposition}\label{prop:em-adjunction}
\index{Eilenberg--Moore category!adjunction}
The pair $F^T\dashv U^T$ is an adjunction, and the monad arising
from this adjunction (via Proposition~\ref{prop:adj-monad}) is
precisely~$(T,\eta,\mu)$.
\end{proposition}

\begin{proof}
\textit{$F^T(A) = (TA,\mu_A)$ is a $T$-algebra:}
The unit axiom is $\mu_A\circ\eta_{TA} = \id_{TA}$ (left unit of the
monad).  The associativity axiom is $\mu_A\circ\mu_{TA} = \mu_A\circ T\mu_A$
(monad associativity).

\textit{Adjunction:} For $(A,a)\in\mathcal{C}^T$ and $B\in\mathcal{C}$,
define
\[
    \Phi\colon\Hom_{\mathcal{C}^T}(F^TB,\,(A,a))
    \longrightarrow \Hom_{\mathcal{C}}(B,\,U^T(A,a) = A)
\]
by $\Phi(f) = f\circ\eta_B$.  The inverse is
$\Psi(g) = a\circ Tg$ for $g\colon B\to A$.
One verifies that $\Psi(g)$ is a $T$-algebra morphism and that
$\Phi,\Psi$ are mutually inverse using the algebra axioms and monad axioms.

\textit{Recovery of the monad:}
$U^TF^T(A) = U^T(TA,\mu_A) = TA = T(A)$, so $U^TF^T = T$.
The unit of the adjunction $F^T\dashv U^T$ is $\eta$.
The counit at $(A,a)$ is the algebra map
$\varepsilon_{(A,a)} = a\colon(T A,\mu_A)\to(A,a)$,
which is a $T$-algebra morphism by the associativity axiom for~$(A,a)$.
The induced multiplication is $U^T\varepsilon_{F^T}\cdot\eta = \mu$,
as desired.
\end{proof}

\begin{example}\label{ex:em-algebras}
\index{Eilenberg--Moore algebra!examples}
\begin{enumerate}[label=(\roman*)]
\item \textbf{Free monoid monad} (Example~\ref{ex:monad-free-monoid}):
    $T$-algebras are monoids.  A structure map $a\colon T(X)\to X$
    takes a word $(x_1,\ldots,x_n)$ and produces a single element:
    this is the multiplication.  The algebra axioms force associativity
    and unitality.
    \index{monoid!as Eilenberg--Moore algebra}

\item \textbf{Free group monad} (Example~\ref{ex:monad-free-group}):
    $T$-algebras are groups.  More precisely,
    $\Set^T\simeq\Grp$.
    \index{group!as Eilenberg--Moore algebra}

\item \textbf{Free module monad} (Example~\ref{ex:monad-free-module}):
    $T$-algebras are $R$-modules, i.e.\ $\Set^T\simeq R\text{-}\mathsf{Mod}$.
    \index{module!as Eilenberg--Moore algebra}

\item \textbf{Power set monad} (Example~\ref{ex:monad-powerset}):
    $T$-algebras are complete sup-lattices with sup-preserving maps
    as morphisms.
    \index{complete lattice!as Eilenberg--Moore algebra}

\item \textbf{Closure operators} (Example~\ref{ex:monad-closure}):
    the algebras for a closure operator $c$ on a poset~$P$ are exactly
    the \textbf{closed elements}, i.e.\ those $x$ with $c(x) = x$.
    \index{closed element}
\end{enumerate}
\end{example}


%% ============================================================
\section{Free algebras}\label{sec:free-algebras}
%% ============================================================

\begin{definition}[Free $T$-algebra]\label{def:free-algebra}
\index{free algebra}
\index{$T$-algebra!free}
A $T$-algebra $(A,a)$ is \textbf{free} if it is isomorphic to
$F^T(X) = (TX,\mu_X)$ for some $X\in\mathcal{C}$.
\end{definition}

\begin{proposition}\label{prop:free-algebra-universal}
\index{free algebra!universal property}
The free $T$-algebra on~$X$ satisfies the universal property:
for every $T$-algebra $(A,a)$ and every morphism $f\colon X\to A$
in~$\mathcal{C}$, there is a unique $T$-algebra morphism
$\bar{f}\colon(TX,\mu_X)\to(A,a)$ with $\bar{f}\circ\eta_X = f$.
Explicitly, $\bar{f} = a\circ Tf$.
\end{proposition}

\begin{proof}
This is simply the adjunction $F^T\dashv U^T$ applied to the pair
$(X,(A,a))$: $\Hom_{\mathcal{C}^T}(F^TX,(A,a))\cong\Hom_{\mathcal{C}}(X,A)$.
\end{proof}


%% ============================================================
\section{The Kleisli category}\label{sec:kleisli}
%% ============================================================

\begin{definition}[Kleisli category]\label{def:kleisli}
\index{Kleisli category}
\index{$\mathcal{C}_T$}
Let $(T,\eta,\mu)$ be a monad on~$\mathcal{C}$.
The \textbf{Kleisli category} $\mathcal{C}_T$ has:
\begin{itemize}
\item \textbf{Objects:} the same as~$\mathcal{C}$,
\item \textbf{Morphisms:} $\Hom_{\mathcal{C}_T}(A,B) = \Hom_{\mathcal{C}}(A,TB)$,
\item \textbf{Identity:} $\id_A = \eta_A\colon A\to TA$,
\item \textbf{Composition:} given $f\colon A\to TB$ and $g\colon B\to TC$
    (in~$\mathcal{C}$), their Kleisli composite is
    \[
        g\circ_T f \;=\; \mu_C\circ Tg\circ f \;\colon\; A\to TC.
    \]
\end{itemize}
\end{definition}

\begin{proposition}\label{prop:kleisli-category}
\index{Kleisli category!well-defined}
The Kleisli category~$\mathcal{C}_T$ is well-defined, i.e.\
composition is associative and $\eta_A$ is a two-sided identity.
\end{proposition}

\begin{proof}
\textit{Left identity:}
$\id_B\circ_T f = \mu_B\circ T\eta_B\circ f = \id_{TB}\circ f = f$
(right unit law of the monad).

\textit{Right identity:}
$f\circ_T\id_A = \mu_B\circ Tf\circ\eta_A = f$
(naturality of~$\eta$: $Tf\circ\eta_A = \eta_{TB}\circ f$, then
$\mu_B\circ\eta_{TB} = \id_{TB}$).

\textit{Associativity:}
Let $f\colon A\to TB$, $g\colon B\to TC$, $h\colon C\to TD$.  Then:
\begin{align*}
    h\circ_T(g\circ_T f)
    &= \mu_D\circ Th\circ\mu_C\circ Tg\circ f, \\
    (h\circ_T g)\circ_T f
    &= \mu_D\circ T(\mu_D\circ Th\circ g)\circ f
    = \mu_D\circ T\mu_D\circ T^2h\circ Tg\circ f.
\end{align*}
By naturality of~$\mu$ at~$h$: $\mu_D\circ Th = \mu_D\circ Th$.
More precisely, we use $\mu_D\circ T\mu_D = \mu_D\circ\mu_{TD}$
(associativity of the monad) and naturality of~$\mu$ at~$h$:
$\mu_D\circ Th = Th \circ \mu_C$ is not quite right; let us be
more careful.  By naturality of~$\mu$ at $h\colon C\to TD$,
$\mu_D\circ T^2h = Th\circ\mu_C$... actually this requires
$h\colon TC\to TD$.  Instead we use:
\begin{align*}
    \mu_D\circ Th\circ\mu_C
    &= \mu_D\circ\mu_{TD}\circ T^2h \quad\text{(naturality of $\mu$ at $h$: applied as $\mu_D\circ Th = \ldots$)}
\end{align*}
Let us argue differently.  By naturality of~$\mu$ applied to~$Th\colon TC\to T^2D$:
\[
    \mu_D\circ T\mu_D = \mu_D\circ\mu_{TD}
    \quad\text{(monad associativity)}.
\]
And by naturality of~$\mu$ at the morphism $h\colon C\to TD$:
$\mu_{TD}\circ T^2 h = Th\circ\mu_C$... No, $h\colon C\to TD$,
so $Th\colon TC\to T^2D$ and $T^2h\colon T^2C\to T^3D$.
Naturality of $\mu$ gives $\mu_D\circ Th = h^*\circ\mu_C$... this is
getting confused.

The cleanest argument: define $f^{\sharp} = \mu_B\circ Tf$ for
$f\colon A\to TB$, so $g\circ_T f = g^{\sharp}\circ f$.  Then
associativity $(h\circ_T g)\circ_T f = h\circ_T(g\circ_T f)$
is equivalent to $(h^{\sharp}\circ g)^{\sharp} = h^{\sharp}\circ g^{\sharp}$,
i.e.\ $\mu\circ T(h^{\sharp}\circ g) = h^{\sharp}\circ\mu\circ Tg$,
which is $\mu\circ Th^{\sharp}\circ Tg = h^{\sharp}\circ\mu\circ Tg$.
This reduces to $\mu\circ Th^{\sharp} = h^{\sharp}\circ\mu$,
i.e.\ $\mu\circ T(\mu\circ Th) = \mu\circ Th\circ\mu$, which follows
from monad associativity and naturality of~$\mu$.
\end{proof}

\begin{proposition}\label{prop:kleisli-adjunction}
\index{Kleisli category!adjunction}
There is an adjunction $F_T\dashv G_T$ with
$F_T\colon\mathcal{C}\to\mathcal{C}_T$ and
$G_T\colon\mathcal{C}_T\to\mathcal{C}$ such that $G_TF_T = T$
and the monad arising from $F_T\dashv G_T$ is~$(T,\eta,\mu)$.
Explicitly:
\begin{itemize}
\item $F_T(A) = A$ on objects, $F_T(f) = \eta_B\circ f$ on morphisms
    $f\colon A\to B$,
\item $G_T(A) = TA$ on objects, $G_T(g) = \mu_B\circ Tg$ for
    $g\colon A\to TB$ (a Kleisli morphism $A\to B$).
\end{itemize}
\end{proposition}

\begin{proof}
The adjunction is verified directly:
$\Hom_{\mathcal{C}_T}(F_TA,B) = \Hom_{\mathcal{C}}(A,TB) = \Hom_{\mathcal{C}}(A,G_TB)$.
The unit is~$\eta$, the counit at~$B$ is $\id_{TB}\colon TB\to TB$
viewed as a Kleisli morphism $F_TG_TB = F_T(TB)\to B$.
One checks that $G_TF_T = T$ and the induced monad is $(T,\eta,\mu)$.
\end{proof}

\begin{remark}\label{rem:kleisli-em-universal}
\index{Kleisli category!as initial resolution}
\index{Eilenberg--Moore category!as terminal resolution}
The Kleisli category is the \textbf{initial} object and the
Eilenberg--Moore category is the \textbf{terminal} object in the
category of adjunctions giving rise to the monad~$T$.
More precisely, if $F\dashv G\colon\mathcal{D}\to\mathcal{C}$ is
any adjunction with associated monad~$T$, then there exist unique
(up to isomorphism) comparison functors
\[
    \mathcal{C}_T \longrightarrow \mathcal{D}
    \longrightarrow \mathcal{C}^T
\]
compatible with the adjunctions.  The comparison functor
$K\colon\mathcal{D}\to\mathcal{C}^T$ sends $D$ to $(GD,G\varepsilon_D)$.
\end{remark}

\begin{example}[Kleisli for the power set monad]\label{ex:kleisli-powerset}
\index{Kleisli category!power set monad}
\index{relation!as Kleisli morphism}
The Kleisli category for the power set monad~$\mathcal{P}$ on $\Set$
has sets as objects and $\Hom(A,B) = \Hom(A,\mathcal{P}(B)) \cong
\{R\subseteq A\times B\}$ as morphisms, i.e.\ morphisms are
\textbf{relations}.  Kleisli composition corresponds to composition
of relations.
\end{example}

\begin{example}[Kleisli for the Maybe monad]\label{ex:kleisli-maybe}
\index{Kleisli category!maybe monad}
The Kleisli category for the Maybe monad $T(X)=X\sqcup\{*\}$ has
morphisms $A\to B\sqcup\{*\}$, i.e.\ partial functions.  This
captures the idea of computations that may fail.
\end{example}


%% ============================================================
\section{Beck's monadicity theorem}\label{sec:beck}
%% ============================================================

The central question is: given an adjunction $F\dashv G$, when is the
comparison functor $K\colon\mathcal{D}\to\mathcal{C}^T$ an equivalence?

\begin{definition}[Comparison functor]\label{def:comparison-functor}
\index{comparison functor}
\index{$K$}
Let $F\dashv G\colon\mathcal{D}\to\mathcal{C}$ be an adjunction with
associated monad $(T,\eta,\mu)$ on~$\mathcal{C}$ (where $T = GF$).
The \textbf{comparison functor}
$K\colon\mathcal{D}\to\mathcal{C}^T$
is defined by:
\begin{itemize}
\item On objects: $K(D) = (GD,\, G\varepsilon_D)$, where
    $\varepsilon_D\colon FGD\to D$ is the counit.
    (This is a $T$-algebra: the unit and associativity axioms follow
    from the triangle identities and naturality of~$\varepsilon$.)
\item On morphisms: $K(f) = Gf$ for $f\colon D\to D'$.
    (This is a $T$-algebra morphism by naturality of~$\varepsilon$.)
\end{itemize}
\end{definition}

\begin{definition}[Monadic functor]\label{def:monadic}
\index{monadic functor}
\index{monadic adjunction}
An adjunction $F\dashv G$ (equivalently, the right adjoint~$G$) is
\textbf{monadic} if the comparison functor $K$ is an equivalence of
categories.  It is \textbf{strictly monadic} if $K$ is an isomorphism
of categories.
\end{definition}

\begin{definition}[$G$-split coequaliser]\label{def:g-split-coeq}
\index{coequaliser!$G$-split}
\index{split coequaliser}
\index{contractible coequaliser}
Let $G\colon\mathcal{D}\to\mathcal{C}$ be a functor.
A parallel pair $f,g\colon D\rightrightarrows D'$ in~$\mathcal{D}$ is
called a \textbf{$G$-split pair} if the image
$Gf,Gg\colon GD\rightrightarrows GD'$ admits a \textbf{split coequaliser}
in~$\mathcal{C}$: a diagram
\[
\begin{tikzcd}
GD \arrow[r, shift left, "Gf"] \arrow[r, shift right, "Gg"'] &
GD' \arrow[r, "e"] \arrow[l, bend left=30, "s"'] &
C \arrow[l, bend left=30, "t"']
\end{tikzcd}
\]
with $e\circ Gf = e\circ Gg$, $e\circ s = \id_C$, $Gf\circ t = s\circ e$,
and $Gg\circ t = \id_{GD'}$.

A \textbf{$G$-split coequaliser} of $f,g$ is a coequaliser of $f,g$
in~$\mathcal{D}$ whose image under~$G$ is this split coequaliser.
\end{definition}

\begin{remark}\label{rem:split-coeq-absolute}
\index{split coequaliser!absolute}
Split coequalisers are \textbf{absolute}: they are preserved by any
functor whatsoever.  Indeed, if $(e,s,t)$ is a split coequaliser as
above, then applying any functor~$H$ preserves all the relevant
equations, so $He$ is the coequaliser of $HGf,HGg$.  This crucial
property is what makes Beck's theorem work.
\end{remark}

\begin{lemma}\label{lem:split-coeq}
\index{split coequaliser!is a coequaliser}
A split coequaliser is a coequaliser.
\end{lemma}

\begin{proof}
Let $Gf,Gg\colon A\rightrightarrows B$ with split coequaliser $(e,s,t)$
as in Definition~\ref{def:g-split-coeq} (with $A=GD$, $B=GD'$).
First, $e$ coequalises: $e\circ Gf = e\circ Gg$ by assumption.
For universality, let $h\colon B\to Z$ with $h\circ Gf = h\circ Gg$.
Define $u = h\circ t\colon C\to Z$.  Then
$u\circ e = h\circ t\circ e = h\circ Gf\circ s = h\circ Gg\circ s$...
wait, we need: $u\circ e = h\circ t\circ e$.  Using $Gf\circ t = s\circ e$:
$h\circ s\circ e = h\circ Gf\circ t$... let us recompute.

We have $e\circ s = \id_C$, so $u\circ e = h\circ t\circ e$.
Now $t\circ e = Gf\circ t$... no, $Gf\circ t = s\circ e$.
Actually: from $Gg\circ t = \id_B$, we get
\[
    u\circ e = h\circ t\circ e = h\circ\id_B = h?
\]
No. Let us be careful.  We need:
From $Gg\circ t = \id_B$: this gives $t\colon C\to A$ is a section of
$Gg$... no, $t\colon C\to B$ and $Gg\colon A\to B$, so this does
not type-check.  Let me reconsider.

In the standard formulation, a split coequaliser of $f_1,f_2\colon A\rightrightarrows B$
is a diagram $A\rightrightarrows B\xrightarrow{e} C$ with
$s\colon C\to B$ and $t\colon B\to A$ such that
$e\circ f_1 = e\circ f_2$, $e\circ s = \id_C$,
$f_1\circ t = s\circ e$, $f_2\circ t = \id_B$.

Then for universality: let $h\colon B\to Z$ with $h\circ f_1 = h\circ f_2$.
Define $u = h\circ s\colon C\to Z$.  Then:
$u\circ e = h\circ s\circ e = h\circ f_1\circ t = h\circ f_2\circ t
= h\circ\id_B^{} = h$... wait: $f_2\circ t = \id_B$, so
$h\circ f_2\circ t = h$.  And $h\circ f_1\circ t = h\circ f_2\circ t = h$
since $h\circ f_1 = h\circ f_2$.  So $u\circ e = h\circ s\circ e
= h\circ f_1\circ t = h$.

Uniqueness: if $u'\circ e = h$, then $u' = u'\circ e\circ s = h\circ s = u$.
\end{proof}

Now we can state Beck's theorem.

\begin{theorem}[Beck's Monadicity Theorem]\label{thm:beck}
\index{Beck's monadicity theorem}
\index{monadicity theorem}
\index{monadic functor!characterisation}
Let $F\dashv G$ with $G\colon\mathcal{D}\to\mathcal{C}$.
The following are equivalent:
\begin{enumerate}[label=(\roman*)]
\item $G$ is monadic (the comparison functor $K$ is an equivalence).
\item $G$ reflects isomorphisms and $\mathcal{D}$ has coequalisers of
    all $G$-split pairs, and $G$ preserves them.
\end{enumerate}
\end{theorem}

\begin{proof}
The proof is substantial.  We break it into parts.

\medskip\noindent
\textbf{(i) $\Rightarrow$ (ii):}

Assume $K\colon\mathcal{D}\to\mathcal{C}^T$ is an equivalence.
Since $U^T\circ K = G$ and $K$ is an equivalence, it suffices to
show that $U^T\colon\mathcal{C}^T\to\mathcal{C}$ reflects isomorphisms
and that $\mathcal{C}^T$ has coequalisers of $U^T$-split pairs.

\emph{$U^T$ reflects isomorphisms:}
Let $f\colon(A,a)\to(B,b)$ be a $T$-algebra morphism such that
$f\colon A\to B$ is an isomorphism in~$\mathcal{C}$.
We claim $f^{-1}\colon(B,b)\to(A,a)$ is a $T$-algebra morphism.
We need $f^{-1}\circ b = a\circ Tf^{-1}$.  Since $f\circ a = b\circ Tf$
(as $f$ is a $T$-algebra morphism), we get
$a = f^{-1}\circ b\circ Tf$, hence
$a\circ Tf^{-1} = f^{-1}\circ b\circ Tf\circ Tf^{-1}
= f^{-1}\circ b\circ T(f\circ f^{-1}) = f^{-1}\circ b$.

\emph{$\mathcal{C}^T$ has coequalisers of $U^T$-split pairs:}
Let $f,g\colon(A,a)\rightrightarrows(B,b)$ be $T$-algebra morphisms
such that $f,g\colon A\rightrightarrows B$ admits a split coequaliser
$(e\colon B\to C,\, s\colon C\to B,\, t\colon B\to A)$ in~$\mathcal{C}$.
We must endow~$C$ with a $T$-algebra structure.

Define $c\colon TC\to C$ by
\[
    c \;=\; e\circ b\circ Ts \;\colon\; TC\to C.
\]
\emph{$c$ is well-defined as a structure map:}
We verify the algebra axioms.  Unit: $c\circ\eta_C = e\circ b\circ Ts\circ\eta_C
= e\circ b\circ\eta_B\circ s = e\circ s = \id_C$
(using naturality of~$\eta$ and the algebra unit for~$(B,b)$).
Associativity: $c\circ\mu_C = e\circ b\circ Ts\circ\mu_C
= e\circ b\circ T(b\circ Ts) = e\circ b\circ Tb\circ T^2s$...
We use the algebra associativity for $(B,b)$: $b\circ Tb = b\circ\mu_B$.
So $c\circ\mu_C = e\circ b\circ\mu_B\circ T^2s = e\circ b\circ Ts\circ T(e\circ b\circ Ts)$...
let us use $\mu_C = T(s)\circ... $ no. Actually:
\begin{align*}
    c\circ Tc &= (e\circ b\circ Ts)\circ T(e\circ b\circ Ts) \\
    &= e\circ b\circ Ts\circ Te\circ Tb\circ T^2s \\
    &= e\circ b\circ T(s\circ e)\circ Tb\circ T^2s \\
    &= e\circ b\circ T(f\circ t)\circ Tb\circ T^2s
    \quad\text{(since $s\circ e = f\circ t$... no, $f_1\circ t = s\circ e$)}
\end{align*}
We have $f\circ t = s\circ e$ from the split coequaliser.
So $T(s\circ e) = Tf\circ Tt$.  Then:
\begin{align*}
    c\circ Tc &= e\circ b\circ Tf\circ Tt\circ Tb\circ T^2s \\
    &= e\circ f\circ a\circ Tt\circ Tb\circ T^2s
    \quad\text{($b\circ Tf = f\circ a$ since $f$ is a $T$-morphism)}\\
    &= e\circ g\circ a\circ Tt\circ Tb\circ T^2s
    \quad\text{($e\circ f = e\circ g$)}\\
    &= e\circ b\circ Tg\circ Tt\circ Tb\circ T^2s
    \quad\text{($b\circ Tg = g\circ a$ since $g$ is a $T$-morphism)}\\
    &= e\circ b\circ T(g\circ t)\circ Tb\circ T^2s \\
    &= e\circ b\circ T(\id_B)\circ Tb\circ T^2s
    \quad\text{($g\circ t = \id_B$)} \\
    &= e\circ b\circ Tb\circ T^2s \\
    &= e\circ b\circ\mu_B\circ T^2s
    \quad\text{(associativity for $(B,b)$)} \\
    &= e\circ b\circ Ts\circ\mu_C
    \quad\text{(naturality of $\mu$ at $s$: $\mu_C = Ts... $ no)}
\end{align*}
Naturality of~$\mu$ at $s\colon C\to B$ gives
$\mu_B\circ T^2s = Ts\circ\mu_C$.  Hence:
\[
    c\circ Tc = e\circ b\circ Ts\circ\mu_C = c\circ\mu_C.
\]
So $(C,c)$ is indeed a $T$-algebra.

\emph{$e$ is a $T$-algebra morphism:}
We check $e\circ b = c\circ Te$.
$c\circ Te = e\circ b\circ Ts\circ Te = e\circ b\circ T(s\circ e)
= e\circ b\circ Tf\circ Tt = e\circ f\circ a\circ Tt$.
But also $e\circ b = e\circ b\circ T(g\circ t) = e\circ g\circ a\circ Tt
= e\circ f\circ a\circ Tt$... this is not right either.
We need $c\circ Te = e\circ b$ directly.
$c\circ Te = e\circ b\circ T(s\circ e) = e\circ b\circ T(f\circ t)
= e\circ b\circ Tf\circ Tt = e\circ f\circ a\circ Tt$.
And $e\circ b = e\circ b\circ\id_{TB} = e\circ b\circ T(g\circ t)
= e\circ g\circ a\circ Tt = e\circ f\circ a\circ Tt$.
Here we used $g\circ t = \id_B$ and $e\circ g = e\circ f$.
So indeed $c\circ Te = e\circ b$.

\emph{$(C,c)$ is the coequaliser:}
Suppose $h\colon(B,b)\to(Z,z)$ is a $T$-algebra morphism with
$h\circ f = h\circ g$.  Since the underlying diagram is a split
coequaliser, there is a unique $u\colon C\to Z$ in~$\mathcal{C}$ with
$u\circ e = h$.  We must show $u$ is a $T$-algebra morphism, i.e.\
$u\circ c = z\circ Tu$.
\begin{align*}
    u\circ c &= u\circ e\circ b\circ Ts = h\circ b\circ Ts
    = z\circ Th\circ Ts = z\circ T(h\circ s)
    = z\circ T(u\circ e\circ s) = z\circ Tu.
\end{align*}
(Using $h\circ b = z\circ Th$ since $h$ is a $T$-morphism, and
$e\circ s = \id_C$.)

\medskip\noindent
\textbf{(ii) $\Rightarrow$ (i):}

Assume $G$ reflects isomorphisms and $\mathcal{D}$ has coequalisers
of $G$-split pairs, preserved by~$G$.  We show $K$ is an equivalence
by constructing a quasi-inverse.

\emph{$K$ is essentially surjective:}
Let $(A,a)\in\mathcal{C}^T$.  Consider the pair
$FTA \rightrightarrows FA$ in~$\mathcal{D}$ given by
\[
    Fa \;\text{ and }\; F\mu_A \quad (\text{no: }
    \varepsilon_{FA}\colon FGFA = FTA\to FA).
\]
More precisely, consider the canonical pair:
$F\mu_A, \varepsilon_{FTA}\colon FT^2A = FGFTA\rightrightarrows FTA$...

Let us use the standard construction.  For $(A,a)\in\mathcal{C}^T$,
consider the parallel pair in~$\mathcal{D}$:
\[
    FTGA \;=\; FTA \underset{\varepsilon_{FA}}{\overset{Fa}{\rightrightarrows}} FA.
\]
Wait---$a\colon TA\to A$ is a morphism in~$\mathcal{C}$, so $Fa\colon FTA\to FA$.
And $\varepsilon_{FA}\colon FGFA = FTA\to FA$.  So we have
$Fa,\varepsilon_{FA}\colon FTA\rightrightarrows FA$.

Applying~$G$:
$GFa,G\varepsilon_{FA}\colon GFTA = T^2A\rightrightarrows TA = GFA$.
That is, $Ta,\mu_A\colon T^2A\rightrightarrows TA$.
This pair has a split coequaliser in~$\mathcal{C}$ with
$e = a\colon TA\to A$, $s = \eta_A\colon A\to TA$,
$t = T\eta_A... $ let us verify:
\begin{itemize}
\item $e\circ Ta = a\circ Ta$ and $e\circ\mu_A = a\circ\mu_A$.
    The algebra associativity says $a\circ Ta = a\circ\mu_A$. \checkmark
\item $e\circ s = a\circ\eta_A = \id_A$. \checkmark\ (algebra unit)
\item $Ta\circ t = s\circ e$: We need $t\colon TA\to T^2A$ with
    $Ta\circ t = \eta_A\circ a$ and $\mu_A\circ t = \id_{TA}$.
    Take $t = T\eta_A$.  Then $\mu_A\circ T\eta_A = \id_{TA}$
    (monad right unit). \checkmark
\item $Ta\circ T\eta_A = T(a\circ\eta_A) = T(\id_A) = \id_{TA}$.
    But we need $Ta\circ t = s\circ e = \eta_A\circ a$.
    We have $Ta\circ T\eta_A = \id_{TA}\neq\eta_A\circ a$ in general.
\end{itemize}

So let us take $t = \eta_{TA}\colon TA\to T^2A$ instead.
\begin{itemize}
\item $\mu_A\circ\eta_{TA} = \id_{TA}$. \checkmark\ (monad left unit)
\item $Ta\circ\eta_{TA} = \eta_A\circ a$: this is naturality of~$\eta$
    at~$a$. \checkmark
\end{itemize}
So the split coequaliser data is: $e = a$, $s = \eta_A$, $t = \eta_{TA}$,
with the coequalised pair being $(Ta,\mu_A)$.

So the pair $(Fa,\varepsilon_{FA})$ is a $G$-split pair.
By hypothesis, $\mathcal{D}$ has a coequaliser
$q\colon FA\to D$ of $(Fa,\varepsilon_{FA})$, and $G$ preserves it.
Since $G$ preserves this coequaliser and the split coequaliser of
$(Ta,\mu_A)$ is $a\colon TA\to A$, we get $Gq = a$ (up to the
canonical isomorphism).  More precisely, $GD\cong A$ via the
isomorphism coming from the fact that $Gq\colon GFA = TA\to GD$
is the coequaliser of $(Ta,\mu_A)$ in~$\mathcal{C}$, which is~$a$
with target~$A$.  So $GD\cong A$.

Now $K(D) = (GD, G\varepsilon_D)$.
Under the identification $GD\cong A$, we get a $T$-algebra structure
on~$A$ that agrees with~$a$.  Hence $K(D)\cong(A,a)$, and $K$ is
essentially surjective.

\emph{$K$ is fully faithful:}
For $D,D'\in\mathcal{D}$, the map
$K_{D,D'}\colon\Hom_{\mathcal{D}}(D,D')\to
\Hom_{\mathcal{C}^T}(K(D),K(D'))$
sends $f\mapsto Gf$.  Since morphisms of $T$-algebras are just
morphisms in~$\mathcal{C}$ satisfying a compatibility condition,
$\Hom_{\mathcal{C}^T}(K(D),K(D')) \subseteq \Hom_{\mathcal{C}}(GD,GD')$.

\emph{Faithful:}
If $Gf = Gf'$, we want $f = f'$.  Consider the coequalisers
$q\colon FA\to D$ from the essential surjectivity argument.
Since $Gf = Gf'$ and $G$ reflects the relevant coequalisers, one
shows $f = f'$.  Alternatively: for any $D\in\mathcal{D}$,
the counit $\varepsilon_D\colon FGD\to D$ is a (regular) epimorphism
(it is the coequaliser of a $G$-split pair).
If $Gf = Gf'$, then $FGf = FGf'$, so
$f\circ\varepsilon_D = \varepsilon_{D'}\circ FGf
= \varepsilon_{D'}\circ FGf' = f'\circ\varepsilon_D$.
Since $\varepsilon_D$ is an epimorphism, $f = f'$.

\emph{Full:}
Let $\varphi\colon GD\to GD'$ be a $T$-algebra morphism
$K(D)\to K(D')$.  Consider the diagram in~$\mathcal{D}$:
\[
\begin{tikzcd}
FGFGD \arrow[r, shift left, "F(G\varepsilon_D)"]
       \arrow[r, shift right, "\varepsilon_{FGD}"'] &
FGD \arrow[r, "\varepsilon_D"] &
D
\end{tikzcd}
\]
and the corresponding one for~$D'$.  The map $F\varphi\colon FGD\to FGD'$
satisfies the required compatibility (since $\varphi$ is a $T$-algebra
morphism).  Since $\varepsilon_D$ is the coequaliser and
$\varepsilon_{D'}\circ F\varphi$ coequalises the pair for~$D$,
there exists a unique $f\colon D\to D'$ with
$f\circ\varepsilon_D = \varepsilon_{D'}\circ F\varphi$.
Applying~$G$: $Gf\circ G\varepsilon_D = G\varepsilon_{D'}\circ GF\varphi
= G\varepsilon_{D'}\circ T\varphi$.  But for a $T$-algebra morphism,
$\varphi\circ G\varepsilon_D = G\varepsilon_{D'}\circ T\varphi$, so
$Gf\circ G\varepsilon_D = \varphi\circ G\varepsilon_D$.
Since $G\varepsilon_D = a$ is an epimorphism (it is a split epimorphism),
$Gf = \varphi$.

Hence $K$ is fully faithful and essentially surjective, so it is an
equivalence.
\end{proof}

\begin{remark}\label{rem:beck-crude}
\index{crude monadicity theorem}
There is a \textbf{crude monadicity theorem} (also called the
\textbf{tripleability theorem} in older terminology): $G$ is monadic
if and only if $G$ has a left adjoint, $G$ reflects isomorphisms,
and $\mathcal{D}$ has and $G$ preserves coequalisers of
\emph{reflexive} $G$-split pairs (a reflexive pair is one admitting
a common section).  This is sometimes easier to verify in practice.
\end{remark}

\begin{corollary}\label{cor:monadic-reflects}
\index{monadic functor!reflects isomorphisms}
A monadic functor reflects isomorphisms and creates coequalisers of
$G$-split pairs.
\end{corollary}

\begin{proof}
This follows immediately from the (i) $\Rightarrow$ (ii) direction
of Beck's theorem.
\end{proof}


%% ============================================================
\section{Applications of monadicity}\label{sec:monadicity-apps}
%% ============================================================

\begin{example}[Algebraic categories are monadic over $\Set$]%
\label{ex:monadic-algebraic}
\index{monadic functor!algebraic categories}
\index{algebraic category}
The forgetful functors from $\Grp$, $\Ab$, $\mathsf{Ring}$,
$R\text{-}\mathsf{Mod}$, and more generally any variety of
algebras in the sense of universal algebra, to~$\Set$ are monadic.
We verify the conditions of Beck's theorem:
\begin{enumerate}[label=(\alph*)]
\item The forgetful functor $U$ has a left adjoint (the free functor).
\item $U$ reflects isomorphisms: a bijective homomorphism is an
    isomorphism.
\item $\Set$ has all coequalisers, and the free-forgetful adjunction
    creates them via the algebra structure.
\end{enumerate}
\end{example}

\begin{example}[Compact Hausdorff spaces]\label{ex:monadic-comphaus}
\index{monadic functor!compact Hausdorff spaces}
\index{compact Hausdorff space}
The forgetful functor $U\colon\mathsf{CompHaus}\to\Set$ is monadic.
The associated monad is the ultrafilter monad~$\beta$, whose
algebras are compact Hausdorff spaces.  This is a deep theorem of
Manes (1969).
\index{ultrafilter monad}
\index{Manes' theorem}
\end{example}

\begin{example}[Non-monadic functors]\label{ex:non-monadic}
\index{monadic functor!non-examples}
\begin{enumerate}[label=(\roman*)]
\item The forgetful functor $U\colon\Top\to\Set$ is \emph{not} monadic.
    It has a left adjoint (discrete topology) but does not reflect
    isomorphisms: a continuous bijection need not be a homeomorphism.

\item The inclusion $\Ab\hookrightarrow\Grp$ is not monadic:
    it does not reflect isomorphisms in general, and the associated
    monad (abelianisation) does not recover~$\Ab$ as the category of
    algebras (the Eilenberg--Moore category is equivalent to~$\Ab$,
    but the inclusion is not the comparison functor from~$\Grp$).
    Actually, the functor $\Ab\to\Grp$ does reflect isomorphisms
    (it is fully faithful), so the issue is that coequalisers of
    split pairs in~$\Grp$ need not land in~$\Ab$.
\end{enumerate}
\end{example}

\begin{proposition}\label{prop:monadic-creates-limits}
\index{monadic functor!creates limits}
A monadic functor creates all limits that exist in the base category.
In particular, if $G\colon\mathcal{D}\to\mathcal{C}$ is monadic
and $\mathcal{C}$ is complete, then so is~$\mathcal{D}$.
\end{proposition}

\begin{proof}
Since $G$ is monadic, $\mathcal{D}\simeq\mathcal{C}^T$ for a
monad~$T$ on~$\mathcal{C}$.  It suffices to show $U^T$ creates limits.
Let $D\colon\mathcal{J}\to\mathcal{C}^T$ be a diagram with
$D(j) = (A_j,a_j)$, and suppose $L = \varprojlim_j A_j$ exists
in~$\mathcal{C}$ with projections $\pi_j\colon L\to A_j$.
Then $TL = T(\varprojlim A_j)$ and we need a $T$-algebra structure
on~$L$.  Since $T$ (being a right adjoint composite $G\circ F$... no,
$T$ need not preserve limits in general).

For each~$j$, the maps $a_j\circ T\pi_j\colon TL\to A_j$ form a cone
over the diagram $\{A_j\}$ (using compatibility of the $a_j$ with
the algebra morphisms).  By the universal property of~$L$, there is a
unique $\ell\colon TL\to L$ with $\pi_j\circ\ell = a_j\circ T\pi_j$
for all~$j$.  One verifies that $(L,\ell)$ is a $T$-algebra and is
the limit of~$D$ in~$\mathcal{C}^T$.  The algebra axioms for~$(L,\ell)$
follow from those for each~$(A_j,a_j)$ and the fact that the $\pi_j$
jointly detect equality (since $L$ is their limit).
\end{proof}


%% ============================================================
\section{Exercises}\label{sec:monad-exercises}
%% ============================================================

\begin{exercise}\label{exer:monad-triangle}
\index{monad!from adjunction}
Complete the details of Proposition~\ref{prop:adj-monad}:
verify each monad axiom carefully using the triangle identities.
\end{exercise}

\begin{exercise}\label{exer:monad-list}
\index{monad!list monad}
\index{list monad}
Let $T\colon\Set\to\Set$ be the ``list'' functor with
$T(X) = \coprod_{n\geq 0}X^n$.  Define $\eta$ and $\mu$ explicitly
and verify the monad axioms.  Show that $T$-algebras are monoids.
\end{exercise}

\begin{exercise}\label{exer:em-adjoint}
\index{Eilenberg--Moore category!forgetful-free adjunction}
Verify all details of Proposition~\ref{prop:em-adjunction}: show that
$F^T\dashv U^T$ is an adjunction and that the associated monad is
$(T,\eta,\mu)$.
\end{exercise}

\begin{exercise}\label{exer:kleisli-composition}
\index{Kleisli category!composition}
For the free group monad $T$ on $\Set$, describe Kleisli composition
explicitly.  Show that a Kleisli morphism $A\to B$ is a function
$A\to F(B)$ (where $F(B)$ is the free group on~$B$) and that
composition corresponds to substitution followed by reduction.
\end{exercise}

\begin{exercise}\label{exer:idempotent-monad}
\index{monad!idempotent}
\index{idempotent monad}
A monad $(T,\eta,\mu)$ is \textbf{idempotent} if
$\mu\colon T^2\Rightarrow T$ is a natural isomorphism.
\begin{enumerate}[label=(\alph*)]
\item Show that the following are equivalent: (i) $T$ is idempotent;
    (ii) $T\eta = \eta T$; (iii) every $T$-algebra is free.
\item Show that the monad arising from a reflective subcategory
    (i.e.\ a full subcategory whose inclusion has a left adjoint)
    is idempotent.
\item Deduce that the category of algebras for an idempotent monad
    is equivalent to the full subcategory of~$\mathcal{C}$ consisting
    of objects of the form~$TA$.
\end{enumerate}
\end{exercise}

\begin{exercise}\label{exer:monad-distributive-law}
\index{monad!distributive law}
\index{distributive law}
Let $(S,\eta^S,\mu^S)$ and $(T,\eta^T,\mu^T)$ be monads on~$\mathcal{C}$.
A \textbf{distributive law} of $S$ over~$T$ is a natural transformation
$\lambda\colon TS\Rightarrow ST$ satisfying four axioms (compatibility
with units and multiplications of both monads).
\begin{enumerate}[label=(\alph*)]
\item Write down the four axiom diagrams.
\item Show that given a distributive law, the composite $ST$ carries a
    monad structure with unit $\eta^S\eta^T$ and multiplication
    $\mu^S\cdot S\lambda T\cdot\mu^T$.
\item Show that the ring axiom ``multiplication distributes over
    addition'' is a distributive law for the free monoid monad over
    the free abelian group monad.
\end{enumerate}
\end{exercise}

\begin{exercise}\label{exer:beck-groups}
\index{Beck's monadicity theorem!applied to groups}
Use Beck's monadicity theorem to give an alternative proof that the
forgetful functor $U\colon\Grp\to\Set$ is monadic.
\emph{Hint:} Verify that $U$ reflects isomorphisms and that $\Grp$
has coequalisers of $U$-split pairs.
\end{exercise}

\begin{exercise}\label{exer:monadic-top}
\index{monadic functor!$\Top$ is not monadic over $\Set$}
Show that the forgetful functor $U\colon\Top\to\Set$ is not monadic by
exhibiting a continuous bijection that is not a homeomorphism.
What is the monad $T = UF$ (where $F$ is the discrete topology functor)?
\end{exercise}

\begin{exercise}\label{exer:comonad-coalgebra}
\index{comonad!coalgebra}
\index{coalgebra!over a comonad}
Dualise the theory of this chapter: define a \textbf{coalgebra} for a
comonad $(S,\varepsilon,\delta)$ on~$\mathcal{C}$ and formulate the
dual of Beck's theorem for comonadic functors.
\end{exercise}

\begin{exercise}\label{exer:monad-maps}
\index{monad!morphism}
A \textbf{morphism of monads} $\varphi\colon(S,\eta^S,\mu^S)\to
(T,\eta^T,\mu^T)$ on the same category~$\mathcal{C}$ is a natural
transformation $\varphi\colon S\Rightarrow T$ compatible with units and
multiplications:
$\varphi\circ\eta^S = \eta^T$ and
$\varphi\circ\mu^S = \mu^T\circ\varphi\varphi$
(where $\varphi\varphi = \varphi T\circ S\varphi = T\varphi\circ\varphi S$).
\begin{enumerate}[label=(\alph*)]
\item Show that a monad morphism $\varphi\colon S\to T$ induces a functor
    $\varphi^*\colon\mathcal{C}^T\to\mathcal{C}^S$ (restriction of scalars).
\item Show that $\varphi^*$ has a left adjoint (extension of scalars).
\end{enumerate}
\end{exercise}

\begin{exercise}\label{exer:monad-coequaliser-detail}
\index{Beck's monadicity theorem!coequaliser construction}
In the proof of Beck's theorem (direction (i) $\Rightarrow$ (ii)),
verify directly that the candidate $T$-algebra structure
$c = e\circ b\circ Ts$ on~$C$ satisfies the associativity axiom, filling
in any steps that were left implicit.
\end{exercise}

\index{monad|)}
% Chapter 7-8: Abelian Categories and Derived Categories
% Continuation -- no preamble

\chapter{Abelian Categories and Exact Sequences}\label{ch:abelian}
\minitoc

\section{Preadditive and Additive Categories}\label{sec:preadditive}

The passage from general category theory to homological algebra requires
enriching the categorical framework with algebraic structure on morphism sets.
This section introduces the hierarchy of additive structures on categories.

\begin{definition}[Preadditive category]\label{def:preadditive-cat}
\index{category!preadditive}\index{preadditive category}
A \textbf{preadditive category} (or \textbf{$\Ab$-category}) is a category $\mathcal{A}$
such that:
\begin{enumerate}
    \item For every pair of objects $A, B \in \Ob(\mathcal{A})$, the set $\Hom_{\mathcal{A}}(A,B)$
    carries the structure of an abelian group, with addition denoted $+$ and zero element
    denoted $0_{A,B}$.
    \item Composition is bilinear: for all morphisms $f, f' \colon A \to B$ and
    $g, g' \colon B \to C$,
    \[
        g \circ (f + f') = g \circ f + g \circ f', \qquad
        (g + g') \circ f = g \circ f + g' \circ f.
    \]
\end{enumerate}
Equivalently, $\mathcal{A}$ is a category enriched over $(\Ab, \otimes_{\mathbb{Z}}, \mathbb{Z})$.
\end{definition}

\begin{remark}\label{rem:preadditive-zero}
\index{zero morphism}
In a preadditive category, for any objects $A, B$, the zero element $0_{A,B} \in \Hom(A,B)$
satisfies $0_{B,C} \circ f = 0_{A,C}$ and $g \circ 0_{A,B} = 0_{A,C}$ for all
$f \colon A \to B$ and $g \colon B \to C$. Thus the zero morphisms form a compatible system.
If $\mathcal{A}$ has a zero object $0$, then $0_{A,B}$ is the composite $A \to 0 \to B$.
\end{remark}

\begin{example}\label{ex:preadditive-examples}
\index{Ab@$\Ab$!as preadditive category}\index{Mod@$\Mod$!as preadditive category}
\begin{enumerate}
    \item The category $\Ab$ of abelian groups is preadditive, with pointwise addition
    of homomorphisms.
    \item For any ring $R$, the category $\Mod_R$ of right $R$-modules is preadditive.
    \item The category $\Set$ is \emph{not} preadditive: there is no natural abelian group
    structure on $\Hom_{\Set}(X,Y)$ for general sets $X, Y$.
    \item Any ring $R$ can be viewed as a preadditive category with a single object $*$,
    where $\Hom(*, *) = R$.
\end{enumerate}
\end{example}

\begin{definition}[Additive functor]\label{def:additive-functor}
\index{functor!additive}\index{additive functor}
Let $\mathcal{A}, \mathcal{B}$ be preadditive categories. A functor
$F \colon \mathcal{A} \to \mathcal{B}$ is \textbf{additive} if for all objects $A, B$
the map
\[
    F_{A,B} \colon \Hom_{\mathcal{A}}(A,B) \longrightarrow \Hom_{\mathcal{B}}(FA, FB)
\]
is a group homomorphism: $F(f + g) = F(f) + F(g)$.
\end{definition}

\begin{definition}[Biproduct]\label{def:biproduct}
\index{biproduct}\index{direct sum!in additive category}
Let $\mathcal{A}$ be a preadditive category and let $A_1, \ldots, A_n$ be objects.
A \textbf{biproduct} of $A_1, \ldots, A_n$ is an object $A_1 \oplus \cdots \oplus A_n$
together with morphisms
\[
    \iota_k \colon A_k \to A_1 \oplus \cdots \oplus A_n, \qquad
    \pi_k \colon A_1 \oplus \cdots \oplus A_n \to A_k, \qquad 1 \leq k \leq n,
\]
satisfying:
\begin{enumerate}
    \item $\pi_k \circ \iota_j = \delta_{jk} \cdot \id_{A_k}$ for all $j, k$,
    \item $\sum_{k=1}^n \iota_k \circ \pi_k = \id_{A_1 \oplus \cdots \oplus A_n}$.
\end{enumerate}
\end{definition}

\begin{proposition}\label{prop:biproduct-universal}
\index{biproduct!universal property}
A biproduct $(A_1 \oplus \cdots \oplus A_n, \iota_k, \pi_k)$ is simultaneously
a product and a coproduct. More precisely:
\begin{enumerate}
    \item $(A_1 \oplus \cdots \oplus A_n, (\pi_k))$ is a product of $A_1, \ldots, A_n$.
    \item $(A_1 \oplus \cdots \oplus A_n, (\iota_k))$ is a coproduct of $A_1, \ldots, A_n$.
\end{enumerate}
\end{proposition}

\begin{proof}
We prove the product property; the coproduct property is dual. Given morphisms
$f_k \colon B \to A_k$ for $1 \leq k \leq n$, define
$f = \sum_{k=1}^n \iota_k \circ f_k \colon B \to A_1 \oplus \cdots \oplus A_n$.
Then
\[
    \pi_j \circ f = \sum_{k=1}^n \pi_j \circ \iota_k \circ f_k
    = \sum_{k=1}^n \delta_{jk} f_k = f_j.
\]
For uniqueness, if $g \colon B \to A_1 \oplus \cdots \oplus A_n$ satisfies
$\pi_k \circ g = f_k$ for all $k$, then
\[
    g = \Bigl(\sum_{k=1}^n \iota_k \circ \pi_k\Bigr) \circ g
    = \sum_{k=1}^n \iota_k \circ f_k = f. \qedhere
\]
\end{proof}

\begin{definition}[Additive category]\label{def:additive-cat}
\index{category!additive}\index{additive category}
An \textbf{additive category} is a preadditive category $\mathcal{A}$ that:
\begin{enumerate}
    \item has a zero object $0$,
    \item has finite biproducts: for any objects $A, B \in \Ob(\mathcal{A})$,
    the biproduct $A \oplus B$ exists.
\end{enumerate}
\end{definition}

\begin{proposition}\label{prop:additive-criterion}
\index{additive category!characterization}
Let $\mathcal{A}$ be a preadditive category with a zero object and finite products
(or finite coproducts). Then $\mathcal{A}$ is additive.
\end{proposition}

\begin{proof}
Suppose $\mathcal{A}$ has finite products. Let $A_1, A_2 \in \Ob(\mathcal{A})$,
and let $(A_1 \times A_2, \pi_1, \pi_2)$ be their product. Define
$\iota_1 \colon A_1 \to A_1 \times A_2$ as the unique morphism with
$\pi_1 \circ \iota_1 = \id_{A_1}$ and $\pi_2 \circ \iota_1 = 0$,
and similarly for $\iota_2$. One checks that $\iota_1 \circ \pi_1 + \iota_2 \circ \pi_2 = \id$
using the universal property and the fact that
\[
    \pi_j \circ (\iota_1 \circ \pi_1 + \iota_2 \circ \pi_2) = \delta_{1j}\pi_1 + \delta_{2j}\pi_2 = \pi_j
\]
for $j = 1, 2$. Thus $(A_1 \times A_2, \iota_k, \pi_k)$ is a biproduct.
\end{proof}

\begin{example}\label{ex:biproduct-matrix}
\index{biproduct!matrix description}
In $\Mod_R$, the biproduct $M \oplus N$ is the usual direct sum. The morphism set
$\Hom(M_1 \oplus M_2, N_1 \oplus N_2)$ can be identified with $2 \times 2$ matrices:
\[
    \begin{pmatrix} f_{11} & f_{12} \\ f_{21} & f_{22} \end{pmatrix},
    \qquad f_{ij} \colon M_j \to N_i,
\]
with composition given by matrix multiplication. This matrix calculus is a key
feature of additive categories.
\end{example}


\section{Kernels, Cokernels, and Abelian Categories}\label{sec:abelian-def}

\begin{definition}[Kernel and cokernel]\label{def:kernel-cokernel}
\index{kernel!in additive category}\index{cokernel}
Let $f \colon A \to B$ be a morphism in a preadditive category $\mathcal{A}$.
\begin{enumerate}
    \item A \textbf{kernel} of $f$ is an equalizer of $f$ and $0_{A,B}$: it is a morphism
    $k \colon K \to A$ such that $f \circ k = 0$ and for every $g \colon X \to A$ with
    $f \circ g = 0$, there exists a unique $\bar{g} \colon X \to K$ with $k \circ \bar{g} = g$.
    % Diagram
    \[
    \begin{tikzcd}
        X \arrow[rd, "g"'] \arrow[rrd, "0", bend left=20] \arrow[d, dashed, "\exists!\bar{g}"'] \\
        K \arrow[r, "k"'] & A \arrow[r, "f"'] & B
    \end{tikzcd}
    \]
    We write $\Ker(f) = (K, k)$ or simply $\Ker(f) = K$.
    \item A \textbf{cokernel} of $f$ is a coequalizer of $f$ and $0_{A,B}$: it is a morphism
    $c \colon B \to C$ such that $c \circ f = 0$ and for every $h \colon B \to Y$ with
    $h \circ f = 0$, there exists a unique $\bar{h} \colon C \to Y$ with $\bar{h} \circ c = h$.
    We write $\coker(f) = (C, c)$.
\end{enumerate}
\end{definition}

\begin{definition}[Image and coimage]\label{def:image-coimage}
\index{image!categorical}\index{coimage}
Let $f \colon A \to B$ be a morphism in a preadditive category.
\begin{enumerate}
    \item The \textbf{image} of $f$ is $\im(f) = \Ker(\coker(f))$.
    \item The \textbf{coimage} of $f$ is $\operatorname{coim}(f) = \coker(\Ker(f))$.
\end{enumerate}
There is a canonical morphism $\overline{f} \colon \operatorname{coim}(f) \to \im(f)$
induced by the universal properties.
\end{definition}

\begin{definition}[Abelian category]\label{def:abelian-cat}
\index{abelian category}\index{category!abelian}
An \textbf{abelian category} is an additive category $\mathcal{A}$ satisfying:
\begin{enumerate}
    \item[(AB1)] Every morphism has a kernel and a cokernel.
    \item[(AB2)] For every morphism $f$, the canonical morphism
    $\overline{f} \colon \operatorname{coim}(f) \to \im(f)$ is an isomorphism.
\end{enumerate}
\end{definition}

\begin{remark}\label{rem:abelian-equiv}
\index{abelian category!equivalent formulations}
Condition (AB2) can equivalently be stated: every monomorphism is a kernel and
every epimorphism is a cokernel. That is, every mono $m \colon K \hookrightarrow A$ is
the kernel of some morphism out of $A$, and dually for epimorphisms.
\end{remark}

\begin{example}\label{ex:abelian-examples}
\begin{enumerate}
    \item\index{Ab@$\Ab$!as abelian category}
    $\Ab$ is an abelian category. Kernels and cokernels are the usual ones.
    \item\index{Mod@$\Mod$!as abelian category}
    $\Mod_R$ is abelian for any ring $R$. This is the prototypical example.
    \item\index{sheaf!category of sheaves}
    For a topological space $X$, the category $\operatorname{Sh}(X, \Ab)$ of sheaves of
    abelian groups on $X$ is abelian.
    \item The category of finitely generated abelian groups is abelian.
    \item\index{category!not abelian}
    The category of free abelian groups is additive but \emph{not} abelian: the cokernel
    of $\mathbb{Z} \xrightarrow{2} \mathbb{Z}$ would be $\mathbb{Z}/2\mathbb{Z}$, which
    is not free.
\end{enumerate}
\end{example}

\begin{definition}[Grothendieck axioms]\label{def:grothendieck-axioms}
\index{Grothendieck axioms}\index{AB3}\index{AB4}\index{AB5}\index{AB3*}\index{AB4*}\index{AB5*}
Let $\mathcal{A}$ be an abelian category. Grothendieck introduced additional axioms:
\begin{enumerate}
    \item[(AB3)] $\mathcal{A}$ has arbitrary (set-indexed) coproducts.
    \item[(AB3*)] $\mathcal{A}$ has arbitrary (set-indexed) products.
    \item[(AB4)] (AB3) holds and coproducts are exact: if $f_i \colon A_i \to B_i$ is a
    family of monomorphisms, then $\bigoplus f_i \colon \bigoplus A_i \to \bigoplus B_i$
    is a monomorphism.
    \item[(AB4*)] (AB3*) holds and products are exact.
    \item[(AB5)] (AB3) holds and filtered colimits are exact: filtered colimits commute
    with finite limits.
    \item[(AB5*)] (AB3*) holds and filtered limits (cofiltered limits) are exact.
\end{enumerate}
A \textbf{Grothendieck category}\index{Grothendieck category} is an (AB5) abelian category
with a generator.
\end{definition}

\begin{example}\label{ex:grothendieck-cat}
\begin{enumerate}
    \item $\Mod_R$ satisfies (AB3)--(AB5) and (AB3*)--(AB4*). It is a Grothendieck category
    with generator $R$. It does \emph{not} satisfy (AB5*) in general.
    \item $\operatorname{Sh}(X, \Ab)$ is a Grothendieck category.
    \item The category $\Ab^{\mathrm{f}}$ of finite abelian groups satisfies (AB1)--(AB2)
    but not (AB3): infinite coproducts do not exist.
\end{enumerate}
\end{example}

\begin{theorem}[Freyd--Mitchell embedding]\label{thm:freyd-mitchell}
\index{Freyd--Mitchell embedding theorem}\index{embedding theorem}
Let $\mathcal{A}$ be a small abelian category. Then there exists a ring $R$ and an
exact, fully faithful functor $F \colon \mathcal{A} \hookrightarrow \Mod_R$.
\end{theorem}

\begin{remark}\label{rem:freyd-mitchell}
The Freyd--Mitchell theorem justifies ``diagram chasing'' in abstract abelian categories:
any statement about finite diagrams involving kernels, cokernels, exactness, etc., that
holds in $\Mod_R$ for all rings $R$ automatically holds in any abelian category.
We shall freely use element-wise arguments in proofs below, understood via this embedding.
\end{remark}


\section{Exact Sequences}\label{sec:exact-seq}

\begin{definition}[Exact sequence]\label{def:exact-seq}
\index{exact sequence}\index{sequence!exact}
Let $\mathcal{A}$ be an abelian category. A sequence of morphisms
\[
    \cdots \longrightarrow A_{n+1} \xrightarrow{f_{n+1}} A_n \xrightarrow{f_n} A_{n-1}
    \longrightarrow \cdots
\]
is \textbf{exact at $A_n$} if $\im(f_{n+1}) = \Ker(f_n)$ (as subobjects of $A_n$).
Equivalently, $f_n \circ f_{n+1} = 0$ and the induced morphism
$\operatorname{coim}(f_{n+1}) \to \Ker(f_n)$ is an isomorphism.
The sequence is \textbf{exact} if it is exact at every term.
\end{definition}

\begin{definition}[Short exact sequence]\label{def:ses}
\index{short exact sequence}\index{exact sequence!short}
A \textbf{short exact sequence} (SES) in $\mathcal{A}$ is an exact sequence
\[
    0 \longrightarrow A \xrightarrow{\;\;f\;\;} B \xrightarrow{\;\;g\;\;} C \longrightarrow 0.
\]
This means:
\begin{itemize}
    \item $f$ is a monomorphism (exactness at $A$),
    \item $g$ is an epimorphism (exactness at $C$),
    \item $\im(f) = \Ker(g)$ (exactness at $B$).
\end{itemize}
\end{definition}

\begin{example}\label{ex:ses-examples}
\index{short exact sequence!examples}
\begin{enumerate}
    \item In $\Ab$: $0 \to \mathbb{Z} \xrightarrow{\times n} \mathbb{Z} \to \mathbb{Z}/n\mathbb{Z} \to 0$
    for any $n \geq 1$.
    \item In $\Mod_R$: for any submodule $N \subseteq M$, we have
    $0 \to N \hookrightarrow M \twoheadrightarrow M/N \to 0$.
    \item A SES $0 \to A \to B \to C \to 0$ \textbf{splits}\index{split exact sequence}
    if there exists $s \colon C \to B$ with $g \circ s = \id_C$, or equivalently
    $r \colon B \to A$ with $r \circ f = \id_A$. In this case $B \cong A \oplus C$.
\end{enumerate}
\end{example}

\begin{proposition}[Splitting lemma]\label{prop:splitting-lemma}
\index{splitting lemma}
For a short exact sequence $0 \to A \xrightarrow{f} B \xrightarrow{g} C \to 0$
in an abelian category, the following are equivalent:
\begin{enumerate}
    \item There exists $s \colon C \to B$ with $g \circ s = \id_C$ (right splitting).
    \item There exists $r \colon B \to A$ with $r \circ f = \id_A$ (left splitting).
    \item $B \cong A \oplus C$ via an isomorphism compatible with $f$ and $g$.
\end{enumerate}
\end{proposition}

\begin{proof}
$(3) \Rightarrow (1)$ and $(3) \Rightarrow (2)$ are clear from the biproduct structure.

$(1) \Rightarrow (3)$: Given $s \colon C \to B$ with $g \circ s = \id_C$, define
$\varphi \colon A \oplus C \to B$ by $\varphi = f \circ \pi_A + s \circ \pi_C$.
Then $g \circ \varphi \circ \iota_C = g \circ s = \id_C$ and
$g \circ \varphi \circ \iota_A = g \circ f = 0$. We claim $\varphi$ is an isomorphism.
By the five lemma (Theorem~\ref{thm:five-lemma} below), or by direct argument:
for $b \in B$, write $b = f(a) + s(g(b))$ where $a$ is the unique element with
$f(a) = b - s(g(b))$ (which lies in $\Ker(g) = \im(f)$).

$(2) \Rightarrow (3)$: Dual argument.
\end{proof}


\section{Exact Functors}\label{sec:exact-functors}

\begin{definition}[Exact functor]\label{def:exact-functor}
\index{functor!exact}\index{exact functor}\index{functor!left exact}\index{functor!right exact}
Let $F \colon \mathcal{A} \to \mathcal{B}$ be an additive functor between abelian categories.
\begin{enumerate}
    \item $F$ is \textbf{left exact} if for every short exact sequence
    $0 \to A \to B \to C \to 0$, the sequence $0 \to FA \to FB \to FC$ is exact.
    \item $F$ is \textbf{right exact} if for every SES, $FA \to FB \to FC \to 0$ is exact.
    \item $F$ is \textbf{exact} if it is both left and right exact, i.e., it preserves
    short exact sequences.
\end{enumerate}
\end{definition}

\begin{example}\label{ex:exact-functors}
\begin{enumerate}
    \item\index{Hom@$\Hom$!left exactness}
    For any $R$-module $M$, the functor $\Hom_R(M, -)$ is left exact (covariant).
    The functor $\Hom_R(-, M)$ is left exact (contravariant, i.e., exact as a functor
    $\Mod_R^{\op} \to \Ab$).
    \item\index{tensor product!right exactness}
    The tensor product functor $M \otimes_R -$ is right exact for any right $R$-module $M$.
    \item The forgetful functor $\Mod_R \to \Ab$ is exact.
    \item\index{flat module}
    An $R$-module $M$ is \textbf{flat} if and only if $M \otimes_R -$ is exact.
\end{enumerate}
\end{example}

\begin{proposition}\label{prop:exact-functor-les}
\index{exact functor!characterizations}
Let $F \colon \mathcal{A} \to \mathcal{B}$ be an additive functor. Then:
\begin{enumerate}
    \item $F$ is left exact if and only if $F$ preserves kernels.
    \item $F$ is right exact if and only if $F$ preserves cokernels.
    \item $F$ is exact if and only if $F$ preserves short exact sequences.
\end{enumerate}
\end{proposition}


\section{The Snake Lemma}\label{sec:snake-lemma}

The snake lemma is one of the fundamental tools of homological algebra. It
produces long exact sequences from short ones.

\begin{theorem}[Snake lemma]\label{thm:snake-lemma}
\index{snake lemma}\index{connecting morphism}
Consider a commutative diagram in an abelian category $\mathcal{A}$ with exact rows:
\[
\begin{tikzcd}
    & A' \arrow[r, "f'"] \arrow[d, "a"] & B' \arrow[r, "g'"] \arrow[d, "b"] & C' \arrow[r] \arrow[d, "c"] & 0 \\
    0 \arrow[r] & A \arrow[r, "f"'] & B \arrow[r, "g"'] & C &
\end{tikzcd}
\]
Then there exists an exact sequence
\[
    \Ker(a) \xrightarrow{\;\Ker(f')\;} \Ker(b) \xrightarrow{\;\Ker(g')\;} \Ker(c)
    \xrightarrow{\;\;\delta\;\;} \coker(a) \xrightarrow{\;\coker(f)\;} \coker(b)
    \xrightarrow{\;\coker(g)\;} \coker(c),
\]
where $\delta$ is the \textbf{connecting morphism}\index{connecting morphism!snake lemma}.
Moreover:
\begin{itemize}
    \item If $f'$ is a monomorphism, so is $\Ker(a) \to \Ker(b)$.
    \item If $g$ is an epimorphism, so is $\coker(b) \to \coker(c)$.
\end{itemize}
\end{theorem}

\begin{proof}
We give a complete proof using element-chasing (justified by the Freyd--Mitchell theorem).

\medskip\noindent
\textbf{Step 1: The induced morphisms.}
The morphisms $\Ker(a) \to \Ker(b)$ and $\Ker(b) \to \Ker(c)$ are induced by $f'$ and $g'$
respectively. Indeed, if $x \in \Ker(a)$, then $b(f'(x)) = f(a(x)) = f(0) = 0$, so
$f'(x) \in \Ker(b)$. Similarly for $g'$. The morphisms $\coker(a) \to \coker(b) \to \coker(c)$
are induced dually.

\medskip\noindent
\textbf{Step 2: Construction of $\delta$.}
Let $x \in \Ker(c) \subseteq C'$. Since $g'$ is surjective (onto $C'$), choose
$y \in B'$ with $g'(y) = x$. Consider $b(y) \in B$. Then
\[
    g(b(y)) = c(g'(y)) = c(x) = 0,
\]
so $b(y) \in \Ker(g) = \im(f)$. Choose $z \in A$ with $f(z) = b(y)$ (unique since $f$
is injective). Define
\[
    \delta(x) = [z] \in \coker(a) = A / \im(a).
\]

\medskip\noindent
\textbf{Step 3: Well-definedness of $\delta$.}
If $y'$ is another choice with $g'(y') = x$, then $g'(y - y') = 0$, so $y - y' = f'(w)$
for some $w \in A'$. Then $b(y) - b(y') = b(f'(w)) = f(a(w))$, so $z - z' = a(w)$
since $f$ is injective. Thus $[z] = [z']$ in $\coker(a)$.

\medskip\noindent
\textbf{Step 4: Exactness at $\Ker(b)$.}
The composition $\Ker(a) \to \Ker(b) \to \Ker(c)$ is $g' \circ f' = 0$ (since the top
row is a complex). Conversely, let $y \in \Ker(b)$ with $g'(y) \in \Ker(c)$ mapping to $0$
in $\Ker(c)$---that is, $g'(y) = 0$. Then $y \in \Ker(g') = \im(f')$, so $y = f'(x)$
for some $x \in A'$. Now $f(a(x)) = b(f'(x)) = b(y) = 0$, and since $f$ is injective,
$a(x) = 0$, so $x \in \Ker(a)$.

\medskip\noindent
\textbf{Step 5: Exactness at $\Ker(c)$.}
The composition $\Ker(b) \to \Ker(c) \xrightarrow{\delta} \coker(a)$ is zero: if $y \in \Ker(b)$
and $x = g'(y) \in \Ker(c)$, then in the construction of $\delta$, we may choose $y$
itself, and $b(y) = 0$, so $z = 0$, giving $\delta(x) = 0$.

Conversely, let $x \in \Ker(c)$ with $\delta(x) = 0$. With notation as in Step~2,
$[z] = 0$ means $z = a(w)$ for some $w \in A'$. Then $b(y) = f(z) = f(a(w)) = b(f'(w))$,
so $y - f'(w) \in \Ker(b)$ and $g'(y - f'(w)) = g'(y) - 0 = x$.

\medskip\noindent
\textbf{Step 6: Exactness at $\coker(a)$.}
The composition $\delta$ followed by $\coker(a) \to \coker(b)$ is zero: in the
notation above, $\delta(x) = [z]$ and the image in $\coker(b)$ is $[f(z)] = [b(y)] = 0$.

Conversely, let $[z] \in \coker(a)$ with $[f(z)] = 0$ in $\coker(b)$. Then
$f(z) = b(y)$ for some $y \in B'$. Now $c(g'(y)) = g(b(y)) = g(f(z)) = 0$, so
$g'(y) \in \Ker(c)$, and $\delta(g'(y)) = [z]$ by construction.

\medskip\noindent
\textbf{Step 7: Exactness at $\coker(b)$.}
Similar argument, left as an exercise.

\medskip\noindent
\textbf{Step 8: Additional exactness.}
If $f'$ is mono, then $\Ker(a) \to \Ker(b)$ (given by $f'|_{\Ker(a)}$) is mono.
If $g$ is epi, then $\coker(b) \to \coker(c)$ is epi by a dual argument.
\end{proof}

The snake lemma is often remembered via the following diagram, where the ``snake''
traces the path through the kernels and cokernels:

\[
\begin{tikzcd}[column sep=1.2em, row sep=2.5em]
    & \Ker(a) \arrow[r] \arrow[d, hookrightarrow] & \Ker(b) \arrow[r] \arrow[d, hookrightarrow]
    & \Ker(c) \arrow[d, hookrightarrow] \arrow[dll, "\delta"', dashed,
    rounded corners, to path={
        -- ([xshift=2em]\tikztostart.east)
        -- ([xshift=2em, yshift=-2em]\tikztostart.east)
        -- ([xshift=-2em, yshift=-2em]\tikztotarget.west)
        -- ([xshift=-2em]\tikztotarget.west)
        -- (\tikztotarget.west)}] \\
    & A' \arrow[r, "f'"] \arrow[d, "a"] & B' \arrow[r, "g'"] \arrow[d, "b"]
    & C' \arrow[r] \arrow[d, "c"] & 0 \\
    0 \arrow[r] & A \arrow[r, "f"'] \arrow[d, twoheadrightarrow]
    & B \arrow[r, "g"'] \arrow[d, twoheadrightarrow] & C \arrow[d, twoheadrightarrow] \\
    & \coker(a) \arrow[r] & \coker(b) \arrow[r] & \coker(c)
\end{tikzcd}
\]


\section{The Five Lemma}\label{sec:five-lemma}

\begin{theorem}[Five lemma]\label{thm:five-lemma}
\index{five lemma}
Consider a commutative diagram with exact rows in an abelian category:
\[
\begin{tikzcd}
    A_1 \arrow[r, "f_1"] \arrow[d, "\alpha_1"'] & A_2 \arrow[r, "f_2"] \arrow[d, "\alpha_2"']
    & A_3 \arrow[r, "f_3"] \arrow[d, "\alpha_3"'] & A_4 \arrow[r, "f_4"] \arrow[d, "\alpha_4"']
    & A_5 \arrow[d, "\alpha_5"'] \\
    B_1 \arrow[r, "g_1"'] & B_2 \arrow[r, "g_2"'] & B_3 \arrow[r, "g_3"']
    & B_4 \arrow[r, "g_4"'] & B_5
\end{tikzcd}
\]
\begin{enumerate}
    \item If $\alpha_1$ is an epimorphism and $\alpha_2, \alpha_4$ are monomorphisms,
    then $\alpha_3$ is a monomorphism.
    \item If $\alpha_5$ is a monomorphism and $\alpha_2, \alpha_4$ are epimorphisms,
    then $\alpha_3$ is an epimorphism.
    \item If $\alpha_1$ is an epimorphism, $\alpha_5$ is a monomorphism, and
    $\alpha_2, \alpha_4$ are isomorphisms, then $\alpha_3$ is an isomorphism.
\end{enumerate}
\end{theorem}

\begin{proof}
We use element chasing (justified by the Freyd--Mitchell embedding).

\medskip\noindent
\textbf{Part (1):} Suppose $\alpha_3(a_3) = 0$ for some $a_3 \in A_3$. Then
$\alpha_4(f_3(a_3)) = g_3(\alpha_3(a_3)) = 0$. Since $\alpha_4$ is mono,
$f_3(a_3) = 0$. By exactness at $A_3$, there exists $a_2 \in A_2$ with $f_2(a_2) = a_3$.
Now $g_2(\alpha_2(a_2)) = \alpha_3(f_2(a_2)) = \alpha_3(a_3) = 0$. By exactness at $B_2$,
there exists $b_1 \in B_1$ with $g_1(b_1) = \alpha_2(a_2)$. Since $\alpha_1$ is epi,
$b_1 = \alpha_1(a_1)$ for some $a_1 \in A_1$. Then
\[
    \alpha_2(f_1(a_1)) = g_1(\alpha_1(a_1)) = g_1(b_1) = \alpha_2(a_2).
\]
Since $\alpha_2$ is mono, $f_1(a_1) = a_2$, so
$a_3 = f_2(a_2) = f_2(f_1(a_1)) = 0$ by exactness at $A_2$ (since $\im(f_1) = \Ker(f_2)$).
Thus $\alpha_3$ is mono.

\medskip\noindent
\textbf{Part (2):} Let $b_3 \in B_3$. Consider $g_3(b_3) \in B_4$. Since $\alpha_4$
is epi, there exists $a_4 \in A_4$ with $\alpha_4(a_4) = g_3(b_3)$. Now
$\alpha_5(f_4(a_4)) = g_4(\alpha_4(a_4)) = g_4(g_3(b_3)) = 0$ by exactness at $B_4$.
Since $\alpha_5$ is mono, $f_4(a_4) = 0$.
By exactness at $A_4$, there exists $a_3 \in A_3$ with $f_3(a_3) = a_4$.
Consider $b_3 - \alpha_3(a_3)$:
\[
    g_3(b_3 - \alpha_3(a_3)) = g_3(b_3) - g_3(\alpha_3(a_3))
    = g_3(b_3) - \alpha_4(f_3(a_3)) = g_3(b_3) - \alpha_4(a_4) = 0.
\]
By exactness at $B_3$, there exists $b_2 \in B_2$ with $g_2(b_2) = b_3 - \alpha_3(a_3)$.
Since $\alpha_2$ is epi, $b_2 = \alpha_2(a_2)$ for some $a_2 \in A_2$. Then
\[
    b_3 = \alpha_3(a_3) + g_2(\alpha_2(a_2)) = \alpha_3(a_3) + \alpha_3(f_2(a_2))
    = \alpha_3(a_3 + f_2(a_2)).
\]
Thus $\alpha_3$ is epi.

\medskip\noindent
\textbf{Part (3):} Immediate from (1) and (2).
\end{proof}

\begin{corollary}[Short five lemma]\label{cor:short-five-lemma}
\index{five lemma!short}
Given a commutative diagram with exact rows
\[
\begin{tikzcd}
    0 \arrow[r] & A' \arrow[r] \arrow[d, "\alpha"'] & B' \arrow[r] \arrow[d, "\beta"'] & C' \arrow[r] \arrow[d, "\gamma"'] & 0 \\
    0 \arrow[r] & A \arrow[r] & B \arrow[r] & C \arrow[r] & 0
\end{tikzcd}
\]
if $\alpha$ and $\gamma$ are isomorphisms, then $\beta$ is an isomorphism.
\end{corollary}

\begin{proof}
Apply the five lemma with $A_1 = B_1 = 0$ and $A_5 = B_5 = 0$, taking
$\alpha_1 = \alpha_5 = 0$, $\alpha_2 = \alpha$, $\alpha_3 = \beta$, $\alpha_4 = \gamma$.
\end{proof}


\section{Injective and Projective Objects}\label{sec:inj-proj}

\begin{definition}[Projective object]\label{def:projective}
\index{projective object}\index{object!projective}
An object $P$ in an abelian category $\mathcal{A}$ is \textbf{projective} if the
functor $\Hom_{\mathcal{A}}(P, -)$ is exact. Equivalently, for every epimorphism
$p \colon B \twoheadrightarrow C$ and every morphism $f \colon P \to C$, there exists
a morphism $\tilde{f} \colon P \to B$ with $p \circ \tilde{f} = f$:
\[
\begin{tikzcd}
    & P \arrow[d, "f"] \arrow[dl, dashed, "\tilde{f}"'] \\
    B \arrow[r, "p"', twoheadrightarrow] & C \arrow[r] & 0
\end{tikzcd}
\]
\end{definition}

\begin{definition}[Injective object]\label{def:injective}
\index{injective object}\index{object!injective}
An object $I$ in $\mathcal{A}$ is \textbf{injective} if $\Hom_{\mathcal{A}}(-, I)$ is exact.
Equivalently, for every monomorphism $i \colon A \hookrightarrow B$ and every
$f \colon A \to I$, there exists $\tilde{f} \colon B \to I$ with $\tilde{f} \circ i = f$:
\[
\begin{tikzcd}
    0 \arrow[r] & A \arrow[r, hookrightarrow, "i"] \arrow[d, "f"'] & B \arrow[dl, dashed, "\tilde{f}"] \\
    & I &
\end{tikzcd}
\]
\end{definition}

\begin{example}\label{ex:proj-inj-examples}
\begin{enumerate}
    \item\index{projective module}
    In $\Mod_R$, free modules are projective. A module is projective if and only if it
    is a direct summand of a free module.
    \item\index{injective!abelian group}
    In $\Ab$, the group $\mathbb{Q}$ is injective. More generally, $\mathbb{Q}/\mathbb{Z}$
    is injective.
    \item In a semisimple abelian category (every SES splits), every object is both
    projective and injective.
\end{enumerate}
\end{example}

\begin{definition}[Enough projectives/injectives]\label{def:enough-proj-inj}
\index{enough projectives}\index{enough injectives}
An abelian category $\mathcal{A}$ has \textbf{enough projectives} if for every object $A$,
there exists an epimorphism $P \twoheadrightarrow A$ with $P$ projective. It has
\textbf{enough injectives} if for every $A$, there exists a monomorphism $A \hookrightarrow I$
with $I$ injective.
\end{definition}

\begin{theorem}[Baer's criterion]\label{thm:baer-criterion}
\index{Baer's criterion}
Let $R$ be a ring and $I$ an $R$-module. Then $I$ is injective if and only if for every
left ideal $J \subseteq R$ and every homomorphism $f \colon J \to I$, there exists an
extension $\tilde{f} \colon R \to I$ with $\tilde{f}|_J = f$.
\end{theorem}

\begin{theorem}\label{thm:grothendieck-injectives}
\index{Grothendieck category!enough injectives}
Every Grothendieck category has enough injectives. In particular, $\Mod_R$ and
$\operatorname{Sh}(X, \Ab)$ have enough injectives.
\end{theorem}


\section{Resolutions}\label{sec:resolutions}

\begin{definition}[Projective resolution]\label{def:proj-resolution}
\index{projective resolution}\index{resolution!projective}
A \textbf{projective resolution} of an object $A$ in an abelian category $\mathcal{A}$
is an exact sequence
\[
    \cdots \longrightarrow P_2 \xrightarrow{d_2} P_1 \xrightarrow{d_1} P_0
    \xrightarrow{\varepsilon} A \longrightarrow 0
\]
where each $P_n$ is projective. We write this as $P_\bullet \xrightarrow{\varepsilon} A$.
\end{definition}

\begin{definition}[Injective resolution]\label{def:inj-resolution}
\index{injective resolution}\index{resolution!injective}
An \textbf{injective resolution} of $A$ is an exact sequence
\[
    0 \longrightarrow A \xrightarrow{\eta} I^0 \xrightarrow{d^0} I^1
    \xrightarrow{d^1} I^2 \longrightarrow \cdots
\]
where each $I^n$ is injective. We write this as $A \xrightarrow{\eta} I^\bullet$.
\end{definition}

\begin{proposition}\label{prop:resolution-existence}
If $\mathcal{A}$ has enough projectives, then every object admits a projective resolution.
If $\mathcal{A}$ has enough injectives, then every object admits an injective resolution.
\end{proposition}

\begin{proof}
We construct a projective resolution inductively. Choose an epi $P_0 \twoheadrightarrow A$
with $P_0$ projective and let $K_0 = \Ker(P_0 \to A)$. Then choose $P_1 \twoheadrightarrow K_0$
with $P_1$ projective, set $d_1 \colon P_1 \twoheadrightarrow K_0 \hookrightarrow P_0$,
let $K_1 = \Ker(P_1 \to K_0)$, and continue.
\end{proof}

\begin{theorem}[Comparison theorem]\label{thm:comparison}
\index{comparison theorem!for resolutions}
Let $P_\bullet \xrightarrow{\varepsilon} A$ be a projective resolution and
$Q_\bullet \xrightarrow{\varepsilon'} B$ be any resolution (not necessarily projective).
Then any morphism $f \colon A \to B$ lifts to a chain map
$f_\bullet \colon P_\bullet \to Q_\bullet$:
\[
\begin{tikzcd}
    \cdots \arrow[r] & P_2 \arrow[r, "d_2"] \arrow[d, dashed, "f_2"'] & P_1 \arrow[r, "d_1"] \arrow[d, dashed, "f_1"'] & P_0 \arrow[r, "\varepsilon"] \arrow[d, dashed, "f_0"'] & A \arrow[r] \arrow[d, "f"] & 0 \\
    \cdots \arrow[r] & Q_2 \arrow[r, "d_2'"'] & Q_1 \arrow[r, "d_1'"'] & Q_0 \arrow[r, "\varepsilon'"'] & B \arrow[r] & 0
\end{tikzcd}
\]
Moreover, any two such lifts are chain homotopic.
\end{theorem}

\begin{proof}
\textit{Existence.} We construct $f_n$ inductively. The composite $f \circ \varepsilon \colon P_0 \to B$
factors through $\varepsilon' \colon Q_0 \to B$ since $P_0$ is projective and $\varepsilon'$
is epi. This gives $f_0$. For the inductive step, given $f_{n-1}$, we need $f_n$ making the
square commute. The morphism $f_{n-1} \circ d_n \colon P_n \to Q_{n-1}$ satisfies
$d'_{n-1} \circ f_{n-1} \circ d_n = f_{n-2} \circ d_{n-1} \circ d_n = 0$. Thus
$f_{n-1} \circ d_n$ factors through $\Ker(d'_{n-1}) = \im(d'_n)$, and projectivity of $P_n$
gives $f_n$.

\textit{Uniqueness up to homotopy.} Suppose $f_\bullet, g_\bullet$ are two lifts.
We construct a chain homotopy $s_n \colon P_n \to Q_{n+1}$ with
$f_n - g_n = d'_{n+1} s_n + s_{n-1} d_n$ inductively, using projectivity at each step.
\end{proof}


\section{Exercises}\label{sec:abelian-exercises}

\begin{exercise}\label{exo:preadditive-zero}
\index{zero object!in preadditive category}
Let $\mathcal{A}$ be a preadditive category. Show that $\mathcal{A}$ has a zero object if and
only if it has an object $Z$ such that $\id_Z = 0_{Z,Z}$.
\end{exercise}

\begin{exercise}\label{exo:additive-functor-zero}
Show that an additive functor preserves zero objects, biproducts, kernels, and cokernels
(when they exist).
\end{exercise}

\begin{exercise}\label{exo:idempotent-split}
\index{idempotent}\index{Karoubi envelope}
Let $\mathcal{A}$ be an additive category and $e \colon A \to A$ an idempotent ($e^2 = e$).
We say $e$ \textbf{splits} if there exist morphisms $p \colon A \to B$ and $i \colon B \to A$
with $p \circ i = \id_B$ and $i \circ p = e$.
\begin{enumerate}
    \item Show that in an abelian category, every idempotent splits.
    \item Show that $B = \Ker(\id_A - e)$ works.
\end{enumerate}
\end{exercise}

\begin{exercise}\label{exo:nine-lemma}
\index{nine lemma}\index{$3 \times 3$ lemma}
(The $3 \times 3$ lemma.) Consider a commutative diagram with exact columns and
the middle and bottom rows exact:
\[
\begin{tikzcd}
    & 0 \arrow[d] & 0 \arrow[d] & 0 \arrow[d] & \\
    0 \arrow[r] & A' \arrow[r] \arrow[d] & B' \arrow[r] \arrow[d] & C' \arrow[r] \arrow[d] & 0 \\
    0 \arrow[r] & A \arrow[r] \arrow[d] & B \arrow[r] \arrow[d] & C \arrow[r] \arrow[d] & 0 \\
    0 \arrow[r] & A'' \arrow[r] \arrow[d] & B'' \arrow[r] \arrow[d] & C'' \arrow[r] \arrow[d] & 0 \\
    & 0 & 0 & 0 &
\end{tikzcd}
\]
Prove that the top row is also exact.
\end{exercise}

\begin{exercise}\label{exo:horseshoe}
\index{horseshoe lemma}
(Horseshoe lemma.) Let $0 \to A' \to A \to A'' \to 0$ be a short exact sequence in an
abelian category with enough projectives. Given projective resolutions
$P'_\bullet \to A'$ and $P''_\bullet \to A''$, show there exists a projective resolution
$P_\bullet \to A$ with $P_n = P'_n \oplus P''_n$ fitting into a short exact sequence
of chain complexes $0 \to P'_\bullet \to P_\bullet \to P''_\bullet \to 0$.
\end{exercise}

\begin{exercise}\label{exo:injective-divisible}
\index{divisible group}\index{injective!abelian group}
Show that an abelian group $G$ is injective in $\Ab$ if and only if it is divisible
(for every $g \in G$ and $n \geq 1$, there exists $h \in G$ with $nh = g$).
\end{exercise}

\begin{exercise}\label{exo:projective-characterization}
\index{projective module!characterization}
Let $R$ be a ring and $P$ an $R$-module. Prove the equivalence:
\begin{enumerate}
    \item $P$ is projective.
    \item Every short exact sequence $0 \to A \to B \to P \to 0$ splits.
    \item $P$ is a direct summand of a free module.
\end{enumerate}
\end{exercise}

\begin{exercise}\label{exo:abelian-pullback}
Let $\mathcal{A}$ be an abelian category. Show that the pullback of an epimorphism is
an epimorphism, and the pushout of a monomorphism is a monomorphism.
\end{exercise}


%% ========================================================================
\chapter{Derived Categories --- Introduction}\label{ch:derived}
\minitoc

\section{Chain Complexes and Their Morphisms}\label{sec:chain-complexes}

\begin{definition}[Chain complex]\label{def:chain-complex}
\index{chain complex}\index{complex!chain}
Let $\mathcal{A}$ be an abelian category. A \textbf{chain complex} (or simply \textbf{complex})
in $\mathcal{A}$ is a sequence of objects and morphisms
\[
    A^\bullet \colon \quad \cdots \longrightarrow A^{n-1} \xrightarrow{d^{n-1}} A^n
    \xrightarrow{d^n} A^{n+1} \longrightarrow \cdots
\]
such that $d^n \circ d^{n-1} = 0$ for all $n \in \mathbb{Z}$. The morphisms $d^n$ are
called \textbf{differentials}\index{differential}. We use cohomological (superscript)
indexing throughout.
\end{definition}

\begin{definition}[Cohomology]\label{def:cohomology-complex}
\index{cohomology!of a complex}
The $n$-th \textbf{cohomology} of a complex $A^\bullet$ is
\[
    H^n(A^\bullet) = \Ker(d^n) / \im(d^{n-1}) = \Ker(d^n \colon A^n \to A^{n+1}) /
    \im(d^{n-1} \colon A^{n-1} \to A^n).
\]
A complex is \textbf{acyclic}\index{acyclic complex} (or \textbf{exact}) if $H^n(A^\bullet) = 0$
for all $n$.
\end{definition}

\begin{definition}[Category of complexes]\label{def:complex-category}
\index{Ch(A)@$\operatorname{Ch}(\mathcal{A})$}
\index{category!of complexes}
The \textbf{category of chain complexes} $\operatorname{Ch}(\mathcal{A})$ has:
\begin{itemize}
    \item Objects: chain complexes in $\mathcal{A}$.
    \item Morphisms: a \textbf{chain map}\index{chain map} $f \colon A^\bullet \to B^\bullet$
    is a collection of morphisms $f^n \colon A^n \to B^n$ commuting with differentials:
    $d_B^n \circ f^n = f^{n+1} \circ d_A^n$ for all $n$.
\end{itemize}
We also define the bounded variants:
\begin{align*}
    \operatorname{Ch}^+(\mathcal{A}) &= \{ A^\bullet \mid A^n = 0 \text{ for } n \ll 0 \}
    &&\text{(bounded below)}, \\
    \operatorname{Ch}^-(\mathcal{A}) &= \{ A^\bullet \mid A^n = 0 \text{ for } n \gg 0 \}
    &&\text{(bounded above)}, \\
    \operatorname{Ch}^b(\mathcal{A}) &= \operatorname{Ch}^+(\mathcal{A}) \cap \operatorname{Ch}^-(\mathcal{A})
    &&\text{(bounded)}.
\end{align*}
\end{definition}

\begin{remark}\label{rem:ch-abelian}
\index{Ch(A)@$\operatorname{Ch}(\mathcal{A})$!abelian structure}
The category $\operatorname{Ch}(\mathcal{A})$ is abelian when $\mathcal{A}$ is abelian,
with kernels, cokernels, and exact sequences computed degreewise.
\end{remark}

\begin{definition}[Shift functor]\label{def:shift-functor}
\index{shift functor}\index{translation functor}
The \textbf{shift} (or \textbf{translation}) functor $[1] \colon \operatorname{Ch}(\mathcal{A})
\to \operatorname{Ch}(\mathcal{A})$ is defined by
\[
    (A[1])^n = A^{n+1}, \qquad d_{A[1]}^n = -d_A^{n+1}.
\]
More generally, $A[k]^n = A^{n+k}$ with $d_{A[k]}^n = (-1)^k d_A^{n+k}$.
For a chain map $f \colon A^\bullet \to B^\bullet$, we set $f[1]^n = f^{n+1}$.
\end{definition}

\begin{definition}[Mapping cone]\label{def:mapping-cone}
\index{mapping cone}\index{cone!of a morphism}
Let $f \colon A^\bullet \to B^\bullet$ be a chain map. The \textbf{mapping cone}
$\operatorname{Cone}(f)$ is the complex defined by
\[
    \operatorname{Cone}(f)^n = A^{n+1} \oplus B^n, \qquad
    d_{\operatorname{Cone}(f)}^n = \begin{pmatrix} -d_A^{n+1} & 0 \\ f^{n+1} & d_B^n \end{pmatrix}.
\]
One verifies $d^2 = 0$ using $d_B \circ f = f \circ d_A$.
There is a natural short exact sequence of complexes:
\[
    0 \longrightarrow B^\bullet \xrightarrow{\;\iota\;} \operatorname{Cone}(f)
    \xrightarrow{\;\pi\;} A[1]^\bullet \longrightarrow 0.
\]
\end{definition}

\begin{proposition}\label{prop:cone-les}
\index{mapping cone!long exact sequence}
For any chain map $f \colon A^\bullet \to B^\bullet$, the mapping cone gives a long exact
sequence in cohomology:
\[
    \cdots \to H^n(A^\bullet) \xrightarrow{H^n(f)} H^n(B^\bullet) \to
    H^n(\operatorname{Cone}(f)) \to H^{n+1}(A^\bullet) \xrightarrow{H^{n+1}(f)} \cdots
\]
In particular, $f$ is a quasi-isomorphism if and only if $\operatorname{Cone}(f)$ is acyclic.
\end{proposition}


\section{Chain Homotopies and the Homotopy Category}\label{sec:homotopy-cat}

\begin{definition}[Chain homotopy]\label{def:chain-homotopy}
\index{chain homotopy}\index{homotopy!of chain maps}
Let $f, g \colon A^\bullet \to B^\bullet$ be chain maps. A \textbf{chain homotopy}
from $f$ to $g$ is a collection of morphisms $s^n \colon A^n \to B^{n-1}$ such that
\[
    f^n - g^n = d_B^{n-1} \circ s^n + s^{n+1} \circ d_A^n \quad \text{for all } n.
\]
\[
\begin{tikzcd}[column sep=2em]
    \cdots \arrow[r] & A^{n-1} \arrow[r, "d^{n-1}"] \arrow[d, shift left=1, "f^{n-1}"]
    \arrow[d, shift right=1, "g^{n-1}"']
    & A^n \arrow[r, "d^n"] \arrow[d, shift left=1, "f^n"]
    \arrow[d, shift right=1, "g^n"'] \arrow[ld, dashed, "s^n"']
    & A^{n+1} \arrow[r] \arrow[d, shift left=1, "f^{n+1}"]
    \arrow[d, shift right=1, "g^{n+1}"'] \arrow[ld, dashed, "s^{n+1}"']
    & \cdots \\
    \cdots \arrow[r] & B^{n-1} \arrow[r, "d^{n-1}"'] & B^n \arrow[r, "d^n"'] & B^{n+1} \arrow[r] & \cdots
\end{tikzcd}
\]
We write $f \sim g$ and say $f$ and $g$ are \textbf{homotopic}.
If $f \sim 0$, we say $f$ is \textbf{null-homotopic}\index{null-homotopic}.
\end{definition}

\begin{proposition}\label{prop:homotopy-equiv-rel}
Chain homotopy is an equivalence relation on $\Hom_{\operatorname{Ch}(\mathcal{A})}(A^\bullet, B^\bullet)$,
compatible with composition: if $f \sim g$ then $h \circ f \sim h \circ g$ and
$f \circ k \sim g \circ k$ for any compatible chain maps $h, k$.
\end{proposition}

\begin{proof}
Reflexivity: take $s = 0$. Symmetry: negate $s$. Transitivity: add homotopies.
For compatibility, if $s$ is a homotopy from $f$ to $g$, then $h \circ s$ is a homotopy
from $h \circ f$ to $h \circ g$, and $s \circ k$ is a homotopy from $f \circ k$ to $g \circ k$.
\end{proof}

\begin{proposition}\label{prop:homotopic-cohomology}
\index{chain homotopy!and cohomology}
If $f \sim g \colon A^\bullet \to B^\bullet$, then $H^n(f) = H^n(g)$ for all $n$.
\end{proposition}

\begin{proof}
For $[a] \in H^n(A^\bullet)$ with $d_A^n(a) = 0$:
$f^n(a) - g^n(a) = d_B^{n-1}(s^n(a)) + s^{n+1}(d_A^n(a)) = d_B^{n-1}(s^n(a))$,
which is a coboundary. Thus $[f^n(a)] = [g^n(a)]$.
\end{proof}

\begin{definition}[Homotopy category]\label{def:homotopy-category}
\index{homotopy category}\index{K(A)@$K(\mathcal{A})$}
The \textbf{homotopy category} $K(\mathcal{A})$ has:
\begin{itemize}
    \item Objects: chain complexes in $\mathcal{A}$ (same as $\operatorname{Ch}(\mathcal{A})$).
    \item Morphisms: $\Hom_{K(\mathcal{A})}(A^\bullet, B^\bullet) =
    \Hom_{\operatorname{Ch}(\mathcal{A})}(A^\bullet, B^\bullet) / {\sim}$,
    i.e., chain maps modulo homotopy.
\end{itemize}
The bounded variants $K^+(\mathcal{A})$, $K^-(\mathcal{A})$, $K^b(\mathcal{A})$ are
defined similarly.
\end{definition}

\begin{remark}\label{rem:homotopy-not-abelian}
\index{homotopy category!not abelian}
The homotopy category $K(\mathcal{A})$ is additive but \emph{not} abelian in general.
It does, however, carry a \textbf{triangulated structure}\index{triangulated category}
(see Section~\ref{sec:triangulated}).
\end{remark}

\begin{definition}[Homotopy equivalence]\label{def:homotopy-equivalence}
\index{homotopy equivalence}
A chain map $f \colon A^\bullet \to B^\bullet$ is a \textbf{homotopy equivalence} if
it becomes an isomorphism in $K(\mathcal{A})$, i.e., there exists $g \colon B^\bullet \to A^\bullet$
with $g \circ f \sim \id_A$ and $f \circ g \sim \id_B$.
\end{definition}


\section{Quasi-Isomorphisms and the Derived Category}\label{sec:derived-cat}

\begin{definition}[Quasi-isomorphism]\label{def:quasi-iso}
\index{quasi-isomorphism}
A chain map $f \colon A^\bullet \to B^\bullet$ is a \textbf{quasi-isomorphism} if the induced
maps $H^n(f) \colon H^n(A^\bullet) \to H^n(B^\bullet)$ are isomorphisms for all $n \in \mathbb{Z}$.
We denote the class of quasi-isomorphisms by $\operatorname{qis}$.
\end{definition}

\begin{remark}\label{rem:homotopy-vs-quasi}
Every homotopy equivalence is a quasi-isomorphism (by Proposition~\ref{prop:homotopic-cohomology}),
but the converse is false in general. The derived category is obtained by formally
inverting all quasi-isomorphisms.
\end{remark}

\begin{example}\label{ex:quasi-iso-not-homotopy}
\index{quasi-isomorphism!vs.\ homotopy equivalence}
In $\Ab$, consider the complexes $A^\bullet \colon 0 \to \mathbb{Z} \xrightarrow{2} \mathbb{Z} \to 0$
(in degrees 0 and 1) and $B^\bullet \colon 0 \to \mathbb{Z}/2\mathbb{Z} \to 0$ (in degree 1).
The natural projection $A^\bullet \to B^\bullet$ is a quasi-isomorphism but not a homotopy
equivalence (since $\Hom(\mathbb{Z}/2\mathbb{Z}, \mathbb{Z}) = 0$, there is no map backward
in degree 1).
\end{example}

\subsection{Localization of categories}\label{subsec:localization}

\begin{definition}[Localization]\label{def:localization-cat}
\index{localization!of categories}\index{category!localization}
Let $\mathcal{C}$ be a category and $S$ a class of morphisms in $\mathcal{C}$.
The \textbf{localization} $\mathcal{C}[S^{-1}]$ is a category equipped with a functor
$Q \colon \mathcal{C} \to \mathcal{C}[S^{-1}]$ such that:
\begin{enumerate}
    \item $Q(s)$ is an isomorphism for every $s \in S$.
    \item $Q$ is universal with this property: for any functor $F \colon \mathcal{C} \to \mathcal{D}$
    sending $S$ to isomorphisms, there exists a unique $\bar{F} \colon \mathcal{C}[S^{-1}] \to \mathcal{D}$
    with $\bar{F} \circ Q = F$.
\end{enumerate}
\[
\begin{tikzcd}
    \mathcal{C} \arrow[r, "Q"] \arrow[rd, "F"'] & \mathcal{C}[S^{-1}] \arrow[d, dashed, "\exists!\bar{F}"] \\
    & \mathcal{D}
\end{tikzcd}
\]
\end{definition}

\begin{remark}\label{rem:localization-construction}
\index{localization!Gabriel--Zisman}
The localization always exists (Gabriel--Zisman): objects of $\mathcal{C}[S^{-1}]$ are those
of $\mathcal{C}$, and morphisms are equivalence classes of ``zigzag'' chains
\[
    A = C_0 \xleftarrow{s_1} C_1 \xrightarrow{f_1} C_2 \xleftarrow{s_2}
    C_3 \xrightarrow{f_2} \cdots \to B,
\]
where the backward arrows $s_i$ belong to $S$, subject to appropriate relations.
In general, $\Hom$ sets may fail to be sets (set-theoretic issues). However,
when $S$ admits a \textbf{calculus of fractions}\index{calculus of fractions},
morphisms simplify to ``roofs'' and the issues become manageable.
\end{remark}

\begin{definition}[Right/left roof]\label{def:roof}
\index{roof}\index{calculus of fractions}
When $S$ admits a \textbf{calculus of right fractions}, every morphism in $\mathcal{C}[S^{-1}]$
from $A$ to $B$ can be represented by a \textbf{right roof}:
\[
\begin{tikzcd}
    & C \arrow[dl, "s"'] \arrow[dr, "f"] & \\
    A & & B
\end{tikzcd}
\]
where $s \in S$. This represents the ``fraction'' $f \circ s^{-1}$.
Two roofs $(C, s, f)$ and $(C', s', f')$ represent the same morphism if and only if
they can be ``completed'' to a common refinement.
\end{definition}

\subsection{The derived category}\label{subsec:derived-cat-def}

\begin{definition}[Derived category]\label{def:derived-cat}
\index{derived category}\index{D(A)@$D(\mathcal{A})$}
Let $\mathcal{A}$ be an abelian category. The \textbf{derived category} $D(\mathcal{A})$
is the localization of $K(\mathcal{A})$ (the homotopy category) at the class of
quasi-isomorphisms:
\[
    D(\mathcal{A}) = K(\mathcal{A})[\operatorname{qis}^{-1}].
\]
The bounded variants are:
\begin{align*}
    D^+(\mathcal{A}) &= K^+(\mathcal{A})[\operatorname{qis}^{-1}], \\
    D^-(\mathcal{A}) &= K^-(\mathcal{A})[\operatorname{qis}^{-1}], \\
    D^b(\mathcal{A}) &= K^b(\mathcal{A})[\operatorname{qis}^{-1}].
\end{align*}
\end{definition}

\begin{remark}\label{rem:why-two-steps}
\index{derived category!construction}
The construction proceeds in two steps: $\operatorname{Ch}(\mathcal{A}) \to K(\mathcal{A}) \to D(\mathcal{A})$.
The first step (modding out homotopies) ensures that quasi-isomorphisms satisfy the Ore
conditions (calculus of fractions), making the second localization well-behaved.
If one tried to localize $\operatorname{Ch}(\mathcal{A})$ directly at quasi-isomorphisms,
the calculus of fractions would fail.
\end{remark}

\begin{theorem}\label{thm:derived-cat-roofs}
\index{derived category!morphisms as roofs}
The class of quasi-isomorphisms in $K(\mathcal{A})$ admits both a left and a right calculus
of fractions. Therefore, morphisms in $D(\mathcal{A})$ from $A^\bullet$ to $B^\bullet$ are
represented by roofs
\[
\begin{tikzcd}
    & C^\bullet \arrow[dl, "s"', "\sim"] \arrow[dr, "f"] & \\
    A^\bullet & & B^\bullet
\end{tikzcd}
\]
where $s$ is a quasi-isomorphism in $K(\mathcal{A})$ and $f$ is any morphism in $K(\mathcal{A})$.
\end{theorem}

\begin{proposition}\label{prop:derived-cat-embedding}
\index{derived category!embedding of $\mathcal{A}$}
There is a fully faithful functor $\mathcal{A} \hookrightarrow D(\mathcal{A})$ sending
an object $A$ to the complex concentrated in degree zero:
\[
    A \mapsto (\cdots \to 0 \to A \to 0 \to \cdots).
\]
Under this embedding, $\Hom_{D(\mathcal{A})}(A, B) \cong \Hom_{\mathcal{A}}(A, B)$
for objects $A, B \in \mathcal{A}$.
\end{proposition}

\begin{theorem}\label{thm:injective-resolution-equiv}
\index{derived category!via injective resolutions}
Let $\mathcal{A}$ be an abelian category with enough injectives, and let
$K^+(\mathcal{I})$ be the full subcategory of $K^+(\mathcal{A})$ consisting of
complexes of injective objects. Then the natural functor
\[
    K^+(\mathcal{I}) \xrightarrow{\;\sim\;} D^+(\mathcal{A})
\]
is an equivalence of categories. Dually, if $\mathcal{A}$ has enough projectives,
$K^-(\mathcal{P}) \xrightarrow{\sim} D^-(\mathcal{A})$.
\end{theorem}

\begin{proof}[Proof sketch]
The functor is the composition $K^+(\mathcal{I}) \hookrightarrow K^+(\mathcal{A})
\xrightarrow{Q} D^+(\mathcal{A})$.

\textit{Essential surjectivity:} Every bounded below complex $A^\bullet$ admits a
quasi-isomorphism $A^\bullet \xrightarrow{\sim} I^\bullet$ to a complex of injectives
(by inductively resolving, using an injective Cartan--Eilenberg resolution).

\textit{Fully faithfulness:} For complexes of injectives $I^\bullet$, $J^\bullet$,
a quasi-isomorphism $s \colon I^\bullet \xrightarrow{\sim} J^\bullet$ is already a
homotopy equivalence (a standard argument using injectivity at each degree).
Thus roofs reduce to ordinary morphisms in $K^+(\mathcal{I})$.
\end{proof}


\section{Triangulated Structure}\label{sec:triangulated}

\begin{definition}[Triangulated category]\label{def:triangulated-cat}
\index{triangulated category}\index{distinguished triangle}
A \textbf{triangulated category} is an additive category $\mathcal{T}$ equipped with:
\begin{enumerate}
    \item An automorphism $[1] \colon \mathcal{T} \to \mathcal{T}$ (the \textbf{shift functor}).
    \item A class of \textbf{distinguished triangles} (or \textbf{exact triangles})
    \[
        A \xrightarrow{f} B \xrightarrow{g} C \xrightarrow{h} A[1],
    \]
\end{enumerate}
satisfying axioms (TR1)--(TR4):
\begin{enumerate}
    \item[(TR1)] (a) For every $A$, the triangle $A \xrightarrow{\id} A \to 0 \to A[1]$ is
    distinguished. (b) Every morphism $f \colon A \to B$ can be completed to a distinguished
    triangle $A \xrightarrow{f} B \to C \to A[1]$. (c) A triangle isomorphic to a
    distinguished triangle is distinguished.
    \item[(TR2)] (\textbf{Rotation})\index{rotation of triangles}
    $A \xrightarrow{f} B \xrightarrow{g} C \xrightarrow{h} A[1]$ is distinguished
    if and only if $B \xrightarrow{g} C \xrightarrow{h} A[1] \xrightarrow{-f[1]} B[1]$ is.
    \item[(TR3)] (\textbf{Morphism of triangles}) Given distinguished triangles and morphisms
    $\alpha, \beta$ making the left square commute, there exists (not necessarily unique)
    $\gamma$ completing the morphism of triangles:
    \[
    \begin{tikzcd}
        A \arrow[r, "f"] \arrow[d, "\alpha"'] & B \arrow[r, "g"] \arrow[d, "\beta"']
        & C \arrow[r, "h"] \arrow[d, dashed, "\gamma"'] & A{[1]} \arrow[d, "\alpha{[1]}"'] \\
        A' \arrow[r, "f'"'] & B' \arrow[r, "g'"'] & C' \arrow[r, "h'"'] & A'{[1]}
    \end{tikzcd}
    \]
    \item[(TR4)] (\textbf{Octahedron axiom})\index{octahedron axiom}
    Given composable morphisms $A \xrightarrow{f} B \xrightarrow{g} C$ and distinguished
    triangles completing $f$, $g$, and $g \circ f$, there exists a distinguished triangle
    relating the three cones, making all faces commute.
\end{enumerate}
\end{definition}

\begin{theorem}\label{thm:K-A-triangulated}
\index{homotopy category!triangulated structure}
The homotopy category $K(\mathcal{A})$ is triangulated with:
\begin{itemize}
    \item Shift: $A^\bullet \mapsto A[1]^\bullet$ (Definition~\ref{def:shift-functor}).
    \item Distinguished triangles: those isomorphic in $K(\mathcal{A})$ to
    \[
        A^\bullet \xrightarrow{f} B^\bullet \xrightarrow{\iota} \operatorname{Cone}(f)
        \xrightarrow{\pi} A[1]^\bullet
    \]
    for some chain map $f$.
\end{itemize}
\end{theorem}

\begin{theorem}\label{thm:D-A-triangulated}
\index{derived category!triangulated structure}
The derived category $D(\mathcal{A})$ inherits a triangulated structure from $K(\mathcal{A})$
via the localization functor $Q \colon K(\mathcal{A}) \to D(\mathcal{A})$: a triangle in
$D(\mathcal{A})$ is distinguished if it is isomorphic to the image of a distinguished
triangle in $K(\mathcal{A})$.
\end{theorem}

\begin{proposition}\label{prop:ses-to-triangle}
\index{short exact sequence!and distinguished triangles}
A short exact sequence of complexes $0 \to A^\bullet \xrightarrow{f} B^\bullet
\xrightarrow{g} C^\bullet \to 0$ gives rise to a distinguished triangle
\[
    A^\bullet \xrightarrow{f} B^\bullet \xrightarrow{g} C^\bullet
    \xrightarrow{\delta} A[1]^\bullet
\]
in $D(\mathcal{A})$, where $\delta$ is the connecting morphism. In particular, the
long exact sequence in cohomology associated to a short exact sequence of complexes
is a consequence of the triangulated structure.
\end{proposition}

\begin{proposition}\label{prop:cohomological-functor}
\index{cohomological functor}
Let $\mathcal{T}$ be a triangulated category and $\mathcal{A}$ an abelian category.
An additive functor $H \colon \mathcal{T} \to \mathcal{A}$ is \textbf{cohomological} if
for every distinguished triangle $A \to B \to C \to A[1]$, the sequence
\[
    H(A) \longrightarrow H(B) \longrightarrow H(C)
\]
is exact. In this case, writing $H^n = H \circ [n]$, every distinguished triangle
gives a long exact sequence
\[
    \cdots \to H^n(A) \to H^n(B) \to H^n(C) \to H^{n+1}(A) \to \cdots
\]
\end{proposition}

\begin{example}\label{ex:cohomological-functors}
\begin{enumerate}
    \item The cohomology functor $H^0 \colon D(\mathcal{A}) \to \mathcal{A}$ is cohomological.
    \item For any object $X$, $\Hom_{D(\mathcal{A})}(X, -)$ is cohomological with values
    in $\Ab$.
\end{enumerate}
\end{example}


\section{Derived Functors}\label{sec:derived-functors}

\begin{definition}[Right derived functor]\label{def:right-derived-functor}
\index{derived functor!right}\index{RF@$RF$}
Let $F \colon \mathcal{A} \to \mathcal{B}$ be a left exact functor between abelian
categories, with $\mathcal{A}$ having enough injectives. The \textbf{right derived functor}
$RF \colon D^+(\mathcal{A}) \to D^+(\mathcal{B})$ is defined as follows. For
$A^\bullet \in D^+(\mathcal{A})$:
\begin{enumerate}
    \item Choose a quasi-isomorphism $A^\bullet \xrightarrow{\sim} I^\bullet$ with
    $I^\bullet \in K^+(\mathcal{I})$ (injective resolution).
    \item Define $RF(A^\bullet) = F(I^\bullet)$ (apply $F$ degreewise).
\end{enumerate}
By Theorem~\ref{thm:injective-resolution-equiv}, this is well-defined (independent
of the choice of resolution, up to canonical isomorphism in $D^+(\mathcal{B})$).

The \textbf{classical right derived functors}\index{derived functor!classical} are
\[
    R^n F(A) = H^n(RF(A)) \quad \text{for } A \in \mathcal{A}, \; n \geq 0.
\]
Note $R^0 F \cong F$ since $F$ is left exact.
\end{definition}

\begin{definition}[Left derived functor]\label{def:left-derived-functor}
\index{derived functor!left}\index{LF@$LF$}
Let $G \colon \mathcal{A} \to \mathcal{B}$ be a right exact functor with $\mathcal{A}$
having enough projectives. The \textbf{left derived functor}
$LG \colon D^-(\mathcal{A}) \to D^-(\mathcal{B})$ is defined by:
\begin{enumerate}
    \item Choose a quasi-isomorphism $P^\bullet \xrightarrow{\sim} A^\bullet$ with
    $P^\bullet \in K^-(\mathcal{P})$.
    \item Define $LG(A^\bullet) = G(P^\bullet)$.
\end{enumerate}
The classical left derived functors are $L_n G(A) = H^{-n}(LG(A))$ for $n \geq 0$,
with $L_0 G \cong G$.
\end{definition}

\begin{remark}\label{rem:derived-universal}
\index{derived functor!universal property}
The right derived functor $RF$ is characterized by the universal property: it is the
``closest approximation'' to $F$ that is an exact functor (i.e., preserves distinguished
triangles) on the derived category. More precisely, $RF$ is the right Kan extension
of $Q_{\mathcal{B}} \circ F$ along $Q_{\mathcal{A}}$.
\end{remark}

\begin{theorem}[Derived functor of a composition]\label{thm:grothendieck-spectral}
\index{Grothendieck spectral sequence}\index{spectral sequence!Grothendieck}
Let $F \colon \mathcal{A} \to \mathcal{B}$ and $G \colon \mathcal{B} \to \mathcal{C}$
be left exact functors between abelian categories with enough injectives. If $F$ sends
injective objects of $\mathcal{A}$ to $G$-acyclic objects (i.e., $R^n G(F(I)) = 0$ for
$n > 0$ and $I$ injective), then
\[
    R(G \circ F) \cong RG \circ RF.
\]
At the level of classical derived functors, there is a \textbf{Grothendieck spectral sequence}:
\[
    E_2^{p,q} = R^p G(R^q F(A)) \Longrightarrow R^{p+q}(G \circ F)(A).
\]
\end{theorem}


\section{Ext and Tor}\label{sec:ext-tor}

\begin{definition}[Ext groups]\label{def:ext}
\index{Ext}\index{Ext@$\operatorname{Ext}^n$}
Let $\mathcal{A}$ be an abelian category with enough injectives (or enough projectives).
For objects $A, B \in \mathcal{A}$, the \textbf{Ext groups} are
\[
    \operatorname{Ext}^n_{\mathcal{A}}(A, B) = R^n \Hom_{\mathcal{A}}(A, -)(B)
    = H^n(R\Hom(A, B)).
\]
Equivalently, if $I^\bullet$ is an injective resolution of $B$:
\[
    \operatorname{Ext}^n(A, B) \cong H^n(\Hom(A, I^\bullet)).
\]
If $\mathcal{A}$ has enough projectives, using a projective resolution $P_\bullet \to A$:
\[
    \operatorname{Ext}^n(A, B) \cong H^n(\Hom(P_\bullet, B)).
\]
\end{definition}

\begin{proposition}\label{prop:ext-properties}
\index{Ext!properties}
For a module category $\Mod_R$:
\begin{enumerate}
    \item $\operatorname{Ext}^0_R(M, N) \cong \Hom_R(M, N)$.
    \item $\operatorname{Ext}^n_R(M, N) = 0$ for all $n > 0$ if $M$ is projective or $N$
    is injective.
    \item The two definitions (via injective or projective resolutions) agree: this is
    the \textbf{balancing of Ext}\index{Ext!balancing}.
    \item $\operatorname{Ext}^n_R(M, N)$ classifies $n$-fold extensions of $M$ by $N$
    (Yoneda's interpretation)\index{Yoneda Ext}.
\end{enumerate}
\end{proposition}

\begin{example}\label{ex:ext-computations}
\begin{enumerate}
    \item\index{Ext!of abelian groups}
    In $\Ab$: $\operatorname{Ext}^1(\mathbb{Z}/m\mathbb{Z}, \mathbb{Z}/n\mathbb{Z})
    \cong \mathbb{Z}/\gcd(m,n)\mathbb{Z}$.
    Indeed, the projective resolution $0 \to \mathbb{Z} \xrightarrow{m} \mathbb{Z} \to \mathbb{Z}/m\mathbb{Z} \to 0$
    gives $\operatorname{Ext}^1 \cong \mathbb{Z}/n\mathbb{Z} \big/ m(\mathbb{Z}/n\mathbb{Z})
    \cong \mathbb{Z}/\gcd(m,n)\mathbb{Z}$.
    \item $\operatorname{Ext}^n_{\mathbb{Z}}(A, B) = 0$ for all $n \geq 2$ and all
    abelian groups $A, B$ (since $\mathbb{Z}$ has global dimension 1).
\end{enumerate}
\end{example}

\begin{definition}[Tor groups]\label{def:tor}
\index{Tor}\index{Tor@$\operatorname{Tor}_n$}
For a ring $R$ and modules $M \in \Mod_{R^{\op}}$, $N \in \Mod_R$, the \textbf{Tor groups} are
\[
    \operatorname{Tor}_n^R(M, N) = L_n(M \otimes_R -)(N) = H^{-n}(M \otimes_R^L N).
\]
If $P_\bullet \to N$ is a projective resolution:
\[
    \operatorname{Tor}_n^R(M, N) \cong H_n(M \otimes_R P_\bullet).
\]
\end{definition}

\begin{proposition}\label{prop:tor-properties}
\index{Tor!properties}
\begin{enumerate}
    \item $\operatorname{Tor}_0^R(M, N) \cong M \otimes_R N$.
    \item $\operatorname{Tor}_n^R(M, N) = 0$ for all $n > 0$ if $M$ or $N$ is flat.
    \item Tor is balanced: $L_n(M \otimes_R -)(N) \cong L_n(- \otimes_R N)(M)$.
    \item $M$ is flat if and only if $\operatorname{Tor}_1^R(M, N) = 0$ for all $N$.
\end{enumerate}
\end{proposition}

\begin{example}\label{ex:tor-computations}
\begin{enumerate}
    \item\index{Tor!of abelian groups}
    $\operatorname{Tor}_1^{\mathbb{Z}}(\mathbb{Z}/m\mathbb{Z}, \mathbb{Z}/n\mathbb{Z})
    \cong \mathbb{Z}/\gcd(m,n)\mathbb{Z}$.
    Using $0 \to \mathbb{Z} \xrightarrow{m} \mathbb{Z} \to \mathbb{Z}/m\mathbb{Z} \to 0$:
    $\operatorname{Tor}_1 = \Ker(\mathbb{Z}/n\mathbb{Z} \xrightarrow{m} \mathbb{Z}/n\mathbb{Z})
    \cong \mathbb{Z}/\gcd(m,n)\mathbb{Z}$.
    \item $\operatorname{Tor}_n^{\mathbb{Z}}(A, B) = 0$ for $n \geq 2$.
\end{enumerate}
\end{example}


\section{Long Exact Sequences}\label{sec:long-exact-derived}

\begin{theorem}\label{thm:les-derived}
\index{long exact sequence!from derived functors}
Let $F \colon \mathcal{A} \to \mathcal{B}$ be a left exact functor and
$0 \to A' \to A \to A'' \to 0$ a short exact sequence in $\mathcal{A}$.
Then there is a long exact sequence:
\[
    0 \to FA' \to FA \to FA'' \xrightarrow{\delta^0} R^1 FA' \to R^1 FA \to R^1 FA''
    \xrightarrow{\delta^1} R^2 FA' \to \cdots
\]
Dually, for a right exact functor $G$ and $0 \to A' \to A \to A'' \to 0$:
\[
    \cdots \to L_2 GA'' \xrightarrow{\delta_2} L_1 GA' \to L_1 GA \to L_1 GA''
    \xrightarrow{\delta_1} GA' \to GA \to GA'' \to 0.
\]
\end{theorem}

\begin{proof}
By Proposition~\ref{prop:ses-to-triangle}, the short exact sequence gives a distinguished
triangle $A' \to A \to A'' \to A'[1]$ in $D^+(\mathcal{A})$. Applying $RF$ (which
preserves distinguished triangles) gives a distinguished triangle
\[
    RF(A') \to RF(A) \to RF(A'') \to RF(A')[1]
\]
in $D^+(\mathcal{B})$. The cohomological functor $H^0$ (Proposition~\ref{prop:cohomological-functor})
then yields the long exact sequence, since $H^n(RF(A)) = R^n F(A)$.
\end{proof}

\begin{corollary}\label{cor:ext-les}
\index{Ext!long exact sequence}
For any short exact sequence $0 \to A' \to A \to A'' \to 0$ in $\Mod_R$ and any module $M$:
\begin{enumerate}
    \item (Covariant) There is a long exact sequence
    \[
        0 \to \Hom(M, A') \to \Hom(M, A) \to \Hom(M, A'') \to
        \operatorname{Ext}^1(M, A') \to \cdots
    \]
    \item (Contravariant) There is a long exact sequence
    \[
        0 \to \Hom(A'', M) \to \Hom(A, M) \to \Hom(A', M) \to
        \operatorname{Ext}^1(A'', M) \to \cdots
    \]
\end{enumerate}
\end{corollary}

\begin{corollary}\label{cor:tor-les}
\index{Tor!long exact sequence}
For any short exact sequence $0 \to A' \to A \to A'' \to 0$ in $\Mod_R$ and any module $M$:
\[
    \cdots \to \operatorname{Tor}_1(M, A') \to \operatorname{Tor}_1(M, A) \to
    \operatorname{Tor}_1(M, A'') \to M \otimes A' \to M \otimes A \to M \otimes A'' \to 0.
\]
\end{corollary}

\begin{example}\label{ex:les-application}
\index{long exact sequence!computation example}
Consider the exact sequence $0 \to \mathbb{Z} \xrightarrow{n} \mathbb{Z} \to \mathbb{Z}/n\mathbb{Z} \to 0$
in $\Ab$. Applying $\Hom(-, A)$ for an abelian group $A$:
\[
    0 \to \Hom(\mathbb{Z}/n, A) \to \Hom(\mathbb{Z}, A) \xrightarrow{n}
    \Hom(\mathbb{Z}, A) \to \operatorname{Ext}^1(\mathbb{Z}/n, A) \to 0.
\]
This gives $\Hom(\mathbb{Z}/n, A) \cong A[n] = \{a \in A : na = 0\}$ and
$\operatorname{Ext}^1(\mathbb{Z}/n, A) \cong A/nA$.
\end{example}


\section{Derived Category in Practice}\label{sec:derived-practice}

\begin{definition}[RHom and derived tensor]\label{def:rhom-ltensor}
\index{RHom}\index{derived tensor product}\index{tensor product!derived}
For $\mathcal{A} = \Mod_R$:
\begin{enumerate}
    \item $R\Hom_R(M, N) \in D^+(\Ab)$ is the right derived functor of $\Hom_R(M, -)$.
    Its cohomology gives $H^n(R\Hom(M, N)) = \operatorname{Ext}^n_R(M, N)$.
    \item $M \otimes_R^L N \in D^-(\Ab)$ is the left derived functor of $M \otimes_R -$.
    Its cohomology gives $H^{-n}(M \otimes_R^L N) = \operatorname{Tor}_n^R(M, N)$.
\end{enumerate}
These extend to functors on derived categories:
\[
    R\Hom \colon D^-(\Mod_R)^{\op} \times D^+(\Mod_R) \to D^+(\Ab),
\]
\[
    \otimes_R^L \colon D^-(\Mod_{R^{\op}}) \times D^-(\Mod_R) \to D^-(\Ab).
\]
\end{definition}

\begin{proposition}[Derived Hom--Tensor adjunction]\label{prop:derived-adjunction}
\index{adjunction!derived Hom--Tensor}
There is a natural isomorphism in $D(\Ab)$:
\[
    R\Hom_R(M \otimes_S^L N, \; P) \cong R\Hom_S(N, \; R\Hom_R(M, P))
\]
for appropriate bimodule structures.
\end{proposition}

\begin{proposition}[Derived category of a hereditary category]\label{prop:hereditary-derived}
\index{hereditary category}\index{derived category!hereditary case}
Let $\mathcal{A}$ be a hereditary abelian category (i.e., $\operatorname{Ext}^n = 0$ for $n \geq 2$).
Then every object of $D^b(\mathcal{A})$ is isomorphic to the direct sum of its cohomology
objects (each placed in its own degree):
\[
    A^\bullet \cong \bigoplus_{n \in \mathbb{Z}} H^n(A^\bullet)[-n] \quad \text{in } D^b(\mathcal{A}).
\]
\end{proposition}

\begin{remark}\label{rem:derived-cat-philosophy}
\index{derived category!philosophy}
The philosophy of derived categories can be summarized as follows:
\begin{enumerate}
    \item Objects of $\mathcal{A}$ are replaced by complexes, which carry ``higher-order''
    information through their cohomology in all degrees.
    \item Functors that are only partially exact (left or right exact) become exact on derived
    categories, at the cost of working with complexes.
    \item The derived category provides a unified framework for cohomological computations:
    spectral sequences, universal coefficient theorems, K\"unneth formulas, etc.\
    all arise from the triangulated structure.
\end{enumerate}
\end{remark}

\begin{notation}\label{not:derived-summary}
\index{derived category!notation summary}
We summarize the key notation:
\begin{center}
\renewcommand{\arraystretch}{1.3}
\begin{tabular}{ll}
\hline
$\operatorname{Ch}(\mathcal{A})$ & Category of chain complexes \\
$K(\mathcal{A})$ & Homotopy category \\
$D(\mathcal{A})$ & Derived category \\
$D^+, D^-, D^b$ & Bounded variants \\
$RF$ & Right derived functor \\
$LG$ & Left derived functor \\
$R\Hom$ & Derived Hom \\
$\otimes^L$ & Derived tensor product \\
$\operatorname{Ext}^n(A,B)$ & $= H^n(R\Hom(A,B))$ \\
$\operatorname{Tor}_n(M,N)$ & $= H^{-n}(M \otimes^L N)$ \\
\hline
\end{tabular}
\end{center}
\end{notation}


\section{Exercises}\label{sec:derived-exercises}

\begin{exercise}\label{exo:acyclic-cone}
\index{mapping cone!and quasi-isomorphism}
Let $f \colon A^\bullet \to B^\bullet$ be a chain map. Show that $f$ is a quasi-isomorphism
if and only if $\operatorname{Cone}(f)$ is acyclic.
\end{exercise}

\begin{exercise}\label{exo:homotopy-category-additive}
Show that the homotopy category $K(\mathcal{A})$ is additive. Verify that it is not
abelian by finding an explicit example (e.g., in $K(\Ab)$) of a morphism that has no
kernel in $K(\Ab)$.
\end{exercise}

\begin{exercise}\label{exo:shift-exact}
\index{shift functor!exactness}
Show that the shift functor $[1] \colon D(\mathcal{A}) \to D(\mathcal{A})$ is an
exact functor of triangulated categories (sends distinguished triangles to
distinguished triangles).
\end{exercise}

\begin{exercise}\label{exo:ext-computation}
\index{Ext!computation}
Let $k$ be a field. Compute $\operatorname{Ext}^n_{k[x]}(k, k)$ for all $n \geq 0$,
where $k$ is viewed as a $k[x]$-module via $x \mapsto 0$. (Use the projective resolution
$0 \to k[x] \xrightarrow{x} k[x] \to k \to 0$.)
\end{exercise}

\begin{exercise}\label{exo:tor-flat}
\index{Tor!and flatness}
Let $R$ be a commutative ring and $M$ an $R$-module.
\begin{enumerate}
    \item Show that $M$ is flat if and only if $\operatorname{Tor}_1^R(M, R/I) = 0$ for
    every ideal $I \subseteq R$.
    \item Compute $\operatorname{Tor}_n^{\mathbb{Z}}(\mathbb{Q}, \mathbb{Z}/m\mathbb{Z})$
    for all $n, m$.
\end{enumerate}
\end{exercise}

\begin{exercise}\label{exo:derived-ses}
\index{derived functor!long exact sequence}
Let $0 \to M' \to M \to M'' \to 0$ be an exact sequence of $R$-modules and $N$ any
$R$-module. Write down explicitly the long exact sequences in $\operatorname{Ext}$ and
$\operatorname{Tor}$ obtained by applying $\Hom_R(N, -)$, $\Hom_R(-, N)$, and
$N \otimes_R -$.
\end{exercise}

\begin{exercise}\label{exo:global-dimension}
\index{global dimension}\index{projective dimension}
The \textbf{projective dimension} $\operatorname{pd}(M)$ of a module $M$ is the infimum of
lengths of projective resolutions of $M$. The \textbf{global dimension} of $R$ is
$\operatorname{gldim}(R) = \sup_M \operatorname{pd}(M)$.
\begin{enumerate}
    \item Show that $\operatorname{pd}(M) \leq n$ if and only if $\operatorname{Ext}^{n+1}_R(M, N) = 0$
    for all $N$.
    \item Show that $\operatorname{gldim}(\mathbb{Z}) = 1$.
    \item Show that $\operatorname{gldim}(k[x_1, \ldots, x_n]) = n$ for a field $k$
    (this is the Hilbert syzygy theorem\index{Hilbert syzygy theorem}).
\end{enumerate}
\end{exercise}

\begin{exercise}\label{exo:triangulated-hom}
\index{triangulated category!Hom long exact sequence}
Let $\mathcal{T}$ be a triangulated category and $A \xrightarrow{f} B \xrightarrow{g} C
\xrightarrow{h} A[1]$ a distinguished triangle. For any object $X$, show there are long
exact sequences:
\begin{align*}
    \cdots &\to \Hom(X, A) \to \Hom(X, B) \to \Hom(X, C) \to \Hom(X, A[1]) \to \cdots \\
    \cdots &\to \Hom(C, X) \to \Hom(B, X) \to \Hom(A, X) \to \Hom(C[-1], X) \to \cdots
\end{align*}
\end{exercise}

\begin{exercise}\label{exo:derived-cat-Hom-Ext}
\index{derived category!Hom and Ext}
Show that in $D(\mathcal{A})$, for objects $A, B \in \mathcal{A}$ (viewed as complexes
concentrated in degree 0):
\[
    \Hom_{D(\mathcal{A})}(A, B[n]) \cong \operatorname{Ext}^n_{\mathcal{A}}(A, B)
    \quad \text{for all } n \geq 0.
\]
This is one of the key motivations for derived categories: the Ext groups are simply
Hom spaces with shifts.
\end{exercise}

\begin{exercise}\label{exo:verdier-quotient}
\index{Verdier quotient}\index{triangulated category!Verdier quotient}
(Verdier quotient.) Let $\mathcal{T}$ be a triangulated category and
$\mathcal{N} \subseteq \mathcal{T}$ a \textbf{thick subcategory} (a full triangulated
subcategory closed under direct summands). Define the Verdier quotient $\mathcal{T}/\mathcal{N}$
as the localization at morphisms whose cone lies in $\mathcal{N}$.
\begin{enumerate}
    \item Show that $D(\mathcal{A}) \cong K(\mathcal{A}) / \operatorname{Ac}(\mathcal{A})$,
    where $\operatorname{Ac}(\mathcal{A})$ is the thick subcategory of acyclic complexes.
    \item Show that the Verdier quotient inherits a triangulated structure.
\end{enumerate}
\end{exercise}
%%%%%%%%%%%%%%%%%%%%%%%%%%%%%%%%%%%%%%%%%%%%%%%%%%%%%%%%%%%%%
%% CATEGORY THEORY — Chapters 9–11 + End Matter
%% Monoidal Categories / Grothendieck Toposes / ∞-Categories
%%%%%%%%%%%%%%%%%%%%%%%%%%%%%%%%%%%%%%%%%%%%%%%%%%%%%%%%%%%%%

%%%%%%%%%%%%%%%%%%%%%%%%%%%%%%%%%%%%%%%%%%%%%%%%%%%%%%%%%%%%
%%         CHAPTER 9: MONOIDAL AND BRAIDED CATEGORIES     %%
%%%%%%%%%%%%%%%%%%%%%%%%%%%%%%%%%%%%%%%%%%%%%%%%%%%%%%%%%%%%

\chapter{Monoidal and Braided Categories}\label{ch:monoidal}
\minitoc

\vspace{1em}

In many areas of mathematics and theoretical physics, one encounters
categories equipped with a ``tensor product''---a bifunctor that combines
objects in a way that is associative and unital up to coherent isomorphism.
Such categories are called \emph{monoidal}.  They provide the natural
setting for studying algebras, coalgebras, and their modules in a
category-theoretic framework.  When the tensor product is commutative
up to isomorphism, one obtains \emph{braided} and \emph{symmetric} monoidal
categories, which underpin the theory of quantum groups and topological
quantum field theories.  In this chapter we develop the foundations of
monoidal category theory, state Mac~Lane's coherence theorem, study
enriched and closed categories, and touch upon the Curry--Howard--Lambek
correspondence.

\index{monoidal category|(}

%% ============================================================
\section{Monoidal categories}\label{sec:monoidal-cats}
%% ============================================================

\begin{definition}[Monoidal category]\label{def:monoidal-cat}
\index{monoidal category!definition}
\index{tensor product!in a monoidal category}
\index{associator}
\index{unitor!left}
\index{unitor!right}
\index{unit object}
A \textbf{monoidal category} is a tuple
$(\mathcal{C},\, \otimes,\, I,\, \alpha,\, \lambda,\, \rho)$ where:
\begin{enumerate}[label=(\roman*)]
\item $\mathcal{C}$ is a category;
\item $\otimes \colon \mathcal{C}\times\mathcal{C}\to\mathcal{C}$ is a
      bifunctor called the \textbf{tensor product};
\item $I \in \Ob(\mathcal{C})$ is an object called the \textbf{unit object}
      (or \textbf{monoidal unit});
\item $\alpha_{A,B,C}\colon (A\otimes B)\otimes C
      \xrightarrow{\;\sim\;} A\otimes(B\otimes C)$
      is a natural isomorphism called the \textbf{associator};
\item $\lambda_A \colon I\otimes A \xrightarrow{\;\sim\;} A$
      is a natural isomorphism called the \textbf{left unitor};
\item $\rho_A \colon A\otimes I \xrightarrow{\;\sim\;} A$
      is a natural isomorphism called the \textbf{right unitor};
\end{enumerate}
subject to the \emph{pentagon axiom} and the \emph{triangle axiom} below.
\end{definition}

\begin{definition}[Pentagon axiom]\label{def:pentagon}
\index{pentagon axiom}
For all objects $A,B,C,D$ of~$\mathcal{C}$, the following diagram commutes:
\[
\begin{tikzcd}[column sep=small]
& (A\otimes B)\otimes(C\otimes D) \arrow[dr, "\alpha_{A,B,C\otimes D}"] & \\
((A\otimes B)\otimes C)\otimes D
  \arrow[ur, "\alpha_{A\otimes B,C,D}"]
  \arrow[d, "\alpha_{A,B,C}\otimes\id_D"'] & &
A\otimes(B\otimes(C\otimes D)) \\
(A\otimes(B\otimes C))\otimes D
  \arrow[rr, "\alpha_{A,B\otimes C,D}"'] & &
A\otimes((B\otimes C)\otimes D)
  \arrow[u, "\id_A\otimes\alpha_{B,C,D}"']
\end{tikzcd}
\]
\end{definition}

\begin{definition}[Triangle axiom]\label{def:triangle}
\index{triangle axiom}
For all objects $A,B$ of~$\mathcal{C}$, the following diagram commutes:
\[
\begin{tikzcd}
(A\otimes I)\otimes B
  \arrow[rr, "\alpha_{A,I,B}"]
  \arrow[dr, "\rho_A\otimes\id_B"'] & &
A\otimes(I\otimes B)
  \arrow[dl, "\id_A\otimes\lambda_B"] \\
& A\otimes B &
\end{tikzcd}
\]
\end{definition}

\begin{definition}[Strict monoidal category]\label{def:strict-monoidal}
\index{monoidal category!strict}
A monoidal category is called \textbf{strict} if $\alpha$, $\lambda$,
and $\rho$ are all identity natural transformations, i.e.\
$(A\otimes B)\otimes C = A\otimes (B\otimes C)$,
$I\otimes A = A$, and $A\otimes I = A$ on the nose.
\end{definition}

\begin{example}\label{ex:monoidal-examples}
\index{monoidal category!examples}
The following are important examples of monoidal categories.
\begin{enumerate}[label=(\roman*)]
\item $(\Set, \times, \{*\})$: the category of sets with cartesian product.
      \index{Set@$\Set$!monoidal structure}
\item $(\Ab, \otimes_\Z, \Z)$: abelian groups with the tensor product over~$\Z$.
      \index{Ab@$\Ab$!monoidal structure}
\item $(\Vect_\K, \otimes_\K, \K)$: vector spaces over a field~$\K$
      with the usual tensor product.
      \index{Vect@$\Vect$!monoidal structure}
\item $(R\text{-}\Mod, \otimes_R, R)$ for a commutative ring~$R$.
      \index{Mod@$\Mod$!monoidal structure}
\item $(\Cat, \times, \mathbf{1})$: the category of small categories
      with the cartesian product.
      \index{Cat@$\Cat$!monoidal structure}
\item $(\mathrm{End}(\mathcal{C}), \circ, \id_\mathcal{C})$: endofunctors
      of a category~$\mathcal{C}$ with composition.  This is a strict
      monoidal category.
      \index{endofunctor!monoidal category of}
\item $(\Set, \sqcup, \varnothing)$: sets with disjoint union.
\item The category of chain complexes $\mathrm{Ch}(R)$ with the
      tensor product of complexes.
      \index{chain complex!monoidal structure}
\end{enumerate}
\end{example}

\begin{remark}\label{rem:monoidal-functor-notation}
We often write just $(\mathcal{C}, \otimes, I)$ for a monoidal category,
leaving the structural isomorphisms implicit.
\end{remark}


%% ============================================================
\section{Mac Lane's coherence theorem}\label{sec:coherence}
%% ============================================================

\index{coherence theorem|(}
\index{Mac Lane!coherence theorem}

The pentagon and triangle axioms may appear insufficient to guarantee that
\emph{all} diagrams built from $\alpha$, $\lambda$, $\rho$ and their
inverses commute.  Mac~Lane's celebrated coherence theorem asserts that
this is indeed the case.

\begin{theorem}[Mac Lane's coherence theorem]\label{thm:maclane-coherence}
\index{coherence theorem!Mac Lane}
Every monoidal category is monoidally equivalent to a strict monoidal
category.  Equivalently, every diagram in a monoidal category whose
morphisms are composites of instances of $\alpha$, $\alpha^{-1}$,
$\lambda$, $\lambda^{-1}$, $\rho$, $\rho^{-1}$, identities, and tensor
products thereof, commutes.
\end{theorem}

\begin{proof}[Proof sketch]
One constructs a strict monoidal category $\mathcal{C}_s$ whose objects
are finite words in $\Ob(\mathcal{C})$ (including the empty word, which
serves as the strict unit) and whose morphisms are induced by those
of~$\mathcal{C}$.  The multiplication is word concatenation.
An explicit monoidal equivalence $\mathcal{C}_s \simeq \mathcal{C}$ is
built by iterating the tensor product.  The commutativity of arbitrary
constraint diagrams then reduces to the trivially commuting diagrams
in~$\mathcal{C}_s$.  For the full proof, see Mac~Lane~\cite{MacLane1963}
or~\cite{MacLaneCWM}.
\end{proof}

\begin{remark}\label{rem:coherence-practical}
\index{coherence theorem!practical consequence}
The practical consequence of the coherence theorem is that, for
calculations, one may pretend that every monoidal category is strict:
one simply omits all parentheses and all instances of the
associator and unitors.  We shall make frequent use of this
simplification in what follows.
\end{remark}

\index{coherence theorem|)}


%% ============================================================
\section{Braided and symmetric monoidal categories}\label{sec:braided}
%% ============================================================

\index{braided monoidal category|(}

\begin{definition}[Braided monoidal category]\label{def:braided}
\index{braided monoidal category!definition}
\index{braiding}
A \textbf{braided monoidal category} is a monoidal category
$(\mathcal{C}, \otimes, I)$ equipped with a natural isomorphism
\[
   \sigma_{A,B} \colon A\otimes B \xrightarrow{\;\sim\;} B\otimes A
\]
called the \textbf{braiding}, satisfying the two \emph{hexagon axioms}:
\[
\begin{tikzcd}[column sep=1.8em]
(A\otimes B)\otimes C
  \arrow[r, "\alpha"]
  \arrow[d, "\sigma_{A,B}\otimes\id"'] &
A\otimes(B\otimes C)
  \arrow[r, "\sigma_{A,B\otimes C}"] &
(B\otimes C)\otimes A
  \arrow[d, "\alpha"] \\
(B\otimes A)\otimes C
  \arrow[r, "\alpha"'] &
B\otimes(A\otimes C)
  \arrow[r, "\id\otimes\sigma_{A,C}"'] &
B\otimes(C\otimes A)
\end{tikzcd}
\]
and
\[
\begin{tikzcd}[column sep=1.8em]
A\otimes(B\otimes C)
  \arrow[r, "\alpha^{-1}"]
  \arrow[d, "\id\otimes\sigma_{B,C}"'] &
(A\otimes B)\otimes C
  \arrow[r, "\sigma_{A\otimes B,C}"] &
C\otimes(A\otimes B)
  \arrow[d, "\alpha^{-1}"] \\
A\otimes(C\otimes B)
  \arrow[r, "\alpha^{-1}"'] &
(A\otimes C)\otimes B
  \arrow[r, "\sigma_{A,C}\otimes\id"'] &
(C\otimes A)\otimes B
\end{tikzcd}
\]
\end{definition}

\begin{definition}[Symmetric monoidal category]\label{def:symmetric}
\index{symmetric monoidal category}
\index{braiding!symmetric}
A braided monoidal category is \textbf{symmetric} if
$\sigma_{B,A}\circ \sigma_{A,B} = \id_{A\otimes B}$ for all $A,B$.
\end{definition}

\begin{example}\label{ex:braided-symmetric}
\begin{enumerate}[label=(\roman*)]
\item $(\Set, \times, \{*\})$, $(\Ab, \otimes_\Z, \Z)$, and
      $(\Vect_\K, \otimes_\K, \K)$ are all symmetric monoidal categories
      with the evident swap maps.
      \index{symmetric monoidal category!examples}
\item The category of $R$-modules with $\otimes_R$ is symmetric monoidal.
\item The category of graded $R$-modules with the Koszul sign rule
      $\sigma(a\otimes b) = (-1)^{|a||b|}\, b\otimes a$ is symmetric monoidal.
      \index{Koszul sign rule}
\item The category of representations of a quantum group is braided
      monoidal but typically not symmetric.
      \index{quantum group}
\end{enumerate}
\end{example}

\index{braided monoidal category|)}


%% ============================================================
\section{Monoid objects}\label{sec:monoid-objects}
%% ============================================================

\index{monoid object|(}

\begin{definition}[Monoid object]\label{def:monoid-object}
\index{monoid object!definition}
\index{multiplication morphism}
\index{unit morphism}
Let $(\mathcal{C}, \otimes, I)$ be a monoidal category.
A \textbf{monoid object} (or \textbf{monoid}) in~$\mathcal{C}$ is a
triple $(M, \mu, \eta)$ where $M\in\Ob(\mathcal{C})$,
$\mu\colon M\otimes M \to M$ is the \textbf{multiplication}, and
$\eta\colon I\to M$ is the \textbf{unit}, satisfying associativity
and unitality:
\[
\begin{tikzcd}
(M\otimes M)\otimes M \arrow[r, "\alpha"] \arrow[d, "\mu\otimes\id"'] &
M\otimes(M\otimes M) \arrow[r, "\id\otimes\mu"] &
M\otimes M \arrow[d, "\mu"] \\
M\otimes M \arrow[rr, "\mu"'] && M
\end{tikzcd}
\qquad
\begin{tikzcd}
I\otimes M \arrow[r, "\eta\otimes\id"] \arrow[dr, "\lambda"'] &
M\otimes M \arrow[d, "\mu"] &
M\otimes I \arrow[l, "\id\otimes\eta"'] \arrow[dl, "\rho"] \\
& M &
\end{tikzcd}
\]
A \textbf{morphism of monoid objects} $f\colon (M,\mu,\eta)\to (M',\mu',\eta')$
is a morphism $f\colon M\to M'$ in~$\mathcal{C}$ with
$f\circ\mu = \mu'\circ(f\otimes f)$ and $f\circ\eta = \eta'$.
The category of monoid objects in~$\mathcal{C}$ is denoted
$\mathrm{Mon}(\mathcal{C})$.
\index{Mon@$\mathrm{Mon}(\mathcal{C})$}
\end{definition}

\begin{example}\label{ex:monoid-objects}
\index{monoid object!examples}
\begin{enumerate}[label=(\roman*)]
\item A monoid object in $(\Set, \times, \{*\})$ is a monoid in the
      usual sense.
\item A monoid object in $(\Ab, \otimes_\Z, \Z)$ is a ring (not
      necessarily commutative).
      \index{ring!as monoid object}
\item A monoid object in $(\Vect_\K, \otimes_\K, \K)$ is a
      $\K$-algebra.
      \index{algebra!as monoid object}
\item A monoid object in $(\mathrm{End}(\mathcal{C}), \circ, \id_\mathcal{C})$
      is a monad on~$\mathcal{C}$.
      \index{monad!as monoid object}
\item A \textbf{commutative monoid object} in a symmetric monoidal
      category $(\mathcal{C}, \otimes, I, \sigma)$ is a monoid $(M,\mu,\eta)$
      with $\mu \circ \sigma_{M,M} = \mu$.
      \index{monoid object!commutative}
\end{enumerate}
\end{example}

\begin{definition}[Comonoid object]\label{def:comonoid}
\index{comonoid object}
Dually, a \textbf{comonoid object} in $(\mathcal{C}, \otimes, I)$ is
a triple $(C, \Delta, \varepsilon)$ where
$\Delta\colon C\to C\otimes C$ (the \emph{comultiplication}) and
$\varepsilon\colon C\to I$ (the \emph{counit}) satisfy the duals of
the associativity and unitality diagrams.
\end{definition}

\index{monoid object|)}


%% ============================================================
\section{Enriched categories}\label{sec:enriched}
%% ============================================================

\index{enriched category|(}

Ordinary category theory is ``enriched over~$\Set$'': for each pair of
objects $A,B$, the hom $\Hom(A,B)$ is a \emph{set}.  But in practice
these hom-sets often carry additional structure---they may be abelian
groups ($\Ab$-enriched), topological spaces ($\Top$-enriched), or
chain complexes (dg-enriched).

\begin{definition}[$\mathcal{V}$-enriched category]\label{def:enriched-cat}
\index{enriched category!definition}
\index{V-category@$\mathcal{V}$-category}
Let $(\mathcal{V}, \otimes, I)$ be a monoidal category.
A \textbf{$\mathcal{V}$-enriched category} (or \textbf{$\mathcal{V}$-category})
$\mathcal{A}$ consists of:
\begin{enumerate}[label=(\roman*)]
\item a class $\Ob(\mathcal{A})$ of objects;
\item for each pair $A,B\in\Ob(\mathcal{A})$, an object
      $\mathcal{A}(A,B)\in\Ob(\mathcal{V})$ (the \emph{hom-object});
\item for each triple $A,B,C$, a \emph{composition morphism}
      \[
          c_{A,B,C}\colon \mathcal{A}(B,C)\otimes\mathcal{A}(A,B)
          \longrightarrow \mathcal{A}(A,C) \quad\text{in }\mathcal{V};
      \]
\item for each object~$A$, an \emph{identity morphism}
      $j_A\colon I\to\mathcal{A}(A,A)$ in~$\mathcal{V}$;
\end{enumerate}
subject to associativity and unitality axioms expressed as commutative
diagrams in~$\mathcal{V}$.
\end{definition}

\begin{example}\label{ex:enriched-examples}
\index{enriched category!examples}
\begin{enumerate}[label=(\roman*)]
\item A $\Set$-enriched category is an ordinary (locally small) category.
\item An $\Ab$-enriched category is a \emph{preadditive category}:
      hom-sets are abelian groups and composition is bilinear.
      \index{preadditive category}
\item A $\Top$-enriched category is a \emph{topological category}:
      hom-sets are topological spaces and composition is continuous.
\item A category enriched over the monoidal category $([0,\infty]^{\op}, +, 0)$
      (extended non-negative reals, reversed order, addition) is a
      (generalised) \emph{metric space} in the sense of Lawvere.
      \index{Lawvere metric space}
      \index{metric space!as enriched category}
\item A category enriched over chain complexes is a \emph{dg-category}.
      \index{dg-category}
\end{enumerate}
\end{example}

\begin{definition}[$\mathcal{V}$-functor]\label{def:enriched-functor}
\index{enriched category!functor}
\index{V-functor@$\mathcal{V}$-functor}
A \textbf{$\mathcal{V}$-functor} $F\colon\mathcal{A}\to\mathcal{B}$
between $\mathcal{V}$-categories consists of a map
$F\colon\Ob(\mathcal{A})\to\Ob(\mathcal{B})$ and morphisms
$F_{A,B}\colon\mathcal{A}(A,B)\to\mathcal{B}(FA,FB)$ in~$\mathcal{V}$
compatible with composition and identities.
\end{definition}

\begin{definition}[$\mathcal{V}$-natural transformation]\label{def:enriched-nat}
\index{enriched category!natural transformation}
A \textbf{$\mathcal{V}$-natural transformation}
$\eta\colon F\Rightarrow G$ between $\mathcal{V}$-functors
$F,G\colon\mathcal{A}\to\mathcal{B}$ consists of a family of morphisms
$\eta_A\colon I\to\mathcal{B}(FA,GA)$ in~$\mathcal{V}$, one for each
$A\in\Ob(\mathcal{A})$, satisfying the enriched naturality condition.
\end{definition}

\index{enriched category|)}


%% ============================================================
\section{Closed monoidal and cartesian closed categories}%
\label{sec:closed-monoidal}
%% ============================================================

\index{closed monoidal category|(}
\index{cartesian closed category|(}

\begin{definition}[Closed monoidal category]\label{def:closed-monoidal}
\index{closed monoidal category!definition}
\index{internal hom}
A monoidal category $(\mathcal{C}, \otimes, I)$ is
\textbf{(left) closed} if for each object~$B$ the functor
$(-)\otimes B\colon\mathcal{C}\to\mathcal{C}$ has a right adjoint,
denoted $[B, -]$ or $\underline{\Hom}(B,-)$:
\[
    \Hom_\mathcal{C}(A\otimes B,\, C)
    \;\cong\;
    \Hom_\mathcal{C}(A,\, [B,C])
    \quad\text{naturally in } A \text{ and } C.
\]
The object $[B,C]$ is called the \textbf{internal hom}
(or \textbf{exponential object}).
\end{definition}

\begin{definition}[Cartesian closed category]\label{def:cartesian-closed}
\index{cartesian closed category!definition}
\index{exponential object}
A category $\mathcal{C}$ with finite products is \textbf{cartesian closed}
if it is closed monoidal with respect to the cartesian monoidal structure
$(\mathcal{C}, \times, 1)$.  That is, for every object~$B$ the functor
$(-)\times B$ has a right adjoint $(-)^B$:
\[
    \Hom(A\times B,\, C) \;\cong\; \Hom(A,\, C^B).
\]
\end{definition}

\begin{example}\label{ex:cartesian-closed}
\index{cartesian closed category!examples}
\begin{enumerate}[label=(\roman*)]
\item $\Set$ is cartesian closed with $C^B = \Hom_{\Set}(B,C)$.
      \index{Set@$\Set$!cartesian closed}
\item The category $\Cat$ of small categories is cartesian closed with
      the exponential $\mathcal{D}^\mathcal{C} = \Fun(\mathcal{C},\mathcal{D})$.
\item Any elementary topos is cartesian closed (Theorem~\ref{thm:topos-ccc}).
\item $\Top$ is \emph{not} cartesian closed in general; one must pass to
      a convenient subcategory such as compactly generated Hausdorff spaces.
      \index{Top@$\Top$!not cartesian closed}
\item $\Ab$ is not cartesian closed (the cartesian product is the direct
      sum, and $\Hom_\Ab(A\oplus B, C)\not\cong\Hom_\Ab(A, C^B)$ in general),
      but it is closed monoidal with $\otimes_\Z$ and internal hom $\Hom_\Z(B,C)$.
\end{enumerate}
\end{example}

\begin{proposition}[Evaluation and coevaluation]\label{prop:eval-coeval}
\index{evaluation morphism}
\index{coevaluation morphism}
In a closed monoidal category, there are natural morphisms:
\begin{enumerate}[label=(\alph*)]
\item the \textbf{evaluation}: $\mathrm{ev}_{B,C}\colon [B,C]\otimes B\to C$,
      corresponding to $\id_{[B,C]}$ under the adjunction;
\item the \textbf{coevaluation}: $\mathrm{coev}_{A,B}\colon A\to [B, A\otimes B]$,
      corresponding to $\id_{A\otimes B}$ under the adjunction.
\end{enumerate}
\end{proposition}

\begin{proof}
Apply the adjunction isomorphism
$\Hom(A\otimes B, C)\cong\Hom(A,[B,C])$
to the appropriate identity morphisms.
\end{proof}

\index{closed monoidal category|)}
\index{cartesian closed category|)}


%% ============================================================
\section{The Curry--Howard--Lambek correspondence}%
\label{sec:curry-howard-lambek}
%% ============================================================

\index{Curry--Howard--Lambek correspondence|(}
\index{lambda calculus}
\index{intuitionistic logic}

The Curry--Howard--Lambek correspondence is a remarkable three-way
dictionary between logic, computation, and category theory.  We give a
brief account aimed at illustrating the categorical perspective.

\begin{center}
\renewcommand{\arraystretch}{1.3}
\begin{tabular}{lll}
\hline
\textbf{Logic} & \textbf{Type theory} & \textbf{Category theory} \\
\hline
Proposition $A$ & Type $A$ & Object $A$ \\
Proof of $A$ & Term of type $A$ & Morphism $1\to A$ \\
$A\Rightarrow B$ & Function type $A\to B$ & Exponential $B^A$ \\
$A\wedge B$ & Product type $A\times B$ & Product $A\times B$ \\
$A\vee B$ & Sum type $A + B$ & Coproduct $A + B$ \\
$\top$ (true) & Unit type $1$ & Terminal object \\
$\bot$ (false) & Empty type $0$ & Initial object \\
Modus ponens & Function application & Evaluation $B^A\times A\to B$ \\
\hline
\end{tabular}
\end{center}

\begin{theorem}[Lambek]\label{thm:lambek}
\index{Lambek!correspondence}
\index{cartesian closed category!and lambda calculus}
The simply-typed lambda calculus (with products) is the internal language
of cartesian closed categories.  Specifically:
\begin{enumerate}[label=(\alph*)]
\item Every cartesian closed category gives rise to a simply-typed
      lambda calculus whose types are the objects and whose terms are
      the morphisms.
\item Every simply-typed lambda calculus (modulo $\beta\eta$-equivalence)
      gives rise to a cartesian closed category, the \emph{syntactic category}.
\item These two constructions are mutually inverse (up to equivalence).
\end{enumerate}
\end{theorem}

\begin{remark}\label{rem:curry-howard-extensions}
\index{Curry--Howard--Lambek correspondence!extensions}
The correspondence extends far beyond the simply-typed case:
\begin{itemize}
\item Dependent types correspond to locally cartesian closed categories.
\item Linear logic corresponds to $*$-autonomous categories (a form of
      closed monoidal category).
      \index{linear logic}
      \index{*-autonomous category@$*$-autonomous category}
\item Homotopy type theory corresponds to $(\infty,1)$-toposes.
      \index{homotopy type theory}
\end{itemize}
\end{remark}

\index{Curry--Howard--Lambek correspondence|)}


%% ============================================================
\section{Exercises}\label{sec:monoidal-exercises}
%% ============================================================

\begin{exercise}\label{ex:strict-is-monoidal}
Show that every strict monoidal category is a monoidal category
(the axioms are trivially satisfied with identity natural transformations).
\end{exercise}

\begin{exercise}\label{ex:monoidal-unique-unit}
\index{unit object!uniqueness}
Let $(\mathcal{C}, \otimes, I)$ be a monoidal category.  Show that the
unit object~$I$ is unique up to isomorphism, in the sense that if
$(I', \lambda', \rho')$ also satisfies the axioms, then $I\cong I'$.
\end{exercise}

\begin{exercise}\label{ex:commutative-monoid-object}
Let $(\mathcal{C}, \otimes, I, \sigma)$ be a symmetric monoidal category.
Show that the category $\mathrm{CMon}(\mathcal{C})$ of commutative monoid
objects in~$\mathcal{C}$ admits a monoidal structure inherited from that
of~$\mathcal{C}$.
\end{exercise}

\begin{exercise}\label{ex:comonoid-in-set}
\index{comonoid object!in $\Set$}
Show that every set $S$ has a unique comonoid structure in
$(\Set, \times, \{*\})$, given by $\Delta(s) = (s,s)$ and
$\varepsilon(s) = *$.
\end{exercise}

\begin{exercise}\label{ex:enriched-underlying}
\index{underlying category}
Let $\mathcal{A}$ be a $\mathcal{V}$-enriched category where
$\mathcal{V}$ is closed monoidal.  Define the \emph{underlying ordinary
category} $\mathcal{A}_0$ of~$\mathcal{A}$ and show it is indeed a category.
\textit{Hint:} take $\Hom_{\mathcal{A}_0}(A,B) = \Hom_\mathcal{V}(I, \mathcal{A}(A,B))$.
\end{exercise}

\begin{exercise}\label{ex:day-convolution}
\index{Day convolution}
Let $(\mathcal{C}, \otimes, I)$ be a small monoidal category.
Show that the category of presheaves $[\mathcal{C}^{\op}, \Set]$ carries
a monoidal structure given by \emph{Day convolution}:
\[
    (F\star G)(C)
    = \int^{A,B\in\mathcal{C}} \Hom_\mathcal{C}(A\otimes B, C)
      \times F(A)\times G(B).
\]
Identify the unit object of this monoidal structure.
\end{exercise}

\begin{exercise}\label{ex:ccc-lambda}
Verify explicitly that in $\Set$, the adjunction
$\Hom(A\times B, C)\cong\Hom(A, C^B)$ corresponds to currying of functions:
$f\colon A\times B\to C$ is sent to $\hat{f}\colon A\to C^B$ defined by
$\hat{f}(a)(b) = f(a,b)$.
\end{exercise}

\index{monoidal category|)}


%%%%%%%%%%%%%%%%%%%%%%%%%%%%%%%%%%%%%%%%%%%%%%%%%%%%%%%%%%%%
%%    CHAPTER 10: GROTHENDIECK TOPOSES — INTRODUCTION     %%
%%%%%%%%%%%%%%%%%%%%%%%%%%%%%%%%%%%%%%%%%%%%%%%%%%%%%%%%%%%%

\chapter{Grothendieck Toposes---Introduction}\label{ch:toposes}
\minitoc

\vspace{1em}

Topos theory, introduced by Grothendieck and his school in the 1960s
for the purposes of algebraic geometry, has grown into one of the deepest
and most far-reaching branches of category theory.  A \emph{Grothendieck
topos} is a category of sheaves on a site; it simultaneously generalises
the category of sheaves on a topological space and the category of
presheaves on a small category.  Elementary toposes, introduced by
Lawvere and Tierney, isolate the key categorical properties that make
a category ``behave like the category of sets.''

In this chapter we give an introduction to the subject, defining sites,
Grothendieck topologies, sheaves, and the resulting toposes.
We discuss elementary toposes, subobject classifiers, and the internal
logic, then state Giraud's theorem characterising Grothendieck toposes.

\index{topos|(}
\index{Grothendieck topos|(}

%% ============================================================
\section{Sites and Grothendieck topologies}\label{sec:sites}
%% ============================================================

\index{site|(}
\index{Grothendieck topology|(}

\begin{definition}[Sieve]\label{def:sieve}
\index{sieve}
Let $\mathcal{C}$ be a category and $C\in\Ob(\mathcal{C})$.
A \textbf{sieve on~$C$} is a collection~$S$ of morphisms with
codomain~$C$ such that if $f\colon D\to C$ is in~$S$ and
$g\colon E\to D$ is any morphism, then $f\circ g\in S$.
Equivalently, a sieve on~$C$ is a subfunctor of the representable
presheaf $\Hom(-,C)$.
\end{definition}

\begin{definition}[Grothendieck topology]\label{def:grothendieck-topology}
\index{Grothendieck topology!definition}
\index{covering sieve}
A \textbf{Grothendieck topology}~$J$ on a small category~$\mathcal{C}$
assigns to each object~$C$ a collection $J(C)$ of sieves on~$C$, called
\textbf{covering sieves}, satisfying:
\begin{enumerate}[label=(\roman*)]
\item \textbf{Maximal sieve:} The maximal sieve
      $t_C = \{f \mid \mathrm{cod}(f) = C\}$ is in $J(C)$.
\item \textbf{Stability:} If $S\in J(C)$ and $h\colon D\to C$ is any
      morphism, then the pullback sieve
      $h^*S = \{g\colon E\to D \mid h\circ g\in S\}$ is in $J(D)$.
\item \textbf{Transitivity:} If $S\in J(C)$ and $R$ is a sieve on~$C$
      such that for every $f\colon D\to C$ in~$S$ the pullback
      $f^*R\in J(D)$, then $R\in J(C)$.
\end{enumerate}
\end{definition}

\begin{definition}[Site]\label{def:site}
\index{site!definition}
A \textbf{site} is a pair $(\mathcal{C}, J)$ consisting of a small
category~$\mathcal{C}$ and a Grothendieck topology~$J$ on~$\mathcal{C}$.
\end{definition}

\begin{definition}[Coverage]\label{def:coverage}
\index{coverage}
\index{covering family}
In practice, Grothendieck topologies are often specified via
\textbf{coverages}: for each $C\in\Ob(\mathcal{C})$, one gives a
collection of families $\{f_i\colon C_i\to C\}_{i\in I}$, called
\textbf{covering families}, satisfying a stability condition under
pullback.  The covering families generate sieves, hence a Grothendieck
topology.
\end{definition}

\begin{example}\label{ex:topologies}
\index{Grothendieck topology!examples}
\begin{enumerate}[label=(\roman*)]
\item \textbf{Trivial topology.}
      \index{trivial topology}
      Only the maximal sieves cover.  Then every presheaf is a sheaf.
\item \textbf{Discrete topology.}
      \index{discrete topology}
      Every sieve is a covering sieve.  The only sheaf is the terminal
      presheaf.
\item \textbf{Canonical topology on $\Top$.}
      \index{canonical topology}
      For an open set~$U$, a covering sieve is generated by an
      open cover $\{U_i\}$ of~$U$.
\item \textbf{Zariski topology.}
      \index{Zariski topology}
      On the category of commutative rings (with the opposite orientation),
      covering families correspond to collections of localisations
      $\{R\to R[f_i^{-1}]\}$ where $(f_1,\ldots,f_n) = R$.
\item \textbf{\'Etale topology.}
      \index{etale topology@\'etale topology}
      On the category of schemes, covering families are jointly surjective
      families of \'etale morphisms.
\item \textbf{fppf and fpqc topologies.}
      \index{fppf topology}
      \index{fpqc topology}
      Finer topologies on schemes used in algebraic geometry.
\end{enumerate}
\end{example}

\index{Grothendieck topology|)}
\index{site|)}


%% ============================================================
\section{Sheaves on a site}\label{sec:sheaves-site}
%% ============================================================

\index{sheaf|(}

\begin{definition}[Presheaf]\label{def:presheaf-reminder}
\index{presheaf}
A \textbf{presheaf} on a small category~$\mathcal{C}$ is a functor
$F\colon\mathcal{C}^{\op}\to\Set$.  The category of presheaves is
$\mathrm{PSh}(\mathcal{C}) = [\mathcal{C}^{\op}, \Set]$.
\index{PSh@$\mathrm{PSh}(\mathcal{C})$}
\end{definition}

\begin{definition}[Sheaf on a site]\label{def:sheaf-site}
\index{sheaf!on a site}
\index{sheaf condition}
Let $(\mathcal{C}, J)$ be a site.  A presheaf $F\colon\mathcal{C}^{\op}\to\Set$
is a \textbf{sheaf} (for the topology~$J$) if for every covering sieve
$S\in J(C)$, the natural map
\[
    F(C) \longrightarrow \varprojlim_{(f\colon D\to C)\in S} F(D)
\]
is a bijection.  Equivalently, for every covering family
$\{f_i\colon C_i\to C\}$ generating a covering sieve, the diagram
\[
\begin{tikzcd}
F(C) \arrow[r] &
\displaystyle\prod_i F(C_i)
  \arrow[r, shift left=0.5ex]
  \arrow[r, shift right=0.5ex] &
\displaystyle\prod_{i,j} F(C_i\times_C C_j)
\end{tikzcd}
\]
is an equaliser.
\end{definition}

\begin{definition}[Category of sheaves]\label{def:sheaf-cat}
\index{Sh@$\mathrm{Sh}(\mathcal{C}, J)$}
\index{topos!Grothendieck}
The full subcategory of $\mathrm{PSh}(\mathcal{C})$ consisting of
sheaves for the topology~$J$ is denoted
$\mathrm{Sh}(\mathcal{C}, J)$.  It is called a
\textbf{Grothendieck topos}.
\end{definition}

\begin{theorem}[Sheafification]\label{thm:sheafification}
\index{sheafification}
\index{associated sheaf}
The inclusion $i\colon\mathrm{Sh}(\mathcal{C},J)\hookrightarrow
\mathrm{PSh}(\mathcal{C})$ has a left adjoint
$a\colon\mathrm{PSh}(\mathcal{C})\to\mathrm{Sh}(\mathcal{C},J)$
called \textbf{sheafification} (or the \textbf{associated sheaf functor}).
Moreover:
\begin{enumerate}[label=(\alph*)]
\item $a$ is left exact (preserves finite limits);
\item $a\circ i \cong \id$;
\item the adjunction $a\dashv i$ is a reflection.
\end{enumerate}
\end{theorem}

\begin{proof}[Proof sketch]
The sheafification $aF$ is constructed by the ``plus construction''
applied twice: $(F^+)^+$.  Given $F$, one defines
\[
    F^+(C) = \varinjlim_{S\in J(C)} \mathrm{Match}(S, F)
\]
where $\mathrm{Match}(S,F)$ is the set of matching families for the
sieve~$S$.  One application of $(-)^+$ makes~$F$ separated; a second
application yields a sheaf.  Left exactness follows from the fact that
the colimit over covering sieves is filtered.
\end{proof}

\begin{example}\label{ex:sheaves-on-space}
\index{sheaf!on a topological space}
If $X$ is a topological space and $\mathcal{O}(X)$ is the poset of
open subsets (viewed as a category), then the open-cover topology makes
$(\mathcal{O}(X), J)$ a site and $\mathrm{Sh}(\mathcal{O}(X), J)$ is
the classical category of sheaves on~$X$.
\end{example}

\index{sheaf|)}


%% ============================================================
\section{Elementary toposes}\label{sec:elementary-topos}
%% ============================================================

\index{elementary topos|(}

\begin{definition}[Elementary topos]\label{def:elementary-topos}
\index{elementary topos!definition}
\index{subobject classifier}
An \textbf{elementary topos} is a category $\mathcal{E}$ satisfying:
\begin{enumerate}[label=(\roman*)]
\item $\mathcal{E}$ has all finite limits;
\item $\mathcal{E}$ has all finite colimits;
\item $\mathcal{E}$ is cartesian closed;
\item $\mathcal{E}$ has a \textbf{subobject classifier}.
\end{enumerate}
\end{definition}

\begin{definition}[Subobject classifier]\label{def:subobject-classifier}
\index{subobject classifier!definition}
\index{Omega@$\Omega$}
\index{true morphism}
A \textbf{subobject classifier} in a category~$\mathcal{E}$ with
finite limits is a monomorphism $\mathrm{true}\colon 1\to\Omega$
such that for every monomorphism $m\colon S\rightarrowtail A$
there exists a unique morphism $\chi_m\colon A\to\Omega$ making the
following square a pullback:
\[
\begin{tikzcd}
S \arrow[r] \arrow[d, tail, "m"'] & 1 \arrow[d, "\mathrm{true}"] \\
A \arrow[r, "\chi_m"'] & \Omega
\end{tikzcd}
\]
The morphism $\chi_m$ is called the \textbf{characteristic morphism}
(or \textbf{classifying morphism}) of the subobject~$m$.
\index{characteristic morphism}
\index{classifying morphism}
\end{definition}

\begin{example}\label{ex:subobject-classifier}
\index{subobject classifier!examples}
\begin{enumerate}[label=(\roman*)]
\item In $\Set$, $\Omega = \{0,1\}$ with $\mathrm{true}(*)=1$.
      The classifying morphism of a subset $S\subseteq A$ is
      the characteristic function $\chi_S$.
      \index{Set@$\Set$!subobject classifier}
\item In the topos of sheaves $\mathrm{Sh}(X)$ on a topological
      space~$X$, $\Omega$ is the sheaf of open subsets:
      $\Omega(U) = \{V\subseteq U \mid V\text{ open}\}$.
\item In a presheaf topos $[\mathcal{C}^{\op}, \Set]$,
      $\Omega(C) = \{\text{sieves on } C\}$.
      \index{presheaf topos!subobject classifier}
\end{enumerate}
\end{example}

\begin{theorem}\label{thm:topos-ccc}
\index{elementary topos!is cartesian closed}
Every elementary topos is cartesian closed.  In particular, every
Grothendieck topos is cartesian closed.
\end{theorem}

\begin{proof}[Proof sketch]
The power object $P(B) = \Omega^B$ exists by assumption in an elementary
topos (cartesian closedness is part of the definition).  For a
Grothendieck topos, cartesian closedness follows from the fact that
$\mathrm{Sh}(\mathcal{C},J)$ is a reflective subcategory of the
cartesian closed presheaf category, and the reflection preserves
finite products.
\end{proof}

\begin{proposition}\label{prop:topos-properties}
\index{elementary topos!properties}
Every elementary topos $\mathcal{E}$ satisfies:
\begin{enumerate}[label=(\alph*)]
\item $\mathcal{E}$ is balanced (every monic epic is an isomorphism);
\item Epimorphisms in $\mathcal{E}$ are coequalizers;
\item Every morphism $f\colon A\to B$ factors as a regular epimorphism
      followed by a monomorphism (image factorisation).
\end{enumerate}
\end{proposition}

\index{elementary topos|)}


%% ============================================================
\section{Internal logic}\label{sec:internal-logic}
%% ============================================================

\index{internal logic|(}
\index{Mitchell--B\'enabou language}

Every topos has an \emph{internal logic} that allows one to reason about
objects and morphisms using a language that resembles set-theoretic
reasoning, but is \emph{intuitionistic}: the law of excluded middle need
not hold.

\begin{definition}[Internal logic of a topos]\label{def:internal-logic}
\index{internal logic!definition}
The \textbf{internal logic} (or \textbf{Mitchell--B\'enabou language})
of a topos~$\mathcal{E}$ interprets:
\begin{itemize}
\item \emph{types} as objects of~$\mathcal{E}$;
\item \emph{terms} of type~$A$ in context~$\Gamma$ as morphisms
      $\Gamma\to A$;
\item \emph{propositions} in context~$\Gamma$ as morphisms $\Gamma\to\Omega$
      (subobjects of~$\Gamma$);
\item logical connectives via operations on~$\Omega$:
      $\wedge$, $\vee$, $\Rightarrow$, $\neg$, $\forall$, $\exists$.
\end{itemize}
\end{definition}

\begin{remark}\label{rem:internal-logic-intuitionistic}
\index{intuitionistic logic!in a topos}
\index{law of excluded middle}
\index{Boolean topos}
The internal logic of a topos is \textbf{intuitionistic higher-order logic}.
The law of excluded middle $\phi\vee\neg\phi$ holds internally if and only
if $\Omega\cong 1\sqcup 1$, in which case $\mathcal{E}$ is called a
\textbf{Boolean topos}.  The topos $\Set$ is Boolean; sheaf toposes on
non-trivial spaces are typically non-Boolean.
\end{remark}

\begin{remark}\label{rem:kripke-joyal}
\index{Kripke--Joyal semantics}
The precise semantics for interpreting formulae of the internal language
in a topos is called \textbf{Kripke--Joyal semantics} (or \emph{forcing}).
It provides the tool for transferring set-theoretic proofs to arbitrary
toposes.
\end{remark}

\index{internal logic|)}


%% ============================================================
\section{Geometric morphisms}\label{sec:geometric-morphisms}
%% ============================================================

\index{geometric morphism|(}

\begin{definition}[Geometric morphism]\label{def:geometric-morphism}
\index{geometric morphism!definition}
\index{direct image functor}
\index{inverse image functor}
A \textbf{geometric morphism} between toposes
$f\colon\mathcal{E}\to\mathcal{F}$ is an adjunction $f^*\dashv f_*$
where:
\begin{itemize}
\item $f^*\colon\mathcal{F}\to\mathcal{E}$ is called the \textbf{inverse
      image functor};
\item $f_*\colon\mathcal{E}\to\mathcal{F}$ is called the \textbf{direct
      image functor};
\item $f^*$ is required to be left exact (preserve finite limits).
\end{itemize}
\end{definition}

\begin{example}\label{ex:geometric-morphisms}
\index{geometric morphism!examples}
\begin{enumerate}[label=(\roman*)]
\item A continuous map $f\colon X\to Y$ of topological spaces induces
      a geometric morphism $f\colon\mathrm{Sh}(X)\to\mathrm{Sh}(Y)$
      with $f_*$ the direct image of sheaves and $f^*$ the inverse
      image.
\item For any topos $\mathcal{E}$, there is an essentially unique
      geometric morphism $\gamma\colon\mathcal{E}\to\Set$ given by
      $\gamma^* = \Delta$ (constant sheaf) and $\gamma_* = \Gamma$
      (global sections).
      \index{global sections}
      \index{constant sheaf functor}
\item The sheafification adjunction $a\dashv i$ for a site $(\mathcal{C},J)$
      is a geometric morphism
      $\mathrm{Sh}(\mathcal{C},J)\to\mathrm{PSh}(\mathcal{C})$.
\end{enumerate}
\end{example}

\begin{definition}[Categories of geometric morphisms]\label{def:geom-morph-cats}
\index{geometric morphism!natural transformation}
A \textbf{2-cell} between geometric morphisms $f,g\colon\mathcal{E}\to\mathcal{F}$
is a natural transformation $\alpha\colon f^*\Rightarrow g^*$ (equivalently,
$\beta\colon g_*\Rightarrow f_*$ by the mate correspondence).  This gives
a 2-category $\mathfrak{Top}$ of (Grothendieck) toposes.
\index{Top (2-category of toposes)@$\mathfrak{Top}$}
\end{definition}

\index{geometric morphism|)}


%% ============================================================
\section{Giraud's theorem}\label{sec:giraud}
%% ============================================================

\index{Giraud's theorem|(}

Giraud's theorem gives a purely categorical characterisation of
Grothendieck toposes, without reference to sites.

\begin{theorem}[Giraud]\label{thm:giraud}
\index{Giraud's theorem!statement}
A category $\mathcal{E}$ is a Grothendieck topos (i.e.\ equivalent to
$\mathrm{Sh}(\mathcal{C},J)$ for some site $(\mathcal{C},J)$)
if and only if it satisfies the following conditions:
\begin{enumerate}[label=(\roman*)]
\item $\mathcal{E}$ has all small colimits;
\item colimits in $\mathcal{E}$ are universal (stable under pullback);
      \index{universal colimit}
\item coproducts in $\mathcal{E}$ are disjoint;
      \index{disjoint coproduct}
\item every equivalence relation in $\mathcal{E}$ is effective
      (is the kernel pair of its coequaliser);
      \index{effective equivalence relation}
\item $\mathcal{E}$ has a small set of generators.
      \index{generator!small set of}
\end{enumerate}
\end{theorem}

\begin{remark}\label{rem:giraud-consequences}
\index{Giraud's theorem!consequences}
Giraud's theorem implies, in particular, that every Grothendieck topos:
\begin{itemize}
\item is an elementary topos;
\item has all small limits and colimits;
\item is a locally presentable category;
\item satisfies the axiom of choice internally if and only if every
      epimorphism splits.
      \index{axiom of choice!in a topos}
\end{itemize}
\end{remark}

\index{Giraud's theorem|)}


%% ============================================================
\section{Exercises}\label{sec:topos-exercises}
%% ============================================================

\begin{exercise}\label{ex:sieves-subfunctors}
Show that sieves on an object~$C$ in a category~$\mathcal{C}$ are in
bijection with subfunctors of the representable presheaf $\Hom(-,C)$.
\end{exercise}

\begin{exercise}\label{ex:presheaf-topos}
\index{presheaf topos}
Let $\mathcal{C}$ be a small category.  Show that $[\mathcal{C}^{\op}, \Set]$
is a Grothendieck topos (for the trivial topology).
Identify the subobject classifier explicitly.
\end{exercise}

\begin{exercise}\label{ex:omega-internal-heyting}
\index{Heyting algebra}
\index{subobject classifier!Heyting algebra structure}
Let $\mathcal{E}$ be an elementary topos.  Show that the subobject
classifier $\Omega$ carries the structure of an \emph{internal Heyting
algebra}: there exist morphisms
$\wedge, \vee, \Rightarrow\colon\Omega\times\Omega\to\Omega$
satisfying the Heyting algebra axioms internally.
\end{exercise}

\begin{exercise}\label{ex:geometric-morphism-compose}
Show that geometric morphisms compose: if
$f\colon\mathcal{E}\to\mathcal{F}$ and $g\colon\mathcal{F}\to\mathcal{G}$
are geometric morphisms, then so is $g\circ f$ with
$(g\circ f)^* = f^*\circ g^*$ and $(g\circ f)_* = g_*\circ f_*$.
\end{exercise}

\begin{exercise}\label{ex:sheaf-condition-equaliser}
\index{sheaf condition!equaliser form}
Let $(\mathcal{C},J)$ be a site and $\{f_i\colon C_i\to C\}_{i\in I}$
a covering family.  Show that a presheaf~$F$ satisfies the sheaf condition
for this covering if and only if the diagram
\[
\begin{tikzcd}
F(C) \arrow[r] &
\displaystyle\prod_{i\in I} F(C_i)
  \arrow[r, shift left=0.5ex, "d_0"]
  \arrow[r, shift right=0.5ex, "d_1"'] &
\displaystyle\prod_{(i,j)\in I\times I} F(C_i\times_C C_j)
\end{tikzcd}
\]
is an equaliser, where $d_0$ and $d_1$ are induced by the two projections
from the fibre product.
\end{exercise}

\begin{exercise}\label{ex:boolean-topos}
\index{Boolean topos}
Show that a topos $\mathcal{E}$ is Boolean if and only if every subobject
has a complement, i.e.\ for every mono $m\colon S\rightarrowtail A$
there exists $m'\colon S'\rightarrowtail A$ with $S\sqcup S'\cong A$
and $S\cap S'\cong 0$.
\end{exercise}

\begin{exercise}\label{ex:etale-site}
\index{etale topology@\'etale topology}
Describe the \'etale site of $\mathrm{Spec}(\Z)$ and explain why
the resulting topos is related to the absolute Galois group
$\mathrm{Gal}(\overline{\Q}/\Q)$.
\end{exercise}

\index{Grothendieck topos|)}
\index{topos|)}


%%%%%%%%%%%%%%%%%%%%%%%%%%%%%%%%%%%%%%%%%%%%%%%%%%%%%%%%%%%%
%%   CHAPTER 11: ∞-CATEGORIES — OVERVIEW                  %%
%%%%%%%%%%%%%%%%%%%%%%%%%%%%%%%%%%%%%%%%%%%%%%%%%%%%%%%%%%%%

\chapter{$\infty$-Categories---Overview}\label{ch:infinity-cats}
\minitoc

\vspace{1em}

Classical category theory considers objects, morphisms between objects,
and equalities between morphisms.  But in many contexts---homotopy theory,
homological algebra, derived algebraic geometry---one needs morphisms
between morphisms (2-cells), morphisms between those (3-cells), and so on,
ad infinitum, with all compositions associative and unital only up to
coherent higher-dimensional cells.  This is the domain of
\emph{higher category theory} and, in particular,
\emph{$(\infty,1)$-categories}.

This chapter provides a survey of the main ideas, definitions, and
results.  Proofs are mostly omitted; we aim to give the reader a map
of the landscape and precise pointers to the literature.

\index{infinity-category@$\infty$-category|(}
\index{higher category theory|(}

%% ============================================================
\section{Motivation}\label{sec:infty-motivation}
%% ============================================================

\index{infinity-category@$\infty$-category!motivation}

The need for higher categories arises from several independent sources.

\begin{enumerate}[label=(\roman*)]
\item \textbf{Homotopy theory.}
\index{homotopy theory}
The homotopy category $\mathrm{Ho}(\Top)$ of topological spaces
loses too much information: mapping spaces are truncated to sets.
One wants a ``category'' where $\Hom(X,Y)$ is a space (or a
simplicial set), composition is associative up to homotopy, and
all higher coherences are recorded.

\item \textbf{Derived categories.}
\index{derived category}
The derived category $D(\mathcal{A})$ of an abelian category~$\mathcal{A}$
is obtained by formally inverting quasi-isomorphisms, but
the resulting category is poorly behaved (e.g.\ it lacks functorial
cones).  The $\infty$-categorical enhancement---the
\emph{derived $\infty$-category}---remedies these defects.

\item \textbf{Stacks and higher stacks.}
\index{stack!higher}
In algebraic geometry, stacks are ``sheaves of groupoids.''
Higher stacks are sheaves of $\infty$-groupoids, and their natural
home is the $\infty$-topos.

\item \textbf{Topological quantum field theory.}
\index{topological quantum field theory}
The cobordism hypothesis of Baez--Dolan (proved by Lurie) requires
the language of $(\infty,n)$-categories.
\end{enumerate}


%% ============================================================
\section{Quasi-categories}\label{sec:quasi-categories}
%% ============================================================

\index{quasi-category|(}
\index{Joyal}
\index{Boardman--Vogt}

The most developed model for $(\infty,1)$-categories is that of
\emph{quasi-categories}, introduced by Boardman--Vogt and extensively
developed by Joyal and Lurie.

\begin{definition}[Simplicial set]\label{def:simplicial-set-reminder}
\index{simplicial set}
\index{simplex category}
Let $\Delta$ denote the simplex category (finite non-empty ordinals and
order-preserving maps).  A \textbf{simplicial set} is a functor
$X\colon\Delta^{\op}\to\Set$.  We write $X_n = X([n])$ for the set of
$n$-simplices.  The category of simplicial sets is
$\mathrm{sSet} = [\Delta^{\op}, \Set]$.
\index{sSet@$\mathrm{sSet}$}
\end{definition}

\begin{definition}[Horn]\label{def:horn}
\index{horn}
\index{horn!inner}
\index{horn!outer}
For $0\le k\le n$, the \textbf{$k$-th horn} $\Lambda^n_k\subset\Delta^n$
is the simplicial subset obtained by removing the interior of~$\Delta^n$
and the face opposite vertex~$k$.  The horn is called \textbf{inner} if
$0 < k < n$ and \textbf{outer} otherwise.
\end{definition}

\begin{definition}[Quasi-category]\label{def:quasi-category}
\index{quasi-category!definition}
\index{inner horn!filler}
\index{weak Kan complex}
A simplicial set $X$ is a \textbf{quasi-category} (or \textbf{$\infty$-category},
or \textbf{weak Kan complex}) if every inner horn has a filler:
for every $0 < k < n$ and every map $\Lambda^n_k\to X$, there exists
an extension $\Delta^n\to X$:
\[
\begin{tikzcd}
\Lambda^n_k \arrow[r] \arrow[d, hook] & X \\
\Delta^n \arrow[ur, dashed] &
\end{tikzcd}
\]
Note that the filler is \emph{not required to be unique}.
\end{definition}

\begin{remark}\label{rem:kan-vs-quasi}
\index{Kan complex}
\index{infinity-groupoid@$\infty$-groupoid}
If one requires fillers for \emph{all} horns (inner and outer), one
obtains a \textbf{Kan complex}, which models an \emph{$\infty$-groupoid}:
every 1-simplex is invertible up to higher homotopy.  A quasi-category
is to an $\infty$-category what a Kan complex is to an $\infty$-groupoid.
\end{remark}

\begin{example}\label{ex:quasi-cat-examples}
\index{quasi-category!examples}
\begin{enumerate}[label=(\roman*)]
\item The nerve $N(\mathcal{C})$ of an ordinary category~$\mathcal{C}$ is
      a quasi-category in which inner horn fillers are \emph{unique}.
      \index{nerve!as quasi-category}
\item The singular simplicial set $\mathrm{Sing}(X)$ of a topological
      space~$X$ is a Kan complex, hence a quasi-category.
      \index{singular simplicial set}
\item The \textbf{dg-nerve} of a dg-category is a quasi-category.
      \index{dg-nerve}
\item The \textbf{homotopy coherent nerve} of a simplicial category is a
      quasi-category.
      \index{homotopy coherent nerve}
\end{enumerate}
\end{example}

\begin{definition}[Homotopy category of a quasi-category]\label{def:homotopy-cat-qcat}
\index{homotopy category!of a quasi-category}
Given a quasi-category~$X$, its \textbf{homotopy category} $\mathrm{h}X$
is the ordinary category with:
\begin{itemize}
\item objects: the vertices ($0$-simplices) of~$X$;
\item morphisms $x\to y$: homotopy classes of edges ($1$-simplices)
      from~$x$ to~$y$, where two edges are homotopic if they are
      related by a $2$-simplex with degenerate third edge;
\item composition: given by choosing fillers for inner $2$-horns.
\end{itemize}
\end{definition}

\index{quasi-category|)}


%% ============================================================
\section{Models for $(\infty,1)$-categories}\label{sec:models-infty1}
%% ============================================================

\index{(infinity,1)-category@$(\infty,1)$-category|(}

There are several equivalent models for $(\infty,1)$-categories.

\begin{definition}[$(\infty,1)$-category---informal]\label{def:infty1-informal}
\index{(infinity,1)-category@$(\infty,1)$-category!informal definition}
An \textbf{$(\infty,1)$-category} is a ``category enriched in spaces''---a
structure with objects, morphisms, homotopies between morphisms,
homotopies between homotopies, etc., where all $k$-morphisms for $k\ge 2$
are invertible (up to higher morphisms).
\end{definition}

\begin{remark}\label{rem:models-infty1}
\index{(infinity,1)-category@$(\infty,1)$-category!models}
The principal models, all Quillen equivalent, include:
\begin{enumerate}[label=(\roman*)]
\item \textbf{Quasi-categories} (simplicial sets with inner horn fillers)
      --- Joyal, Lurie.
      \index{quasi-category}
\item \textbf{Complete Segal spaces} --- Rezk.
      \index{complete Segal space}
      \index{Rezk}
\item \textbf{Segal categories} --- Hirschowitz--Simpson.
      \index{Segal category}
\item \textbf{Simplicial categories} ($\mathrm{sSet}$-enriched categories)
      --- Bergner.
      \index{simplicial category}
      \index{Bergner}
\item \textbf{Relative categories} (categories with weak equivalences)
      --- Barwick--Kan.
      \index{relative category}
\end{enumerate}
The fact that these models are all equivalent is a deep theorem
(the ``comparison theorem''), proved using Quillen model structures.
\index{comparison theorem}
\end{remark}

\index{(infinity,1)-category@$(\infty,1)$-category|)}


%% ============================================================
\section{The $\infty$-Yoneda lemma}\label{sec:infty-yoneda}
%% ============================================================

\index{Yoneda lemma!infinity version@$\infty$-version|(}

The Yoneda lemma generalises to the $\infty$-categorical setting.

\begin{definition}[$\infty$-presheaf]\label{def:infty-presheaf}
\index{infinity-presheaf@$\infty$-presheaf}
Let $\mathcal{C}$ be a small $(\infty,1)$-category.
The \textbf{$\infty$-category of presheaves} on~$\mathcal{C}$ is
$\mathcal{P}(\mathcal{C}) = \Fun(\mathcal{C}^{\op}, \mathcal{S})$,
where $\mathcal{S}$ denotes the $\infty$-category of spaces
(i.e.\ $\infty$-groupoids, modelled by Kan complexes).
\index{S@$\mathcal{S}$!infinity-category of spaces@$\infty$-category of spaces}
\end{definition}

\begin{theorem}[$\infty$-Yoneda lemma]\label{thm:infty-yoneda}
\index{Yoneda lemma!infinity version@$\infty$-version}
\index{Yoneda embedding!infinity version@$\infty$-version}
Let $\mathcal{C}$ be a small $(\infty,1)$-category.  The Yoneda embedding
\[
    \mathbf{y}\colon \mathcal{C}\hookrightarrow\mathcal{P}(\mathcal{C}),
    \quad C\mapsto\mathrm{Map}_\mathcal{C}(-,C)
\]
is fully faithful.  Moreover, for any presheaf
$F\in\mathcal{P}(\mathcal{C})$ and any $C\in\mathcal{C}$,
\[
    \mathrm{Map}_{\mathcal{P}(\mathcal{C})}(\mathbf{y}(C), F)
    \;\simeq\; F(C)
\]
as objects of~$\mathcal{S}$ (an equivalence of spaces).
\end{theorem}

\begin{remark}\label{rem:infty-yoneda-universal}
\index{free cocompletion!infinity-categorical@$\infty$-categorical}
The $\infty$-category $\mathcal{P}(\mathcal{C})$ is the
\textbf{free cocompletion} of~$\mathcal{C}$ under small colimits,
just as in the 1-categorical case.  Every $\infty$-presheaf is a
colimit of representables.
\end{remark}

\index{Yoneda lemma!infinity version@$\infty$-version|)}


%% ============================================================
\section{$\infty$-adjunctions and $\infty$-limits}\label{sec:infty-adj}
%% ============================================================

\index{infinity-adjunction@$\infty$-adjunction|(}

\begin{definition}[$\infty$-adjunction]\label{def:infty-adjunction}
\index{infinity-adjunction@$\infty$-adjunction!definition}
An \textbf{$\infty$-adjunction} between $(\infty,1)$-categories
$\mathcal{C}$ and $\mathcal{D}$ is a pair of functors
$F\colon\mathcal{C}\to\mathcal{D}$, $G\colon\mathcal{D}\to\mathcal{C}$
together with an equivalence of mapping spaces
\[
    \mathrm{Map}_\mathcal{D}(FC, D)
    \;\simeq\;
    \mathrm{Map}_\mathcal{C}(C, GD)
\]
natural in~$C$ and~$D$.
\end{definition}

\begin{remark}\label{rem:infty-adjunction-properties}
\index{infinity-adjunction@$\infty$-adjunction!properties}
The theory of $\infty$-adjunctions parallels the 1-categorical theory:
\begin{enumerate}[label=(\roman*)]
\item $\infty$-adjoints are unique up to contractible choice;
\item an $\infty$-left adjoint preserves $\infty$-colimits;
\item an $\infty$-right adjoint preserves $\infty$-limits;
\item the $\infty$-adjoint functor theorem holds (under appropriate
      presentability hypotheses).
      \index{adjoint functor theorem!infinity version@$\infty$-version}
\end{enumerate}
\end{remark}

\begin{definition}[$\infty$-limit]\label{def:infty-limit}
\index{infinity-limit@$\infty$-limit}
\index{homotopy limit}
Let $F\colon\mathcal{J}\to\mathcal{C}$ be a diagram in an
$(\infty,1)$-category~$\mathcal{C}$.  An \textbf{$\infty$-limit}
of~$F$ is an object $L\in\mathcal{C}$ together with equivalences
\[
    \mathrm{Map}_\mathcal{C}(X, L)
    \;\simeq\;
    \lim_{j\in\mathcal{J}}\, \mathrm{Map}_\mathcal{C}(X, Fj)
\]
natural in~$X$, where the right-hand side is a homotopy limit of spaces.
Dually for $\infty$-colimits.
\end{definition}

\index{infinity-adjunction@$\infty$-adjunction|)}


%% ============================================================
\section{Higher toposes}\label{sec:higher-toposes}
%% ============================================================

\index{infinity-topos@$\infty$-topos|(}

\begin{definition}[$\infty$-topos]\label{def:infty-topos}
\index{infinity-topos@$\infty$-topos!definition}
An \textbf{$\infty$-topos} is an $(\infty,1)$-category that is an
accessible left exact localisation of an $\infty$-presheaf category
$\mathcal{P}(\mathcal{C})$ for some small $(\infty,1)$-category~$\mathcal{C}$.
\end{definition}

\begin{remark}\label{rem:infty-topos-giraud}
\index{infinity-topos@$\infty$-topos!Giraud-type characterisation}
Lurie proves an $\infty$-categorical analogue of Giraud's theorem:
an $(\infty,1)$-category $\mathcal{E}$ is an $\infty$-topos if and only if:
\begin{enumerate}[label=(\roman*)]
\item $\mathcal{E}$ is presentable (cocomplete and accessible);
\item colimits in $\mathcal{E}$ are universal;
\item groupoid objects in $\mathcal{E}$ are effective.
\end{enumerate}
Note that the disjointness condition for coproducts, present in Giraud's
classical theorem, is automatic in the $\infty$-setting.
\end{remark}

\begin{example}\label{ex:infty-toposes}
\index{infinity-topos@$\infty$-topos!examples}
\begin{enumerate}[label=(\roman*)]
\item The $\infty$-category $\mathcal{S}$ of spaces (Kan complexes) is
      the initial $\infty$-topos, playing the role of~$\Set$.
\item For a topological space~$X$, the $\infty$-category
      $\mathrm{Sh}_\infty(X)$ of $\infty$-sheaves (sheaves of spaces)
      on~$X$ is an $\infty$-topos.
      \index{infinity-sheaf@$\infty$-sheaf}
\item Derived algebraic geometry uses $\infty$-toposes as the ambient
      framework for ``derived stacks.''
      \index{derived algebraic geometry}
\end{enumerate}
\end{example}

\begin{remark}\label{rem:hott}
\index{homotopy type theory}
\index{univalence axiom}
$\infty$-toposes provide the semantic models for \textbf{homotopy type
theory} (HoTT).  The univalence axiom of Voevodsky holds in every
$\infty$-topos, and the internal language of an $\infty$-topos is
(conjecturally) homotopy type theory.  This is an active area of research
connecting category theory, homotopy theory, and the foundations of
mathematics.
\end{remark}

\index{infinity-topos@$\infty$-topos|)}


%% ============================================================
\section{Exercises}\label{sec:infty-exercises}
%% ============================================================

\begin{exercise}\label{ex:nerve-quasi}
\index{nerve!inner horn condition}
Let $\mathcal{C}$ be an ordinary category.  Show that the nerve
$N(\mathcal{C})$ is a quasi-category and that inner horn fillers
are unique.  Deduce that $N$ is fully faithful as a functor
$\Cat\to\mathrm{sSet}$.
\end{exercise}

\begin{exercise}\label{ex:kan-groupoid}
\index{Kan complex!groupoid property}
Show that a quasi-category~$X$ is a Kan complex if and only if every
morphism (1-simplex) in~$X$ is an equivalence (invertible up to a
2-simplex).
\end{exercise}

\begin{exercise}\label{ex:homotopy-cat-nerve}
\index{homotopy category!of a nerve}
Let $\mathcal{C}$ be an ordinary category.  Show that
$\mathrm{h}(N(\mathcal{C})) \cong \mathcal{C}$.
\end{exercise}

\begin{exercise}\label{ex:infty-presheaf-colim}
\index{infinity-presheaf@$\infty$-presheaf!as colimit of representables}
Let $\mathcal{C}$ be a small $(\infty,1)$-category and
$F\in\mathcal{P}(\mathcal{C})$ an $\infty$-presheaf.  Using the
$\infty$-Yoneda lemma, show that~$F$ is a colimit of representable
presheaves.
\end{exercise}

\begin{exercise}\label{ex:infty-adj-unique}
Show that if $F\colon\mathcal{C}\to\mathcal{D}$ is a functor of
$(\infty,1)$-categories that admits a right adjoint~$G$, then $G$ is
unique up to a contractible space of choices.
\textit{Hint:} use the $\infty$-Yoneda lemma to express $G(D)$ via
the presheaf $C\mapsto\mathrm{Map}_\mathcal{D}(FC, D)$.
\end{exercise}

\begin{exercise}\label{ex:infty-topos-slice}
\index{infinity-topos@$\infty$-topos!slice}
Let $\mathcal{E}$ be an $\infty$-topos and $X\in\mathcal{E}$ an object.
Show that the slice $\infty$-category $\mathcal{E}_{/X}$ is again an
$\infty$-topos.
\end{exercise}

\begin{exercise}[Descent]\label{ex:infty-descent}
\index{descent}
\index{infinity-topos@$\infty$-topos!descent}
Let $\mathcal{E}$ be an $\infty$-topos and $f\colon X\to Y$ a morphism
in~$\mathcal{E}$.  Show that the pullback functor
$f^*\colon\mathcal{E}_{/Y}\to\mathcal{E}_{/X}$ preserves all colimits.
This is the $\infty$-categorical formulation of \textbf{descent}.
\end{exercise}

\index{infinity-category@$\infty$-category|)}
\index{higher category theory|)}


%%%%%%%%%%%%%%%%%%%%%%%%%%%%%%%%%%%%%%%%%%%%%%%%%%%%%%%%%%%%
%%                  END MATTER                            %%
%%%%%%%%%%%%%%%%%%%%%%%%%%%%%%%%%%%%%%%%%%%%%%%%%%%%%%%%%%%%

\appendix

\chapter{Review of Algebra and Topology}\label{app:review}
\minitoc

\vspace{1em}

We collect here the basic definitions from algebra and topology that
are used throughout the text.  The reader familiar with these notions
may safely skip this appendix and refer back to it as needed.

%% ============================================================
\section{Groups, rings, and modules}\label{sec:app-algebra}
%% ============================================================

\begin{definition}[Group]\label{def:app-group}
\index{group!definition}
A \textbf{group} is a set~$G$ equipped with a binary operation
$\cdot\colon G\times G\to G$, an identity element $e\in G$, and an
inverse map $(-)^{-1}\colon G\to G$, satisfying associativity,
identity, and inverse laws.  A group is \textbf{abelian} if
$a\cdot b = b\cdot a$ for all $a,b\in G$.
\end{definition}

\begin{definition}[Ring]\label{def:app-ring}
\index{ring!definition}
A \textbf{ring} $(R, +, \cdot, 0, 1)$ is an abelian group $(R,+,0)$
equipped with an associative multiplication with unit~$1$, such that
multiplication distributes over addition.  A ring is \textbf{commutative}
if $ab = ba$ for all $a,b\in R$.
\end{definition}

\begin{definition}[Module]\label{def:app-module}
\index{module!definition}
Let $R$ be a ring.  A \textbf{left $R$-module} is an abelian group $(M,+)$
equipped with a scalar multiplication $R\times M\to M$ satisfying the
usual axioms: $r(m+n) = rm+rn$, $(r+s)m = rm+sm$, $(rs)m = r(sm)$,
$1m = m$.
\end{definition}

\begin{definition}[Field]\label{def:app-field}
\index{field!definition}
A \textbf{field} is a commutative ring~$\K$ in which every non-zero
element is invertible, i.e.\ $\K^\times = \K\setminus\{0\}$.
\end{definition}

\begin{definition}[Algebra over a field]\label{def:app-algebra}
\index{algebra!over a field}
A \textbf{$\K$-algebra} is a ring~$A$ that is also a $\K$-vector space,
with multiplication $\K$-bilinear.  Equivalently, it is a monoid object
in $(\Vect_\K, \otimes_\K, \K)$.
\end{definition}


%% ============================================================
\section{Topological spaces}\label{sec:app-topology}
%% ============================================================

\begin{definition}[Topological space]\label{def:app-top-space}
\index{topological space!definition}
A \textbf{topological space} is a pair $(X, \tau)$ where $X$ is a set
and $\tau\subseteq\mathcal{P}(X)$ is a collection of subsets (called
\emph{open sets}) satisfying: (i)~$\varnothing, X\in\tau$;
(ii)~$\tau$ is closed under arbitrary unions; (iii)~$\tau$ is closed
under finite intersections.
\end{definition}

\begin{definition}[Continuous map]\label{def:app-continuous}
\index{continuous map}
A map $f\colon X\to Y$ between topological spaces is \textbf{continuous}
if $f^{-1}(U)\in\tau_X$ for every $U\in\tau_Y$.
\end{definition}

\begin{definition}[Hausdorff space]\label{def:app-hausdorff}
\index{Hausdorff space}
A topological space~$X$ is \textbf{Hausdorff} (or $T_2$) if for every
pair of distinct points $x\ne y$ there exist disjoint open sets
$U\ni x$ and $V\ni y$.
\end{definition}

\begin{definition}[Compact space]\label{def:app-compact}
\index{compact space}
A topological space~$X$ is \textbf{compact} if every open cover has a
finite subcover.
\end{definition}

\begin{definition}[Connected space]\label{def:app-connected}
\index{connected space}
A topological space~$X$ is \textbf{connected} if it cannot be written
as a disjoint union of two non-empty open sets.
\end{definition}


%% ============================================================
\section{Simplicial sets}\label{sec:app-simplicial}
%% ============================================================

\begin{definition}[Simplex category]\label{def:app-simplex-cat}
\index{simplex category!definition}
The \textbf{simplex category}~$\Delta$ has as objects the finite
ordinals $[n] = \{0, 1, \ldots, n\}$ for $n\ge 0$, and as morphisms
the order-preserving maps.
\end{definition}

\begin{definition}[Simplicial set]\label{def:app-sset}
\index{simplicial set!definition}
A \textbf{simplicial set} is a functor $X\colon\Delta^{\op}\to\Set$.
It is specified by sets $X_n$ (the $n$-simplices) together with
face maps $d_i\colon X_n\to X_{n-1}$ and degeneracy maps
$s_i\colon X_n\to X_{n+1}$ satisfying the simplicial identities.
\end{definition}

\begin{definition}[Geometric realisation]\label{def:app-realisation}
\index{geometric realisation}
The \textbf{geometric realisation} of a simplicial set~$X$ is
\[
    |X| = \left(\coprod_{n\ge 0} X_n\times\Delta^n_{\mathrm{top}}\right)
    \Big/ {\sim}
\]
where $\Delta^n_{\mathrm{top}}$ is the topological $n$-simplex
and the equivalence relation is generated by the face and degeneracy maps.
\end{definition}


%% ============================================================
\section{Basic homological algebra}\label{sec:app-homological}
%% ============================================================

\begin{definition}[Chain complex]\label{def:app-chain-complex}
\index{chain complex!definition}
A \textbf{chain complex} over a ring~$R$ is a sequence of $R$-modules
and homomorphisms
\[
    \cdots \xrightarrow{d_{n+1}} C_n \xrightarrow{d_n} C_{n-1}
    \xrightarrow{d_{n-1}} \cdots
\]
with $d_n\circ d_{n+1} = 0$ for all~$n$.
The \textbf{homology} is $H_n(C_\bullet) = \ker d_n / \mathrm{im}\, d_{n+1}$.
\index{homology}
\end{definition}

\begin{definition}[Exact sequence]\label{def:app-exact}
\index{exact sequence}
A chain complex is \textbf{exact} at~$C_n$ if $H_n = 0$, i.e.\
$\ker d_n = \mathrm{im}\, d_{n+1}$.  A \textbf{short exact sequence}
is an exact sequence $0\to A\to B\to C\to 0$.
\end{definition}

\begin{definition}[Tensor product of modules]\label{def:app-tensor}
\index{tensor product!of modules}
For a commutative ring~$R$ and $R$-modules $M,N$, the \textbf{tensor
product} $M\otimes_R N$ is the $R$-module generated by symbols
$m\otimes n$ ($m\in M$, $n\in N$) subject to bilinearity relations.
It is characterised by the universal property:
$\Hom_R(M\otimes_R N, P) \cong \mathrm{Bilin}_R(M\times N, P)$.
\end{definition}

\begin{definition}[Projective and injective modules]\label{def:app-proj-inj}
\index{projective module}
\index{injective module}
An $R$-module~$P$ is \textbf{projective} if $\Hom_R(P,-)$ is exact.
An $R$-module~$Q$ is \textbf{injective} if $\Hom_R(-,Q)$ is exact.
\end{definition}


%% ============================================================
\section{Posets and lattices}\label{sec:app-posets}
%% ============================================================

\begin{definition}[Partially ordered set]\label{def:app-poset}
\index{partially ordered set}
\index{poset|see{partially ordered set}}
A \textbf{partially ordered set} (or \textbf{poset}) is a set~$P$
equipped with a relation~$\le$ that is reflexive, antisymmetric, and
transitive.  A poset may be viewed as a category with objects~$P$
and a unique morphism $x\to y$ whenever $x\le y$.
\end{definition}

\begin{definition}[Lattice]\label{def:app-lattice}
\index{lattice}
A \textbf{lattice} is a poset in which every pair of elements $\{a,b\}$
has a meet $a\wedge b$ (greatest lower bound) and a join $a\vee b$
(least upper bound).  A \textbf{complete lattice} has meets and joins
for all subsets.
\index{complete lattice}
\end{definition}

\begin{definition}[Heyting algebra]\label{def:app-heyting}
\index{Heyting algebra!definition}
A \textbf{Heyting algebra} is a lattice~$H$ with a binary operation
$\Rightarrow$ (implication) satisfying: $a\wedge b\le c$ if and only
if $a\le (b\Rightarrow c)$.  Every Boolean algebra is a Heyting algebra,
but not conversely.
\end{definition}


%% ============================================================
%% BIBLIOGRAPHY
%% ============================================================

\begin{thebibliography}{99}
\bibitem{MacLaneCWM} Mac\,Lane, S., \emph{Categories for the Working Mathematician}, 2nd ed., Springer, 1998.
\bibitem{MacLane1963} Mac\,Lane, S., ``Natural associativity and commutativity,'' \emph{Rice Univ. Studies}, 49(4), 28--46, 1963.
\bibitem{freyd1964} Freyd, P., \emph{Abelian Categories: An Introduction to the Theory of Functors}, Harper \& Row, 1964.
\bibitem{mitchell1965} Mitchell, B., \emph{Theory of Categories}, Academic Press, 1965.
\bibitem{weibel1994} Weibel, C.A., \emph{An Introduction to Homological Algebra}, Cambridge University Press, 1994.
\bibitem{borceux} Borceux, F., \emph{Handbook of Categorical Algebra} (3 volumes), Cambridge University Press, 1994.
\bibitem{riehl} Riehl, E., \emph{Category Theory in Context}, Dover, 2016.
\bibitem{awodey} Awodey, S., \emph{Category Theory}, 2nd ed., Oxford University Press, 2010.
\bibitem{kashiwara} Kashiwara, M. and Schapira, P., \emph{Categories and Sheaves}, Springer, 2006.
\end{thebibliography}

%% ============================================================
%% INDEX AND CLOSING
%% ============================================================

\cleardoublepage
\phantomsection
\addcontentsline{toc}{chapter}{Index}
\printindex

\end{document}
